% !TEX program = lualatex
%
% AI-Ready Finance Whitepaper (日本語)
% Compile: lualatex whitepaper.tex (2回実行してTOC/参照を確定させる)
%
\documentclass[11pt,a4paper]{ltjsarticle}

% ---------- フォント ----------
\usepackage{luatexja-fontspec}
\setmainjfont{Hiragino Mincho ProN}
\setsansjfont{Hiragino Kaku Gothic Pro}
\setmainfont{Times New Roman}
\setsansfont{Helvetica Neue}

% ---------- レイアウト ----------
\usepackage[a4paper,top=25mm,bottom=28mm,left=22mm,right=22mm]{geometry}
\usepackage{setspace}
\onehalfspacing
\usepackage{parskip}
\setlength{\parskip}{0.6\baselineskip}

% ---------- 色 ----------
\usepackage{xcolor}
\definecolor{brandnavy}{HTML}{0B2545}
\definecolor{brandblue}{HTML}{1565C0}
\definecolor{brandgray}{HTML}{4A4A4A}
\definecolor{brandlight}{HTML}{EEF3FA}
\definecolor{brandaccent}{HTML}{C9972C}

% ---------- 見出し ----------
\usepackage{titlesec}
\titleformat{\section}
  {\normalfont\Large\bfseries\color{brandnavy}}
  {第\thesection 部}{1em}{}
\titlespacing*{\section}{0pt}{3.2ex plus 1ex minus .2ex}{2ex}

\titleformat{\subsection}
  {\normalfont\large\bfseries\color{brandblue}}
  {\thesubsection}{0.8em}{}
\titlespacing*{\subsection}{0pt}{2.6ex plus 1ex minus .2ex}{1.2ex}

\titleformat{\subsubsection}
  {\normalfont\normalsize\bfseries\color{brandgray}}
  {\thesubsubsection}{0.6em}{}
\titlespacing*{\subsubsection}{0pt}{2ex plus .8ex minus .2ex}{1ex}

% ---------- ヘッダー/フッター ----------
\usepackage{fancyhdr}
\pagestyle{fancy}
\fancyhf{}
\renewcommand{\headrulewidth}{0.4pt}
\renewcommand{\footrulewidth}{0pt}
\fancyhead[L]{\small\color{brandgray} AI-Ready Finance ホワイトペーパー}
\fancyhead[R]{\small\color{brandgray} \leftmark}
\fancyfoot[C]{\small\color{brandgray} \thepage}

\fancypagestyle{plain}{
  \fancyhf{}
  \fancyfoot[C]{\small\color{brandgray} \thepage}
  \renewcommand{\headrulewidth}{0pt}
}

% ---------- 表・図 ----------
\usepackage{booktabs}
\usepackage{longtable}
\usepackage{array}
\usepackage{graphicx}
\usepackage{tabularx}
\renewcommand{\arraystretch}{1.25}

% ---------- 図解 (TikZ / pgfplots) ----------
\usepackage{tikz}
\usetikzlibrary{arrows.meta,positioning,shapes.geometric,fit,backgrounds,decorations.pathreplacing,calc}
\usepackage{pgfplots}
\pgfplotsset{compat=1.18}
% 図の共通カラーパレット（本文の brandnavy/brandblue/brandaccent と統一）
\colorlet{figblue}{brandblue}
\colorlet{fignavy}{brandnavy}
\colorlet{figaccent}{brandaccent}
\colorlet{figgray}{brandgray}
\colorlet{figlight}{brandlight}

% ---------- 図版の共通スタイル（統一された洗練された見た目） ----------
% すべての図はこのスタイル群を土台に組む：ソフトな影・角丸・統一された矢印。
% 個々の図では text width / minimum height / align などを追記して調整する。
\usetikzlibrary{shadows.blur}
\tikzset{
  % 既定の矢先を上品な Stealth に統一
  >={Stealth[length=2.4mm,width=2.2mm]},
  % ソフトなドロップシャドウ（浮き上がる立体感）
  figshadow/.style={blur shadow={shadow blur steps=6, shadow blur radius=2pt,
    shadow xshift=0pt, shadow yshift=-1pt, shadow opacity=20}},
  % --- ノード（枠）: 用途別の 5 種 ---
  fnode/.style={draw=fignavy!85, line width=0.7pt, rounded corners=4pt,
    fill=figlight, figshadow, align=center, inner sep=2.8mm, font=\small},
  fnode blue/.style={draw=figblue, line width=0.7pt, rounded corners=4pt,
    fill=figblue!7, figshadow, align=center, inner sep=2.8mm, font=\small},
  fnode accent/.style={draw=figaccent!85!black, line width=0.8pt, rounded corners=4pt,
    fill=white, figshadow, align=center, inner sep=2.8mm, font=\small},
  fnode dark/.style={draw=fignavy, line width=0.8pt, rounded corners=4pt,
    fill=fignavy, text=white, figshadow, align=center, inner sep=2.8mm, font=\small},
  fnode muted/.style={draw=figgray!70, line width=0.6pt, rounded corners=4pt,
    fill=figgray!8, figshadow, align=center, inner sep=2.8mm, font=\small},
  % --- 見出しチップ（列ヘッダ等） ---
  fchip/.style={rounded corners=3pt, figshadow, align=center,
    font=\bfseries\sffamily\small, text=white, inner sep=2mm},
  % --- 矢印: 用途別 ---
  farr/.style={-{Stealth[length=2.6mm,width=2.4mm]}, line width=0.9pt, draw=fignavy!70},
  farr blue/.style={-{Stealth[length=2.6mm,width=2.4mm]}, line width=0.9pt, draw=figblue},
  farr accent/.style={-{Stealth[length=2.6mm,width=2.4mm]}, line width=0.9pt, draw=figaccent},
  farr dash/.style={-{Stealth[length=2.4mm,width=2.2mm]}, line width=0.8pt, draw=figgray, dashed},
  fbiarr/.style={{Stealth[length=2.4mm,width=2.2mm]}-{Stealth[length=2.4mm,width=2.2mm]},
    line width=0.9pt, draw=figaccent},
}

% ---------- 箇条書き ----------
\usepackage{enumitem}
\setlist[itemize]{leftmargin=1.6em, itemsep=0.35em, topsep=0.3em}
\setlist[enumerate]{leftmargin=1.8em, itemsep=0.35em, topsep=0.3em}

% ---------- ボックス ----------
\usepackage[most]{tcolorbox}
\tcbset{
  casebox/.style={
    enhanced, breakable, colback=brandlight, colframe=brandblue,
    boxrule=0.6pt, arc=1.5mm, left=3mm, right=3mm, top=2mm, bottom=2mm,
    fonttitle=\bfseries\sffamily, coltitle=white,
    colbacktitle=brandblue, title={#1}
  },
  insightbox/.style={
    enhanced, breakable, colback=white, colframe=brandaccent,
    boxrule=0.7pt, arc=1mm, left=3mm, right=3mm, top=2mm, bottom=2mm,
    fonttitle=\bfseries\sffamily, coltitle=brandnavy,
    title={\textcolor{brandaccent}{\textbullet}\ #1}
  },
}
\newtcolorbox{casestudy}[1]{casebox={#1}}
\newtcolorbox{insight}[1]{insightbox={#1}}

% コラム（初学者向けの概念解説ボックス）: casestudy(青)/insight(金)と視覚的に区別する
% ため、濃紺の見出し帯＋淡い背景＋左端の金色ストライプで構成する。用途は本文で
% 初めて登場する重要概念を、図解付きで自己完結的に噛み砕くこと。
\tcbset{
  conceptbox/.style={
    enhanced, breakable, colback=brandlight!45, colframe=brandnavy,
    boxrule=0.7pt, arc=1.5mm, left=3mm, right=3mm, top=2mm, bottom=2mm,
    fonttitle=\bfseries\sffamily, coltitle=white,
    colbacktitle=brandnavy, title={コラム｜#1},
    borderline west={2.4pt}{0pt}{brandaccent},
  },
}
\newtcolorbox{concept}[1]{conceptbox={#1}}

% ---------- リンク ----------
\usepackage{xurl}
\usepackage[bookmarksnumbered,hidelinks]{hyperref}
\hypersetup{
  colorlinks=true, linkcolor=brandblue, urlcolor=brandblue, citecolor=brandblue,
  pdftitle={AI-Ready 金融データ基盤とAI-Native投資運用の最前線},
  pdfauthor={株式会社INDX 伊藤克哉 (k1ito)},
}

% ---------- 参考文献 (biblatex + biber) ----------
\usepackage[backend=biber,style=numeric,sorting=nyt,maxbibnames=99,urldate=iso]{biblatex}
\addbibresource{references.bib}

% ---------- 目次の見た目 ----------
\usepackage{tocloft}
\renewcommand{\cftsecfont}{\bfseries\color{brandnavy}}
\renewcommand{\cftsecpagefont}{\color{brandnavy}}

\begin{document}

% ============================================================
% 表紙
% ============================================================
\begin{titlepage}
\begin{center}
\vspace*{18mm}
{\color{brandaccent}\rule{\textwidth}{1.2pt}}\\[6mm]
{\sffamily\color{brandgray}\large AI-READY FINANCE RESEARCH WHITEPAPER}\\[10mm]
{\sffamily\bfseries\color{brandnavy}\LARGE AI-Ready金融データ基盤と\\[2mm] AI-Native投資運用の最前線}\\[6mm]
{\large\color{brandgray} 学術研究・大手金融機関・スタートアップ事例に見る\\ LLM／AIエージェントによる投資意思決定の変革}\\[10mm]
{\color{brandaccent}\rule{\textwidth}{1.2pt}}\\[14mm]

\begin{minipage}{0.8\textwidth}
\centering
\small\color{brandgray}
本書は、440件を超える一次情報（学術論文・企業発表・ニュース・仮説的考察）を\\
体系的に精査し、AI-ready金融データ基盤、AI-nativeヘッジファンド、\\
投資意思決定AIエージェント、LLMによる金融データ分析の4領域について\\
事例・数値・示唆を網羅的に整理したリサーチホワイトペーパーである。\\[2mm]
金融は主役であると同時に、より広いAI業界の技術波の上にある。本書は各領域で、\\
リーガル・医療・ソフトウェア開発など周辺分野の同型の技術・動向を対比し、\\
金融が学べる示唆・借用できる設計・先取りすべきリスクを併記する。
\end{minipage}

\vfill

{\small\color{brandgray} 株式会社INDX　伊藤克哉 (k1ito)}\\[1mm]
{\small\color{brandgray} 2026年7月}

\end{center}
\end{titlepage}

\newpage
\tableofcontents
\newpage

% ============================================================
% 本文（別ファイルからinput）
% ============================================================
\section{エグゼクティブサマリー}
\label{sec:exec-summary}

本ホワイトペーパーは、2026年前半にかけて収集された440件を超える一次情報――学術論文（arXiv等）、大手金融機関・スタートアップの公式発表、業界ニュース、および調査チームによる仮説的考察（Takes）――を体系的に精査し、「金融データのAI-Ready化」と「AI-Nativeな投資運用」という一体の潮流を、4つの視点から統合的に描き出したものである。数百件の周辺情報を漏れなく列挙することよりも、追跡可能な一次資料に裏付けられた、驚きのある・示唆に富む事例を厳選して深く掘り下げることを編集方針としている。

\subsection*{本書の4つの柱}

\begin{enumerate}
  \item \textbf{AI-Ready金融データ基盤}: 決算資料・SEC提出書類・アナリストレポート・市場データ・ニュース・PDF・表・スプレッドシート・社内文書といった大量の構造化・非構造化データを、LLM／AIエージェント／RAG／自動分析ワークフローが確実に利用できる形式へ変換する技術と実装。
  \item \textbf{AI-Nativeヘッジファンド・投資会社}: Bridgewater、Man Group、Renaissance Technologies、Two Sigma、Citadel、D.E.\ Shaw、Point72/Cubist、Voleon、Numerai、QRT、WorldQuantなど、LLM・AIエージェント・機械学習を投資リサーチ・ポートフォリオ管理・リスク管理・トレーディングの中核に据える機関の実例。
  \item \textbf{投資意思決定のためのAIエージェント}: マルチエージェント・アーキテクチャによる財務分析・投資アイデア創出・デューデリジェンス・ポートフォリオモニタリング・リスク分析の支援・自動化。
  \item \textbf{LLM／AIエージェントによる金融データ分析}: 構造化・非構造化データを問わず、財務データを抽出・要約・分類・構造化・比較・推論する技術とその限界。
\end{enumerate}

4つの柱は独立した並列トピックではなく、下図のように一方向に積み上がる依存関係を成す。データ基盤の整備水準がファンド・エージェントの運用可能性を規定し、そこで得られた実例が分析技術の妥当性検証の場となり、最終的に4つの柱を横断する構造的示唆（第\ref{sec:cross-cutting}部）へ収束する。

\begin{center}
\begin{tikzpicture}[
  pillar/.style={fnode, text width=4.0cm, minimum height=1.8cm, font=\small},
  synth/.style={fnode accent, text width=8.6cm, minimum height=1.4cm, font=\small},
  arr/.style={farr blue}
]
\node[pillar] (data)    at (0,3)    {\textbf{(1)\,データ基盤}\\金融文書をAIが確実に利用できる形へ変換};
\node[pillar] (funds)   at (6.6,3)  {\textbf{(2)\,AI-Nativeファンド}\\AIを運用の中核に据える機関の実例};
\node[pillar] (agents)  at (0,0)    {\textbf{(3)\,投資AIエージェント}\\リサーチ・DD・モニタリングの自動化};
\node[pillar] (analysis) at (6.6,0) {\textbf{(4)\,LLM分析技術}\\抽出・要約・推論とその限界};
\node[synth] (synth) at (3.3,-2.6) {\textbf{柱を横断する構造的示唆}（第\ref{sec:cross-cutting}部）\\コスト構造・規制・標準化が実装の遅速を左右する};

\draw[arr] (data) -- (funds);
\draw[arr] (data) -- (agents);
\draw[arr] (funds) -- (analysis);
\draw[arr] (agents) -- (analysis);
\draw[arr] (agents.south) -- ++(0,-0.5) -| (synth.north);
\draw[arr] (analysis.south) -- ++(0,-0.5) -| (synth.north);
\end{tikzpicture}
\end{center}
{\small\color{brandgray}\textbf{図: 本書の4つの柱と全体構造}\ データ基盤がファンド・エージェントの実装水準を規定し、そこで得られた実例が分析技術の限界を検証する場となる。4つの柱の交点に現れる構造的示唆を第\ref{sec:cross-cutting}部で統合する。}

\subsection*{主要な発見}

\begin{itemize}
  \item \textbf{ボトルネックはモデル能力ではなくデータ準備}: GPT-4oはネスト表を含む金融PDFで前処理なしでは56\%の精度に留まるが、専用の前処理パイプラインを挟むと76\%まで改善する（\cite{multistructured2025}、\S\ref{sec:data-infra}）。FinanceBenchではGPT-4-Turbo＋検索でも81\%の質問が誤答・回答拒否となり\cite{financebench2023}、AUDITFLOWではシンボリック検証層を除去すると精度が82\%から18\%へ崩壊する\cite{auditflow2026}。
  \item \textbf{検索アーキテクチャが精度の主要な律速要因}: Fin-RATEベンチマークでは、企業横断比較タスクにおいてファインチューン済み金融LLMの精度が23.91\%から1.27\%へ崩壊する事例が確認されており\cite{finrate2026}、階層型・グラフ型の検索設計がフラットなベクトル検索を上回ることを複数の研究が裏付けている。
  \item \textbf{実運用に到達したAI-Nativeファンドが複数存在}: Bridgewater AIA Labsは実資本を運用し初年度11.9\%のリターンを記録\cite{bridgewater_aialabs_nd}、Renaissance TechnologiesのMedallion Fundは設立来年率39.9\%の純リターンを継続\cite{institutionalinvestor2025renaissance}、QRT（Qube Research \& Technologies）は2025年に主力ファンドで約30\%のリターンを達成する\cite{bloomberg_qube2026}など、AIを核とした運用体制が既に大規模な実績を残している（\S\ref{sec:ai-native-funds}）。
  \item \textbf{ベンチマーク上の「アルファ」の多くは方法論的に脆弱}: 77件のLLMトレーディング研究の系統的レビューでは、取引コストを明示的にモデル化していたのはわずか1件、再現可能で時間整合的な評価プロトコルを用いていたのはわずか2件にとどまる\cite{agentictrading2026}。Look-Ahead-Benchでは、評価期間をモデルの学習カットオフ以降に移すだけで年率アルファが+20.73\%から-1.04\%へほぼ消失する例が報告されている\cite{lookaheadbench2026}（\S\ref{sec:llm-analysis}）。
  \item \textbf{コスト構造がアーキテクチャ選択を左右する}: 反射型自己修正アーキテクチャ（F1 0.943、\$0.430/文書）に対し、最適化された階層型supervisor-workerアーキテクチャはF1 0.924を\$0.148/文書で達成する\cite{benchmarkingmultiagent2026}。100万文書規模では\$282,000のコスト差に相当し、リーダーボード精度だけに基づく調達判断は体系的に誤りうる。
  \item \textbf{規制・ガバナンスが実装の遅速を左右する新たな変数に}: EU AI Actの高リスク分類（与信スコアリング・保険料算定AI）は当初2026年8月2日の完全適用が予定されていたが、Digital Omnibus（欧州議会承認2026年6月）により信用・保険AI義務の拘束的期限は2027年12月2日へ約16カ月延期された\cite{eba2025aiact,x07digitalomnibus2026}。ただし非準拠に最大3,000万ユーロまたは世界売上高6\%の罰金が科される点、適合性評価・技術文書・データガバナンス・人間監督等を追加要求する点は変わらない。米国ではSR 11-7がリスク階層型ガイダンスへ置き換えられ\cite{databricksndmodelrisk}、これまでLLMコパイロットの本番投入を躊躇させてきた曖昧さが緩和されつつある（\S\ref{sec:cross-cutting}、\S\ref{sec:open-gaps}）。
  \item \textbf{価値はデータそのものから「解釈アーキテクチャ」へ移動}: 差別化の焦点は、埋め込み・検索・オーケストレーションといった\emph{解釈}層へ移りつつある。金融ドメインでは素朴なBag-of-Wordsが密な汎用埋め込みを上回る事例すら報告されており、埋め込み選択は精度を左右する「沈黙のボトルネック」になっている（\S\ref{sec:data-infra}）。同時に、JPMorgan（LLM Suiteが23万人超に到達）・Goldman（GS AIプラットフォーム）・BlackRock（Aladdin Copilot）に見られる\textbf{全社的democratization型}の展開は、アルファ創出型・完全自律型に次ぐ第3のAI-Native類型として立ち現れている（\S\ref{sec:ai-native-funds}）。
  \item \textbf{データ接続性はコモディティ化へ向かう}: Bloomberg、LSEG、FactSet、Daloopaが2025〜2026年に相次いでMCP（Model Context Protocol）サーバーを提供開始し\cite{factset_mcp_nd}、標準化されたデータ接続の限界価値はゼロに近づきつつある。差別化の焦点は「データを持っていること」から「データをどう解釈し文脈化し行動に結びつけるか」へ移行している。
  \item \textbf{エージェント標準は接続から協調・決済へ拡張}: OpenAI Responses APIのremote MCP対応\cite{x01openai2025responsesmcp}、Google A2A\cite{x01google2025a2a}、AP2\cite{x01google2025ap2}、Linux FoundationでのA2A普及\cite{x01linuxfoundation2026a2a}は、金融AIの標準化が「データを読む」段階から「複数エージェントが協調し、権限付きで経済行為を行う」段階へ移ったことを示す。
  \item \textbf{信頼性の鍵はグラフ・記号・棄権を組み込んだ検証層}: FinChainは金融推論を実行可能な記号テンプレートで評価し\cite{x05finchain2025}、GBFRは金融メトリック知識グラフ・経路検証・安全な回答棄権を組み合わせる\cite{x05gbfr2026}。単なる自然言語回答精度ではなく、計算経路・根拠・棄権判断を同時に評価する方向へベンチマークが移行している。
  \item \textbf{標準化・集中リスク・業界統制が未解決ギャップの中心へ}: FSBの2026年AI sound practices\cite{x07fsbsound2026}、米財務省の金融AI報告\cite{x07treasuryai2024}、Cyber Risk Instituteの金融AIリスク管理フレームワーク\cite{x07crifsairmf2026}はいずれも、モデル単体よりもプロバイダ集中、情報共有、データ標準、統制証跡を金融AIの主要リスクとして扱い始めている。
  \item \textbf{金融は業界横断のAI技術波の「速い追随者」}: AI-ready／AI-nativeへの移行は金融固有の現象ではなく経済全体の潮流であり、採用曲線ではソフトウェア開発（GitHub Copilotは2,000万ユーザー超\cite{y01techcrunch2025copilot20m}）やエンタープライズ検索が金融の買い手側に先行している。MCPはLinux Foundation傘下の業界横断標準として採用され\cite{y01linuxfoundation2025aaif}、金融は接続層を「作る」側ではなく「乗る」側になっている。Stanford AI Index\cite{y01stanfordhai2025aiindex}のような外部の物差しで自らの現在地を測れることは、投資判断のタイミングを計るうえで実務的な含意を持つ（\S\ref{sec:intro-cross-domain}）。
  \item \textbf{文書AIとエージェント評価は周辺分野が先行し、忠実性の教訓を借用できる}: 契約解析（CUAD\cite{y02cuad2021}）・GraphRAG\cite{y02graphrag2024}・エンタープライズ検索（Glean\cite{y02glean2025}）、およびSWE-bench Verified\cite{y04swebenchverified2024}に象徴される厳密なエージェント評価は、金融よりも文書AI・自律エージェントの成熟が進んでいる領域である。医療・法律の忠実性研究――Med-PaLM 2\cite{y05medpalm2_2025}や法的ハルシネーションの体系的実証\cite{y05dahl2024legalfictions}――は、金融が評価設計・検証層・自律度の段階付けをそのまま借用できる先行事例を提供する（\S\ref{sec:data-infra}、\S\ref{sec:investment-agents}、\S\ref{sec:llm-analysis}）。
  \item \textbf{金融は汎用AI課題が最も早く鋭く顕在化する「ストレステスト領域」}: 評価クライシス、エージェント安全性、基盤モデル集中（algorithmic monoculture\cite{y07monoculture2022}）、水平的規制（EU AI Actは特定業界向けではなく横断的に設計\cite{y07euaiacthorizontal}）はいずれも金融固有ではない汎用AIの未解決課題である。しかし実資本の毀損リスク・執行／決済権限・13F型の公開可測性・既存の規制密度ゆえに、これらは他業界に先んじて金融で最も先鋭化する（\S\ref{sec:open-gaps}）。
\end{itemize}

\subsection*{本書の読み方}

第\ref{sec:data-infra}部から第\ref{sec:llm-analysis}部までが4つの柱に対応する本論であり、各部は具体的な事例（企業名・論文名・数値・出典）を中心に構成される。第\ref{sec:cross-cutting}部では柱を横断する構造的な示唆を、第\ref{sec:open-gaps}部では現時点で未解決の研究・実装ギャップを整理し、第\ref{sec:conclusion}部で総括と示唆を述べる。巻末に参考文献（出典URL）の一覧を付す。各部の本文は、網羅的な事例集としてではなく、\textbf{casestudy}（検証済みの具体事例）と\textbf{insight}（そこから導かれる仮説）を対にして提示する構成を基本としており、量よりも一次情報への遡及可能性を優先している。

\section{はじめに}
\label{sec:introduction}

\subsection{背景と目的}

投資運用業界は、ここ数年で二つの制約に同時に直面してきた。第一に、分析対象となる情報量の爆発的増加である。SEC提出書類、決算説明会トランスクリプト、アナリストレポート、ニュース、オルタナティブデータ、社内メモは、人間のアナリストが処理できる量をとうに超えている。第二に、その情報の大半が非構造化・半構造化の形式――PDF、スキャン画像、ネストした表、自由記述テキスト――で存在し、伝統的な定量分析パイプラインに直接投入できない。

大規模言語モデル（LLM）とAIエージェントは、この二つの制約を同時に解消しうる技術として急速に投資業界に浸透しつつある。しかし「LLMを使えば金融データがAI-readyになる」という単純な図式は成り立たない。本書で精査した一次情報は、モデル単体の能力よりも、前処理・検索アーキテクチャ・シンボリック検証層・組織的なガバナンス設計こそが実運用の成否を分けていることを繰り返し示している。

\begin{concept}{LLM・生成AI・AIエージェント}
生成AI（LLM: 大規模言語モデル）は、膨大な文章を学習し、次に来る語を確率的に予測して文章を紡ぐモデルである。同じ入力でも出力が揺らぎうる点で、通常のプログラムの決定論的な計算とは性質が異なる。決算短信や有価証券報告書のような自由記述テキストを読ませ、要約・抽出・分類させられる点が実務での価値となる。一方AIエージェントは、このLLMを「頭脳」として、検索・計算・社内API・データベースといった外部ツールを呼び出し、記憶を保持しながら「観測→思考→行動」の自律ループを回す仕組みである。一問一答で終わるチャットボットと異なり、「決算書を取得し、表を抽出し、計算し、レポート化する」といった多工程のタスクを自分の判断で最後まで進められる点が本質的な違いである。
\begin{center}
\begin{tikzpicture}[node distance=6mm, every node/.style={font=\scriptsize}]
\node[fchip] (t1) {LLM（生成AI）};
\node[fnode, below=5mm of t1, text width=20mm, align=center] (in) {テキスト入力};
\node[fnode blue, below=6mm of in, text width=22mm, align=center] (pred) {次の語を\\確率的に予測};
\node[fnode, below=6mm of pred, text width=20mm, align=center] (out) {テキスト出力};
\draw[farr] (in) -- (pred);
\draw[farr] (pred) -- (out);
\node[fchip, right=34mm of t1] (t2) {AIエージェント};
\node[fnode dark, below=8mm of t2, text width=22mm, align=center] (brain) {LLM＝頭脳\\（思考）};
\node[fnode muted, right=8mm of brain, text width=20mm, align=center] (tool) {ツール\\（検索・計算・API）};
\node[fnode muted, below=8mm of brain, text width=16mm, align=center] (mem) {記憶};
\draw[farr accent] (brain) to[bend left=22] node[midway, above, font=\tiny] {行動} (tool);
\draw[farr accent] (tool) to[bend left=22] node[midway, below, font=\tiny] {観測} (brain);
\draw[fbiarr] (brain) -- (mem);
\end{tikzpicture}
\end{center}
{\small\color{brandgray}\textbf{図: LLMとAIエージェントの違い}\ LLMは入力から次の語を予測して出力するだけだが、AIエージェントはLLMを頭脳にツールと記憶を使い、観測と行動のループを自律的に回す。}
このループにより、単発の応答では終わらない「取得→抽出→計算→報告」のような多段階業務を、人手を介さず一気通貫で任せられるようになる。
\end{concept}

\begin{center}
\begin{tikzpicture}[
  font=\small,
  constraint/.style={fnode, align=left, text width=70mm, minimum height=27mm},
  merge/.style={fnode dark, text width=92mm, minimum height=10mm, font=\bfseries\small},
  factor/.style={fnode accent, text width=30mm, minimum height=17mm, font=\footnotesize},
  arr/.style={farr}
]

\node[constraint] (c1) {\textbf{制約1: 情報量の爆発}\\[1mm]
  SEC提出書類、決算説明会トランスクリプト、アナリストレポート、ニュース、オルタナティブデータ、社内メモが人間のアナリストの処理能力を超過};
\node[constraint, right=4mm of c1] (c2) {\textbf{制約2: 非構造化・半構造化データ}\\[1mm]
  PDF、スキャン画像、ネストした表、自由記述テキストは伝統的な定量分析パイプラインに直接投入できない};

\node[merge, below=9mm of c1.south, xshift=37mm] (llm) {LLM／AIエージェントが両制約を同時に解消しうる技術として浸透};

\node[factor, below=8mm of llm, xshift=-50mm] (f1) {前処理};
\node[factor, right=3mm of f1] (f2) {検索\\アーキテクチャ};
\node[factor, right=3mm of f2] (f3) {シンボリック\\検証層};
\node[factor, right=3mm of f3] (f4) {組織的な\\ガバナンス設計};

\draw[arr] (c1.south) -- ++(0,-4mm) -| (llm.north);
\draw[arr] (c2.south) -- ++(0,-4mm) -| (llm.north);
\draw[arr] (llm.south) -- ($(f1.north)!0.5!(f4.north)$);
\end{tikzpicture}
\end{center}
{\small\color{brandgray}\textbf{図: 二つの制約とLLM/AIエージェントの位置づけ}\ 情報量の爆発と非構造化データという二つの制約をLLM／AIエージェントが同時に解消しうる一方、実運用の成否は前処理・検索アーキテクチャ・シンボリック検証層・組織的ガバナンス設計という4つの要因に左右される。}

本ホワイトペーパーは、この分野の一次情報を網羅的に精査し、「何が既に実運用に到達しているか」「どこに測定可能な効果があるか」「どこに未解決の課題が残るか」を、誇張や憶測を排して整理することを目的とする。

\subsection{なぜ「いま」なのか――市場的背景}
\label{sec:intro-market}

生成AIの投資業界への浸透は、もはや一部の先進的なクオンツファンドに限られた現象ではない。2025年に相次いで公表された複数の独立した調査は、採用がすでに\textbf{ほぼ普遍的な水準}に達し、争点が「導入するか否か」から「どこまで運用の中核に据えるか」へと移行したことを一致して示している。McKinseyの年次調査『The State of AI』（2025年11月版、105か国・1,993名を対象に2025年6〜7月に実施）では、AIを少なくとも一つの業務機能で利用する組織は78\%に達し、AIエージェントを本番展開（scaling）している組織が23\%、実験段階が39\%と、エージェント技術が2025年最大のトレンドとして立ち上がったことが確認されている\cite{x01mckinsey2025stateofai}。買い手側に限れば採用率はさらに高い。AIMAの2025年調査（運用資産7,880億ドルを代表する運用会社150社を対象）ではファンドマネジャーの95\%がすでに生成AIを業務に用いており（2023年の86\%から上昇）、生成AIが投資意思決定でより大きな役割を果たすと見込む割合は2023年の20\%から58\%へと急増した\cite{x01aima2025}。EYの『GenAI in Wealth \& Asset Management Survey 2025』も、95\%の企業が複数ユースケースへ展開済み、78\%がエージェント型AIを検討中と報告している\cite{x01ey2025genai}。

\begin{table}[htbp]
\centering
\small
\caption{買い手側・エンタープライズにおける生成AI採用の主要指標（2025年）}
\label{tab:intro-adoption}
\begin{tabular}{@{}p{34mm}p{58mm}p{34mm}@{}}
\toprule
\textbf{調査主体} & \textbf{主要な発見} & \textbf{対象・時期} \\
\midrule
McKinsey \cite{x01mckinsey2025stateofai} & AI利用組織78\%／エージェント本番展開23\%・実験39\%／高成果企業（EBIT寄与5\%超）は6\%のみ & 105か国・1,993名／2025年6〜7月 \\
AIMA \cite{x01aima2025} & 生成AI利用95\%（'23年86\%）／投資判断での役割拡大を見込む58\%（'23年20\%） & 運用会社150社・7,880億ドル／2025年9月 \\
EY \cite{x01ey2025genai} & 複数ユースケース展開済み95\%／エージェント型AI検討中78\%／実質的な事業インパクトを実感27\% & 資産・ウェルス運用会社／2025年9月 \\
Coalition Greenwich \cite{altdata2025} & オルタナティブデータ支出の増額を計画する投資家63\% & 買い手側56社／2025年9〜10月 \\
\bottomrule
\end{tabular}
\end{table}

一方で、採用の普遍化は成果の普遍化を意味しない。同じMcKinsey調査は、企業レベルで何らかの財務的インパクトを報告した組織が39\%にとどまり、EBIT寄与5\%超の「高成果企業」はわずか6\%であることを示している\cite{x01mckinsey2025stateofai}。この「採用は進むが規模化した価値は捉えきれていない」というギャップこそが、本書が\S\ref{sec:intro-definition}で論じる\textbf{AI-Ready／AI-Nativeという概念区分}の実務的な出発点となる。

市場規模の予測もまた、この時期の期待の急速な再評価を映し出している。Bloomberg Intelligenceの独自モデルは、生成AI市場が2022年の約400億ドルから2032年には1.3兆ドル（年率約43\%）へ拡大すると2023年に予測し、2026年の改訂では2.3兆ドルへと上方修正した\cite{x01bi2023genaimarket}。金融分野に限っても、McKinseyは生成AIが銀行業に年間2,000〜3,400億ドル（営業利益の9〜15\%相当）の価値をもたらしうると見積もり\cite{x01mckinsey2023banking}、AI全般の金融市場はMarketsandMarketsの推計で2024年の383.6億ドルから2030年には1,903.3億ドル（年率30.6\%）へ成長するとされる\cite{x01marketsandmarkets2024aifinance}。これらの数値は算定基準が調査機関ごとに異なるため額面比較には注意を要するが、いずれも二桁後半〜四十%台という高いCAGRで一致しており、投資判断・データ処理を担う中核業務が生成AIの主要な受益領域と目されている点が本書の関心と重なる。

\subsection{「AI-Ready」と「AI-Native」――定義の明確化}
\label{sec:intro-definition}

本書のタイトルに含まれる二つの鍵概念は、しばしば混同されるが、対象も水準も異なる。両者を明示的に区別することが、本書全体の議論を貫く軸となる。

\textbf{AI-Ready（AIレディ）}は、まず\emph{データ}に関する概念である。Gartnerの定義によれば、AI-Readyデータとは品質・網羅性・関連性・多様性・想定ユースケースへの整合性・倫理／ガバナンス上の健全性という要件を満たし、AIでの利用に適することが検証されたデータを指す。決定的に重要なのは、従来のデータ品質基準で「高品質」と判定されたデータが自動的にAI-Readyになるわけではないという点である\cite{x01gartner2025aiready}。Gartnerは、AI-Readyデータに支えられていないAIプロジェクトの60\%が2026年までに放棄されると予測しており\cite{x01gartner2025aiready}、さらに2026年の調査では、AIで成果を上げている組織はデータ品質・ガバナンス・人材・変革管理といった基盤領域に（売上比で）最大4倍多く投資し、成熟度の高い組織は最大65\%高い事業成果を得ていると報告している\cite{x01gartner2026foundations}。すなわち、成否を分けるのはモデルそのものよりも\textbf{データ基盤への投資の深さ}である。

\textbf{AI-Native（AIネイティブ）}は、\emph{組織・システムのアーキテクチャ}に関する概念である。AIを既存の製品やワークフローに後付けした「AI-Enabled（AI対応）」との対比で語られ、最も簡潔な判定基準は「AIを取り除いたときにその機能が\emph{劣化}するにとどまるか、それとも\emph{完全に停止}するか」である\cite{x01crv2025ainative}。AI-Nativeなシステムは確率的な出力・反復・適応を前提に設計されるため、多段階の業務プロセスが単一のプロンプトとエージェントの推論列へ折り畳まれる。後付け型アーキテクチャの能力上限は固定されるのに対し、AI-Native型の上限はモデルの改良ごとに上昇するため、その差は3〜5年で埋め難いものになると論じられている\cite{x01crv2025ainative}。

\begin{concept}{AI-Ready と AI-Native}
「AI-Ready」と「AI-Native」は似た言葉だが層が違う。AI-Ready（AIレディ）はデータ層の状態を指し、品質・網羅性・メタデータ・ガバナンスが整い、AIが確実に読み解ける形になっていることを意味する。単なる「高品質データ」とは別物である。AI-Native（AIネイティブ）は組織・システムのアーキテクチャの話であり、AIを前提に設計され、AIを外すと機能そのものが止まる状態を指す。後付けでAIを足しただけの「AI-Enabled（AI対応）」との違いは、次の一問で見分けられる：「AIを取り除いたら、機能は劣化にとどまるか、それとも停止するか」。
\begin{center}
\begin{tikzpicture}[node distance=8mm, every node/.style={font=\small}]
\node[fnode dark, text width=30mm] (sys) {対象システム};
\node[fnode muted, above=9mm of sys, text width=30mm] (data) {データ層: AI-Ready\\（品質・網羅性・ガバナンス）};
\node[fnode, right=24mm of sys, yshift=17mm, text width=34mm] (enabled) {AIを外す $\rightarrow$ \textbf{劣化}にとどまる\\＝AI-Enabled（後付け）};
\node[fnode accent, right=24mm of sys, yshift=-17mm, text width=34mm] (native) {AIを外す $\rightarrow$ \textbf{停止}する\\＝AI-Native（設計層）};
\draw[farr] (data) -- (sys);
\draw[farr] (sys.east) -- (enabled.west);
\draw[farr accent] (sys.east) -- (native.west);
\end{tikzpicture}
\end{center}
{\small\color{brandgray}\textbf{図: 「AIを外すテスト」による直感的な判定}\ AIを取り除いたときに劣化で済むならAI-Enabled、停止するならAI-Nativeである。}
このテストは、見た目にAIらしい機能があるかどうかではなく、AIが構造上どれだけ不可欠かを問う点が要である。
\end{concept}

\begin{table}[htbp]
\centering
\small
\caption{AI-Ready と AI-Native の対比}
\label{tab:intro-ready-native}
\begin{tabular}{@{}p{26mm}p{57mm}p{57mm}@{}}
\toprule
 & \textbf{AI-Ready} & \textbf{AI-Native} \\
\midrule
主たる対象 & データ（および前処理・検索基盤） & 組織・システムのアーキテクチャ \\
問い & このデータはAIが確実に利用できる形か & AIを外すと業務は停止するか、単に劣化するか \\
成否の律速 & データ品質・ガバナンス・メタデータ・可観測性 \cite{x01gartner2025aiready} & 確率的出力を前提とした設計・エージェント化の深さ \cite{x01crv2025ainative} \\
本書での対応 & 第\ref{sec:data-infra}部（データ基盤） & 第\ref{sec:ai-native-funds}部（AI-Nativeファンド）ほか \\
\bottomrule
\end{tabular}
\end{table}

\begin{insight}{示唆: データ接続のコモディティ化と差別化の移動}
Bloomberg・LSEG・FactSet・DaloopaがMCP（Model Context Protocol）サーバーを相次いで提供し\cite{factset_mcp_nd}、標準化されたデータ\emph{接続}の限界価値がゼロに近づくと、買い手側の持続的な差別化は「データを保有すること（AI-Readyの達成）」から「AI-Nativeなアーキテクチャでデータを解釈・文脈化し行動に結びつけること」へと移動する。実際、採用率が95\%に達した以上、採用それ自体はもはや差別化要因たりえない\cite{x01aima2025,x01ey2025genai}。ただしこれは仮説であり、反証材料も存在する。Citadelのケン・グリフィンは2025年10月時点で生成AIは生産性を高めるがアルファ創出には「力不足」と述べ、2026年半ばに評価を一変させた\cite{griffin2026}。AI-Nativeの優位は実在しうるが、その到来はベンダーの物語が示唆するよりも遅く、かつ不均等でありうる。
\end{insight}

\subsection{標準化とタイムライン――2023→2026の変曲点}
\label{sec:intro-timeline}

生成AIが投資業界で「実験」から「実装」へ移行した背景には、モデル能力の向上だけでなく、\textbf{データ接続の標準化}と\textbf{規制の明確化}という二つの制度的変化がある。とりわけAnthropicが2024年11月に公開したMCP（Model Context Protocol）は、AIエージェントとデータソース／ツールをつなぐ「N×M問題」を解消する共通規格として急速に普及し\cite{x01anthropic2024mcp}、OpenAIも2025年5月にResponses APIでリモートMCPサーバー対応を追加し、MCPをエージェント開発APIの標準的な道具立てへ組み込んだ\cite{x01openai2025responsesmcp}。2025年12月にはMCPがLinux Foundation傘下のAgentic AI Foundation（Anthropic・Block・OpenAIのプロジェクトを初期貢献とし、AWS・Bloomberg・Cloudflare・Google・Microsoft等が参加）へ寄贈され、ベンダー主導の仕様から中立的な共通インフラへと移行した\cite{x01mcp2026foundation}。この標準化こそが、\S\ref{sec:intro-definition}で述べた「データ接続のコモディティ化」の技術的な裏づけである。

ただし、2025〜2026年の標準化はMCPだけで完結しない。Googleが2025年4月に発表したAgent2Agent（A2A）は、エージェント同士がベンダーやフレームワークをまたいで通信・情報交換・協調行動を行うためのプロトコルであり、50社超の技術・サービスパートナーとともに公表された\cite{x01google2025a2a}。さらに2025年9月のAgent Payments Protocol（AP2）は、A2AとMCPを拡張する形で、エージェントがユーザーの委任に基づいて支払いを開始する際の認可・真正性・責任分界を扱う\cite{x01google2025ap2}。Linux Foundationは2026年4月時点で、A2Aが150超の組織に支持され、Google・Microsoft・AWSの主要クラウドに深く統合され、金融サービスを含む複数業界で本番利用が始まったと発表している\cite{x01linuxfoundation2026a2a}。したがって、投資業界にとっての変曲点は「LLMが賢くなった」ことではなく、\textbf{データに触る、他のエージェントと協調する、経済行為を実行する}という三層の相互運用性が同時に整備され始めた点にある。

\begin{table}[htbp]
\centering
\small
\caption{2024〜2026年に立ち上がったエージェント標準化スタック}
\label{tab:intro-agent-standards}
\begin{tabular}{@{}p{22mm}p{34mm}p{50mm}p{34mm}@{}}
\toprule
\textbf{層} & \textbf{代表標準} & \textbf{解く問題} & \textbf{金融での意味} \\
\midrule
データ／ツール接続 & MCP \cite{x01anthropic2024mcp,x01openai2025responsesmcp} & エージェントが外部データ、API、業務ツールへ共通インターフェースで接続する & 市場データ、文書、社内APIへの接続コストを低下させる \\
エージェント間協調 & A2A \cite{x01google2025a2a,x01linuxfoundation2026a2a} & 異なるベンダー・フレームワークのエージェントが互いを発見し、タスクを委譲・調整する & リサーチ、リスク、コンプライアンスの複数エージェントを横断編成できる \\
経済行為・支払い & AP2 \cite{x01google2025ap2,x01linuxfoundation2026a2a} & ユーザー委任、真正性、責任分界を保ちながらエージェント主導の支払いを扱う & 決済・証券・保険で、エージェントが「助言」から「執行」に近づく際の統制点になる \\
\bottomrule
\end{tabular}
\end{table}

\begin{concept}{MCP（Model Context Protocol）}
MCPとは、AIエージェントと外部のデータ源・ツールを繋ぐための共通規格（プロトコル）である。Anthropicが2024年11月に公開し\cite{x01anthropic2024mcp}、その後業界標準化が進んだ。素朴には、AIアプリ$M$個とデータ源$N$個を個別配線で繋ぐと最大$M\times N$本の専用コネクタが必要になり、組合せが増えるほど接続の管理が破綻する（「N×M問題」）。これは家電の電源プラグが会社ごとにバラバラだと変換器が無数に要るのと同じ状況である。MCPという共通の「差込口」を間に挟めば、AI側・データ側それぞれが一度MCPに対応するだけで済み、必要な接続本数は$N+M$へと激減する。金融領域ではBloomberg・LSEG・FactSet等が既にMCPサーバーを提供し始めており\cite{factset_mcp_nd}、データへの「接続」そのものがコモディティ化しつつある。
\begin{center}
\begin{tikzpicture}[
  node distance=6mm,
  every node/.style={font=\footnotesize},
  aiN/.style={fnode blue, text width=16mm, minimum height=7mm},
  dsN/.style={fnode muted, text width=16mm, minimum height=7mm},
  hub/.style={fnode dark, text width=16mm, minimum height=9mm, font=\bfseries\footnotesize}
]
\node[fchip] (t1) {MCPなし：N\texttimes M本};
\node[aiN, below=8mm of t1, xshift=-12mm] (a1) {AI1};
\node[aiN, below=4mm of a1] (a2) {AI2};
\node[aiN, below=4mm of a2] (a3) {AI3};
\node[dsN, right=20mm of a2] (d1) {データ1};
\node[dsN, above=4mm of d1] (d0) {データ2};
\node[dsN, below=4mm of d1] (d2) {データ3};
\foreach \a in {a1,a2,a3}{
  \foreach \d in {d0,d1,d2}{
    \draw[farr dash] (\a) -- (\d);
  }
}
\node[fchip, right=42mm of t1] (t2) {MCPあり：N+M本};
\node[aiN, below=8mm of t2, xshift=-12mm] (b1) {AI1};
\node[aiN, below=4mm of b1] (b2) {AI2};
\node[aiN, below=4mm of b2] (b3) {AI3};
\node[hub, right=14mm of b2] (hub) {MCP};
\node[dsN, right=14mm of hub, yshift=8mm] (e1) {データ2};
\node[dsN, right=14mm of hub] (e2) {データ1};
\node[dsN, right=14mm of hub, yshift=-8mm] (e3) {データ3};
\draw[farr] (b1) -- (hub);
\draw[farr] (b2) -- (hub);
\draw[farr] (b3) -- (hub);
\draw[farr blue] (hub) -- (e1);
\draw[farr blue] (hub) -- (e2);
\draw[farr blue] (hub) -- (e3);
\end{tikzpicture}
\end{center}
{\small\color{brandgray}\textbf{図: N×M問題の解消}\ 個別配線（左）では接続数が掛け算で増えるが、MCPという共通ハブ（右）を挟むと足し算で済む。}
\end{concept}

\begin{center}
\begin{tikzpicture}[
  font=\footnotesize,
  milestone/.style={fnode, align=left, text width=30mm, minimum height=20mm},
  reg/.style={fnode accent, align=left, text width=30mm, minimum height=20mm},
  arr/.style={farr}
]
\node[milestone] (m1) {\textbf{2023}\\[1mm]金融特化型LLMの登場（BloombergGPT）\cite{bloomberggpt2023}};
\node[milestone, right=5mm of m1] (m2) {\textbf{2024.11}\\[1mm]MCP公開、データ接続の標準化が始動\cite{x01anthropic2024mcp}};
\node[milestone, right=5mm of m2] (m3) {\textbf{2025}\\[1mm]A2A/AP2により協調・支払いの標準化も始動\cite{x01google2025a2a,x01google2025ap2}};
\node[reg, right=5mm of m3] (m4) {\textbf{2026}\\[1mm]A2Aが150超組織へ拡大、AI規制も本格化\cite{x01linuxfoundation2026a2a,eba2025aiact}};

\draw[arr] (m1.east) -- (m2.west);
\draw[arr] (m2.east) -- (m3.west);
\draw[arr] (m3.east) -- (m4.west);
\end{tikzpicture}
\end{center}
{\small\color{brandgray}\textbf{図: 2023→2026の主要な変曲点}\ モデル能力の向上（金融特化型LLM）に、データ接続の標準化（MCP）と規制の明確化（EU AI Act）という制度的変化が重なり、投資業界は生成AIの実装フェーズへ移行した。金色のノードは規制上の節目を示す。}

規制面では、EU AI Actの高リスク分類（与信スコアリング・保険料算定AI等）の完全適用が当初2026年8月2日に予定されていたが、Digital Omnibus（欧州議会承認2026年6月）により信用・保険AI義務の拘束的期限は2027年12月2日へ延期された\cite{eba2025aiact,x07digitalomnibus2026}。非準拠には最大3,000万ユーロまたは世界売上高6\%の罰金が科される点は変わらない。規制・ガバナンスは、これまでLLMコパイロットの本番投入を躊躇させてきた曖昧さを緩和する一方で、実装の遅速を左右する新たな変数として立ち現れている（詳細は第\ref{sec:cross-cutting}部）。

\subsection{周辺分野から見た位置づけ：業界横断のAI移行と金融}
\label{sec:intro-cross-domain}

ここまで見てきた採用率・標準化のうねりは、投資運用業界に固有の現象ではない。むしろ金融は、IT・ソフトウェア産業を筆頭とする\textbf{経済全体のAI移行という、より大きな潮流の上に乗っている}。Stanford大学HAI（Human-Centered Artificial Intelligence Institute）の『2025 AI Index Report』によれば、少なくとも一つの業務機能でAIを利用する組織の割合は、業種を問わない経済全体の横断調査で2023年の55\%から2024年には78\%へ達しており\cite{y01stanfordhai2025aiindex}、この水準は本書\S\ref{sec:intro-market}で見た金融買い手側の高採用率（AIMA調査でファンドマネジャーの95\%\cite{x01aima2025}等）が確立する前後の時期に、業種を問わず既に到達していた基準点である。すなわち金融の高い採用率は、金融固有の突出した先進性というより、経済全体で既に進行していた拡散曲線に金融が追いついた結果と解釈する方が整合的である。本節では、AI-Ready／AI-Nativeという概念、そしてMCP等の標準化がいずれも金融発ではなく業界横断的に立ち上がった経緯を、金融のプロが自らの採用段階を測るための「ものさし」として提示する。

\begin{casestudy}{ソフトウェア開発: AIツール普及曲線のペースセッター}
生成AIの企業導入が最も早く・最も深く進んだ分野はソフトウェア開発である。GitHub Copilotは2025年7月に累計ユーザー数\textbf{2,000万人}を突破し（同年4月の1,500万人から3か月で500万人増）、\textbf{Fortune 100の90\%}が導入済みである\cite{y01techcrunch2025copilot20m}。競合のCursor（開発元Anysphere）は、2025年1月に年換算収益（ARR）\textbf{1億ドル}、6月に\textbf{5億ドル}、11月に\textbf{10億ドル}（同時に290億ドル評価額でシリーズDを実施）、2026年2〜3月には\textbf{20億ドル}へと、13か月で20倍という、Slack・Zoom・Snowflakeを上回るペースでの成長を遂げた\cite{y01techcrunch2025cursorarr,y01businesswire2025cursorseriesd}。Gartnerは、企業のソフトウェアエンジニアのうちAIコーディング支援ツールを利用する割合が、2023年初頭の10\%未満から2028年には75\%に達すると予測している\cite{y01gartner2024aicodeassist}。これらはいずれも金融とは無関係に進行した、AIネイティブなツール導入の教科書的な普及曲線である。
\end{casestudy}

\begin{insight}{金融への示唆: 自らの採用段階を外部のものさしで測る}
ソフトウェア開発におけるCopilot／Cursorの普及曲線は、\S\ref{sec:intro-market}で見た金融買い手側の採用曲線（AIMA調査: 生成AI利用86\%→95\%、2023→2025年）と比較するための「外部のものさし」を提供する。ソフトウェア開発では、単なるツール導入率（3分の2の企業が生成AIツールを導入済み\cite{y01bain2025techreport}）と実際の生産性押し上げ効果には大きな差があり、基礎的なAIコーディング支援だけでは10〜15\%の生産性向上にとどまる一方、業務プロセス自体を再設計した企業では25〜30\%に達する\cite{y01bain2025techreport}。この「導入率と業務プロセス変革の間のギャップ」は、金融における「エージェント型AI本番展開23\%・実験段階39\%」というMcKinsey調査の構図（\S\ref{sec:intro-market}）と酷似しており、金融固有の課題ではなく、AIツール導入一般に共通する構造的な現象であることを示唆する。ソフトウェア開発が既に踏んだ「導入は容易だが、業務プロセスの再設計を伴わなければ効果は限定的」という段階を、金融も遅かれ早かれ通過することを見越した投資判断が求められる。
\end{insight}

\begin{table}[htbp]
\centering
\small
\caption{業界横断で見た生成AI／AIエージェント普及状況（2023〜2026年）}
\label{tab:intro-cross-domain-adoption}
\begin{tabular}{@{}p{30mm}p{62mm}p{34mm}@{}}
\toprule
\textbf{分野} & \textbf{採用状況の指標} & \textbf{出典・時期} \\
\midrule
経済全体（業種横断） & 組織のAI利用率 55\%→78\%（2023→2024年） & Stanford AI Index \cite{y01stanfordhai2025aiindex} \\
ソフトウェア開発（AIコーディング支援） & GitHub Copilot 累計2,000万ユーザー、Fortune 100の90\%が導入 & TechCrunch、2025年7月 \cite{y01techcrunch2025copilot20m} \\
ソフトウェア開発（AIネイティブ新興企業） & Cursor/Anysphere: ARR 1億ドル→20億ドル（13か月） & TechCrunch/Business Wire、2025年1月〜2026年3月 \cite{y01techcrunch2025cursorarr,y01businesswire2025cursorseriesd} \\
ソフトウェア産業（予測） & 企業ソフトウェアエンジニアのAIコード支援利用 10\%未満→75\%（2023→2028年予測） & Gartner、2024年4月 \cite{y01gartner2024aicodeassist} \\
データ接続標準（MCP、業種横断） & Agentic AI Foundation設立会員に通信・小売・SaaS企業が多数（金融はBloombergのみ） & Linux Foundation、2025年12月 \cite{y01linuxfoundation2025aaif} \\
データ接続標準（MCP、ソフトウェア産業） & ソフトウェア組織の41\%がMCPを限定的または本格的に本番運用 & Stacklok、2026年1月 \cite{y01stacklok2026mcpsoftware} \\
資産運用・ヘッジファンド（金融、買い手側） & ファンドマネジャーの生成AI利用 86\%→95\%（2023→2025年） & AIMA、2025年9月 \cite{x01aima2025} \\
\bottomrule
\end{tabular}
\end{table}
{\small\color{brandgray}\textbf{図: 業界横断で見た生成AI／AIエージェント普及状況}\ 経済全体・ソフトウェア産業・データ接続標準・金融買い手側の採用指標を並べると、金融の高い採用率（86\%→95\%）は独立した現象ではなく、より大きな業界横断的な普及曲線の一部として位置づけられることが分かる。}

\begin{casestudy}{MCP標準化の主導権: 金融発ではなく業界横断インフラとして立ち上がった}
\S\ref{sec:intro-timeline}で見たMCPは、2025年12月にLinux Foundation傘下の新設「Agentic AI Foundation（AAIF）」へ寄贈された。プラチナ会員はAWS・Anthropic・Block・Bloomberg・Cloudflare・Google・Microsoft・OpenAIの8社、ゴールド会員にはCisco・Ericsson（通信）、Shopify（小売・EC）、Salesforce・SAP・Snowflake・Okta・Datadog・Twilio（エンタープライズSaaS）が名を連ねる\cite{y01linuxfoundation2025aaif}。金融固有の企業はBloombergのみであり、MCPが金融発の技術ではなく、通信・小売・SaaSを含む業界横断のインフラとして立ち上がったことを裏づける。さらにStacklokの2026年調査では、ソフトウェア組織の\textbf{41\%}が既にMCPサーバーを限定的または本格的に本番運用しており、同社は同時期に小売・金融サービス向けの同種調査も並行して公表している\cite{y01stacklok2026mcpsoftware}。
\end{casestudy}

\begin{insight}{金融への示唆: 接続層の標準化は「乗る」ものであり「作る」ものではない}
金融データベンダー（Bloomberg・LSEG・FactSet・Daloopa）がMCPサーバーを提供し始めたことは（\S\ref{sec:intro-definition}）、金融が独自に接続標準を主導した結果ではなく、既に通信・小売・SaaS業界を巻き込んで確立された業界横断インフラに、後から乗った実装である。この事実は\S\ref{sec:intro-definition}で述べた「データ接続のコモディティ化」という示唆を補強する: 接続層の標準化競争において金融が独自の差別化を主張する余地は元々小さく、差別化の焦点は最初から解釈・オーケストレーション層に置くべきだったことになる。同時に、金融固有のガバナンス・監査要件（本節前段のEU AI Act等）は、業界横断インフラの上に金融だけが追加で積む「規制レイヤー」として理解するのが実務的である。
\end{insight}

ただし、金融が経済全体の潮流に単純に追随しているだけというわけでもない。同じMcKinsey調査は、AI導入によって人員を「純減」させると予想する回答者の割合が「変化なし」を上回る業種は他に見当たらない中で、\textbf{金融サービスだけが「人員純減」を「変化なし」より多く見込む唯一の業種}であったと報告している\cite{x01mckinsey2025stateofai}。これは、金融におけるAI導入が単なる生産性向上の一手段ではなく、業務構造そのものの再編を伴う変化として認識されていることを示唆しており、\S\ref{sec:llm-analysis}以降で論じる自動化の深さの議論に接続する。

これらの観察から得られる一つの仮説は、金融は業界横断のAIツール・標準化の波に対しては「速い追随者（fast follower）」であり、先駆者ではないというものである。ソフトウェア開発・エンタープライズSaaSの採用曲線は金融に先行し、MCP・A2Aといった接続・協調標準も金融発ではなく業界横断で確立された後に金融へ波及した。ただしこの仮説には反証材料もある。同一期間に金融買い手側の生成AI利用率は86\%から95\%へ他業種平均を上回るペースで上昇しており（\S\ref{sec:intro-market}）、Bridgewater AIA Labsのように実運用資本で大規模な実績を残すAI-Nativeファンドは、他業種で同等の事例が確立する前から存在した（第\ref{sec:ai-native-funds}部）。したがって「fast follower」という位置づけは、接続・ツール標準といったインフラ層には当てはまるが、金融固有のアルファ創出アプリケーションには必ずしも当てはまらないと、限定的に理解すべきである。

\subsection{調査手法}

本書の元となったのは、学術研究（Foundations）・業界ニュース（Latest）・仮説的考察（Takes）の3種類、合計212件の一次情報である。各項目は実在する論文・公式発表・報道記事に基づいて収集され、以下の観点で精査・分類した。

\begin{enumerate}
  \item \textbf{テーマ分類}: 各項目を「AI-Ready金融データ基盤」「AI-Nativeヘッジファンド・投資会社」「投資意思決定のためのAIエージェント」「LLM／AIエージェントによる金融データ分析」の4テーマ（重複を許容）に整理した。
  \item \textbf{事実確認の原則}: 企業名・論文名・数値・URLは原資料の記載をそのまま用い、本書側での推測・創作は行っていない。仮説的考察（Takes）に基づく記述は、本文中で「仮説」「Rationale」等と明示し、実証済みの事実と区別している。
  \item \textbf{横断的パターンの抽出}: 個別事例をテーマごとに整理するだけでなく、複数の独立した研究・事例に共通して現れる構造的パターン（例: 検索アーキテクチャのボトルネック化、コスト・精度のトレードオフ、規制の非同期な進展）を横断的に抽出した。
\end{enumerate}

\begin{center}
\begin{tikzpicture}[
  font=\small,
  stage/.style={fnode, align=left, text width=37mm, minimum height=27mm},
  stagelast/.style={fnode accent, align=left, text width=37mm, minimum height=27mm},
  arr/.style={farr},
  node distance=13mm
]

\node[stage] (src) {
  \textbf{1. 一次情報の収集}\\[1.5mm]
  \footnotesize
  合計\textbf{212件}\\
  ・学術研究（Foundations）\\
  ・業界ニュース（Latest）\\
  ・仮説的考察（Takes）
};

\node[stage, right=of src] (theme) {
  \textbf{2. 4テーマへの分類}\\[1.5mm]
  \footnotesize
  ・AI-Ready金融データ基盤\\
  ・AI-Nativeファンド\\
  ・投資意思決定エージェント\\
  ・LLMデータ分析
};

\node[stagelast, right=of theme] (cross) {
  \textbf{3. 横断的パターン抽出}\\[1.5mm]
  \footnotesize
  検索ボトルネック化\\
  コスト・精度トレードオフ\\
  規制の非同期な進展\ 等
};

\draw[arr] (src.east) -- (theme.west);
\draw[arr] (theme.east) -- (cross.west);
\end{tikzpicture}
\end{center}
{\small\color{brandgray}\textbf{図: 本書の調査手法}\ 学術研究・業界ニュース・仮説的考察の3種類、合計212件の一次情報を4つのテーマに分類したうえで、個別事例を横断して繰り返し現れる構造的パターンを抽出した。第\ref{sec:data-infra}部から第\ref{sec:llm-analysis}部が「2」の各テーマに、第\ref{sec:cross-cutting}部が「3」に対応する。}

本書の編集方針は、収集した212件をすべて等価に列挙することではない。個々の項目には検証の厚み・再現性・実運用への到達度に大きな差があるため、本書は\textbf{数の網羅性よりも、事例の質と検証可能性を優先}する。具体的には、（i）定量的な裏付けを伴う事例や、複数の独立した研究が同じ結論に収束している場合を重点的に取り上げ、単発的・宣伝色の強い発表は簡潔な言及にとどめる、（ii）実証済みの「事実」と、著者の解釈である「仮説（Takes）」を本文中で明確に区別し、後者は根拠と反証材料を併記する、という2点を徹底している。この方針は、生成AIをめぐる情報の多くが誇張されたマーケティング表現と厳密な実証研究の間を揺れ動いている現状に対する、本書なりの応答である。

\subsection{対象範囲と限界}

本書が扱う事例の多くは2025年後半から2026年前半にかけて公表されたものであり、生成AI・エージェント技術が投資業界で急速に実装フェーズへ移行した時期の「スナップショット」である。以下の点に留意されたい。

\begin{itemize}
  \item \textbf{ベンチマーク上の成果と実運用の成果は区別が必要}: 本書で紹介する精度・Sharpe比・コスト削減率などの数値は、多くの場合バックテストや限定的なパイロット環境で得られたものであり、複数年にわたる実市場での検証を経たものは一部（Bridgewater AIA Labs、Renaissance Technologies等の実運用ファンド）に限られる。
  \item \textbf{再現性の限界}: 後述するとおり（\S\ref{sec:llm-analysis}）、LLMトレーディング研究の多くは取引コストのモデル化や時間整合的な評価プロトコルを欠いており、公表された「アウトパフォーム」を額面通りに受け取ることはできない。
  \item \textbf{情報の非対称性}: 大手ヘッジファンドの多くは技術的詳細を非公開とする傾向が強く、本書で参照できるのは各社が意図的に公開した情報（公式ブログ、プレスリリース、幹部インタビュー）に限られる。したがって「公開されていない＝存在しない」ではなく、「公開されていない＝検証不能」と解釈すべきである。
\end{itemize}

\begin{center}
\begin{tikzpicture}[
  font=\small,
  period/.style={fnode, text width=140mm, minimum height=11mm, font=\bfseries\small},
  caution/.style={fnode accent, align=left, text width=44mm, minimum height=32mm, font=\footnotesize},
  arr/.style={farr}
]

\node[period] (period) {対象期間: 2025年後半〜2026年前半（生成AI・エージェント技術が投資業界で実装フェーズへ移行した時期のスナップショット）};

\node[caution, below=9mm of period.south, xshift=-48mm] (c1) {\textbf{(a) ベンチマーク\,vs.\,実運用}\\[1mm]
  精度・Sharpe比・コスト削減率の多くはバックテストや限定的パイロットの数値。複数年の実市場検証を経たのは一部（Bridgewater AIA Labs、Renaissance Technologies等）に限られる};
\node[caution, right=4mm of c1] (c2) {\textbf{(b) 再現性の限界}\\[1mm]
  LLMトレーディング研究の多くは取引コストのモデル化や時間整合的な評価プロトコルを欠き、公表された「アウトパフォーム」を額面通りに受け取れない（\S\ref{sec:llm-analysis}）};
\node[caution, right=4mm of c2] (c3) {\textbf{(c) 情報の非対称性}\\[1mm]
  大手ヘッジファンドは技術詳細を非公開とする傾向が強く、参照できるのは意図的に公開された情報のみ。「非公開＝検証不能」と解釈すべき};

\draw[arr] (period.south) -- ++(0,-4mm) -| (c1.north);
\draw[arr] (period.south) -- (c2.north);
\draw[arr] (period.south) -- ++(0,-4mm) -| (c3.north);
\end{tikzpicture}
\end{center}
{\small\color{brandgray}\textbf{図: 本書の対象範囲と留意点}\ 2025年後半〜2026年前半のスナップショットである本書は、(a)ベンチマークと実運用の乖離、(b)再現性の限界、(c)情報の非対称性という3つの留意点を前提に読む必要がある。}

\section{AI-Ready金融データ基盤}
\label{sec:data-infra}

\subsection{課題の本質: モデルではなくデータが律速する}

金融データをLLM／AIエージェントが確実に利用できる形式へ変換する取り組みは、本書が精査した212件のうち最大の比重を占めるテーマである。共通して浮かび上がるのは、「フロンティアモデルの能力」よりも「文書準備・前処理・検索アーキテクチャ」がボトルネックになっているという構図である。

\begin{casestudy}{FinanceBench: 企業導入の最低ラインを可視化した基準ベンチマーク}
上場企業の開示資料から作成された10,231問のQAベンチマークにおいて、GPT-4-Turbo＋検索は\textbf{81\%の質問で誤答または回答拒否}となった。マルチエージェントRAGを用いても精度は56\%が上限であり、企業が財務QAシステムを導入する前に越えるべき最低ラインを定量的に示した。（出典: \cite{financebench2023}）
\end{casestudy}

\begin{casestudy}{多構造金融文書の理解可能性: 前処理が精度を決める}
GPT-4oはネストした表を含む金融PDFに対して前処理なしでは\textbf{56\%}の精度に留まるが、産業用途＋オープンソースの前処理パイプラインを組み合わせるとGPT-4は\textbf{76\%}まで改善する。モデル能力ではなく文書準備こそが鍵であることを直接示した。（出典: \cite{multistructured2025}）
\end{casestudy}

\begin{insight}{示唆}
「モデルを高性能なものに差し替える」だけでは精度は頭打ちになる。実務上のROIは、OCR・レイアウト認識・チャンキング・メタデータ付与といった前処理層への投資に強く依存する。
\end{insight}

本節全体を貫く実証パターンとして、モデル自体を変えずに周辺のアーキテクチャ層だけを着脱した3つの独立した研究が、いずれも精度を劇的に変動させている。

\begin{center}
\begin{tikzpicture}
\begin{axis}[
    ybar,
    bar width=16pt,
    width=0.85\textwidth,
    height=6cm,
    ymin=0, ymax=95,
    ylabel={精度（\%）},
    ylabel style={font=\small},
    symbolic x coords={前処理パイプライン,シンボリック検証層,企業横断検索},
    xtick=data,
    x tick label style={font=\small},
    y tick label style={font=\small},
    legend style={at={(0.5,-0.28)}, anchor=north, legend columns=2, font=\small, draw=none},
    nodes near coords,
    nodes near coords style={font=\scriptsize},
    ymajorgrids=true,
    major grid style={draw=figlight},
    axis line style={draw=figgray},
    tick align=outside,
    enlarge x limits=0.3,
]
\addplot[fill=figblue, draw=fignavy] coordinates {(前処理パイプライン,76) (シンボリック検証層,82) (企業横断検索,23.91)};
\addplot[fill=figaccent, draw=fignavy] coordinates {(前処理パイプライン,56) (シンボリック検証層,18) (企業横断検索,1.27)};
\legend{適切なアーキテクチャあり,なし／崩壊}
\end{axis}
\end{tikzpicture}
\end{center}

{\small \textbf{図: アーキテクチャ投資の有無が精度を左右する三事例。}前処理パイプライン（GPT-4o、\cite{multistructured2025}）、シンボリック検証層（AUDITFLOW、\cite{auditflow2026}）、企業横断の階層的検索（Fin-RATE、\cite{finrate2026}）のいずれも、モデル自体は同一のままアーキテクチャ層の有無だけで精度が大きく変動している。「モデル選びよりアーキテクチャ」が本節を貫く中心テーゼである。}

\subsection{ドキュメントインテリジェンスと前処理パイプライン}

金融文書は表・脚注・多言語混在・レイアウト崩れが常態であり、LLMに投入する前段の抽出・構造化工程の質がそのまま下流タスクの精度上限を規定する。

\begin{concept}{データ抽出と前処理パイプライン}
PDFやスキャン画像のまま渡された金融文書は、LLMにとって「読める」形をしていない。前処理パイプラインとは、生文書から機械可読なデータを取り出す一連の工程を指す。OCR（文字を画像から読み取る技術）でテキスト化し、レイアウト解析でどこが見出し・本文・表かを判定し、VLM（画像と文章を同時に理解できるモデル）で図表を解釈し、入れ子になった表や単位・注記を含む数値を正規化する。いわば「手書きの領収書の束を、そのまま会計ソフトに読み込ませられる電子データへ整える下ごしらえ」であり、この下ごしらえが雑だと、どれほど優秀な調理人（LLM）でも良い料理は作れない。
\begin{center}
\begin{tikzpicture}[node distance=8mm, every node/.style={font=\small},
  stg/.style={fnode, text width=22mm, align=center}]
\node[fnode muted, text width=22mm] (raw) {生文書\\(PDF・画像・表)};
\node[stg, right=10mm of raw] (ocr) {OCR\\レイアウト解析};
\node[stg, right=10mm of ocr] (vlm) {VLM抽出\\テーブル認識};
\node[fnode blue, right=10mm of vlm, text width=24mm] (struct) {構造化データ\\(JSON/整形表)};
\draw[farr] (raw) -- (ocr);
\draw[farr] (ocr) -- (vlm);
\draw[farr] (vlm) -- (struct);
\node[fnode accent, below=10mm of vlm, text width=30mm] (llm) {LLM／RAGへ投入};
\draw[farr accent] (struct.south) -- ++(0,-3mm) -| (llm.north);
\end{tikzpicture}
\end{center}
{\small\color{brandgray}\textbf{図: 前処理パイプラインの基本フロー}\ 生文書はOCR・レイアウト解析・VLM抽出を経て構造化データへ変換され、初めてLLM／RAGが安定して利用できる形になる。}
\end{concept}

\begin{itemize}
  \item \textbf{産業KYCワークフロー向け多段階抽出パイプライン}: 画像前処理→多言語OCR→ハイブリッドページレベル検索→VLM抽出の4段階システムを120件の実運用KYC文書に適用し、\textbf{精度87.27\%}を達成。アブレーション分析ではページレベル検索が支配的要因であり、特に非英語財務諸表で効果が顕著だった。（\cite{multistageextraction2026}）
  \item \textbf{本番運用のマイクロサービス型OCR+LLMパイプライン}: 毎時数千件のマルチページ文書を処理する実運用アーキテクチャでは、エンドツーエンドのボトルネックはLLM推論ではなく\textbf{OCRのレイテンシ}であり、システムの飽和点を決めるのはワーカー数ではなくGPU容量であることが判明した。（\cite{documentai2026}）
  \item \textbf{DataClawBench}: 企業・産業・政策ドメインにまたがる492件のマルチステップ分析タスク、206万レコードをノイズ除去せずそのまま用いたエージェントベンチマーク。現行LLMエージェントは探索を増やしても正答に確実に結びつかず、ベンチマーク精度と実運用の“汚い”データとの間に根本的なギャップがあることを示した。（\cite{dataclawbench2026}）
  \item \textbf{連邦準備制度理事会（FRB）のRAG実証研究}: G-SIB（グローバルなシステム上重要な銀行）の開示文書数千ページからの情報抽出において、対話的なRAG利用が人手のみのベースラインに対し抽出速度を最大\textbf{10倍}に加速し精度も改善、1プロジェクトあたり最大\textbf{268時間}の注釈作業を節約できることを統制実験で示した。中央銀行における前処理・抽出支援の実証事例として重要。（\cite{fedreserverag2025}）
  \item \textbf{MacroLens}: 米国株4,416銘柄（2021〜2026年）を対象に、4,680万件のXBRL会計エントリ、約29.6万件のSECファイリング、53のマクロ経済指標、21.5万件超のニュース記事を統合したマルチタスク・ベンチマーク。実データでLLMを評価する際、時間ゲーティング・報告遅延・レジームリーク（regime leakage）が体系的な失敗モードになることを明らかにした。（\cite{macrolens2026}）
  \item \textbf{長文コンテキスト化（DocFinQA）}: FinQAの7,437問に文書全体のコンテキストを付与し、平均文脈長を700語未満から\textbf{約12.3万語}へ拡張した長文脈金融推論データセット。財務プロフェッショナルが実際に扱うのは数百ページの文書であるという前提を反映しており、検索型QAパイプライン・長文脈LLMのいずれも最長文書で顕著に性能が劣化することを示した。抜粋前提の既存ベンチマークが実運用の難度を過小評価していることを定量化した事例である（ACL 2024、\cite{x02docfinqa2024}）。
  \item \textbf{その他の応用}: 200ページ超の政府財政PDFに階層型LLMパイプラインと財政表の合計値構造を数値検証機構として活用し研究可能なDBへ変換する事例（\cite{fiscaldocs2025}）、決算コールのトランスクリプトから階層的トピック分類を抽出し戦略優先事項の四半期変化を追跡するAgentic手法（SIGIR 2025、\cite{earningscalls2025}）など、前処理パイプラインの応用範囲は拡大している。
\end{itemize}

\begin{casestudy}{FinChart-Bench: チャート画像理解は依然として構造的な弱点}
金融文書の情報はテキストと表だけでなく\textbf{チャート画像}にも埋め込まれている。実世界の金融チャートに特化した初のベンチマークFinChart-Bench（2015〜2024年のBloombergニュース・企業プレゼン資料由来の1,200枚のチャート画像、真偽・多肢選択・自由記述の\textbf{7,016問}）で、\textbf{25種類の最新の視覚言語モデル（LVLM）}を評価した結果、(1)オープンソースとクローズドの性能差が縮小している一方、(2)同一ファミリーの新しいモデルでむしろ性能が\textbf{劣化}する事例があり、(3)空間推論・指示追従に重大な限界があること、そして(4)LVLM自身を自動評価器として使うには信頼性が不足することが示された。テキスト・表・チャートの三モダリティを横断する前処理設計が、金融ドキュメントインテリジェンスの次の課題であることを裏づける（\cite{x02finchartbench2025}）。
\end{casestudy}

\begin{center}
\begin{tikzpicture}[
  stage/.style={fnode, text width=2.9cm, minimum height=2.0cm, font=\footnotesize},
  result/.style={fnode accent, text width=13.6cm, minimum height=1.3cm, font=\small},
  arr/.style={farr}
]
\node[stage] (s1) at (0,0) {① 画像前処理};
\node[stage] (s2) [right=5mm of s1] {② 多言語OCR};
\node[stage] (s3) [right=5mm of s2] {③ ハイブリッド\\ページレベル検索};
\node[stage] (s4) [right=5mm of s3] {④ VLM抽出};
\draw[arr] (s1) -- (s2);
\draw[arr] (s2) -- (s3);
\draw[arr] (s3) -- (s4);
\node[fit=(s1)(s4), inner sep=0pt] (row) {};
\node[result, below=7mm of row] (res) {実運用KYC文書120件で\textbf{精度87.27\%}を達成。アブレーションでは③ページレベル検索が支配的要因（特に非英語財務諸表で顕著）};
\draw[arr] (row.south) -- (res.north);
\end{tikzpicture}
\end{center}

{\small\color{brandgray}\textbf{図: 産業KYCワークフロー向け4段階抽出パイプライン}\ 画像前処理からVLM抽出までの4段階を経て精度87.27\%を達成し、ページレベル検索が精度を左右する支配的要因であることが判明した（\cite{multistageextraction2026}）。}

\subsection{RAGアーキテクチャの進化: フラット検索から階層型・グラフ型へ}

RAG（検索拡張生成）は金融文書QAの標準的アプローチだが、単純なベクトル類似度検索（flat dense retrieval）では企業横断・期間横断のタスクで性能が崩壊することが繰り返し確認されている。

\begin{concept}{RAGと検索アーキテクチャ}
RAG（Retrieval-Augmented Generation、検索拡張生成）とは、質問に関連する文書の断片をあらかじめ検索してLLMに手渡してから回答させる仕組みである。LLMの記憶だけに頼ると古い情報や捏造（ハルシネーション）が混じるため、「回答の前に必ず該当ページを開いて確認させる、根拠付きのオープンブック試験」に近い。検索方式にも段階がある。フラット検索は全文書を一つの巨大な本棚から探すため、似た企業名や年度の資料を取り違えやすい。階層型検索はまず「企業」「年度」で書棚を仕分けてから探すため精度が上がる。グラフ型検索はエンティティ間の関係を辿る地図を使い、複数ステップの推論を要する質問にも対応できる。
\begin{center}
\begin{tikzpicture}[node distance=6mm, every node/.style={font=\small}]
\node[fnode] (q) {質問};
\node[fnode blue, right=9mm of q, text width=24mm] (r) {関連文書を検索};
\node[fnode accent, right=9mm of r, text width=20mm] (g) {文脈＋LLM};
\node[fnode dark, right=9mm of g, text width=22mm] (a) {根拠付き回答};
\draw[farr] (q) -- (r);
\draw[farr] (r) -- (g);
\draw[farr] (g) -- (a);
\node[fnode muted, below=10mm of q, text width=20mm, xshift=6mm] (flat) {フラット\\全文書を一括検索};
\node[fnode muted, right=6mm of flat, text width=20mm] (hier) {階層型\\企業/年度で仕分け};
\node[fnode muted, right=6mm of hier, text width=20mm] (graph) {グラフ型\\関係を辿る};
\end{tikzpicture}
\end{center}
{\small\color{brandgray}\textbf{図: RAGの基本フローと検索方式の違い}\ 質問に対し関連文書を検索し文脈として渡す基本フロー（上段）と、フラット・階層型・グラフ型という3つの検索方式の違い（下段）。}
\end{concept}

\begin{casestudy}{Fin-RATEベンチマーク: 検索失敗が精度崩壊の主因}
Goldman Sachs\ +\ Yaleによる17モデル評価で、企業横断比較タスクにおいてファインチューン済み金融LLMの精度が\textbf{23.91\%から1.27\%}へ崩壊し、時系列トラッキングでも\textbf{18.6ポイント}低下することが確認された。専門家パネルは「ボトルネックは検索失敗であり生成品質ではない」「企業・年度で事前にバケット化する階層的検索が結果を劇的に改善する」と結論している。\cite{finrate2026}
\end{casestudy}

\begin{casestudy}{FinReflectKG-MultiHop: ナレッジグラフ検索の優位性}
S\&P 100の10-Kから構築した1,750万トリプレットの知識グラフを用いたKGガイド検索は、フラット検索に対して\textbf{正答率約24\%向上、トークン使用量約84.5\%削減}を達成した（ICAIF 2025）。\cite{finreflectkg2025}
\end{casestudy}

その他、検索アーキテクチャに関する主要な研究・実装のうち、特筆すべき成果を挙げる。

\begin{itemize}
  \item \textbf{FinSage}: マルチモーダル前処理・ハイブリッドスパース密検索・DPOチューニングによるリランキングを組み合わせたマルチアスペクトRAG。財務提出書類QAで\textbf{再現率92.51\%、ベースライン比精度24.06\%改善}を達成。単純なテキストのみRAGはFinanceBenchで19\%に留まる一方、マルチエージェントRAGは56\%を記録した。（\cite{finsage2025}）
  \item \textbf{MimirRAG}: テーブル認識型チャンキングとメタデータ抽出を組み込んだマルチエージェントRAGで、FinanceBenchにおいて\textbf{89.3\%}の精度を達成。（\cite{mimirrag2026}）
  \item \textbf{HybridRAG}: ナレッジグラフ検索（KG-RAG）とベクトル検索（VectorRAG）を統合したハイブリッド手法。決算コールのトランスクリプトを対象に、検索精度・回答品質の両面でいずれの単独手法をも一貫して上回ることを示し、異種金融文書向けハイブリッドRAGの原型となった。（\cite{hybridrag2024}）
  \item \textbf{PDF解析・チャンキングの実証評価}: 複数のPDFパーサとチャンク重複戦略を、FinanceBenchおよび新規ベンチマークTableQuest上で体系的に比較し、表・数式・文書構造を忠実に保持するRAGパイプラインの実務指針を導出した。前処理段の選択がそのまま下流QA精度を左右することを追試している。（\cite{pdfparsingrag2026}）
  \item \textbf{BM25対密検索のベンチマーク}: テキストと表が混在する文書の金融QAクエリ23,088件を対象とした2026年のベンチマークでは、BM25単体がRecall@20を除く全指標で最新の密検索を上回った。検索失敗の\textbf{73\%}は表構造の不整合に起因し、HyDEのようなクエリ拡張はLLMが架空の数値を生成するため逆に性能を悪化させる。密検索を安全な既定値とみなす業界の通念に反証を突きつけている。（\cite{bm25hybrid2026}）
  \item \textbf{金融QA向けRAG検索戦略の最適化}: 3段階RAG（事前検索・ハイブリッド密＋疎検索・後処理再ランキング）がFinanceBench、FinQA、TATQAなど7つの金融QAデータセットで精度を大きく改善する再現可能な設計図を提示（ICLR 2025ワークショップ）。（\cite{ragretrieval2025}）
  \item \textbf{SMARTFinRAG}: 検索・生成コンポーネントを実行時に交換可能なモジュール型オープンソースRAGプラットフォーム。モデル自体を固定してもRAGの構成要素選択だけで性能に大きな差が出ることを実証し、「モデルよりコンポーネント選択」というアーキテクチャ優位のテーゼを補強する。（\cite{smartfinrag2025}）
  \item \textbf{T\textsuperscript{2}-RAGBench}: テキストと表が混在する実文書に対する\textbf{32,908件}の問い–文脈–解答トリプルからなるRAG評価ベンチマーク。既存データセットの多くが正解文脈を明示的に与える「Oracle設定」であるのに対し、本ベンチマークはモデルにまず正しい文脈を\textbf{検索}させたうえで数値推論を要求する。既存データセットを文脈非依存の形式へ変換し（専門家検証で91.3\%が文脈非依存）、データ汚染を排除することで、検索段を含めた現実的なRAG性能を測れるようにした（EACL 2026、\cite{x02t2ragbench2026}）。
  \item \textbf{HiREC（階層的検索＋エビデンス精選）}: 標準化された開示文書に対するオープンドメイン金融QA向けに、階層的検索と生成前のエビデンス精選（evidence curation）を組み合わせ、無関係な文脈を削減する手法。専用ベンチマークLOFin-benchを併せて提案し、長文の標準化文書で検索・QA双方の性能を改善した。「モデル規模ではなく検索アーキテクチャがオープンドメインQAの律速」という本節のテーゼを補強する（\cite{x02hirec2025}）。
  \item \textbf{メタデータ駆動RAG（contextual chunks）}: FinanceBench上でLLM生成メタデータを活用する多段RAGを構築し、最大の精度向上は再ランカーやクエリ拡張ではなく\textbf{チャンクにメタデータをテキストごと埋め込む「文脈化チャンク」}から得られることを実証。強力な再ランカーは精度に必須だが、自作のメタデータ再ランカーが商用品より費用対効果に優れる点も示した（\cite{x02metadatarag2025}）。
  \item \textbf{検索パラダイムの再考（サーベイ）}: 金融ドメインでは古典的なベクトル類似度RAGが、エージェント型・非ベクトル型の推論システムへ移行しつつあるとし、密検索・エージェント判断型検索・SQL／グラフ等の代替機構を切り分けるタクソノミを提示した（\cite{x02rethinkretrieval2025}）。
  \item その他、実務家の略語・簡潔表現を反映した5,703件のクエリ三つ組\textbf{FinDER}（\cite{finder2025}）、クエリ再構成・クロスエンコーダ再ランキングを組み込んだ\textbf{Fintech向けAgentic RAG}（\cite{ragfintech2025}）、NASDAQファンダメンタルズで前処理・照会層を追加し精度78.6\%を達成した\textbf{LLM-RAG金融データ分析システム}（\cite{finllmrag2025}）など、モジュール構成を競う実装が相次いでいる。
\end{itemize}

\begin{table}[htbp]
\centering
\small
\caption{金融RAG／QAベンチマークの比較（本節で参照する主要データセット）}
\label{tab:finrag-benchmarks}
\begin{tabularx}{\textwidth}{@{}l r l X@{}}
\toprule
\textbf{ベンチマーク} & \textbf{規模} & \textbf{焦点} & \textbf{中心的な知見} \\
\midrule
FinanceBench \cite{financebench2023} & 10,231問 & 開示資料QA & GPT-4-Turbo＋検索でも81\%が誤答・拒否 \\
FinQA拡張・DocFinQA \cite{x02docfinqa2024} & 7,437問 & 長文脈推論 & 平均12.3万語、最長文書で性能劣化 \\
Fin-RATE \cite{finrate2026} & 17モデル評価 & 企業横断・時系列 & 検索失敗で精度23.91→1.27\%へ崩壊 \\
T\textsuperscript{2}-RAGBench \cite{x02t2ragbench2026} & 32,908トリプル & テキスト＋表検索 & Oracle文脈を排除、汚染なしの現実的評価 \\
FinChart-Bench \cite{x02finchartbench2025} & 7,016問 & チャート画像理解 & 25 LVLMで空間推論・指示追従に限界 \\
FinBalance \cite{x02finbalance2026} & 710レコード & 多文書照合 & 最良モデルでも貸借一致は46\%止まり \\
\bottomrule
\end{tabularx}
\end{table}

{\small\color{brandgray}\textbf{表\ref{tab:finrag-benchmarks}について:}\ いずれのベンチマークも、モデル自体の性能ではなく検索段・前処理段・数値照合段の設計が精度上限を規定することを、異なる切り口から繰り返し示している。}

\begin{insight}{仮説（Takes）: GraphRAGがフラットRAGを置き換える}
KG-based検索の優位性を示す複数の実証結果を根拠に、知識グラフ拡張検索（GraphRAG）が今後2〜3年でエンタープライズ金融AIにおける企業横断・縦断タスクの標準アーキテクチャとなるという仮説が提示されている。ただしグラフ構築コストとエンティティ抽出パイプラインの整備負担、および事前定義タクソノミが新規のクロスドメインクエリを捕捉できないリスクが反証材料として指摘されている。
\end{insight}

\begin{center}
\begin{tikzpicture}[
  font=\small,
  oldbox/.style={fnode muted, text width=6.4cm, align=left, minimum height=4.4cm, font=\footnotesize},
  newbox/.style={fnode blue, text width=6.4cm, align=left, minimum height=4.4cm, font=\footnotesize},
  arr/.style={farr accent}
]
\node[oldbox] (flat) at (0,0) {
  \textbf{フラット密検索（Flat Dense Retrieval）}\\[2mm]
  全文書を単一ベクトル空間で類似度検索\\[2mm]
  企業横断比較タスク: 精度\textbf{23.91\%\,→\,1.27\%}に崩壊\\[1mm]
  時系列トラッキング: \textbf{-18.6pt}\ \cite{finrate2026}
};
\node[newbox, right=10mm of flat] (hier) {
  \textbf{階層型・グラフ型検索}\\[2mm]
  企業・年度で事前バケット化＋KGガイド検索\\[2mm]
  KG検索（FinReflectKG）: 正答率\textbf{約+24\%}\\トークン使用量\textbf{約-84.5\%}\ \cite{finreflectkg2025}
};
\draw[arr] (flat.east) -- (hier.west) node[midway, above=1mm, font=\footnotesize\bfseries, color=figaccent] {アーキテクチャ進化};
\end{tikzpicture}
\end{center}

{\small\color{brandgray}\textbf{図: フラット検索から階層型・グラフ型検索への進化}\ 同一のクエリ・モデル条件下でも、検索アーキテクチャの選択だけでFin-RATEでは精度が崩壊し、FinReflectKG-MultiHopのKGガイド検索では逆に精度・効率とも改善する。}

\subsection{埋め込みモデルとベクトル検索: 汎用埋め込みの限界とドメイン適応}

RAGの検索精度は最終的に埋め込みモデルの表現力に依存する。ところが、汎用テキスト埋め込みは金融特有の語彙・数値表現・文脈で体系的に劣化することが複数の研究で示され、ドメイン適応済み埋め込みの重要性が浮き彫りになっている。

\begin{concept}{埋め込みとベクトル検索}
埋め込み（embedding）とは、文章を「意味を数値で表した座標点」に変換する技術である。意味が近い文同士は座標空間上でも近い位置に置かれる仕組みで、いわば「意味を地図上の緯度経度に変換する」作業といえる。ベクトル検索は、質問文を同じ座標に変換し、最も近い場所にある文書を探し出す（最近傍探索）ことでRAGの土台を支えている。ただし汎用の埋め込みモデルは、決算用語や数値表現・単位といった金融特有の語彙を「意味の近さ」として正しく捉えられないことが多く、専門分野向けに再調整（ドメイン適応）したモデルが必要になる。
\begin{center}
\begin{tikzpicture}[font=\small]
\node[fnode, text width=20mm] (text) at (0,0) {テキスト};
\node[fnode blue, text width=22mm] (emb) at (3.3,0) {埋め込み\\モデル};
\node[fnode muted, minimum width=36mm, minimum height=28mm,
  label={[font=\footnotesize\color{brandgray}]above:ベクトル空間（意味の地図）}] (space) at (8.7,0) {};
\draw[farr] (text) -- (emb);
\draw[farr] (emb) -- (space.west);
\fill[figgray] (7.7,0.6) circle (1.4pt) node[above=0.4mm, font=\scriptsize] {文書};
\fill[figgray] (9.6,0.7) circle (1.4pt) node[above=0.4mm, font=\scriptsize] {文書};
\fill[figgray] (9.5,-0.7) circle (1.4pt) node[below=0.4mm, font=\scriptsize] {文書};
\fill[figaccent] (8.4,-0.2) circle (1.9pt) node[below=0.5mm, font=\scriptsize, text=figaccent] {クエリ};
\draw[farr accent] (8.4,-0.2) -- (7.7,0.6);
\end{tikzpicture}
\end{center}
{\small\color{brandgray}\textbf{図: 埋め込みとベクトル検索の仕組み}\ テキストを埋め込みモデルでベクトル空間上の点に変換し、クエリに最も近い点（文書）を検索する。}
\end{concept}

\begin{casestudy}{FinMTEB: 汎用ベンチマークの成績は金融タスクを予測しない}
64個の金融ドメイン特化データセット・7タスク（英中両言語、金融ニュース・年次報告書・ESGレポート・規制提出書類・決算コール）で15種類の埋め込みモデルを評価した結果、(1)汎用埋め込みベンチマークの成績は金融ドメインタスクの成績と\textbf{相関が弱く}、(2)ドメイン適応モデルが一貫して汎用モデルを上回り、(3)金融の意味的類似度（STS）タスクでは\textbf{単純なBag-of-Wordsが高度な密埋め込みを上回る}という反直感的な結果が得られた。ペルソナベースの合成データで訓練した金融特化モデルFin-E5は平均\textbf{0.6767}でSOTAを達成し、商用モデルをも凌駕した（EMNLP 2025、\cite{x02finmteb2025}）。
\end{casestudy}

\begin{casestudy}{BAM埋め込み: ドメイン適応がRecall@1を大きく引き上げる}
金融文書は一般テキストとは語彙分布が大きく異なるため埋め込みも適応が必要だとする研究では、1,430万件のクエリ–パッセージ対でファインチューンしたBAM埋め込みが、保留テストセットでRecall@1\ \textbf{62.8\%}を達成した（OpenAIの最良汎用埋め込みは39.2\%）。FinanceBench上のQA精度も\textbf{8\%}改善しており、検索段の埋め込み選択がRAG全体の精度上限を直接押し上げることを定量的に示した（EMNLP 2024、\cite{x02bamembeddings2024}）。
\end{casestudy}

\begin{insight}{示唆: 埋め込み層は「静かな律速」}
検索アーキテクチャの議論はチャンキングや再ランキングに集中しがちだが、その土台となる埋め込みモデルのドメイン適応度こそが検索品質の下限を決める。汎用ベンチマークのスコアで埋め込みを選定すると金融タスクでの実効性能を見誤りやすく、金融STSではBoWのような素朴な手法が密埋め込みを上回る領域さえ残る。埋め込み層は表面に現れにくい「静かな律速」として、前処理・再ランキングと並ぶ投資対象と位置づけるべきである。
\end{insight}

\subsection{財務知識グラフ: 検索を超えた推論の足場}

知識グラフはRAGの検索精度改善だけでなく、エンティティ間関係を明示的に保持することでマルチホップ推論・信用リスク評価・投資判断の足場としても機能し始めている。

\begin{concept}{ナレッジグラフ（知識グラフ）}
ナレッジグラフとは、企業・人物・製品といった「モノ（エンティティ）」を点（ノード）、モノ同士の関係を線（エッジ）で表したデータ構造である。最小単位は「主語―述語―目的語」の三つ組（トリプレット）で、例えば「A社―供給する―B社」のように表現する。表やベクトル検索が「一つの文書の中身」を扱うのに対し、ナレッジグラフは関係そのものを明示的に持つため、「A社のサプライヤーの競合はどこか」のように複数のリンクをたどるマルチホップ推論が可能になる。いわば「人名と会社名だけの名簿」ではなく「誰と誰がどうつながっているかを描いた相関図」である。
\begin{center}
\begin{tikzpicture}[node distance=13mm, every node/.style={font=\small}]
\node[fnode dark, text width=16mm] (a) {A社};
\node[fnode, above right=10mm and 20mm of a, text width=16mm] (ceo) {CEO};
\node[fnode blue, below right=10mm and 20mm of a, text width=16mm] (supplier) {サプライヤーC};
\node[fnode accent, right=32mm of a, text width=18mm] (rival) {競合B社};
\draw[farr] (a) -- node[midway, above, font=\scriptsize] {率いる} (ceo);
\draw[farr] (a) -- node[midway, below, font=\scriptsize] {供給} (supplier);
\draw[farr accent] (supplier) -- node[midway, below, font=\scriptsize] {供給} (rival);
\draw[farr dash] (a) -- node[midway, above, font=\scriptsize] {競合} (rival);
\end{tikzpicture}
\end{center}
{\small\color{brandgray}\textbf{図: ナレッジグラフとマルチホップ推論}\ A社からサプライヤーCを経て競合B社へ至る経路（濃い矢印）のように、複数のエッジをたどることで「サプライヤーの競合」といった多段の問いに答えられる。}
\end{concept}

\begin{itemize}
  \item \textbf{FinDKG}: 金融ニュース・開示文書から動的知識グラフを構築するICKG（ファインチューン済みLLM生成器）、オープンソースデータセット、KGTransformer（アテンション型GNN）を提案。リンク予測で優れた性能を示し、実運用のテーマ型ETFを上回る投資成果を記録（ACM ICAIF 2024）。（\cite{findkg2024}）
  \item \textbf{FinKG-News}: ニュース中心の金融知識グラフを用いて根拠付き信用リスクレポートを生成。ベースライン比で\textbf{品質19〜34\%向上}を達成する一方、自動ハルシネーション検出は信頼性不足であり専門家レビューの必要性を明確に指摘した。（\cite{creditrisk2026}）
  \item \textbf{Agentic GraphRAG}: フラットなベクトルRAGをNeo4jナレッジグラフに置き換え、スイス商業登記の7年分のデータから構築。分析エージェントがグラフ上をクエリし、フラットRAGベースラインを一貫して上回る。（\cite{agenticgraphrag2026}）
\end{itemize}

\begin{center}
\begin{tikzpicture}[
  src/.style={fnode muted, text width=2.4cm, minimum height=1.3cm, font=\scriptsize},
  kg/.style={draw=fignavy, line width=0.8pt, fill=brandlight, figshadow, circle, minimum size=24mm, font=\small\bfseries, align=center, inner sep=1mm},
  outbox/.style={fnode blue, text width=5.6cm, align=left, minimum height=1.5cm, font=\scriptsize},
  arr/.style={farr}
]
\node[src] (news) at (0,2.4) {ニュース};
\node[src] (filing) at (0,0) {開示文書};
\node[src] (registry) at (0,-2.4) {商業登記\\データ};
\node[kg] (kg) at (4.0,0) {財務\\知識\\グラフ};
\node[outbox] (findkg) at (10.2,2.4) {\textbf{FinDKG}: リンク予測に優れ、テーマ型ETFを上回る投資成果\ \cite{findkg2024}};
\node[outbox] (finkgnews) at (10.2,0) {\textbf{FinKG-News}: 根拠付き信用リスクレポートの品質\textbf{+19〜34\%}\ \cite{creditrisk2026}};
\node[outbox] (agentic) at (10.2,-2.4) {\textbf{Agentic GraphRAG}: Neo4j上のクエリでフラットRAGを一貫して上回る\ \cite{agenticgraphrag2026}};

\draw[arr] (news) -- (kg);
\draw[arr] (filing) -- (kg);
\draw[arr] (registry) -- (kg);
\draw[arr] (kg) -- (findkg);
\draw[arr] (kg) -- (finkgnews);
\draw[arr] (kg) -- (agentic);
\end{tikzpicture}
\end{center}

{\small\color{brandgray}\textbf{図: 財務知識グラフのデータフローと応用}\ ニュース・開示文書・登記データを統合した知識グラフが、リンク予測、信用リスク評価、分析エージェントのクエリ基盤という3つの異なる用途で検索を超えた推論の足場として機能する。}

\subsection{数値精度とシンボリック検証: LLM単体の限界}

複数の独立した研究が、LLM単体では正確な会計演算・数値照合を安定して行えないことを示している。この弱点を補うシンボリック検証層の導入が、実運用可否を分ける分岐点になりつつある。

\begin{concept}{ハルシネーションと忠実性}
ハルシネーションとは、LLMがもっともらしい文章や数値を、根拠がないまま作り出してしまう現象である。文章の要約なら多少の言い回しの違いは許容されるが、金融では「売上高を1桁間違える」「存在しない条項を引用する」といった誤りが致命的な損失につながりかねない。忠実性（faithfulness）とは、出力が実際の根拠（検索した文書や計算結果）にどれだけ忠実かを表す指標であり、文章が流暢で自信ありげに見えることと、内容が正確であることは別物である。この弱点を補うのがシンボリック検証層——数式や会計規則に基づいて機械的に計算・照合し直す仕組みで、いわば「作文の採点」とは別に「電卓で検算する係」を置くようなものである。
\begin{center}
\begin{tikzpicture}[node distance=10mm, every node/.style={font=\small}]
\node[fnode, text width=30mm] (llm) {LLMのそれらしい出力};
\node[fnode accent, right=16mm of llm, text width=30mm] (check) {検証層\\数値照合・根拠チェック};
\node[fnode blue, above right=10mm and 12mm of check, text width=22mm] (ok) {合格→採用};
\node[fnode muted, below right=10mm and 12mm of check, text width=22mm] (ng) {不一致→差し戻し};
\draw[farr] (llm) -- (check);
\draw[farr dash] (check) -- (ok);
\draw[farr dash] (check) -- (ng);
\end{tikzpicture}
\end{center}
{\small\color{brandgray}\textbf{図: シンボリック検証層による忠実性チェック}\ LLMの出力を数値照合・根拠チェックにかけ、合格したものだけを採用し不一致は差し戻す。}
\end{concept}

\begin{casestudy}{AUDITFLOW: シンボリック検証層がなければ精度は崩壊する}
米国GAAPタクソノミとXBRL提出事実に基づく監査検証環境で、2ジュニア＋1シニアのマルチエージェント監査フレームワークからシンボリック検証層を除去すると、精度は\textbf{82\%から18\%}へ崩壊した。LLM単体では信頼できる監査検証に不十分であることを定量的に実証した。\cite{auditflow2026}
\end{casestudy}

\begin{casestudy}{FinVerBench: 数値の丸め処理だけで再現率が21ポイント低下}
14種類のLLMをSEC 10-Kの財務諸表検証タスクで評価した結果、現実的な数値の丸め処理を行っただけで最良モデルの再現率が\textbf{100\%から79\%}へ低下した。14モデル中9モデルは、誤りのないクリーンな財務諸表に対して\textbf{95〜100\%という高い偽陽性率}を示した。モデル間の性能差より前処理（数値表記の正規化）による差の方が大きいことを示す一方、過度な正規化は契約条項の閾値判定を覆すリスクもあり、規制・契約上の閾値近傍では人間レビューが必要と結論している。\cite{finverbench2026}
\end{casestudy}

\begin{itemize}
  \item \textbf{FinGround}: 原子的claim検証による金融ハルシネーション検出の3段階パイプライン。既存の検出器は意味的類似性に依存するため金融文書の計算誤りの\textbf{43\%を見逃す}ことを実証。GPT-4o比で\textbf{ハルシネーション率78\%削減}を達成し、蒸留版8Bモデルは1クエリあたり約0.003ドルで稼働する。4週間のアナリストパイロットで実運用適性を検証済み。（\cite{finground2026}）
  \item \textbf{FinSheet-Bench}: プライベートエクイティ・ファンドのスプレッドシートを対象にした合成ベンチマーク。最良モデル（Gemini 3.1 Pro）は単純ファイルで82\%だが、大規模マルチファンド構造では\textbf{49\%}まで急落し、文書理解と決定的計算の分離の必要性を示唆。（\cite{finsheetbench2026}）
  \item \textbf{Finch}: Enronアーカイブと金融機関（2000〜2025年）から収集した1,710個のスプレッドシート（2,700万セル）、384タスク、172の複合ワークフローからなる会計・財務ベンチマーク。最良のGPT-5.1 Proでも平均16.8分のエージェント作業を経て\textbf{ワークフローの38.4\%}しか通過できず、文書理解層と決定的計算エンジンの分離という設計課題を改めて裏づけた。（\cite{finchbench2025}）
  \item \textbf{FinBalance}: ソース文書束を出典付きの仕訳と貸借対照表へ照合させる多文書会計照合ベンチマーク（710レコード、8業種・3期間種別・5難易度、23種の不整合コード）。最良モデルでも\textbf{貸借一致は正答率46\%}に止まり、6モデル中4モデルが自らの仕訳から導かれる貸借対照表と報告値との間に\textbf{26〜41ポイント}の乖離（集約失敗）を示した。さらに勘定科目・金額が正しい仕訳の\textbf{52\%が誤った出典文書を引用}しており、文書理解層と決定的な会計計算エンジンを分離すべきという設計課題を多文書設定で改めて裏づけた（\cite{x02finbalance2026}）。
  \item \textbf{MMFCTUB}: 信用ドメインの5種類の表タイプにまたがる7,600件超のサンプルからなるマルチモーダル信用表理解ベンチマーク。表を跨いだ構造理解・ドメイン知識の適用・数値推論をプロプライエタリ／オープンソース双方のマルチモーダルLLMで評価し、構造化された信用データに対する推論能力の弱点を可視化した。（\cite{mmfctub2026}）
  \item \textbf{FinTagging}: XBRLタグ付け財務情報抽出を数値抽出（FinNI）と概念リンク（FinCL、1万以上のUS GAAP概念空間）に分解して評価。LLMは生テキスト抽出では強いが、細粒度の会計タクソノミへの概念整合は\textbf{約17\%の精度}に留まり、AI-readiness上の重大なギャップとして繰り返し指摘されている。（\cite{fintagging2025}）
\end{itemize}

\begin{center}
\begin{tikzpicture}[
  font=\small,
  layerbad/.style={draw=figgray, thick, rounded corners, fill=figlight, text width=4.8cm, align=center, minimum height=1.8cm, inner sep=2mm, font=\footnotesize},
  layer/.style={draw=figblue, thick, rounded corners, fill=white, text width=4.8cm, align=center, minimum height=1.8cm, inner sep=2mm, font=\footnotesize},
  layertop/.style={draw=figaccent, thick, rounded corners, fill=brandlight, text width=4.8cm, align=center, minimum height=1.6cm, inner sep=2mm, font=\footnotesize},
  lbl/.style={font=\footnotesize\bfseries}
]
\node[layerbad] (bad) at (0,0) {2ジュニア＋1シニアの\\マルチエージェント\\監査フレームワーク};
\node[lbl, color=figgray, above=2mm of bad] {LLM単体（検証層なし）};
\node[lbl, color=figaccent, below=2mm of bad] {精度\ \textbf{18\%}};

\node[layer] (good1) at (8.2,0) {2ジュニア＋1シニアの\\マルチエージェント\\監査フレームワーク};
\node[layertop, above=2mm of good1] (good2) {シンボリック検証層\\(GAAPタクソノミ／XBRL照合)};
\node[lbl, color=fignavy, above=2mm of good2] {LLM\ +\ シンボリック検証};
\node[lbl, color=figblue, below=2mm of good1] {精度\ \textbf{82\%}};
\end{tikzpicture}
\end{center}

{\small\color{brandgray}\textbf{図: シンボリック検証層の有無による精度差（AUDITFLOW）}\ 同一のマルチエージェント監査フレームワークからシンボリック検証層を除去するだけで精度が82\%から18\%へ崩壊し、LLM単体では信頼できる数値照合に不十分であることを示す（\cite{auditflow2026}）。}

\subsection{大規模データプロバイダーのAI-Ready化とMCP標準化}

大手金融データプロバイダーは、自社の情報資産をLLM・エージェントから直接利用可能な形へ再構築する取り組みを2025〜2026年にかけて加速させている。

\begin{casestudy}{FactSet \texttimes\ Snowflake Cortex: エンタープライズ規模のAI-readyコーパス}
FactSetは、グローバル規制開示・決算コール・StreetAccountニュースを標準化・ベクトル化・強化し、Snowflake Cortex Knowledge Extension経由でAWS/Azure/Databricks環境から利用可能なAI-readyコーパスを提供。セマンティック検索とエージェント型モニタリングを実現している。\cite{factset_snowflake2025}
\end{casestudy}

\begin{casestudy}{MCP標準化: データ接続のコモディティ化}
2025〜2026年にBloomberg（ASKB）、FactSet、LSEG/Microsoft、Daloopa、Moody's、S\&P Globalが相次いでMCP（Model Context Protocol）またはMCP互換のエージェント接続を提供開始した。Bloomberg ASKBは数億件規模の文書、800以上のリサーチ提供元、Bloomberg Intelligence等の独自調査に根拠づけて応答し、データ分析時にはBQLコードを提示してExcelやBQuantへ接続できる\cite{x02bloombergai2026}。LSEG/Microsoftは33ペタバイト超のAI-readyコンテンツとタクソノミをMicrosoft 365 Copilot/Copilot Studioに接続し、LSEG管理MCPサーバーで段階的に提供する計画を公表した\cite{x02lsegmicrosoft2025}。標準化されたプロトコルでAIエージェントが構造化金融データにクエリできるようになると、データ接続性そのものの限界価値はゼロに近づき、差別化はデータの解釈・文脈化・活用へシフトするという仮説が提示されている。（出典: \cite{factset_mcp_nd,x02moodysmcp2026,x02spglobalmcp2026}）
\end{casestudy}

\begin{itemize}
  \item \textbf{SageX AI}: NY拠点のスタートアップが、リサーチレポート・決算トランスクリプト・開示資料・メール・オルタナティブデータなど1万種類以上の文書フォーマットを構造化情報に変換する継続型AIデータ変換レイヤーを発表。エンタープライズのデータ処理コストを\textbf{80〜90\%削減}、ノーコードで1日以内に導入可能と主張（2026年4月）。
  \item \textbf{Two Sigma「データをコードとして扱う」＋「Semantic Types」}: データパイプライン自体にバージョン管理・CI/CD・テスト規律を適用する取り組みと、システム間のデータ統合・発見・検証を改善するセマンティックメタデータ型の提案を並行させ、下流のML/LLMシステムが金融データを利用可能にするための基盤整備を進めている。
  \item \textbf{Point72 CSP（Composable Stream Processing）オープンソース化}: ヘッジファンドが社内のストリーム処理フレームワークを初のオープンソースプロジェクトとしてGitHubで公開した異例の事例。同社が持続的な競争優位をデータ配管ではなくAI/MLモデルの質と人材に見出していることを示唆する。
  \item \textbf{ドイツ連邦銀行（Bundesbank）}: 長文の二言語証券目論見書に対する担保適格性チェックを生成AIパイプラインで自動化し、実運用環境で文書レベル精度最大\textbf{91\%}を達成。中央銀行における実装事例として重要。（\cite{llmeligibilitycriteria2026}）
  \item \textbf{S\&P Global「Kensho LLM-ready API」（ベータ）}: S\&P Capital IQ Financials・Compustat・市場データを、Claude／GPT／Geminiなどのモデルを介して自然言語でクエリ可能にするAPI（2024年11月オープンベータ）。生成された関数呼び出しの監査証跡（audit trail）を残し、結果の追跡可能性を担保する。テキスト入出力のLLMと、表形式で高度に関係的な金融データとの構造的ミスマッチを埋める試みである。（\cite{kenshoapi2024}）
  \item \textbf{Benzinga}: リアルタイムニュース・SECファイリング・決算データ・アナリスト見解・株価イベントを対象としたAPI群を、LLM学習パイプラインとRAG本番システム向けに明示的に拡張。AIネイティブなライセンス枠組みと多言語データセットを追加し、金融AI専用のデータレイヤーとして自らを位置づけた（2026年5月）。（\cite{benzingaai2026}）
  \item \textbf{Morgan Stanley（MCP開放）}: 大手ウォール街銀行として初めて、S\&P 500の約40\%・\$1.2兆の資産をカバーする株式報酬管理プラットフォーム（ShareWorks・Equity Edge）を、AnthropicのMCP経由で外部の法人AIエージェントに開放。人手を介さずデータ取得やレポート生成を実行できるようにした（2026年6月）。（\cite{morganstanleymcp2026}）
  \item \textbf{FactSet「Portfolio Analytics MCP」}: 事前計算・ガバナンス済みのパフォーマンス・アトリビューション・リスク分析を、データパイプラインの再構築なしにクライアントのプライベート環境内でLLM／エージェント型ワークフローに開放するMCPサーバーを限定リリース。発表を受けFDS株は約11\%上昇し、MCP対応が市場からも評価される段階に入ったことを示した。（\cite{factset_portfolio_mcp2026}）
  \item \textbf{Moody's「GenAI-Ready Data」MCPサーバー}: 信用格付・企業ファンダメンタルズ・エンティティ連関・市場指標といった自社データと分析ツールを、MCP経由でLLM／エージェントに開放（エンドポイント \texttt{api.moodys.com/genai-ready-data}）。MCPのOAuth仕様を実装し、全てのやり取りを認証・ログ記録・エンタープライズデータポリシー準拠とする、規制業界向けの堅牢なMCP実装である。（\cite{x02moodysmcp2026}）
  \item \textbf{S\&P Global「ChatGPT向けアプリ＋MCPコネクタ」}: ライセンス済みのS\&Pデータに根拠づけて複雑な財務質問に回答する「検証済みS\&P Global MCPコネクタ」をChatGPT上のアプリとして提供（2026年2月）。並行してDatabricks MarketplaceにもMCPサーバーを公開し、顧客が自社データ／AIワークフローとS\&Pの金融・コモディティ情報を接続できるようにした。「検証済みコネクタ」「ライセンスデータ根拠付け」を前面に出した点が特徴。（\cite{x02spglobalmcp2026}）
\end{itemize}

\begin{table}[htbp]
\centering
\small
\caption{主要データプロバイダーのAI-ready接続戦略}
\label{tab:mcp-provider-strategy}
\begin{tabularx}{\textwidth}{@{}l l X X@{}}
\toprule
\textbf{提供者} & \textbf{接続面} & \textbf{AI-ready化の焦点} & \textbf{示唆} \\
\midrule
Bloomberg & ASKB／BQL & 数億件の文書、800以上のリサーチ提供元、BQLコード提示\cite{x02bloombergai2026} & 端末内の独自データを、根拠付き回答と再利用可能な分析コードへ変換する \\
LSEG/Microsoft & LSEG管理MCP & 33ペタバイト超のAI-readyコンテンツとタクソノミをCopilot Studioへ接続\cite{x02lsegmicrosoft2025} & データカタログ、権限、低コードエージェント構築が一体化する \\
FactSet & MCP／Portfolio Analytics MCP & 事前計算済み・承認済みのパフォーマンス、アトリビューション、リスク分析を自然言語で提供\cite{factset_portfolio_mcp2026} & 生成AIの自由回答ではなく、book-of-recordの分析結果へ誘導する \\
Daloopa & MCP構造化DB & 500問の金融検索で構造化DB追加により89〜91\%へ収束\cite{x02daloopabench2026,x02finretrieval2026} & モデル差よりツール・データ構造の差が精度を支配する \\
Moody's／S\&P & MCP／ChatGPTアプリ & 信用格付・企業データ・コモディティ情報を認証・ログ付きで接続\cite{x02moodysmcp2026,x02spglobalmcp2026} & 金融データ接続は監査証跡とライセンス統制を前提に標準化される \\
\bottomrule
\end{tabularx}
\end{table}

\begin{casestudy}{Daloopa: 構造化データがエージェントの検索精度を最大71ポイント押し上げる}
Daloopaのベンチマーク報告（2026年2月10日）は、500件の実務的な財務質問を3種類のエージェント構成——OpenAI Agents SDK、Anthropic Agent SDK、Google ADK——で評価した。公開Webソースの代わりに\textbf{監査可能な構造化金融データベース}を参照させると、精度は\textbf{89〜91\%}へ収束し、Claude Opus構成ではWeb検索のみの19.8〜20\%から構造化API利用時の90.8〜91\%へ改善した。arXiv版FinRetrievalは14構成の応答・ツール呼び出しトレースを公開し、地域差5.6ポイントの一因がモデル能力ではなく会計年度命名規則にあることも示した\cite{x02finretrieval2026}。同社の基盤は5,000社超の上場企業についてSECファイリング・プレスリリース・投資家プレゼン・トランスクリプトを出典追跡付きで構造化し、99\%超の精度・3\%未満のハルシネーション率を主張する。エージェント型金融AIにおいて「構造化データがWebスクレイピングに勝る」というテーゼを定量化した事例である。（\cite{x02daloopabench2026,x02finretrieval2026}）
\end{casestudy}

\begin{center}
\begin{tikzpicture}[
  prov/.style={fnode, text width=4.0cm, align=left, minimum height=1.5cm, font=\scriptsize},
  hub/.style={fnode accent, text width=2.4cm, minimum height=2.4cm, font=\small\bfseries},
  agentn/.style={fnode blue, text width=2.6cm, minimum height=2.4cm, font=\small},
  arr/.style={farr}
]
\node[prov] (bbg) at (0,3) {Bloomberg\\ASKB／BQL};
\node[prov] (fds) at (0,1) {FactSet\\MCP／Portfolio Analytics};
\node[prov] (lseg) at (0,-1) {LSEG/Microsoft\\33ペタバイトのカタログ};
\node[prov] (daloopa) at (0,-3) {Daloopa\\M365 Copilotコネクタ};
\node[hub] (mcp) at (5.6,0) {MCP\\標準\\プロトコル};
\node[agentn] (agents) at (9.6,0) {AI\\エージェント};

\draw[arr] (bbg.east) -- (mcp.west);
\draw[arr] (fds.east) -- (mcp.west);
\draw[arr] (lseg.east) -- (mcp.west);
\draw[arr] (daloopa.east) -- (mcp.west);
\draw[arr] (mcp.east) -- (agents.west);
\end{tikzpicture}
\end{center}

{\small\color{brandgray}\textbf{図: MCP標準化によるデータ接続のコモディティ化}\ 2025〜2026年に主要プロバイダーが相次いでMCPサーバーまたはMCP互換のエージェント接続を提供し、AIエージェントが標準プロトコル経由で構造化金融データにクエリできるようになりつつある（\cite{factset_mcp_nd,x02bloombergai2026,x02lsegmicrosoft2025,x02daloopabench2026}）。}

\subsection{大規模オープンコーパスと財務基盤モデル}

\begin{itemize}
  \item \textbf{BloombergGPT}: Bloombergが3630億トークンの金融コーパス＋3450億トークンの一般コーパスで学習した500億パラメータモデル。金融特化事前学習が金融NLPベンチマークで汎用モデルを大きく上回ることを実証し、ドメイン特化LLMの原型を確立した。（\cite{bloomberggpt2023}）
  \item \textbf{Stanford EDGAR Filings Dataset (SEFD)}: スタンフォード大学が公開した1,520億トークン規模のSECファイリング・オープンコーパス（MultiMarkdown形式に変換）。LLM事前学習と長文脈金融推論ベンチマーク向けの大規模AI-readyインフラ資源。（\cite{stanfordedgar2026}）
  \item \textbf{Kronos}: 45取引所・120億件超のローソク足レコードで事前学習した金融市場専用の基盤モデル。専用のマーケット・トークナイザにより、既存の時系列基盤モデルに対して価格予測・ボラティリティ予測タスクで\textbf{93\%の性能改善}（AAAI 2026）。（\cite{kronos2025}）
  \item \textbf{時系列基盤モデルの金融タスクへの適用}: 汎用予測向けに設計された事前学習済み時系列基盤モデル（TimesFM・Chronos・TTM）を金融固有タスクで評価する動きが加速している。TimesFMをバリュー・アット・リスク（VaR）推定に初めて適用した研究では複数のテール分位で最良〜準最良の成績を示す一方、高リスク用途でのハルシネーションリスクが指摘され（\cite{tsfmvar2024}）、TTM・Chronosを米国債利回り・通貨ボラティリティ・株式スプレッドで検証した研究では、データが限られる状況で\textbf{25〜50\%}高精度、かつ専門特化モデルに追いつくのに必要な履歴が\textbf{3〜10年}少なくて済むと報告された。ただし多くのタスクでは依然として専門特化モデルがゼロショットの基盤モデルに匹敵ないし凌駕しており、金融特有の非定常性に合わせたファインチューニングの必要性が示唆される（\cite{tsfmmultivariate2025}）。
  \item \textbf{Open FinLLM Leaderboard}: Linux Foundation・Hugging Faceと提携したコミュニティ型リーダーボード。36データセット・24タスク（情報抽出、QA、リスク管理、予測、意思決定）を横断し、「financial AI readiness」を単一タスクスコアではなく総合力として評価する枠組みを提示。（\cite{finllmleaderboard2025}）
  \item \textbf{その他の基盤モデル}: 180万件のロイター記事とFinancial PhraseBankで事前学習し財務センチメント分類のデファクトベースラインとなった\textbf{FinBERT}（\cite{finbert2020}）、テキストのみのFinGPT/BloombergGPTを超え市場データ・マクロ指標・決算音声・画像・動画を統合処理するマルチモーダル基盤モデルの動向を整理したコロンビア大学SecureFinAI Labの\textbf{MFFMサーベイ}（ACM ICAIF 2024、\cite{mffm2025}）など、周辺領域の整備も進む。
\end{itemize}

\begin{center}
\begin{tikzpicture}[
  font=\small,
  card/.style={fnode, align=left, text width=3.2cm, minimum height=3.8cm, font=\footnotesize},
  card2/.style={fnode blue, align=left, text width=3.2cm, minimum height=3.8cm, font=\footnotesize},
  node distance=3mm
]
\node[card] (bgpt) {\textbf{BloombergGPT}\\[1.5mm] 500億パラメータ\\金融コーパス3630億トークン\\+一般コーパス3450億トークン\\ドメイン特化LLMの原型};
\node[card2, right=of bgpt] (sefd) {\textbf{Stanford SEFD}\\[1.5mm] SECファイリング\\\textbf{1,520億トークン}の\\オープンコーパス\\長文脈金融推論向け};
\node[card, right=of sefd] (kronos) {\textbf{Kronos}\\[1.5mm] 45取引所\\\textbf{120億件超}の\\ローソク足レコード\\予測精度\textbf{93\%改善}};
\node[card2, right=of kronos] (finbert) {\textbf{FinBERT}\\[1.5mm] ロイター記事\\180万件\\Financial PhraseBank\\センチメント分類の\\デファクト標準};
\end{tikzpicture}
\end{center}

{\small\color{brandgray}\textbf{図: 大規模オープンコーパス・財務基盤モデルの規模比較}\ 学習トークン数・レコード数の規模でみると、テキスト系（BloombergGPT\cite{bloomberggpt2023}・Stanford SEFD\cite{stanfordedgar2026}・FinBERT\cite{finbert2020}）と市場データ系（Kronos\cite{kronos2025}）の両輪でAI-readyコーパスの整備が進んでいる。}

\subsection{多言語・ESG・規制文書のカバレッジギャップ}

英語・米国基準に最適化されたパイプラインは、グローバルな投資運用にとって構造的な死角となりうる。

\begin{itemize}
  \item \textbf{CLEF-2026 FinMMEval}: 金融LLMの初のマルチリンガル・マルチモーダル評価フレームワーク。既存の金融ベンチマークがほぼ英語中心であることを指摘。東京証券取引所の企業開示の\textbf{75\%超が日本語のみ}で公開されており、EU企業もCSRD義務化の開示を現地語で公開しているため、英語中心のAIパイプラインには大きなカバレッジの死角がある。（\cite{finmmeval2026}）
  \item \textbf{ESGReveal}: RAG拡張LLMによるESGレポート構造化抽出フレームワーク。HKEX（香港取引所）166社のESGレポートで、抽出精度\textbf{76.9\%}、開示分析精度\textbf{83.7\%}を達成。（\cite{esgreveal2023}）
  \item \textbf{EulerESG}: デュアルチャネル検索によりSASB整合の指標テーブルで平均精度最大\textbf{0.95}を達成する一方、EU Taxonomy整合の前向き（forward-looking）KPI抽出では精度が大幅に低いという構造的なギャップが指摘されている。米国中心の枠組みに基づく自動ESGスコアが、EUのESGスクリーニングファンドの資本配分に既に影響している点が論点。（\cite{eulereesg2025}）
  \item \textbf{ComplianceNLP}: ナレッジグラフ拡張RAGによる規制コンプライアンスAI。本番環境4か月間の実運用でアナリスト効率\textbf{3.1倍向上}、再現率\textbf{96\%}、適合率\textbf{90.7\%}を達成した、金融コンプライアンスAIとして最も強力な本番検証済み実績の一つ。（\cite{compliancenlp2026}）
\end{itemize}

\begin{center}
\begin{tikzpicture}[
  font=\small,
  gbox/.style={fnode, align=left, text width=4.5cm, minimum height=4.4cm, font=\footnotesize},
  gbox2/.style={fnode blue, align=left, text width=4.5cm, minimum height=4.4cm, font=\footnotesize},
  gbox3/.style={fnode accent, align=left, text width=4.5cm, minimum height=4.4cm, font=\footnotesize},
  node distance=4mm
]
\node[gbox] (lang) {
  \textbf{①言語カバレッジ}\\[1.5mm]
  東証上場企業開示の\\\textbf{75\%超が日本語のみ}\\
  EU企業もCSRD開示を\\現地語で公開\\[1mm]
  → 英語中心パイプラインの\\構造的な死角\ \cite{finmmeval2026}
};
\node[gbox2, right=of lang] (esg) {
  \textbf{②ESG抽出精度}\\[1.5mm]
  ESGReveal（HKEX 166社）:\\抽出\textbf{76.9\%}／\\開示分析\textbf{83.7\%}\ \cite{esgreveal2023}\\[1mm]
  EulerESG: SASB整合\textbf{最大0.95}も\\EU Taxonomy前向きKPIは急落\ \cite{eulereesg2025}
};
\node[gbox3, right=of esg] (comp) {
  \textbf{③規制コンプライアンス}\\[1.5mm]
  ComplianceNLP\\（本番環境4か月）:\\
  効率\textbf{3.1倍}\\再現率\textbf{96\%}\\適合率\textbf{90.7\%}\ \cite{compliancenlp2026}
};
\end{tikzpicture}
\end{center}

{\small\color{brandgray}\textbf{図: 多言語・ESG・規制文書のカバレッジギャップ地図}\ 言語・ESG・規制コンプライアンスの各領域で達成済みの精度・実績と、非英語圏や新規タクソノミへ拡張した際に残る構造的ギャップを並置した。}

\subsection{プライバシー保護協調学習とデータガバナンス}

機密性の高い顧客データを機関横断で共有できないという制約は、金融AIのAI-Ready化におけるもう一つの構造的障壁である。生データを移動させずに勾配更新だけを協調させる連合学習（federated learning）や、データソースの取捨選択・メタデータ整備といったガバナンス層の設計が、モデル選び以上に実運用の可否を左右し始めている。

\begin{itemize}
  \item \textbf{連合学習による金融セキュリティ（サーベイ）}: デジタル金融における連合学習の応用を規制上の露出度で分類し、不正検知・ポートフォリオ最適化・信用モデリングを対象に、差分プライバシーや秘密計算（MPC）といった防御策をコンプライアンス要件へ対応づけた包括的サーベイ（IEEE GCAIoT 2025採録）。（\cite{federatedlearning2025}）
  \item \textbf{連合学習による金融予測}: 市場方向の二値分類を対象としたプライバシー保護型の連合LSTMが、現実的な非IIDデータと差分プライバシー制約の下で、中央集約学習と同等の精度を達成しつつ単一機関モデルを大きく上回ることを実証した（2025年9月）。（\cite{fedforecasting2025}）
  \item \textbf{The Challenger: 新データソースへの切替判断}: 新しいデータソースが既存モデルの退役を正当化する条件を、経済的・統計的に統合した枠組みで定式化。閉形式の切替ルールと先読み型の逐次方策がオラクル性能に迫ることを実際の信用スコアリングデータで検証し、オルタナティブデータ採用の意思決定に定量的な基準を与えた（2025年12月）。（\cite{challenger2025}）
  \item \textbf{合成金融データのプライバシー・トレードオフの計測}: 6つの金融データセット（1千〜15万サンプル、少数派クラス比6.68〜30\%）で合成データ生成器を、データ品質・下流有用性（XGBoost分類）・プライバシー（メンバーシップ推論攻撃、最近傍距離）の3軸で評価。TabDiffは列形状で品質\textbf{0.96超}を示す一方、DP-TVAEは深刻なモード崩壊を起こし少数派クラスの表現を\textbf{0\%}まで消失させる場合があった。メンバーシップ推論の成功率は0.5前後（ランダム推測水準）で、プライバシー評価自体が信頼できないことが多いと指摘。クラスバランス調整は下流精度を最大\textbf{49.6\%}改善した（ACM Web Conference 2026、\cite{x02synthprivacy2026}）。差分プライバシーが希少・高シグナルの少数派パターンを不均衡に破壊するという本節の懸念（\cite{dpsynthetic2026}）を実証データで裏づける。
\end{itemize}

\begin{insight}{仮説（Takes）: ガバナンス層こそが精度の律速}
差分プライバシーを合成データに適用する際、プライバシー予算$\varepsilon \le 2$ではダウンストリームの精度劣化が顕著になるという仮説が提示されている。金融取引データは相手方・期間で著しく不均質であり、DPノイズが不正取引系列や信用悪化シグナルといった希少だが高シグナルの少数派パターンを不均衡に破壊するためである（\cite{dpsynthetic2026}）。同様に、エンタープライズのNL2SQLが本番水準（精度80\%超）に届かない主因はモデル規模ではなくメタデータ文書の不足だとする見方もある。最新LLMはSpider 1.0のような整備済みベンチマークで85〜91\%を達成する一方、簡潔で社内造語的な列名に意味注釈が伴わない実際の金融データベースでは10〜20\%まで崩落する。業務用語集・テーブル関係グラフ・オントロジタグでスキーマ文脈を補強したFI-NL2PY2SQLは、それだけでF1を\textbf{3.7〜6.3ポイント}改善しており、律速がメタデータ品質にあることを示唆する（\cite{nl2sqlmetadata2025}）。
\end{insight}

本節で確認した3つの潮流——フラット検索から階層型・グラフ型RAGへの進化、シンボリック検証層によるLLM単体の限界の補完、大手データプロバイダーによるMCP標準化——は、独立に見えて同じ方向を指している。データ接続そのものはコモディティ化しつつある一方、企業横断・言語横断で崩壊しない検索・検証アーキテクチャの設計は、今なお差別化の源泉であり続けている。

\subsection{周辺分野の文書AIから借りる：リーガル・医療・開発現場の到達点}

「非構造化文書をAIが確実に使える形へ変換する」という本節の課題は、金融に固有のものではない。契約書レビュー、電子カルテからの情報抽出、ソフトウェア開発現場のコードベース検索、企業横断の社内文書検索——いずれも同種の壁（フラット検索の限界・汎用埋め込みの機能不全・チャンキング設計の失敗）に突き当たり、同種の解法（構造を意識したチャンキング・ドメイン適応埋め込み・グラフ型検索）へ収束してきた。以下では、これらの周辺分野のうち\textbf{金融側が実際に転用・警戒すべき知見}のみを厳選して取り上げる。並べること自体を目的とせず、すべて金融への示唆に着地させる。

\textbf{①リーガル：契約書は金融文書と同じ「入れ子構造＋高い正確性要求」を持つ}

契約書レビューは、EDGARの開示資料と同じ文書群（SEC提出書類に添付された契約）を対象にしながら、金融のQAベンチマークにはない「条項単位の構造化抽出」というタスク設計を10年以上かけて磨いてきた分野である。The Atticus Projectの法律専門家が510件の商業契約に41種類の条項カテゴリを付与したCUAD（13,000件超の注釈、NeurIPS 2021）が標準ベンチマークの原型となり、ContractEval（19モデル評価、CUADテストセット4,128件・102契約）やACORD（契約条項の精緻な検索データセット、114クエリ・12.6万件超のクエリ―条項ペア）へと発展してきた。（\cite{y02cuad2021,y02contracteval2025,y02acord2025}）

\begin{casestudy}{LegalOn 2026年契約レビューベンチマーク: アーキテクチャがモデルに勝る}
LegalOn Technologiesが2026年6月3日に公表したベンチマークでは、Anthropic・Google・OpenAIの汎用フロンティアモデルを含む11モデルを、21件の精密なガイドラインに基づく契約条項3,282件のペア比較でLegalOnの専用パイプラインと比較した。独立した審査者は、テストしたすべての条項タイプでLegalOnのレビューを汎用モデルより高品質と評価し、例えばトリプルネットリースの譲渡条項では最良の汎用モデル（GPT-5）に対して\textbf{3.8倍}の頻度で高評価を得た。処理時間もLegalOnの\textbf{2.3秒}に対しClaude Opus 4.6は\textbf{40.4秒}を要した。同社の結論は「契約レビューの精度を左右するのは、フロンティアモデルの生の能力ではなくガイドライン連動型パイプラインという設計そのものである」というものである。（\cite{y02legalonbench2026}）
\end{casestudy}

\begin{insight}{金融への示唆}
これは本節が繰り返し示してきた「モデルよりアーキテクチャ」というテーゼ（AUDITFLOWのシンボリック検証層、Fin-RATEの検索失敗）と完全に同型の結論を、リーガルという全く独立した分野が導き出したものである。与信契約・ISDAマスター契約・目論見書の条項チェックにも、汎用チャットボットではなく、ガイドライン library と検索を密結合させた狭いパイプラインを優先投資すべきという教訓がそのまま転用できる。
\end{insight}

\begin{casestudy}{MLEB: 汎用埋め込みのランキング逆転はリーガルにも存在する}
Isaacusが公開したMassive Legal Embedding Benchmark（MLEB、10データセット・6法域・5文書種別・3タスク種別）では、汎用埋め込みベンチマーク（MTEB）で1位のGemini Embeddingがリーガル特化のMLEBでは\textbf{7位}に後退し、MTEBで23位のVoyage 3.5がMLEBでは\textbf{3位}に浮上する逆転現象が確認された。リーガル特化埋め込みのKanon 2 Embedder（NDCG@10\ \textbf{86.03}）は、汎用埋め込みのGemini Embedding（80.90）やOpenAI Text Embedding 3 Large（78.91）を\textbf{5〜12ポイント}上回る。同社のLegal RAG Benchでも、Kanon 2は主要フロンティアモデル・埋め込みに対し平均\textbf{17ポイント}の精度向上（Text Embedding 3 Large比では正答率+17.5・検索精度+34ポイント）を示した。ベンチマーク設計では章・節の階層構造を保った上で512トークン単位の意味的チャンキングを行っており、「条項こそがリーガル推論の最小単位である」という原則を明示している。（\cite{y02mleb2025,y02legalragbench2026}）
\end{casestudy}

\begin{insight}{金融への示唆}
「汎用埋め込みベンチマークの順位がドメイン性能を予測しない」というFinMTEBの知見（本節既出）と、MLEBのGemini逆転現象はほぼ同一の構造を示している。独立した2つの専門分野が同じ結論に達したことは、これが金融固有の癖ではなく専門テキスト全般に共通する構造的性質であることを裏づける。また「条項が最小単位」という発想は、金融文書における「表の行・脚注のペアこそが最小単位であり、固定トークン長のチャンキングでは崩れる」という、本節既出のBM25対密検索ベンチマークの知見（検索失敗の73\%が表構造の不整合に起因）と直接呼応する。
\end{insight}

\textbf{②医療：構造化スキーマがあっても、AIエージェントはそれを使いこなせるとは限らない}

電子カルテ・臨床ノートは、金融文書と同様に「自由記述テキストの中に構造化されるべき情報が埋め込まれている」典型例である。

\begin{casestudy}{CLEAR: エンティティ主導検索がベクトル検索と全文投入の両方に勝る}
スタンフォード大学のLopez・Swaminathanらが npj Digital Medicine（2025年1月）に発表したClinical Entity Augmented Retrieval（CLEAR）は、6種類のLLM・2万件の臨床ノート・18種類の抽出変数を対象に、(1)臨床エンティティ抽出に基づく検索、(2)通常の埋め込みRAG、(3)全文をそのままコンテキストに投入、の3手法を比較した。平均F1スコアはCLEARが\textbf{0.90}、埋め込みRAGが0.86、全文投入が0.79であり、入力トークン数は全文投入比\textbf{81\%減}・埋め込みRAG比\textbf{71\%減}、処理時間も1ノートあたり4.95秒（埋め込みRAGは17.41秒、全文投入は20.08秒）と、精度・コストの両面で優位性を示した。（\cite{y02clear2025}）
\end{casestudy}

\begin{casestudy}{FHIR-AgentBench: 標準化されたスキーマでもエージェントは苦戦する}
米国の21st Century Cures Act最終規則により、2021年以降すべての認定EHRシステムはHL7 FHIR標準に基づくAPIの提供を義務付けられている。にもかかわらずFHIR-AgentBench（2,931件の実臨床質問をFHIR構造化データに紐づけたベンチマーク）は、直接API呼び出しか専用ツールか、単一ターンか複数ターンか、自然言語かコード生成かといった条件を変えても、現行のLLMエージェントがFHIRのリソースベースのスキーマを確実に検索・推論できないことを示した。（\cite{y02fhiragentbench2025}）
\end{casestudy}

\begin{insight}{金融への示唆}
CLEARの「エンティティ主導検索」は、本節が示すフラット・階層型・グラフ型に続く\textbf{第4の検索様式}として、企業名・勘定科目コード・XBRLコンセプトタグ・カウンターパーティIDといった構造化エンティティを先に抽出してから検索する層を、金融RAGにも追加投資する余地を示唆する。一方でFHIR-AgentBenchは、MCP・XBRLのような標準化されたスキーマを整備すればエージェントの信頼性が自動的に解決するという楽観論への強い警鐘である。医療がFHIRという政府義務化された成熟した標準の上ですら苦戦している以上、金融もMCP標準化（本節既出）を「データ接続の解決」で終わらせず、「XBRL-AgentBench」に相当するエージェント評価基盤への投資を並行させる必要がある。
\end{insight}

\textbf{③ソフトウェア開発：関係性に依存するタスクはベクトル類似度では解けない}

コードベース検索は、単一ファイルの類似度検索では「この関数を呼び出しているすべての箇所」のような多段階の関係性クエリに答えられないという、金融の企業横断検索と同型の壁に直面してきた分野である。

\begin{casestudy}{CodexGraph・Sourcegraph Cody: グラフ構造とエンタープライズ規模の実運用}
CodexGraph（2024年8月、arXiv:2408.03910）は、コードリポジトリの構造（関数・クラス・呼び出し／インポート関係）から抽出したグラフデータベースをLLMエージェントに接続し、埋め込み類似度検索の代わりに構造化グラフクエリでコードを探索させる。著者らは、埋め込み類似度検索が多段階の関係性クエリで再現率が低下することを課題として挙げ、CrossCodeEval・SWE-bench・EvoCodeBenchで既存手法に対する優位性を確認した。実運用面では、Sourcegraph Codyが企業の全リポジトリ・全コードホストを継続的に再インデックスし、ローカルファイル・リポジトリ全体・リモート他リポジトリの文脈をロールベースアクセス制御付きで横断検索できるアーキテクチャを、\textbf{30万件超のリポジトリ}・90GB超のモノレポを持つ顧客で実運用していると公表している。（\cite{y02codexgraph2024,y02sourcegraphcody2024}）
\end{casestudy}

\begin{insight}{金融への示唆}
「意味的類似度だけでは構造的な関係性を捉えられない」というCodexGraphの指摘は、Fin-RATEの検索崩壊やFinReflectKG-MultiHopのKG検索優位性（本節既出）と全く独立に、別分野が同じ結論へ到達した点で説得力を増す。Sourcegraph Codyの実運用規模はさらに、金融機関の内部文書量をはるかに上回るスケールで階層的検索基盤が本番稼働している既成事実を示す。特に転用価値が高いのは、\textbf{アクセス制御を検索層の後付けフィルタではなく一級市民として設計する}という点であり、顧客の守秘義務・インサイダー情報の遮断（チャイニーズウォール）といった多エンティティ・多法域の権限制約を持つ金融RAGにそのまま応用できる設計思想である。
\end{insight}

\textbf{④企業横断検索：グラフ＋RAGは金融外で既に大規模商用化されている}

エンタープライズ検索は、金融が用いるGraphRAG系アーキテクチャの技術的起源であると同時に、それが金融の外側で既に大規模な商業的成功を収めていることを示す領域でもある。

\begin{casestudy}{Glean: ナレッジグラフ＋RAGが年間経常収益2億ドルに到達}
企業内検索・AIエージェント基盤のGleanは、組織のコンテンツ・人物・やり取りからナレッジグラフを構築し、その上にRAGを重ねて出典付きの回答を返すアーキテクチャで、2025年1月期に年間経常収益\textbf{1億ドル}を突破した後、2025年12月には\textbf{2億ドル}に到達（9か月で倍増）。同時期に立ち上げたGlean Agentsは、数か月で年間\textbf{1億件超}のエージェントアクションを処理する規模に達している。Wellington Managementが主導する1.5億ドルのシリーズFで評価額は72億ドルに達した。（\cite{y02glean2025}）
\end{casestudy}

Gleanが用いるグラフ＋RAGアーキテクチャの技術的起源は、Microsoft ResearchのEdgeらが2024年に発表した「From Local to Global」論文（GraphRAG）に遡る。エンティティ知識グラフを構築した上でコミュニティ（関連エンティティ群）ごとの要約を事前生成し、「このデータセット全体の主要テーマは何か」といった大域的な問いにmap-reduce方式で答える手法であり、約100万トークン規模のデータセットで従来の平坦なRAGに対し回答の網羅性・多様性を大きく改善することを示した。（\cite{y02graphrag2024}）

\begin{insight}{金融への示唆}
Gleanは、金融の開示資料のように整形されクリーンな文書群ではなく、社内チケット・チャット・Wikiのような雑多な文書を対象に、グラフ＋RAGが年間2億ドル規模の商用実績を築けることを実証した——本節が示すFinReflectKG-MultiHopやAgentic GraphRAG（いずれも既出）の優位性が、金融の整った文書だからこそ成立する特殊事情ではないことを裏づける。またGraphRAGの原論文が想定する「大域的な意味把握」の問いは、個別文書QA（FinanceBenchが典型）では測れない「今四半期、与信ファイル全体を横断してどのようなリスクテーマが浮上しているか」といったポートフォリオ横断の問いにこそ活きる。金融機関が社内RAGの内製・外部調達を判断する際、GleanのようなドメインニュートラルなグラフRAG基盤は、金融特化ベンダーと比較検討すべき現実的な選択肢になっている。
\end{insight}

\begin{table}[htbp]
\centering
\small
\caption{周辺分野の文書AI到達点と金融への転用可能性}
\label{tab:crossdomain-docai}
\begin{tabularx}{\textwidth}{@{}
  >{\raggedright\arraybackslash}p{1.5cm}
  >{\raggedright\arraybackslash\hsize=1.30\hsize}X
  >{\raggedright\arraybackslash\hsize=0.70\hsize}X
  >{\raggedright\arraybackslash\hsize=1.00\hsize}X @{}}
\toprule
\textbf{分野} & \textbf{検索・埋め込みの到達点} & \textbf{最小単位（チャンキング）} & \textbf{金融への転用点} \\
\midrule
リーガル & ドメイン適応埋め込みが汎用比\,5〜12pt優位（MLEB）\cite{y02mleb2025} & 条項（clause） & 与信契約・ISDA条項抽出への専用パイプライン投資 \\
医療 & エンティティ主導検索がRAG・全文比で高精度・低コスト（CLEAR）\cite{y02clear2025} & 臨床エンティティ言及 & 企業名・勘定科目・XBRLタグ主導の検索層追加 \\
開発現場 & グラフDBが埋め込み類似度の多段階関係性クエリを代替（CodexGraph）\cite{y02codexgraph2024} & 関数・クラス・呼び出し関係 & アクセス制御を検索層の一級市民として設計 \\
企業検索 & ナレッジグラフ＋RAGが年間2億ドル規模で商用実証（Glean）\cite{y02glean2025} & エンティティ・コミュニティ要約 & 大域的・ポートフォリオ横断の問いへの応用 \\
\bottomrule
\end{tabularx}
\end{table}

{\small\color{brandgray}\textbf{表\ref{tab:crossdomain-docai}について:}\ いずれの分野も、金融がFin-RATE・FinMTEB・FinReflectKG-MultiHopで独自に発見してきた知見（検索アーキテクチャの重要性、ドメイン適応埋め込み、チャンキング単位の設計）と構造的に同型の結論に達している。}

\begin{insight}{仮説（Takes）: 金融の文書AIは「特殊」ではなく「後発」である}
リーガル・医療・開発現場・企業検索という4つの独立した分野が、金融と同じ3段階の解法——構造を意識したチャンキング、ドメイン適応埋め込み、グラフ型検索——へ収束している事実は、金融の文書AI課題が独自に困難なのではなく、他分野が既にベンチマーク済み・商用実証済みの「標準レシピ」の採用において1〜2年遅れている可能性を示唆する。Gleanの年間2億ドルのARRやSourcegraphの30万件超リポジトリでの実運用は、このレシピが金融の文書量をはるかに超える規模で既に量産段階にあることを示す。

ただし反証材料もある。4分野を通じて調査した限り、いずれも金融のAUDITFLOW・FinVerBenchが示すような「数値・会計演算の忠実性」を検証する層を持たない。契約条項の有無、臨床エンティティの存在、コード関数の呼び出し関係はいずれも真偽・存在の判定問題であるのに対し、金融は「貸借が一致するか」「注記の単位換算が正しいか」という決定的な数値検証を必要とする。したがって、\textbf{検索層のレシピは借用できる一方、数値検証層は金融が独自に構築し続けなければならない}可能性が高い。
\end{insight}

\section{AI-Nativeヘッジファンド・投資会社}
\label{sec:ai-native-funds}

本節では、LLM・AIエージェント・機械学習を投資リサーチ、ポートフォリオ管理、リスク管理、トレーディングの中核に据える機関の実例を整理する。伝統的クオンツファンドの「AI強化」から、AIエージェントが意思決定の相当部分を担う「AI-Native」設計まで、実装の深度には大きな幅がある。AI活用を標榜するファンド・ベンダーは無数に存在するため、本節は網羅的なリストではなく、実運用資本の規模・具体的な数値開示・意思決定構造の観点で特に示唆に富む旗艦事例とインフラ企業に絞って取り上げる。

\begin{concept}{AI-Nativeヘッジファンドとクオンツ運用}
ヘッジファンドとは、私募形式で資金を集め、相場の上下によらない絶対収益を狙う運用形態を指す。その中でも\textbf{クオンツ運用}は、勘や経験ではなく数理モデル・統計・アルゴリズムを使って売買判断を体系化する手法である。運用成績を測る際は、単なる儲けの大きさではなく「市場平均（ベータ）を上回る超過収益」を意味する\textbf{アルファ（$\alpha$）}や、「リスク1単位あたりどれだけ超過リターンを稼げたか」を示す\textbf{Sharpe比}（超過リターン÷変動の大きさ）が重視され、値が大きいほど効率よく稼げていることを意味する。こうした手法は実運用に投入する前に、過去データで戦略の有効性を検証する\textbf{バックテスト}を経るのが通例である。本節に登場する各社は、いわば「運転支援（人が主役でAIが助手席）」から「自動運転（AIが操縦の大半を担う）」まで幅のあるAI活用度合いに位置づけられる。
\begin{center}
\begin{tikzpicture}[
  stg/.style={fnode, text width=26mm, align=center, font=\scriptsize, minimum height=17mm},
  lbl/.style={font=\tiny, color=figgray, align=center, text width=26mm}
]
\node[stg] (s1) {伝統的\\クオンツ};
\node[stg, fnode blue, right=9mm of s1] (s2) {AI強化／\\コパイロット};
\node[stg, fnode accent, right=9mm of s2] (s3) {アルファ\\生成型};
\node[stg, fnode dark, right=9mm of s3] (s4) {完全\\自律型};
\node[lbl, below=2mm of s1] {数理モデルを\\人が設計・運用};
\node[lbl, below=2mm of s2] {人が主役、\\AIは補助・提案};
\node[lbl, below=2mm of s3] {AIが売買\\シグナルを生成};
\node[lbl, below=2mm of s4] {AIが調査〜\\執行まで担う};
\draw[farr] (s1) -- (s2);
\draw[farr] (s2) -- (s3);
\draw[farr accent] (s3) -- (s4);
\end{tikzpicture}
\end{center}
{\small\color{brandgray}\textbf{図: AI-Nativeファンドの深度スペクトル}\ 右へ行くほどAIが担う意思決定の範囲が広がり、人間の関与は補助から監督へと後退する。}
\end{concept}


\subsection{実運用資本を動かす旗艦事例}

\begin{casestudy}{Bridgewater Associates AIA Labs: 実資本を運用する「人工投資家」}
Co-CIOのGreg Jensen率いる専属AIラボが、LLM・表形式ML・独自マクロ経済データベースを用いた「人工投資家」を構築。実運用資本を大規模に運用し、\textbf{初年度フルイヤーで11.9\%のリターン}を達成した。AIネイティブ投資リサーチシステムが実運用に到達した旗艦事例として位置づけられる。（出典: Bridgewater公式サイト \cite{bridgewater_aialabs_nd}）
\end{casestudy}

\begin{casestudy}{Renaissance Technologies: 設立来年率39.9\%という金字塔}
1988年設立のMedallion Fundは、設立来平均\textbf{年率39.9\%の純リターン}、1日15万〜30万件の自動売買、勝率約50.75\%という実績を誇る。もっとも、2025年10月には創業者James Simons死去後初となる公開ファンド（RIEF、RIDA）が大幅な損失（RIEF -14.39\%、RIDA -15.6\%）を記録し、11月には12.6\%超反発するなど、AI／統計的手法をもってしてもボラティリティから免れないことも同時に示している。（出典: Institutional Investor \cite{institutionalinvestor2025renaissance}）
\end{casestudy}

\begin{casestudy}{QRT (Qube Research \& Technologies): 急成長する新興クオンツ大手}
2016年Credit Suisseのクオンツ部門を起源とし2018年にMBOで独立。2026年1月時点でAUM \textbf{380億ドル}（前年230億ドルから増加）、従業員約1,400名。主力ファンドは2025年に約\textbf{+30\%}（Torus +20.4\%、Prism +7.4\%）と、HFRI Global Hedge Fund Indexの+12.6\%を大きく上回った。2026年にはQRT Labsを設立し、ケンブリッジ・インペリアル・オックスフォード大学で70名超の若手研究者に資金提供している。（出典: Bloomberg \cite{bloomberg_qube2026}）
\end{casestudy}

\begin{casestudy}{Magnetar Capital: 数百のAIボットがアナリスト業務を担う「AI-first」ファンド}
AUM \textbf{180億ドル}の運用会社Magnetar Capitalは、新ファンド向けにアイデア発掘・企業分析・投資推奨・トレンド予測を担う\textbf{数百のAIボット}をNvidia製インフラ上に配備し、従来はアナリストチームが担っていた業務を代替する構想を明らかにした。ただし最終的な発注権限は人間のポートフォリオマネージャーに留保される。数百のエージェントをリサーチの前工程に投入しつつ意思決定の最終責任は人間に残すという設計は、大手ファンドにおける現実的な「AI-first」の一形態を示している。（出典: Hedgeweek \cite{magnetar2026}）
\end{casestudy}

\begin{center}
\begin{tikzpicture}[
  card/.style={fnode, text width=4.5cm, align=left, font=\footnotesize}
]
\node[card] (bw) {\textbf{\color{fignavy}Bridgewater AIA Labs}\\[0.5mm]
Co-CIO Greg Jensen率いる専属AIラボ\\
LLM＋表形式ML＋独自マクロDBで「人工投資家」を構築\\
実運用資本を大規模に運用};
\node[card, right=5mm of bw] (rt) {\textbf{\color{fignavy}Renaissance（Medallion）}\\[0.5mm]
1988年設立\\
1日15万〜30万件を自動売買\\
勝率約50.75\%};
\node[card, right=5mm of rt] (qrt) {\textbf{\color{fignavy}QRT}\\[0.5mm]
2018年MBOで独立\\
AUM 380億ドル（前年230億ドルから増加）\\
従業員約1,400名、QRT Labsで70名超の若手研究者に資金提供};
\end{tikzpicture}
\end{center}

\smallskip
{\small\color{figgray}\textbf{図: 実運用資本を動かす3社の体制比較。}}3社とも実運用資本でAIを投資判断の中核に据えるが、体制・規模の性質は異なる——Bridgewaterは専属AIラボによる「人工投資家」、Renaissanceは長年蓄積された自動売買実績、QRTは急成長するAUMと研究者ネットワークが特徴である。3社を含むリターン水準の比較は本節後半の図を参照。

\subsection{大手クオンツファンドのAI活用スタンス}

老舗クオンツ勢は、完全自律型を掲げるのではなく、人材育成とインフラ投資を両輪とした漸進的なAI導入を選好する傾向が強い。以下ではその代表例のうち、経営陣の発言の転換・戦略の二正面展開・全社的な導入設計といった物語性が際立つCitadel、D.E.\ Shaw、AQR、Balyasnyの4社を独立事例として取り上げ、残りは要点を絞った比較としてまとめる。とりわけAQRとBalyasnyの事例は、2025年が「AIを取り込んだ系統的戦略が実運用で顕著に上回った変曲点」であったことを、懐疑派の翻意と全社展開という異なる角度から裏づけている\cite{x03forbes_billionaires2026}。

\begin{casestudy}{Citadel: Ken Griffinの「AI＝ゴミ」発言からの手のひら返しと社内リサーチAI}
創業者Ken Griffinは2026年1月のダボス会議でAIを「ゴミ」と酷評したが、その後方針を転換し「profoundly more powerful（はるかに強力）」と評価を改めた。2025年12月には、規制フィリング・決算トランスクリプト・ブローカーリサーチ・自社戦略で訓練した社内AIアシスタントを株式リサーチ向けに導入し、従来ファイナンスPhDが数週間を要した作業を数時間に短縮した\cite{griffin2026}。Citadel Securitiesはさらに「The Economics of Intelligence」で、S\&P500決算コールトランスクリプト9,646件の分析を根拠に、AIによる労働代替ではなく人間／AIの補完性を主張している\cite{citadelsecurities2026economics}。創業者本人の公言の急転換という点で、AI懐疑派の翻意を象徴する事例である。
\end{casestudy}

\begin{casestudy}{D.E. Shaw: 独立稼働するMLグループと「人間主導」新ファンドへの二正面投資}
2018年設立のMachine Learning and Research Group（Pedro Domingos率いる、2019年IJCAI John McCarthy Award受賞）が独立稼働する一方、2025年9月には対照的に約50億ドルを調達する裁量型・人間主導の新ファンド「Cogence」を立ち上げた。既存のOculusマクロファンドはAUM 850億ドル超で約+28.2\%のリターンを記録している\cite{bloomberg_deshaw2025}。同一運用会社がAI主導と人間裁量主導を並行して強化するという、二極化する業界の縮図とも言える賭け方である。
\end{casestudy}

\begin{casestudy}{AQR: 「機械にさらに明け渡した」——クオンツ正統派の翻意}
伝統的クオンツの旗手Cliff Asnessは、かつて機械学習がAQRで大きな役割を担うことに懐疑的だったが、方針を転換。AQRは2つのML専用戦略で外部資本を募り、機械学習が\textbf{主力マルチストラテジー・ファンドのシグナルの約5分の1}を担うに至った\cite{x03aqr_bloomberg2025}。運用資産は2025年に約730億ドル増加し\textbf{約1,870億ドル}へ膨張、主力Apexファンドとロング・ショートのDelphiファンドはいずれも2025年に堅調なリターンを記録した\cite{x03forbes_billionaires2026}。Asnessは「機械にさらに（意思決定を）明け渡した」と述べつつ、生成AIを「革命的というより進化的な、高度な統計ツール」と位置づけ、乱高下する相場での解釈可能性の低下という副作用にも言及している\cite{x03aqr_hedgeco2025}。長年の懐疑派が実運用でMLの寄与を認めた点で、業界全体のAI受容が新たな段階に入ったことを象徴する事例である。
\end{casestudy}

\begin{casestudy}{Balyasny: 中央集権AIチームが全投資チームへエージェントを配る「フェデレーテッドAI」}
マルチマネージャー型のBalyasny Asset Managementは、2022年末に約20名の研究者・エンジニア・ドメイン専門家からなる中央集権的な\textbf{Applied AIチーム}を設置。エージェントフレームワーク・ツールチェーン・コンプライアンスガードレールを中核で開発し、資産クラス別にスコープを絞ったエージェントを\textbf{全投資チームの約95\%}へ配布する「フェデレーテッド」展開を採る\cite{x03balyasny_openai2026,x03balyasny_wikipedia2026}。推論エンジンにはOpenAIのGPT-5.4を内製モデルと併用しタスク単位で使い分け、本番投入前に予測精度・数値推論・シナリオ分析・ノイズ耐性など\textbf{12以上の次元}でモデルを評価する\cite{x03balyasny_mlq2026}。SEC提出書類・ブローカーリサーチ・決算／エキスパートコールなど数万件の文書を統合する「Deep Research」エージェントは従来数日を要した調査を数時間に短縮し、中央銀行スピーチ分析エージェントはマクロシナリオ分析を約2日から約30分へ圧縮した\cite{x03balyasny_openai2026}。中核チームがガードレールと評価基盤を一元的に整備しつつポッドの自律性を保つこの設計は、プラットフォーム型ファンドにおける現実的なAI展開の参照アーキテクチャとなりつつある。
\end{casestudy}

\begin{itemize}
  \item \textbf{Man Group / Man AHL}: エージェント型AIパイプラインAlphaTrendが、量的トレンドフォローシグナルの生成・実装・バックテストをエンドツーエンドで担う。統制実験ではClaudeがシグナル提案間で高い相関を示す一方GPT-5はより多様性が高く、「良いアイデア」はベンチマークSharpe比0.95に対し約1.05にクラスタリングした\cite{manndalphatrend}。CIOのRussell Korgaonkarは「包括的なAIエージェント群」の構築を明言する一方\cite{man2025bigimperative}、Man AHL自身は生成AIがクオンツ／PMを代替せず単独で市場を上回る新システムも生んでいないと率直に認め、効果を主に生産性向上（アルファ創出への集中時間の増加）に限定して評価している\cite{man2024bigpicture}。Oxford-Man Instituteとの連携により2016年以降19名のML研究者を追加した\cite{man2026riseofml}。
  \item \textbf{Two Sigma}: CTO Jeff WeckerとChief AI Innovation Officer Matt Greenwoodは、AIを研究ワークフロー全体の「オペレーティングシステム」として位置づけつつ、AI過大評価とバックテスト／レジーム変化リスクに警鐘を鳴らす\cite{twosigma2026aifirst}。社内では「Platform Thinking」としてadvisory・oracle・operational・agenticの4つの役割別AIフレームワークを整理し\cite{twosigmandplatformthinking}、数千のAIエージェントが並列で動く「combinatorial research」構想を提示している\cite{twosigma2026aioutlook2}。LLMプログラミングフレームワークをOSI風の階層モデルに整理する学術的貢献も行っている\cite{twosigmandllmabstractions}。2026年には新共同CEO（Carter Lyons／Scott Hoffman）体制下での戦略見直しの一環として約200名の人員削減に踏み切り、AIを研究の「オペレーティングシステム」に据える方針が組織再編としても具体化しつつある（2025年末AUM約700億ドル）\cite{x03twosigma_hedgeweek2026,x03twosigma_bloomberg2026}。
  \item \textbf{Millennium Management}: 2026年3月時点でAUM約842億ドル。2025年11月には創業者Israel Englander後の体制を見据えた少数株式売却を進めるなど、マルチマネージャー型「プラットフォーム」モデルの制度的定着が進む\cite{x03millennium_bloomberg2025}。自社の生成AI施策の公表は限定的だが、Llamaベースのアルファ探索基盤BRAINを擁するWorldQuant\cite{forbes2025worldquant}が大規模な系統的スリーブを主にMillennium向けに運用しており、AIは自社ブランドよりもサブマネージャーや中央インフラを通じて浸透している。
  \item \textbf{Point72 / Cubist Systematic Strategies}: AUM約400億ドルのうち約17\%をCubistが運用（2026年6月時点で546名の従業員）。2025年9月にWorldQuant元CIOのGeoffrey Lauprete氏が新責任者に就任した\cite{bloomberg_cubistchief2025}。人材育成にも注力しており、Cubist Quant Academy（2020年開始、比較で約25\%のチーム成長）\cite{point72_quantacademy2023,point72_5years2025}、GenAIハッカソン（ワルシャワ拠点、65名・17チーム、2024年3月）\cite{point72_genaipoland2024}、Global Technology Hackathon（5拠点約100名、2025年3月）\cite{point72_globalhackathon2025}など多数のAI/ML人材パイプライン施策を展開している。
  \item \textbf{Voleon Group}: ディープラーニングを唯一の運用手法とする希少な純ML運用会社。2026年1月末時点AUM 234億ドル（収益1億9,180万ドル）。Compositionファンドは2022年+9\%、2024年+13.5\%、2025年4月までのYTDで+12\%\cite{voleon_wikipedia2026}。
  \item \textbf{WorldQuant}: CEO Igor Tulchinskyは、クラウドソース型リサーチ基盤BRAINにLlamaベースのLLMを早期段階で導入し、アルファ発見と非構造化データ構造化を加速させていると証言している\cite{forbes2025worldquant}。
\end{itemize}

\begin{center}
\resizebox{\textwidth}{!}{%
\begin{tikzpicture}
\begin{axis}[
    xbar,
    bar width=6mm,
    width=13.5cm,
    height=6.2cm,
    xlabel={公表リターン (\%、各社の開示ベース)},
    xlabel style={font=\small, color=figgray},
    tick label style={font=\small},
    symbolic y coords={
        Voleon Composition・2024年通期,
        D.E. Shaw Oculus・直近開示,
        QRT主力ファンド計・2025年通期,
        Renaissance Medallion・設立来年率平均,
        Bridgewater AIA Labs・初年度
    },
    ytick=data,
    yticklabel style={align=right, text width=6.0cm},
    nodes near coords={\pgfmathprintnumber[fixed,precision=1]{\pgfplotspointmeta}\%},
    nodes near coords style={font=\small, text=fignavy},
    xmin=0, xmax=46,
    enlarge y limits=0.18,
    ymajorgrids=false,
    xmajorgrids=true,
    major grid style={draw=figlight},
    axis line style={draw=figgray},
    every axis plot/.append style={fill=figblue, draw=fignavy, line width=0.5pt},
]
\addplot coordinates {
    (13.5,Voleon Composition・2024年通期)
    (28.2,D.E. Shaw Oculus・直近開示)
    (30,QRT主力ファンド計・2025年通期)
    (39.9,Renaissance Medallion・設立来年率平均)
    (11.9,Bridgewater AIA Labs・初年度)
};
\end{axis}
\end{tikzpicture}%
}

\smallskip
{\small\color{figgray}\textbf{図: AI-Nativeファンド5社の公表リターン比較（測定条件は不揃い）。}測定期間・戦略・ファンド種別が異なるため数値を直接比較することはできない点に注意されたい——Bridgewaterは新設AI投資家の初年度単年、Renaissanceは1988年設立来の年率平均、QRTは2025年通期（参考: HFRI Global Hedge Fund Indexは+12.6\%）、D.E.\ Shawは直近開示値、Voleonは2024年通期の数値であり、レバレッジ・運用規模・戦略類型も異なる。ここでの棒の長さの差は、AI活用の巧拙の優劣を示すものではない。}
\end{center}

\subsection{運用インフラを支えるAIエージェント・プラットフォーム企業}

大手金融機関・データベンダーの側でも、エージェント型AIを本番の運用インフラに組み込む動きが進んでいる。以下は開示された定量効果の大きさで特に際立つ事例を優先して取り上げる。

\begin{casestudy}{Moody's: 与信メモ作成を「40時間」から「約2分」へ圧縮するマルチエージェントAI}
エンティティ解決・10-K抽出・ピア分析・リスク評価それぞれに特化したエージェントを協調させるマルチエージェントシステムを、5.9億エンティティ規模のデータ基盤上に構築。ローン引受における与信メモ作成を、アナリストが\textbf{40時間}を要していた業務からおよそ\textbf{2分}へと圧縮した（2026年4月、AWS Marketplaceでローンチ）。単一の万能モデルではなく機能分業したエージェント群によって、監査可能性を保ちながら数百倍規模の時間短縮を実現した点が、他の効率化事例と比べても際立つ\cite{moodys2026creditmemo}。
\end{casestudy}

\begin{itemize}
  \item \textbf{BNY Mellon（Eliza 2.0）}: マルチエージェントプラットフォームで130以上の自律型「デジタル従業員」をカストディ・コンプライアンス・トレード決済業務に配備し、累計\textbf{2万件}のAIアシスタント展開規模でコアカストディ取引の単位コストを約5\%削減した\cite{bnymellon2026}。
  \item \textbf{Arcesium Intelligence}: 100超の投資領域特化型エージェントとセルフサービス「Agent Studio」を運用基盤Opterra・データ管理スイートAquataに組み込み、P\&L分解・NAVレビュー・ポートフォリオモニタリングを自動化する\cite{arcesium2026}。
  \item \textbf{Broadridge}: エージェント型AIを40超のクライアント・日次取引高15兆ドル規模の本番環境に展開し、取引失敗解消や口座ワークフロー自動化により初日\textbf{30\%のコスト削減}を掲げる\cite{broadridge2026}。
  \item \textbf{AlphaSense}: カスタムエージェントと25万件のトランスクリプトを用いたAI主導のエキスパートコールでARR6億ドル超・評価額75億ドルに到達し、AIネイティブな投資リサーチ基盤カテゴリーの中心的存在となっている\cite{alphasense2026}。2026年6月にはSEC提出書類・決算コール・ブローカーリサーチ・エキスパートトランスクリプトを含む\textbf{5億件超}の金融文書を24時間365日自律監視し、複数ステップのリサーチワークフローを実行する常時稼働エージェント「SuperAnalyst」を投入。同時にJPMorgan Asset ManagementとD.E.\ Shaw Venturesが出資する\textbf{3.5億ドル}の資金調達を評価額\textbf{75億ドル}で実施した\cite{alphasensesuper2026}。
  \item \textbf{Anthropic（金融サービス向けエージェントテンプレート）}: 2026年5月、ピッチブック作成・KYCスクリーニング・決算レビュー・月次決算締めなどをカバーする\textbf{10種類}の即用可能なClaudeエージェントテンプレートを提供開始。導入企業としてCitadel、Walleye Capital（400名の全従業員に展開）、BNY、Carlyle、FIS、みずほが名を連ねており、AIエージェントが特定ファンドの内製ツールにとどまらず金融機関横断で標準化されつつあることを示す\cite{anthropicfinanceagents2026}。
  \item \textbf{Numerai}: 匿名データサイエンティスト群のモデルをStake-Weighted Meta Modelで統合する分散型ファンド運用モデル。2025年11月、評価額5億ドル（2023年比5倍）でシリーズC 3,000万ドルを調達し、J.P.\ Morganから5億ドルのトレーディング枠を確保した\cite{numerai_seriesc2025}。微分可能な凸ポートフォリオ最適化で勾配寄与度を測るTrue Contribution\cite{numerai_alien_nd}や、ファクター中立化・流動性/ボラティリティ加重によるAlpha scoring\cite{numerai_alpha_nd}など、独自のインセンティブ設計を継続的に高度化している。
\end{itemize}

\begin{center}
\begin{tikzpicture}[
  card/.style={fnode, text width=4.5cm, align=left, font=\footnotesize}
]
\node[card] (moodys) {\textbf{\color{fignavy}Moody's}\\[0.5mm] 与信メモ作成: 40時間\,$\to$\,約2分（マルチエージェント、5.9億エンティティ基盤）};
\node[card, right=5mm of moodys] (bny) {\textbf{\color{fignavy}BNY Mellon（Eliza 2.0）}\\[0.5mm] デジタル従業員130超、累計2万件展開、コア決済単位コスト▲5\%};
\node[card, right=5mm of bny] (arc) {\textbf{\color{fignavy}Arcesium Intelligence}\\[0.5mm] 100超の投資領域特化エージェント、Agent Studio};
\node[card, below=5mm of moodys] (bro) {\textbf{\color{fignavy}Broadridge}\\[0.5mm] 40超クライアント・日次取引高15兆ドル、初日コスト▲30\%};
\node[card, right=5mm of bro] (alp) {\textbf{\color{fignavy}AlphaSense}\\[0.5mm] トランスクリプト25万件、ARR6億ドル超、評価額75億ドル};
\node[card, right=5mm of alp] (num) {\textbf{\color{fignavy}Numerai}\\[0.5mm] 評価額5億ドル（2023年比5倍）、J.P.\ Morganから5億ドルの取引枠};
\end{tikzpicture}
\end{center}

\smallskip
{\small\color{figgray}\textbf{図: 運用インフラを支えるAIエージェント・プラットフォーム企業6社。}}各社が開示する定量効果は業務内容（与信メモ作成、カストディ決済、投資リサーチ、取引執行、資金調達等）により大きく異なるが、いずれもエージェント型AIを本番運用インフラに組み込んでいる点で共通する。

Hebbia（Matrix）の処理ページ数急増やMorgan Stanleyのアドバイザー向けGPT-4展開など、他にも大きな数値を掲げる事例が報じられているが、本節では一次資料で裏付けが取れた範囲に絞って掲載した。

\subsection{大手銀行・資産運用会社の全社的AI展開}

ヘッジファンドと並行して、大手銀行・資産運用会社は\emph{全社基盤}としてのLLM／エージェント展開を急速に進めている。特徴は、単一ツールの導入ではなく、数万〜数十万人規模の従業員に生成AIを行き渡らせる「democratization（民主化）」と、そこから自律エージェントへ移行する段階設計にある。

\begin{casestudy}{JPMorgan LLM Suite: 23万人超に展開する7モデルの社内基盤とエージェント段階への移行}
JPMorgan Chaseの社内基盤「LLM Suite」は、セキュアなインフラ上で\textbf{7つのファインチューン済みLLM}を提供し、2024年夏のリリースから8か月で約20万人を、現在は\textbf{全世界23万人超}の従業員をオンボード（約半数が日次利用）した\cite{x03jpmorgan_ciodive2024,x03jpmorgan_cnbc2025}。プライベートバンクの「Connect Coach」による顧客面談準備の短縮など\textbf{450件超のユースケース}を生み、約\textbf{20億ドル}のAI由来アップサイドを見込む。同行はこれを「完全にAIで接続されたメガバンク」構想の初段と位置づけ、多段タスクを担うエージェント型AIへ移行しつつある\cite{x03jpmorgan_cnbc2025}。
\end{casestudy}

\begin{casestudy}{Goldman Sachs: 全社AIプラットフォームと自律コーダーDevinによる「ハイブリッド労働力」}
Goldman Sachsは2025年6月、GPT-4・Gemini・Llama・Claude・内製モデルを自社ネットワーク内（暗号化・プロンプトフィルタ・RBAC・監査ログ付き）でホストするGS AIプラットフォームを\textbf{46,500名超}の従業員に開放し、2026年に100\%利用を目標に掲げた（現在は月間約100万件のプロンプトが書かれる）。2025年7月にはCognition社の自律ソフトウェアエンジニア\textbf{Devin}のパイロットを発表し、約1万2,000名の人間開発者と並走する形で数百から\textbf{数千のAIエージェント}へ拡大する計画を示した。CTO Marco Argentiは従来ツール比\textbf{3〜4倍}の生産性向上を見込み、エージェントを監督しつつ最終責任を負う「AIネイティブ」人材からなる「ハイブリッド労働力」という組織像を提示している\cite{x03goldman_cnbc2025,x03goldman_ibm2025}。
\end{casestudy}

\begin{casestudy}{BlackRock Aladdin Copilot: 運用プラットフォームに組み込まれたマルチエージェント層}
資産運用最大手BlackRockは、投資運用プラットフォームAladdin全体を横断する生成AI層「Aladdin Copilot」を全Aladdinクライアントへ一般提供した。監督型のマルチエージェント構成（LangChain＋LangGraph）とGPT-4のfunction callingにより、自然言語クエリを\textbf{数百のドメイン特化Aladdin API}への呼び出しへ変換し、回答提示・レポート生成・リサーチ要約・能動的アラートを行う。ハルシネーション抑制やコンテンツフィルタといったガードレールでAladdinの範囲内に限定し、投資助言は行わない設計である\cite{x03blackrock_aladdincopilot,x03blackrock_microsoft2024}。運用会社が日々用いる基盤そのものにエージェント型AIを織り込む動きを示す。
\end{casestudy}

\begin{casestudy}{Morgan Stanley: 評価駆動のAI@Morgan Stanleyから外部エージェント接続へ}
Morgan Stanleyは、AI @ Morgan Stanley Assistant／Debriefをウェルスマネジメントの業務フローに組み込み、Assistantは\textbf{アドバイザリーチームの98\%超}が利用する段階に到達した。OpenAIの公式事例では、回答対象が当初の7,000問から10万文書規模のコーパスへ拡大し、文書アクセス範囲も20\%から80\%へ広がったとされる。展開の中核はモデル単体ではなく、要約・翻訳・検索品質を専門家が採点するeval、日次の回帰テスト、ゼロデータ保持、Debrief出力をアドバイザーが確定前に確認する人間監督である\cite{x03morganstanley_openai2026}。さらに機関投資家向けには、OpenAIのGPT-4でMorgan Stanley Researchの年7万本超のレポートを横断要約するAskResearchGPTを投入し\cite{x04askresearchgpt2024}、株式報酬管理プラットフォームではMCP経由で外部AIエージェントがShareWorks／Equity Edgeに接続できる設計を開いた\cite{morganstanleymcp2026}。同行の事例は、銀行・証券会社におけるAI-Native化が「完全自律売買」ではなく、評価済みの社内知識検索、会議後処理、リサーチ要約、外部エージェント接続という周辺ワークフローから制度化されることを示す。
\end{casestudy}

これらは前節までのヘッジファンドと異なり、\emph{アルファ創出}よりも\emph{全社的な生産性と業務基盤の再編}を主眼とする点に特徴がある。以下に本節で取り上げた主な機関のAIアプローチと根拠を一覧化する。

\begin{center}
\footnotesize
\begin{tabularx}{\textwidth}{@{}>{\raggedright\arraybackslash}p{2.9cm} >{\raggedright\arraybackslash}p{2.6cm} >{\raggedright\arraybackslash}X >{\raggedright\arraybackslash}p{3.0cm}@{}}
\toprule
\textbf{機関} & \textbf{AIアプローチの類型} & \textbf{中核的な取り組み} & \textbf{開示された根拠} \\
\midrule
Bridgewater (AIA Labs) & AIを投資判断の中核 & LLM＋表形式ML＋独自マクロDBで「人工投資家」 & 初年度+11.9\%\cite{bridgewater_aialabs_nd} \\
Renaissance (Medallion) & 自動売買の長期実績 & 1日15万〜30万件の統計的自動売買 & 設立来年率39.9\%\cite{institutionalinvestor2025renaissance} \\
QRT & 中核＋研究者網 & 急成長AUM＋QRT Labs & 2025年主力計約+30\%、AUM 380億ドル\cite{bloomberg_qube2026} \\
AQR & 懐疑派の翻意 & ML＝主力ファンドのシグナル約1/5 & AUM約1,870億ドル\cite{x03aqr_bloomberg2025,x03forbes_billionaires2026} \\
Balyasny & フェデレーテッドAI & 中央Applied AI→全チーム95\%へエージェント & Deep Research：数日→数時間\cite{x03balyasny_openai2026} \\
Man Group & 漸進・生産性重視 & AlphaTrendエージェントパイプライン & Sharpe約1.05へクラスタ\cite{manndalphatrend} \\
Two Sigma & 研究の「OS」化＋再編 & Platform Thinking／数千エージェント構想 & 200名削減、AUM約700億ドル\cite{x03twosigma_hedgeweek2026} \\
Point72 / Cubist & 人材育成両輪 & Quant Academy・GenAIハッカソン & Cubist AUM約400億ドルの約17\%\cite{bloomberg_cubistchief2025} \\
D.E.\ Shaw & 二正面投資 & 独立MLグループ＋人間主導新ファンド & Oculus約+28.2\%\cite{bloomberg_deshaw2025} \\
Citadel & 懐疑から転換 & 社内リサーチAIで数週間→数時間 & Griffin発言の急転換\cite{griffin2026} \\
JPMorgan & 全社基盤の民主化 & LLM Suite（7モデル）→エージェント & 従業員23万人超、450+ユースケース\cite{x03jpmorgan_cnbc2025} \\
Goldman Sachs & 全社基盤＋自律コーダー & GS AI＋Devin「ハイブリッド労働力」 & 46,500名、生産性3〜4倍見込み\cite{x03goldman_cnbc2025} \\
BlackRock & 運用基盤への組込み & Aladdin Copilot（マルチエージェント） & 全クライアント提供、数百API\cite{x03blackrock_aladdincopilot} \\
Morgan Stanley & 評価駆動の業務AI & Assistant／Debrief／AskResearchGPT／MCP接続 & Advisorチーム98\%超、文書アクセス20\%→80\%\cite{x03morganstanley_openai2026} \\
Standard Signal & 完全自律型 & 人間PM不在のエンドツーエンド運用 & ライブSharpe比3超を主張\cite{ycombinator2026standardsignal} \\
\bottomrule
\end{tabularx}
\end{center}

\smallskip
{\small\color{figgray}\textbf{表: 主要機関のAIアプローチ比較（機関×類型×取り組み×根拠）。}開示ベース・測定条件が不揃いのため数値の直接比較はできない。左列上部ほどAIを投資判断の中核に、下部ほど全社基盤・自律型に位置づく傾向がある。}

\subsection{完全自律型AIヘッジファンドという急進モデル}

\begin{casestudy}{Standard Signal (Y Combinator 2026 Spring): 完全自律型を標榜}
AIが調査・意思決定・執行をすべてエンドツーエンドで行う完全自律型AIヘッジファンドを標榜し、強化学習モデルを用いて人間のポートフォリオマネージャーを介さずに取引を実行すると主張。ライブSharpe比3超を主張している（2026年4月）。ただし、後述する再現性・先読みバイアスの問題（\S\ref{sec:llm-analysis}）を踏まえると、こうした主張は独立した検証を経ていない限り慎重に評価すべきである。（出典: \cite{ycombinator2026standardsignal}）
\end{casestudy}

\begin{casestudy}{Lumenai × ETS Asset Management Factory: 「初の機関投資家向け完全エージェント型」を標榜}
Lumenaiは運用会社ETS Asset Management Factoryと提携し、\textbf{完全エージェント型アーキテクチャ}上に構築した「初の機関投資家向けヘッジファンド」と称するファンドの立ち上げを2026年6月に発表した。自律エージェント群がグローバル株式のロング・ショートのアイデアを継続的に生成・評価・運用し、人間はガバナンス上の監督のみを担うとされる。Standard Signalと同様、意思決定を人間PMではなくエージェントに委ねる急進モデルであり、機関投資家向けを標榜する点で完全自律型が個人・スタートアップ領域を越えて機関市場に波及しつつあることを示す一方、その実運用実績と検証可能性は今後の開示を待つ必要がある。（出典: Hedgeweek \cite{lumenai2026}）
\end{casestudy}

\begin{center}
\begin{tikzpicture}[
  spec/.style={font=\footnotesize, align=left, text width=3.0cm, inner sep=1mm},
  marker/.style={circle, fill=figblue, draw=fignavy, figshadow, inner sep=1.6pt},
  leader/.style={draw=figgray, thin}
]
\draw[farr, thick] (0,0) -- (16.0,0);
\node[font=\footnotesize, color=fignavy, anchor=west] at (0,0.35) {\textbf{漸進的・人間主導}};
\node[font=\footnotesize, color=figaccent, anchor=east] at (16.0,0.35) {\textbf{急進的・AI主導}};

\node[marker] (m1) at (1.7,0) {};
\node[marker] (m2) at (4.9,0) {};
\node[marker] (m3) at (8.0,0) {};
\node[marker] (m4) at (11.1,0) {};
\node[marker] (m5) at (14.3,0) {};

\node[spec] (l1) at (1.7,-1.9) {\textbf{Man Group\,/\,Two Sigma\\Point72\,/\,Voleon\\WorldQuant}\\人材育成・インフラ投資で漸進導入};
\node[spec] (l2) at (4.9,1.7) {\textbf{Citadel}\\「AI＝ゴミ」発言を撤回し社内リサーチAIを導入};
\node[spec] (l3) at (8.0,-1.7) {\textbf{D.E.\ Shaw}\\独立MLグループと人間主導新ファンドの二正面投資};
\node[spec] (l4) at (11.1,1.7) {\textbf{Bridgewater\,/\,Renaissance\\QRT}\\AIを中核に据えて実運用資本を稼働};
\node[spec] (l5) at (14.3,-1.7) {\textbf{Standard Signal}\\人間PM不在の完全自律型を標榜};

\draw[leader] (m1) -- (l1.north);
\draw[leader] (m2) -- (l2.south);
\draw[leader] (m3) -- (l3.north);
\draw[leader] (m4) -- (l4.south);
\draw[leader] (m5) -- (l5.north);
\end{tikzpicture}
\end{center}

\smallskip
{\small\color{figgray}\textbf{図: AI活用アプローチの二極化スペクトラム。}}左端の「漸進的・人間主導」から右端の「急進的・AI主導」まで、本節で取り上げた事例を概念的に配置した。位置は各社の開示情報にみられる体制の性質を示す模式図であり、厳密な定量指標に基づく順位付けではない点に注意されたい。

\begin{insight}{示唆: 二極化する実装アプローチ}
大手ヘッジファンド（Man Group、Point72/Cubist、Two Sigma、AQR等）はAI活用を人材育成とインフラ投資の両輪で漸進的に進める一方、Standard SignalのようなYC発スタートアップは完全自律型という急進的モデルを掲げる。さらにこの二極に加え、JPMorgan・Goldman・BlackRock・Morgan Stanleyに代表される大手金融機関は、アルファ創出よりも数万〜数十万人規模での生産性向上と業務基盤の再編を主眼とする「全社基盤としての民主化＋段階的なエージェント移行」という第三の型を形成しつつある。Balyasnyの中央Applied AIチームが全投資チームへエージェントを配る「フェデレーテッドAI」は、この全社基盤型とアルファ創出型を橋渡しする設計として注目に値する。本節前半に示した公表リターン比較図の水準そのものにも表れているように、両者は前提とする戦略・資本規模・検証可能性の水準がそもそも異なり、単純な優劣比較はできない。実証性（特にライブSharpe比のような主張の検証可能性）には注意が必要であり、\S\ref{sec:llm-analysis}で扱う再現性の課題と併せて評価する必要がある。
\end{insight}

\subsection{他業界の「AI-Native組織」と金融ファンドの対比}
\label{sec:ai-native-funds-crossdomain}

AIを運用の中核（operating model）に据える組織づくりは金融の専売特許ではない。ソフトウェア開発企業、コンサルティング・クラウドベンダー、法務・カスタマーサポートといった周辺のIT・AI業界では、金融より一足先に「AI-Native組織」への転換を本番運用の規模で進めている。本項では、これらの周辺分野の事例のうち\textbf{金融の組織設計にそのまま転用・比較できるもの}のみを厳選し、いずれも最終的に金融ファンドへの示唆に着地させる。並べること自体を目的としない。

\begin{casestudy}{ソフトウェア開発組織のAI-Native化: AnthropicとCognition（Devin）の「エージェント既定化」}
AI企業自身が自社の開発組織をAI-Nativeに再設計する事例が現れている。Anthropicは自社のClaude Code／Claude Cowork開発組織において、マネージャーもICから出発するフラットな体制と、独立に裁量を持つ小規模「ポッド」を採用し、従来の半年単位のロードマップを、プロトタイプと社内ユーザーのフィードバックに基づく\textbf{Just-in-Time（JIT）プランニング}に置き換えた。コードレビューはスタイル・lint・テスト・バグをClaudeが担い、人間は法務・セキュリティ・プロダクト判断など専門性が要る領域に集中する。Director of EngineeringのFiona Fungは「デフォルトで、すべてのコミットはClaude支援済み——過去4か月間、Claude支援でないコミットを見た記憶がない」と述べている\cite{y03anthropic2026ainativeorg}。同様に、自律コーディングエージェントDevinを開発するCognitionでは、\textbf{自社コードの89\%}が既にDevin作成であり（2025年12月時点の13\%から急上昇）、2026年5月時点で年換算収益（ARR）は前年の3,700万ドルから\textbf{4億9,200万ドル}へと13倍に拡大した。エンタープライズ顧客にはGoldman Sachs・Citi・Santanderといった金融機関も含まれ、Ita\'uはDevinによりセキュリティ脆弱性の70\%を自動解消、Mercedes-Benzは8か月を要したレガシー刷新を8日間に短縮したと開示している\cite{y03techcrunch2026cognitionraise}。
\end{casestudy}

\begin{insight}{金融への示唆: 「エージェント既定化率」を定性的スローガンでなく実測指標にする}
AnthropicとCognitionの共通点は、「AIをエージェント既定にする」という方針を、コミット比率（89\%、Claude支援コミット比率）という\textbf{検証可能な内部指標}として開示している点にある。本節で紹介したBalyasnyの「全投資チームの約95\%へエージェント配布」\cite{x03balyasny_openai2026}やJPMorgan LLM Suiteの「約半数が日次利用」\cite{x03jpmorgan_cnbc2025}は、金融側でも同種の実測開示が既に始まっていることを示すが、Anthropic・Cognitionほど粒度の細かい生産物ベースの指標（コード生成比率のような「アウトプットのうちAIが作った割合」）を継続的に開示している金融機関はまだ少ない。ファンド・銀行がAI-Native化の進捗を対外的に主張する際は、稼働率のようなインプット指標だけでなく、Cognitionの「コミットの89\%」に相当する\textbf{アウトプットベースの指標}（例: リサーチメモ・与信メモのうちエージェント初稿の割合）を定期開示することが、Standard Signalのような未検証の主張と一線を画す実務的な処方箋になる。
\end{insight}

\begin{casestudy}{全社的AIプラットフォームの「中央集約」型: コンサル・ERPベンダーに見る収斂パターン}
15社のAI先進企業への取材に基づくMcKinseyの「AI-Native企業の7つの運用原則」は、\textbf{エージェントを同僚として扱う}ことが即座に\textbf{build vs buy}の判断を迫り、ナレッジ層の整備・構成可能でガバナンスの効いたアーキテクチャ・段階的な信頼構築・適切な組織設計・そして一過性のIT導入ではなく\textbf{文化的フライホイールとしての定着}へと連鎖することを指摘する\cite{y03mckinsey2026sevenoperatingtruths}。MIT Sloan Management ReviewとBoston Consulting Groupの共同調査（2025年11月）も、エージェント型AIの導入率が既に35\%に達し44\%が導入を計画する一方、\textbf{競争優位はエージェントへのアクセスの早さではなく組織設計そのものから生まれる}と結論づけ、「投資か雇用か」「監督か自律か」など4つの戦略的緊張関係への対応が今後3年の運用モデル・役職・キャリアパスの根本的変化につながると予測する\cite{y03mitsloanbcg2025agenticenterprise}。実装面では、PwCとGoogle Cloudが共同で設置した「Gemini Enterprise Center of Excellence」が好例で、データ系譜管理・ガバナンス・認証に特化した\textbf{30以上の再利用可能なエージェント群}を中央チームが構築したうえで、小売・医療・金融変革など業界横断でクライアントごとにカスタマイズして展開している\cite{y03pwc2026googlecoe}。同様に中央プラットフォームベンダー自身が自社業務でエージェントを実証する例として、Salesforceのエージェント基盤Agentforceは年換算収益（ARR）が\textbf{12億ドル超}（顧客数18,500社超、前年比+205\%）に達し、自社のヘルプデスクhelp.salesforce.comでは430万件の問い合わせのうち\textbf{70\%を人手を介さず自動解決}している\cite{y03salesforce2026agentforcesuccess}。ERP最大手のSAPも2026年、自社スイート全体を「Autonomous Enterprise」として再定義し、業務基盤としての会計・サプライチェーン・人事プロセスをエージェントがエンドツーエンドで担い、人間は例外処理とガバナンスに専念する構想を打ち出した\cite{y03forbes2026sapautonomous}。もっとも、Gartnerは企業アプリの40\%が2026年末までにタスク特化型AIエージェントを組み込むと予測する一方（2025年時点で5\%未満から急上昇）、2027年末までにエージェント型AIプロジェクトの40\%超がコスト・不明瞭な価値・不十分なリスク統制を理由に中止されるとも警告しており、拡大と幻滅が同時進行する段階にあることを示している\cite{y03gartner2025agentapps40pct}。
\end{casestudy}

\begin{insight}{金融への示唆: Balyasny・BlackRockの設計は業界横断で収斂しつつある「型」である}
PwC×Google Cloudの「中央チームが再利用可能なエージェント群を構築し、事業部門・クライアントごとに展開する」という設計は、本節で見たBalyasnyの中央Applied AIチームが全投資チームへエージェントを配布する「フェデレーテッドAI」\cite{x03balyasny_openai2026,x03balyasny_mlq2026}や、BlackRockのAladdin Copilotが全Aladdinクライアントへ横断的に提供される構造\cite{x03blackrock_aladdincopilot}と、構造的にほぼ同型である。これは、金融ファンドが独自に発明した設計ではなく、\textbf{AIを本番運用に組み込む組織一般が収斂する「中央構築・現場配布」パターン}であることを示唆する。同時に、Gartnerが警告する「プロジェクトの40\%超が中止される」という数字は、本節前半で紹介したStandard Signal・Lumenaiのような検証途上の急進的主張（\S\ref{sec:ai-native-funds}）を評価する際の外部的な物差しとなる——中央プラットフォーム型の地道な定着（JPMorgan・Goldman・Balyasny）と、急進的な完全自律の主張とでは、業界横断で見ても前者の方が持続する確率が高いことを示唆している。
\end{insight}

\begin{casestudy}{完全自動化の限界と揺り戻し: Klarnaの巻き戻しとHarveyの「人間監督付きエンドツーエンド」}
カスタマーサポート分野の代表事例であるKlarnaは、2024年にAIエージェントが700人分（後に\textbf{853人分}）の顧客対応業務を代替したと発表し、「AI-Nativeが人員を代替する」旗艦事例として広く引用された\cite{y03cxdive2025klarna853}。しかし2025年から2026年にかけて、顧客満足度の低下や運用上の不具合を受けて人的な顧客対応体制を静かに再構築し、CEOのSebastian Siemiatkowskiは「AI主導の人員削減はやり過ぎだった」と認めた。現在はAIが問い合わせの約3分の2を処理し、エスカレーション・感情的に複雑な案件・判断を要する案件は人間が担う\textbf{ハイブリッド型}に落ち着いている\cite{y03fastcompany2026klarnareversal}。対照的に、法務AIのHarveyは2026年3月に評価額\textbf{110億ドル}（2025年6月の50億ドルから急拡大）を達成し、契約書分類・訴訟ワークフロー・デューデリジェンス追跡といった業務をエージェントがエンドツーエンドで実行する一方、ノーコードのWorkflow Builderにより各ステップに人間の承認・ロールベースの権限設定を組み込む設計を貫いている\cite{y03harvey2026raise11b}。両社の対比は、「AIに業務を委ねる」ことと「人間の関与点を設計から排除する」ことが同義ではないことを示している。
\end{casestudy}

\begin{insight}{金融への示唆: Magnetar型の「委任するが最終権限は人間」設計の妥当性を裏づける}
Klarnaの巻き戻しは、本節で紹介したMagnetar Capitalが数百のAIボットにリサーチ業務を委ねながら\textbf{最終的な発注権限を人間のポートフォリオマネージャーに留保する}設計\cite{magnetar2026}や、Morgan StanleyがDebriefの出力をアドバイザーが確定前に確認する人間監督を組み込む設計\cite{x03morganstanley_openai2026}の妥当性を、他業界の失敗事例から裏づけるものである。逆に、Standard Signalの「人間PM不在の完全自律型」\cite{ycombinator2026standardsignal}やLumenaiの「初の機関投資家向け完全エージェント型」\cite{lumenai2026}は、Klarnaが顧客対応という定型性の高い業務ですら人間の安全弁を再導入せざるを得なかった前例に照らすと、より不確実性の高い投資判断において同種の巻き戻しリスクを内包している可能性がある。他方でHarveyのように、エンドツーエンド実行と各ステップでの人間承認を両立させる設計は、Moody'sの与信メモ生成エージェント\cite{moodys2026creditmemo}やBalyasnyのDeep Researchエージェント\cite{x03balyasny_openai2026}が採る「エージェントが実行、人間がチェックポイントで検証」という金融側の主流パターンと軌を一にしており、これが業界を問わず持続可能な設計であることを補強している。
\end{insight}

\begin{center}
\begin{tikzpicture}[
  card/.style={fnode, text width=3.5cm, align=left, font=\footnotesize}
]
\node[card] (sw) {\textbf{\color{fignavy}ソフトウェア開発}\\[0.5mm] Anthropic：全コミットAI支援既定化\\ Cognition：自社コードの89\%をDevinが作成};
\node[card, fnode blue, right=4mm of sw] (coe) {\textbf{\color{fignavy}コンサル・ERP}\\[0.5mm] PwC×Google Cloud：中央30+エージェント群を業界横断展開\\ SAP：ERP全体を自律化構想};
\node[card, fnode accent, right=4mm of coe] (support) {\textbf{\color{fignavy}カスタマーサポート・法務}\\[0.5mm] Klarna：完全自動化→人間監督つきハイブリッドへ巻き戻し\\ Harvey：エンドツーエンド実行＋各ステップ承認};
\node[card, fnode dark, right=4mm of support] (startup) {\textbf{\color{fignavy}AI-Nativeスタートアップ}\\[0.5mm] 新規ユニコーンの約25\%がAI-Native\\ 「ソロコーン」：一桁人数で評価額10億ドル超};
\end{tikzpicture}
\end{center}

\smallskip
{\small\color{figgray}\textbf{図: 周辺分野4類型におけるAI-Native組織化の到達点。}}\cite{y03anthropic2026ainativeorg,y03techcrunch2026cognitionraise,y03pwc2026googlecoe,y03forbes2026sapautonomous,y03cxdive2025klarna853,y03fastcompany2026klarnareversal,y03harvey2026raise11b,y03crunchbase2025unicorninvestors} いずれも金融ファンド固有ではなく、AIを本番運用に組み込む組織一般に見られる型である点に注意されたい。

\begin{center}
\footnotesize
\begin{tabularx}{\textwidth}{@{}>{\raggedright\arraybackslash}p{2.6cm} >{\raggedright\arraybackslash}p{3.6cm} >{\raggedright\arraybackslash}X >{\raggedright\arraybackslash}p{3.4cm}@{}}
\toprule
\textbf{周辺業界の類型} & \textbf{組織パターン} & \textbf{根拠} & \textbf{対応する金融ファンドの事例} \\
\midrule
ソフトウェア開発 & エージェント支援をデフォルト化し比率で実測開示 & Claude支援コミット「デフォルト」、Devinコード比率89\%\cite{y03anthropic2026ainativeorg,y03techcrunch2026cognitionraise} & Balyasny：全投資チームの約95\%へエージェント配布\cite{x03balyasny_openai2026} \\
コンサル・ERP & 中央チームが再利用可能エージェント群を構築し現場・クライアントへ配布 & PwC×Google Cloud CoE：30+エージェント、業界横断展開\cite{y03pwc2026googlecoe} & BlackRock Aladdin Copilot：全クライアントへの横断提供\cite{x03blackrock_aladdincopilot} \\
カスタマーサポート・法務 & 完全自動化の反省からエンドツーエンド実行＋人間チェックポイントへ収斂 & Klarnaの巻き戻し、Harveyの承認付きワークフロー\cite{y03fastcompany2026klarnareversal,y03harvey2026raise11b} & Magnetar：数百ボット委任、発注権限は人間に留保\cite{magnetar2026} \\
AI-Nativeスタートアップ & 少人数チーム＋エージェントレバレッジで高評価額を実現 & 新規ユニコーンの約25\%がAI-Native、「ソロコーン」現象\cite{y03crunchbase2025unicorninvestors} & Voleon：約234名でAUM 234億ドルを運用\cite{voleon_wikipedia2026} \\
\bottomrule
\end{tabularx}
\end{center}

\smallskip
{\small\color{figgray}\textbf{表: 周辺業界の組織パターンと金融ファンド事例の対応関係。}}測定条件・開示粒度が業界間で不揃いなため、数値の直接比較ではなく\textbf{構造的パターンの対応}として読むべき点に注意されたい。

\begin{insight}{示唆: AI-Native化は金融固有の現象ではなく、組織設計として業界横断で収斂しつつある}
本項で見た4類型はいずれも、程度の差はあれ「①エージェント関与を既定化し実測指標で開示する」「②中央チームが再利用可能な基盤を構築し現場へ federated に展開する」「③完全自動化の主張は人間のチェックポイントを伴う設計に収斂しやすい」という3つの共通パターンへ向かっている。これは、Balyasny・BlackRock・Magnetar・JPMorganといった本節の金融事例が採る組織設計が、金融特有の思いつきではなく、AIを本番運用の中核に据える組織一般が辿り着く合理的な均衡点に近いことを示唆する。逆に言えば、Standard SignalやLumenaiのように人間の関与点を設計から排除する急進モデルは、周辺業界（特にKlarnaの前例）が既に経験した巻き戻しリスクを、独立した検証データなしに繰り返している可能性がある点に、金融機関・投資家は留意すべきである。
\end{insight}

\section{投資意思決定のためのAIエージェント}
\label{sec:investment-agents}

本節では、LLM・AIエージェントが財務分析・投資アイデア創出・デューデリジェンス・ポートフォリオモニタリング・リスク分析・意思決定を支援・自動化する事例を扱う。マルチエージェント・アーキテクチャは急速に高度化している一方、独立したベンチマークは依然として大きな性能ギャップを一貫して報告している。

\begin{concept}{マルチエージェントとオーケストレーション}
1つのLLMエージェントに情報収集からリスク管理・執行まで全部を任せると、判断が偏ったり大事な見落としが起きやすく、たった一人で全部署の仕事をこなそうとする社員のように行き詰まりやすい。マルチエージェントとは、ファンダメンタル分析役・センチメント分析役・リスク管理役・執行役のように役割を分担した複数のエージェントが協働する構成である。この協調をどう制御するかの設計を\textbf{オーケストレーション}と呼び、代表的な型として、監督役が計画し担当者に割り振る\textbf{supervisor-worker}（階層型）、複数の担当が互いに反論し合い誤りを減らすが偽の合意に陥りうる\textbf{debate}、自分の推論の弱点を自己診断して直す\textbf{reflection（内省）}がある。

\begin{center}
\begin{tikzpicture}[font=\small,
  sup/.style={fnode dark, text width=3.0cm, minimum height=1.0cm, font=\bfseries\small},
  wk/.style={fnode blue, text width=2.3cm, minimum height=1.0cm, font=\scriptsize},
  arr/.style={farr blue}
]
\node[sup] (sup) at (4.35,1.6) {オーケストレーター\\（supervisor）};
\node[wk] (a) at (0,0) {ファンダメンタル\\分析役};
\node[wk] (b) at (2.9,0) {センチメント役};
\node[wk] (c) at (5.8,0) {リスク管理役};
\node[wk] (d) at (8.7,0) {執行役};

\draw[arr] (sup.south) -- ++(0,-0.3) -| (a.north);
\draw[arr] (sup.south) -- ++(0,-0.3) -| (b.north);
\draw[arr] (sup.south) -- ++(0,-0.3) -| (c.north);
\draw[arr] (sup.south) -- ++(0,-0.3) -| (d.north);

\node[font=\scriptsize\itshape, text=figgray, align=left, text width=10.3cm] at (4.35,-1.6)
  {他の代表的な型: \textbf{debate}（複数エージェントが反論し合い誤りを減らすが偽の合意に陥る危険もある）／\textbf{reflection}（自分の推論を自己診断し修正するループ）};
\end{tikzpicture}

\vspace{1mm}
{\small\color{brandgray}\textbf{図: マルチエージェント・オーケストレーションの入門地図}\ supervisor-worker型では中央のオーケストレーターが役割分担した複数のworkerに作業を割り振る。debateとreflectionは別の代表的な協調パターンで、詳細は後段の図で扱う。}
\end{center}
\end{concept}

\subsection{マルチエージェント・リサーチワークフロー}

\begin{casestudy}{TradingAgents: ディベート構造を持つマルチエージェント取引フレームワーク}
強気／弱気ディベーター、ファンダメンタル・センチメント・テクニカルアナリスト、専属リスクチームからなるマルチエージェントシステム。シミュレーション市場条件下で、単一エージェント基準に対しSharpe比と累積リターンの改善を実証した。（出典: \cite{tradingagents2024}）
\end{casestudy}

\begin{casestudy}{MarketSenseAI 2.0: 強力な実運用評価を伴う株式分析エージェント}
Chain-of-Agentsアーキテクチャで10-K/10-Q、決算コールトランスクリプト、マクロデータ、価格履歴を処理。S\&P 100バックテスト（2023〜2024年）で\textbf{累積リターン125.9\%}（指数の73.5\%を上回る）、S\&P 500 2024年でSortinoレシオ\textbf{33.8\%改善}という、公表されているLLM投資エージェントの中でも最も強力な実運用評価事例の一つ。（出典: \cite{marketsenseai2025}）
\end{casestudy}

\begin{center}
\begin{tikzpicture}[font=\small,
  agentbox/.style={fnode blue, text width=3.1cm, minimum height=1.4cm},
  debatebox/.style={fnode accent, text width=2.9cm, minimum height=1.2cm},
  riskbox/.style={fnode, text width=7.6cm, minimum height=1.1cm},
  decbox/.style={fnode dark, text width=7.6cm, minimum height=1.3cm, font=\bfseries\small},
  arr/.style={farr blue}
]
\node[agentbox] (fund) at (0,0) {ファンダメンタル\\アナリスト};
\node[agentbox] (sent) at (3.5,0) {センチメント\\アナリスト};
\node[agentbox] (tech) at (7.0,0) {テクニカル\\アナリスト};

\node[debatebox] (bull) at (1.2,-2.0) {強気ディベーター};
\node[debatebox] (bear) at (5.8,-2.0) {弱気ディベーター};

\node[riskbox] (risk) at (3.5,-3.7) {専属リスクチーム};
\node[decbox] (dec) at (3.5,-5.3) {投資意思決定\\（単一エージェント基準比でSharpe比・累積リターンが改善）};

\draw[arr] (fund.south) -- ++(0,-0.6) -| (bull.north);
\draw[arr] (sent.south) -- ++(0,-0.3) -| (bull.north);
\draw[arr] (sent.south) -- ++(0,-0.3) -| (bear.north);
\draw[arr] (tech.south) -- ++(0,-0.6) -| (bear.north);
\draw[fbiarr] (bull.east) -- (bear.west)
      node[midway, above, font=\scriptsize\itshape, text=figaccent] {ディベート};
\draw[arr] (bull.south) -- ++(0,-0.4) -| (risk.north);
\draw[arr] (bear.south) -- ++(0,-0.4) -| (risk.north);
\draw[arr] (risk) -- (dec);
\end{tikzpicture}

\vspace{1mm}
{\small\color{figgray}\textbf{図: マルチエージェント・リサーチワークフローの例（TradingAgents）}\
ファンダメンタル・センチメント・テクニカルの3アナリストが情報を提供し、強気／弱気ディベーターが議論、専属リスクチームを経て投資意思決定に至る構造。}
\end{center}

\begin{itemize}
  \item \textbf{FundaPod}: 異なる投資哲学を体現する複数のAIペルソナが共有ナレッジグラフ上で協働し、根拠に紐づくファンダメンタル・リサーチメモを生成する、人間-AIハイブリッド型の機関投資家向けプラットフォーム（\cite{fundapod2026}）。
  \item \textbf{QRAFTI}: MCPサーバー経由の構造化パネルデータツールと内省型プランニングを統合し、動的コード生成単独よりもツール連鎖呼び出しの方が株式ファクター研究タスクで優位と実証した（\cite{qrafti2026}）。
  \item \textbf{AlphaAgents（BlackRock）}: ファンダメンタル・センチメント・バリュエーションの3マイクロエージェントが、Microsoft AutoGen上のグループチャットで協調して株式ポートフォリオを構築する。意見が分岐した場合はラウンドロビン型の\textbf{マルチエージェント・ディベート}で合意へ収束させて幻覚と説明不能性を抑え、プロンプトエンジニアリングでエージェントごとに\textbf{リスク許容度}（risk-neutral／risk-averse）を付与して実際の投資家プロファイルを模擬する。役割特化エージェント群の成果を既存ベンチマークと比較して批判的に評価している（\cite{alphaagents2025}）。
  \item \textbf{Hubble}: LLM＋ドメイン特化オペレーター言語＋ASTサンドボックス＋デュアルチャネルRAGを組み合わせたエージェント型パイプラインで株式アルファファクターを自動探索。3ラウンド・100候補の評価で実行クラッシュ\textbf{0件}を達成し、ヘッジファンドのクオンツリサーチワークフローを直接標的とする。（\cite{hubblealpha2026}）
  \item \textbf{The Self Driving Portfolio}: Andrew Ang氏らが提案する約50体のエージェント群による機関投資家向け資産運用アーキテクチャ。専門化されたエージェントが資本市場前提を生成し、20種類以上のポートフォリオ構築手法を競わせ投票で出力を決定。メタエージェントが過去の予測誤差からエージェントのコードを継続的に改善し、人間はInvestment Policy Statementの遵守を監督するガバナンス役に後退するという設計思想を提示している。（\cite{selfdrivingportfolio2026}）
  \item \textbf{FinRobot}: 金融アプリケーション向けのオープンソースなAIエージェント・プラットフォーム。エージェント層／LLMアルゴリズム層／LLMOps・DataOps層／マルチソースLLM層の4層アーキテクチャで構成され、プロのアナリストから一般ユーザーまでが独自のエージェント・パイプラインを組み立てられるよう、金融AIエージェントのツールチェーンへのアクセスを民主化する（\cite{finrobot2024}）。
  \item \textbf{HedgeAgents}: 中央のファンドマネージャー・エージェントが、3種類のエージェント間カンファレンスを通じてドメイン特化のヘッジングエージェント群を協調させる「バランス志向」のマルチエージェント取引システム。3年間のシミュレーションで年率\textbf{70\%}・累積\textbf{400\%}のリターンを達成したと報告し、WWW 2025でoral発表された（\cite{hedgeagents2025}）。
  \item \textbf{TradingGPT}: 本系譜の初期（2023年9月）に位置づけられる先駆的フレームワーク。各エージェントの記憶を短期・中期・長期の3層に分け、それぞれに減衰機構を持たせて人間の認知過程に近づける「階層型メモリ」を提案した。各層内の記憶は\textbf{recency（新しさ）・relevancy（関連性）・importance（重要性）}の3指標でランク付けされ、個性を付与されたエージェント同士がディベートを行う。後続のTradingAgents・FinConが採用するメモリ＋ディベート型アーキテクチャの概念的原型である（\cite{x04tradinggpt2023}）。
  \item \textbf{FinCon}: 実在する投資会社の組織構造に着想を得た\textbf{マネージャー-アナリスト階層}を持つマルチエージェントシステム。中核となる\textbf{概念的言語強化（conceptual verbal reinforcement）}では、獲得した投資信念を言語化した教訓として、更新を必要とするエージェントノードにのみ選択的に伝播させ、全対全の冗長な通信を抑制する。さらにリスク制御コンポーネントが挿話的（episodic）に自己批判を起動し、systematicな投資信念を改訂する。単一銘柄取引とポートフォリオ管理の双方で、最先端のLLMベース・DRLベースのエージェントを累積リターン・Sharpe比で上回りつつ、最大ドローダウンを最小水準に抑えたと報告した（NeurIPS 2024、\cite{x04fincon2024}）。
  \item \textbf{FinRobot（エクイティリサーチ版）}: FinRobotプラットフォーム（\cite{finrobot2024}）とは別に提案された、Chain-of-Thoughtを多段で連鎖させるエクイティリサーチ・パイプライン。多源データを統合する\textbf{Data-CoT}、アナリストの推論を模倣する\textbf{Concept-CoT}、投資thesisとレポートに統合する\textbf{Thesis-CoT}の3エージェントからなる。生成レポートは投資銀行アナリストのレビューで正確性・一貫性・叙述の整合性が高く評価され、バリュエーションタスクで約\textbf{75\%}の精度を報告した（\cite{x04finrobotequity2024}）。
\end{itemize}

\begin{insight}{仮説: 構造化された敵対的ディベートによる投資リサーチ}
主張を担うthesisエージェントが、専属の強気・弱気カウンターエージェントと少なくとも3ラウンド対峙してから最終判断に至る構造化された敵対的ディベートは、単一エージェント生成が見落とす反証や隠れた前提を明示的に浮かび上がらせ、アンカリングや確証バイアスの克服に資すると期待される。検証案としてRussell 2000から200銘柄を選び、単一エージェント／単純アンサンブル／3ラウンドの構造化ディベートで投資thesisを生成し、3か月先の方向的正確性と情報係数（IC）で比較する90日間のプロスペクティブ実験が考えられる。ただしディベート参加エージェントが同一の訓練コーパス由来の潜在バイアスを共有する場合、偽の合意に収束しうること、役割を制約するシステムプロンプトの慎重な設計が必要なこと、トークンコストが3〜5倍に増える点が、流動性の高い大型株のルーチンなカバレッジでは正当化されない可能性が交絡・リスクとして挙げられる（出典: \cite{multiagentdebate2025}）。
\end{insight}

\subsection{オーケストレーション・パターンと「調整プリマシー」}

前項までのフレームワークは、表面的には多様に見えて、その内部は少数の反復するオーケストレーション（協調制御）パターンの組み合わせに還元できる。実務上支配的なのは次の3類型である。\textbf{(1) supervisor-worker（監督-作業者）階層}――中央のスーパーバイザーが計画を立て制約を監視し、専門ワーカーに最小限の手順を割り当てる。軌道（trajectory）の安定性に優れる一方、スーパーバイザーが単一障害点かつスループットのボトルネックになる。\textbf{(2) debate（ディベート）}――複数エージェントが共有会話で反論し合い、単一モデル照会に比べ幻覚を減らすが、誤っていても多数派に同調して\emph{偽の合意}を生みやすい。\textbf{(3) reflection／reflexion（内省・自己修正）}――意思決定後に自らの推論の弱点を診断し計画を再生成する批判-改訂ループで、TradingGroup（\cite{x04tradinggroup2025}）やSEP系の研究が採用する。実運用のマルチエージェントシステムの多くは、これらを合成した\textbf{ハイブリッド}（例: 階層の末端チームだけがメッシュ協調する）として構成される。

\begin{center}
\begin{tikzpicture}[font=\small,
  sup/.style={fnode accent, text width=2.2cm, minimum height=0.8cm, font=\scriptsize\bfseries},
  wk/.style={fnode blue, text width=1.7cm, minimum height=0.7cm, font=\scriptsize},
  ag/.style={fnode blue, text width=1.7cm, minimum height=0.7cm, font=\scriptsize},
  rf/.style={fnode, text width=1.9cm, minimum height=0.7cm, font=\scriptsize},
  arr/.style={farr blue}, biarr/.style={fbiarr}
]
% (1) supervisor-worker
\node[font=\scriptsize\bfseries, text=figaccent!70!black] at (1.6,1.2) {(1) supervisor-worker};
\node[sup] (sp) at (1.6,0.4) {スーパーバイザー};
\node[wk] (w1) at (0.2,-1.1) {ワーカーA};
\node[wk] (w2) at (1.6,-1.1) {ワーカーB};
\node[wk] (w3) at (3.0,-1.1) {ワーカーC};
\draw[arr] (sp) -- (w1); \draw[arr] (sp) -- (w2); \draw[arr] (sp) -- (w3);
% (2) debate
\node[font=\scriptsize\bfseries, text=figaccent!70!black] at (6.5,1.2) {(2) debate};
\node[ag] (d1) at (5.7,0.2) {エージェント1};
\node[ag] (d2) at (7.4,0.2) {エージェント2};
\node[ag] (d3) at (6.5,-1.1) {エージェント3};
\draw[biarr] (d1) -- (d2); \draw[biarr] (d1) -- (d3); \draw[biarr] (d2) -- (d3);
% (3) reflection
\node[font=\scriptsize\bfseries, text=figaccent!70!black] at (11.2,1.2) {(3) reflection};
\node[rf] (act) at (11.2,0.3) {行動案の生成};
\node[rf] (crit) at (11.2,-1.2) {自己批判・診断};
\draw[arr] (act.east) .. controls +(1.1,0) and +(1.1,0) .. (crit.east)
      node[midway,right,font=\scriptsize\itshape,text=figgray] {批判};
\draw[arr] (crit.west) .. controls +(-1.1,0) and +(-1.1,0) .. (act.west)
      node[midway,left,font=\scriptsize\itshape,text=figgray] {改訂};
\end{tikzpicture}

\vspace{1mm}
{\small\color{figgray}\textbf{図: 投資エージェントで反復する3つのオーケストレーション・パターン。}\
supervisor-workerは軌道安定性、debateは反証の可視化、reflectionは逐次的自己修正を担う。実運用系はこれらの合成で構成される。}
\end{center}

\begin{insight}{示唆: 性能を左右するのはモデル規模より「調整プロトコル」の設計か}
12のマルチエージェントシステムと2つの単一エージェント基準を横断分析した評価研究は、アーキテクチャ・調整機構・メモリ・ツール統合の\textbf{4次元タクソノミー}を提示し、\textbf{調整プリマシー仮説（Coordination Primacy Hypothesis）}――「エージェント間の調整プロトコル設計は、モデル規模の拡大以上に取引意思決定の質を左右しうる」を反証可能な仮説として定式化した。同研究は、報告リターンの符号すら反転させうる5つの評価上の落とし穴（ルックアヘッド・生存者バイアス・バックテスト過学習・取引コスト無視・レジーム転換への盲目）を列挙し、調整がコスト差し引き後に価値を生むかを検証する\textbf{Coordination Breakeven Spread（CBS）}指標を導入した。この主張を検証するための評価基盤自体がまだ整っていない点が、本分野の根本的なギャップである（\cite{x04coordinationprimacy2026}）。
\end{insight}

\subsection{ユースケース別の実装地図と証拠水準}

投資エージェントを「調査を助けるAI」から「資本配分を動かすAI」へ近づけるほど、求められる証拠水準は非線形に上がる。2026年のagentic financeサーベイは、金融エージェントをデータ知覚・推論エンジン・戦略生成・統制付き執行の4層として整理し、当面の均衡は完全自律ではなく、人間の意思決定過程に埋め込まれた\textbf{bounded autonomy}であると見る（\cite{aiagentsfinmarkets2026}）。別の包括サーベイも、agentic AIの差分を目標志向の自律性・継続学習・マルチエージェント協調に置きつつ、同じ能力が市場安定性・規制遵守・解釈可能性・システミックリスクを新たな制約に変えると整理する（\cite{x04agenticaisurvey2026}）。したがって実務導入では、ユースケースごとに「どの役割分担が必要か」「どこで人間が止めるか」「どの評価を越えたら次段階へ進めるか」を明示した証拠ラダーが要る。

\begingroup
\footnotesize
\setlength{\tabcolsep}{4pt}
\begin{longtable}{@{}>{\bfseries}p{2.5cm}p{3.0cm}p{3.1cm}p{3.5cm}p{2.3cm}@{}}
\toprule
ユースケース & 推奨オーケストレーション & 最低限の評価プロトコル & 主な失敗モード & 自律性の上限\\
\midrule
\endfirsthead
\toprule
ユースケース & 推奨オーケストレーション & 最低限の評価プロトコル & 主な失敗モード & 自律性の上限\\
\midrule
\endhead
投資アイデア創出 & 検索/RAGワーカー＋強気・弱気ディベート & 観測時点で利用可能な文書だけを使う時点固定評価、アナリストによるblind ranking & よく知られた大型株ナラティブの再生、反証不足 & 推奨のみ\\
\addlinespace
デューデリジェンス & supervisor-worker＋抽出検証ワーカー＋欠損フラグ & BankerToolBench型の成果物ルーブリック、数式・出典・仮定の監査 & 値のハードコード、欠損データの幻覚、複数資料間の不整合 & 人間承認必須\\
\addlinespace
ポートフォリオ監視 & 永続メモリ＋reflection＋異常検知ワーカー & 月次/四半期の縦断評価、thesis driftとイベント反応時間の測定 & 古い信念の温存、短期ノイズへの過剰反応 & アラート生成\\
\addlinespace
リスク分析 & リスク専門エージェント＋ストレスシナリオ生成＋反事実レビュー & 相関・流動性・取引コストを含むwalk-forward評価、PortBench/CLQT型診断 & 相関構造の誤読、レジーム転換への盲目、コスト過小評価 & 配分案まで\\
\addlinespace
取引・執行 & constrained executionモジュール＋kill switch＋監査ログ & 学習カットオフ後のforward test、paper trading、制約違反率の監査 & ルックアヘッド、プロンプトインジェクション、流動性枯渇時の連鎖誤発注 & 小口・制約付き執行\\
\bottomrule
\end{longtable}
\endgroup

\vspace{-1mm}
{\small\color{figgray}\textbf{表: 投資エージェントのユースケース別実装地図。}\
自律性を上げるほど、単なる正答率ではなく、時点整合性・監査可能な推論軌跡・制約違反率・取引コスト控除後の成果を同時に満たす必要がある。}

\begin{insight}{示唆: 「研究補助」から「制約付き執行」への証拠ラダー}
本節の事例を横断すると、成功している実装は、Morgan StanleyのAskResearchGPT/Debriefのようなスコープ限定型、MarketSenseAIのような観測時点シグナルに基づく株式分析、FinRobotのような人間アナリストの推論工程を模倣するレポート生成に偏る。一方、DeepFundのforward test、FINSABERの長期・広範なバックテスト、PortBenchの相関考慮型評価、Look-Ahead-Benchの学習カットオフ検査は、完全自律的な資本配分ほど崩れやすいことを示す（\cite{deepfund2025,x04finsaber2025,portbench2026,lookaheadbench2026}）。したがって導入判断は「PoCで良いリターンが出たか」ではなく、\textbf{offline研究品質 → 時点固定バックテスト → walk-forward/forward test → paper trading → 制約付き小口執行}という段階的な証拠ラダーとして扱うべきである。77研究を監査志向でレビューしたAgentic Tradingサーベイが、取引コスト・時間整合的分割・完全再現性の欠如を構造的問題として指摘したことも、この段階設計を要求する根拠になる（\cite{agentictrading2026}）。
\end{insight}

\subsection{主要マルチエージェント・フレームワークの比較}

以下の表は、本節で扱った代表的なマルチエージェント投資フレームワークを、アーキテクチャ・主タスク・報告された結果・注意点の4軸で対比したものである。報告された結果は各文献の自己申告値であり、\S\ref{sec:investment-agents}冒頭で述べた独立ベンチマークの厳しい結果と併せて読む必要がある。

\begingroup
\footnotesize
\setlength{\tabcolsep}{4pt}
\begin{longtable}{@{}>{\bfseries}p{2.0cm}p{2.9cm}p{1.7cm}p{4.3cm}p{3.5cm}@{}}
\toprule
フレームワーク & アーキテクチャ & 主タスク & 報告された結果 & 注意点（caveat）\\
\midrule
\endfirsthead
\toprule
フレームワーク & アーキテクチャ & 主タスク & 報告された結果 & 注意点（caveat）\\
\midrule
\endhead
TradingGPT & 3層メモリ（減衰付）＋個性化＋ディベート & 株式・ファンド取引 & 人間認知に近い記憶想起を提案（初期研究） & 大規模な独立検証は限定的（\cite{x04tradinggpt2023}）\\
\addlinespace
TradingAgents & ディベート構造（強気/弱気＋専属リスクチーム） & 取引意思決定 & 単一エージェント比でSharpe比・累積リターン改善 & シミュレーション条件下の評価（\cite{tradingagents2024}）\\
\addlinespace
FinCon & マネージャー-アナリスト階層＋概念的言語強化 & 単一銘柄取引・ポートフォリオ管理 & SOTA LLM/DRLを累積リターン・Sharpeで上回り最大DD最小級 & 自己申告値。窓が短期（\cite{x04fincon2024}）\\
\addlinespace
AlphaAgents & 3マイクロエージェント＋AutoGenディベート＋リスク許容度 & 株式ポートフォリオ構築 & 役割特化で説明性・幻覚抑制を改善 & 既存ベンチとの比較で批判的評価（\cite{alphaagents2025}）\\
\addlinespace
FinRobot (ER版) & Data/Concept/Thesis の多段CoT & エクイティリサーチ・バリュエーション & IBアナリスト評価で高評価、valuation約75\% & 生成レポートの限定サンプル（\cite{x04finrobotequity2024}）\\
\addlinespace
MarketSenseAI 2.0 & Chain-of-Agents（多源データ統合） & 株式分析・銘柄選択 & S\&P100で累積リターン125.9\%（指数比+73.5pt） & 実運用でも短期窓・特定期間（\cite{marketsenseai2025}）\\
\addlinespace
HedgeAgents & 中央FM＋エージェント間カンファレンス & ヘッジング取引 & 3年シミュレーションで年率70\%・累積400\% & シミュレーション報告値（\cite{hedgeagents2025}）\\
\addlinespace
TradingGroup & 5エージェント＋自己内省＋データ合成 & 取引意思決定 & 5データセットで既存LLM/RL/ML戦略を上回る & 定量指標は要本文確認（\cite{x04tradinggroup2025}）\\
\addlinespace
SHARP & ニューロシンボリック（人間可読ルーブリック） & 取引方策の自己進化 & 小型モデルを平均10〜20ポイント改善、walk-forward検証 & ルール束縛が上振れを制約しうる（\cite{x04sharp2026}）\\
\bottomrule
\end{longtable}
\endgroup

\vspace{-1mm}
{\small\color{figgray}\textbf{表: 代表的マルチエージェント投資フレームワークの比較。}\
「報告された結果」は各文献の自己申告であり、後述の独立ベンチマークが示す性能ギャップとの併読を要する。}

\subsection{デューデリジェンスの自動化}

\begin{casestudy}{BankerToolBench: プロ水準には遠いジュニアバンカー業務の自動化}
502名の投資銀行員と共同開発したオープンソースベンチマーク。財務モデル・ピッチデック・M\&A分析等の現実的なジュニアバンカー業務でAIエージェントを評価した結果、最高性能モデルでも基準の\textbf{約半数}を満たせず、クライアント提出可能な出力は皆無だった。（出典: \cite{bankertoolbench2026}）
\end{casestudy}

\begin{itemize}
  \item \textbf{VCデューデリジェンス向けマルチエージェント・オーケストレーション}: LLMとリアルタイムWeb検索、レイアウト認識OCR（公式企業登記文書向け）を組み合わせたイベント駆動型マルチエージェントシステム。欠損データを幻覚生成せず明示的にフラグする設計を採用し、手動デューデリジェンスの代替可能性を実証した。（\cite{vcduediligence2026}）
  \item \textbf{VCBench}: 創業者経歴・製品説明・求人情報等の非構造化テキストからLLMが投資シグナルを抽出。市場ベースラインの\textbf{6倍の精度}を達成し、スタートアップのカテゴリ埋め込み単独で特徴量重み\textbf{15.6\%}（どの財務指標よりも高い）を占めることを確認した。（\cite{vcbench2025}）
  \item \textbf{MBABench / WorkstreamBench}: エンドツーエンドのスプレッドシート構築タスク評価。エージェントは複雑度の高いマルチシート財務モデルで大きく性能劣化し、透明性のある数式ではなく値をハードコードする傾向があり、DCF・LBO・M\&Aモデリングで求められるプロ水準の出力基準を一貫して満たせないことが確認された。（\cite{workstreambench2026}）
  \item \textbf{IPO Finance Agent（SpaceX/SPCX IPO事例）}: SpaceXのS-1申告書をベースにした1,000問のIPOデューデリジェンス・ベンチマーク。主要LLM金融アナリストエージェントの精度は約\textbf{80\%}が上限であり、AI支援IPO分析の具体的なギャップを明らかにしている。（\cite{ipofinanceagent2026}）
\end{itemize}

\begin{center}
\begin{tikzpicture}[font=\small,
  stage/.style={fnode blue, align=left, text width=13.5cm, minimum height=1.1cm},
  gapstage/.style={fnode accent, align=left, text width=13.5cm, minimum height=1.1cm},
  arr/.style={farr}
]
\node[stage] (s1) at (0,0)
  {\textbf{① 資料収集・OCR}\ ―\ 企業登記文書等の非構造化資料を取り込み、欠損データは幻覚生成せず明示的にフラグ};
\node[stage] (s2) at (0,-1.7)
  {\textbf{② 投資シグナル抽出}\ ―\ 創業者経歴・製品説明・求人情報から、VCBenchは市場ベースラインの\textbf{6倍}の精度を達成};
\node[stage] (s3) at (0,-3.4)
  {\textbf{③ 財務モデル構築}\ ―\ MBABench/WorkstreamBench: DCF・LBO・M\&Aモデリング。複雑なマルチシートモデルで性能が大きく劣化};
\node[gapstage] (s4) at (0,-5.3)
  {\textbf{④ 品質基準チェック}\ ―\ BankerToolBench: 最高性能モデルでも基準の約\textbf{半数}止まり／IPO Finance Agent: 精度上限は約\textbf{80\%}};

\draw[arr] (s1) -- (s2);
\draw[arr] (s2) -- (s3);
\draw[arr] (s3) -- (s4);
\end{tikzpicture}

\vspace{1mm}
{\small\color{figgray}\textbf{図: デューデリジェンス自動化パイプラインの4段階}\
資料収集からシグナル抽出、財務モデル構築を経て品質基準チェックに至る過程で、各段階の代表的ベンチマークが到達した性能水準を示す。いずれの段階にも人間による最終レビューを要するギャップが残る。}
\end{center}

\subsection{ポートフォリオ管理・執行エージェントの評価と限界}

\begin{casestudy}{DeepFund: フロンティアモデルでも純損失}
リアルタイム市場データ（学習カットオフ後）でのマルチエージェント・フォワードテストアリーナ。DeepSeek-V3やClaude-3.7-Sonnet等フロンティアモデルでも\textbf{純損失}を記録し、LLM投資への過度な期待に一石を投じた。（出典: \cite{deepfund2025}）
\end{casestudy}

\begin{casestudy}{PortBench: 相関構造を読み誤るポートフォリオ管理エージェント}
6資産クラス・10年間をカバーする、相関を考慮したLLM駆動型ポートフォリオ管理のフルパイプラインベンチマーク。LLMモデル×リスクプロファイルの組み合わせの\textbf{90\%が単純な均等加重ベースラインを上回れず}、モデルは一貫して相関構造を誤読し、歴史的ストレスシナリオ下で壊滅的なドローダウンを被った。（出典: \cite{portbench2026}）
\end{casestudy}

\begin{casestudy}{KTD-Fin: ティッカー名を匿名化すると消える「アルファ」}
ティッカー名・日付・価格系列をプロンプトとツール呼び出し全体で一貫して匿名化し、訓練データからの記憶済み市場ナラティブを排除するベンチマーク。10種のフロンティアLLMをCSI300（2024〜2026年）で検証した結果、ポートフォリオリターンの大半はパッシブな市場・スタイルエクスポージャーで説明でき、持続的な銘柄選択アルファの証拠は限定的だった。「多くの報告されているLLMトレーディングリターンは、ティッカー名・日付・価格系列を匿名化するとパッシブな市場ファクターエクスポージャーに収束する」ことを裏付け、現行ベンチマークが実際のスキルを体系的に過大評価している可能性を示す、本節の中核テーマ（ルックアヘッド／記憶バイアス）を最も直接的に補強する結果である。（出典: \cite{fromknowingtodoing2026}）
\end{casestudy}

\begin{casestudy}{FINSABER: 評価窓を広げるとLLMの優位が消える}
多源データモジュール・モジュラー戦略ベース・バイアス考慮型2段階バックテストパイプラインの3モジュールからなる評価フレームワーク。2000〜2024年の\textbf{20年超・100銘柄超}という広い横断面・長期窓で検証すると、既報のLLM戦略の優位が大きく減衰した。とりわけLLM戦略は\textbf{強気相場では過度に保守的}でパッシブ基準を下回り、\textbf{弱気相場では過度に攻撃的}で大きな損失を被る傾向が確認され、狭い期間・限られた銘柄集合での評価が生存者バイアス・データスヌーピングにより有効性を過大評価していたことを示した。（出典: \cite{x04finsaber2025}）
\end{casestudy}

\begin{itemize}
  \item \textbf{CLQT}: 5段階の取引サイクル（情報収集・統合・配分・実行・振り返り）でLLMベースのポートフォリオ管理エージェントを監査する診断ベンチマーク。コスト考慮型・戦略一貫性スコアリングとハッシュチェーン化された意思決定記録により、リターンによるランキングだけでなくAIエージェントの失敗箇所を特定できる（\cite{clqt2026}）。
  \item \textbf{Belief at Risk / FinTrace}: 前者はLLMを意味的観測システムとみなしベイズフィルタリングで監査可能な事後信念を維持し金融意思決定のモデルリスクを規制上防御可能な形で定量化する数理フレームワーク（\cite{beliefatrisk2026}）。後者は長期horizonタスクにおけるツール呼び出しの軌道（trajectory）レベル評価で、最終回答の正答率だけでは逐次推論の破綻を捕捉できないことを実証し、実運用ワークフローの失敗がツール呼び出し系列の複合エラーに起因するという仮説の根拠となる（\cite{fintrace2026}）。
  \item \textbf{Absorbing Complexity}: 目標・リスク選好・市場前提などのコンテキスト管理負担をユーザーからシステム側へ移す「interaction-native」ナレッジハーネス。構造化・非構造化データを跨いだ状態保持を可能にする設計（\cite{absorbingcomplexity2026}）。
\end{itemize}

\begin{insight}{仮説: 永続的な多層メモリが長期horizonで複利的優位を生む}
エピソード記憶（過去の分析セッションの想起）・意味記憶（thesis関連事実の蓄積）・時間グラフ記憶を備えたLLM投資リサーチエージェントは、長期horizonで複利的に優位を積み上げうるという仮説である。Mem0やGraphiti／Zepといったメモリ拡張アーキテクチャは、汎用のエージェントベンチマークでフルコンテキスト基準に対し26〜94\%の性能改善と、レイテンシ・トークンコストの約90\%削減を示している。一方、TradingAgentsやMarketSenseAIなど公表されているLLM投資エージェントの評価は数日〜数週間の短い窓に限られ、そこではコンテキストウィンドウ型エージェントと計算的にほぼ等価になる。株式thesisの構築は数十回の決算サイクル・経営戦略の転換・マクロ局面の変化にまたがるため、過去の分析ステップを想起する価値は短窓ベンチマークでは見えない形で複利的に効くはずである。検証案として、永続メモリ搭載型と毎セッション新規コンテキスト型の2変種が同一の中型株20銘柄を追跡し、12回の月次チェックポイントで買い手側アナリストがブラインド評価する6か月の縦断ベンチマークが考えられる。ただしメモリは陳腐化・矛盾した過去信念による品質劣化に直面するため、Graphitiが双時間グラフエッジで行うような時間的なメモリ無効化を明示的に実装しなければ、蓄積された記憶がかえって後続分析を劣化させうる点が交絡・リスクとなる（出典: \cite{mem0_2025}）。
\end{insight}

\begin{center}
\begin{tikzpicture}[font=\small,
  stg/.style={fnode, text width=5.4cm, minimum height=1.1cm},
  gate/.style={fnode accent, text width=5.4cm, minimum height=1.1cm, font=\bfseries\small},
  note/.style={font=\scriptsize, text=figgray, align=left, text width=6.6cm},
  arr/.style={farr blue}
]
\node[stg] (info) at (0,0) {① 情報収集};
\node[stg] (integ) at (0,-1.4) {② 統合};
\node[stg] (alloc) at (0,-2.8) {③ 配分};
\node[gate] (gateN) at (0,-4.2) {人間による承認ゲート\\（Human-in-the-Loop）};
\node[stg] (exec) at (0,-5.9) {④ 執行};
\node[stg] (review) at (0,-7.3) {⑤ 振り返り};

\draw[arr] (info) -- (integ);
\draw[arr] (integ) -- (alloc);
\draw[arr] (alloc) -- (gateN);
\draw[arr] (gateN) -- (exec);
\draw[arr] (exec) -- (review);
\draw[arr] (review.west) .. controls +(-2.2,0) and +(-2.2,-0.5) .. (info.west);

\node[note, anchor=west] at (3.2,-2.8)
  {PortBench: モデル×リスクプロファイルの\textbf{90\%}が均等加重ベースラインを上回れず};
\node[note, anchor=west] at (3.2,-5.9)
  {DeepFund: フロンティアモデルでも\textbf{純損失}／KTD-Fin: ティッカー匿名化でアルファが消失};
\end{tikzpicture}

\vspace{1mm}
{\small\color{figgray}\textbf{図: ポートフォリオ管理エージェントの5段階サイクルと承認ゲート（CLQTに基づく）}\
情報収集・統合・配分を経て執行に移る前に人間による承認ゲートを置く設計。右側の注記は、こうしたゲートを欠く／不十分な場合に実運用・ベンチマークで観測された失敗モード。}
\end{center}

\subsection{ガバナンスと自律性の設計}

\begin{casestudy}{Morgan Stanley: スコープを絞った人間協調型エージェントの大規模実運用}
Morgan StanleyはOpenAIとの協業で2つの生成AIを本番投入した。\textbf{AskResearchGPT}（2024年10月）は、投資銀行・セールス\&トレーディング・リサーチ部門向けに、同社リサーチのレポート群を横断合成し出典付きで複雑な問いに答えるアシスタント。\textbf{AI @ Morgan Stanley Debrief}（2024年6月〜）は、約1.5〜1.6万人のウェルスアドバイザー向けに顧客ミーティングを自動で文字起こし・要約しCRMへのフォローアップ草案まで生成する。Debriefのアドバイザーチーム採用率は\textbf{98\%}超と報告されている。完全自律の投資エージェントが独立ベンチマークで苦戦する一方、\textbf{タスクを厳密にスコープし人間をループに残した}エージェントは広範な実運用採用に到達するという、本節を貫く対比を裏づける事例である。（出典: \cite{x04askresearchgpt2024}, \cite{x04morganstanleyai_nd}）
\end{casestudy}

\begin{casestudy}{Adaptive AI Delegation: 動的な権限委任のガバナンスモデル}
リアルタイムの証拠品質と不確実性に基づき、AIと人間のポートフォリオマネージャー間で逐次的意思決定権限をどれだけ委任するかを動的調整するGovernance-Aware POMDPフレームワーク。5種類の静的ガバナンス代替案を上回る成果を示し、アカウンタブルなAI-Native投資ワークフロー設計の形式的モデルを提供する。（出典: \cite{adaptiveai2026}）
\end{casestudy}

\begin{casestudy}{BondGPT: 完全エージェント化された債券トレーディングと未解決のガードレール課題}
LTXのBondGPTは2026年6月に完全エージェント型AIをローンチし、AIエージェントがリアルタイム市場を自律的に監視し、取引チケットを生成し、ディーラーを選定し、RFQを起動し、価格を受諾できるようになった。ガードレールは現状ポリシー駆動の取引サイズ制限に留まるが、固定収益市場は2020年3月の「dash-for-cash」や2023年3月SVB取り付け騒ぎのように、エピソード的に流動性が枯渇する局面が存在し、取引サイズ閾値だけでは流動性レジームの変化に対応できない可能性が指摘されている。（出典: \cite{prnewswirendbondgpt}）
\end{casestudy}

\begin{casestudy}{Agent-to-Agent Finance: 自律AIエージェント同士の機械媒介型金融レイヤー}
自律AIエージェントが取引相手を発見し、互いに決済し合う「エージェント間金融（agent-to-agent finance）」を、ブロックチェーンベースのアイデンティティ・プログラマブルな決済・境界づけられた自律性（bounded-autonomy）のガバナンス基盤の上に成り立つ機械媒介型の金融レイヤーとして定義する。エージェントが人間の逐次的関与なしに相互取引を行う際に必要な、識別・信頼・決済・統制のインフラ要件を整理しており、自律性を高めたエージェント経済に向けた金融インフラの設計論を提示する。（出典: \cite{agenttoagentfinance2026}）
\end{casestudy}

\begin{insight}{仮説: プロンプトインジェクションという新たな攻撃面}
市場ニュースフィード・クライアントメール・サードパーティのリサーチPDFなど外部の非構造化データを取り込むLLMエージェントワークフローに対し、OWASPがLLM01:2025としてランク付けするプロンプトインジェクションは、敵対的テストデータセットで50〜84\%の成功率を記録している。ニュース取り込み・分析・執行にわたるツール呼び出しを連鎖させるマルチエージェント金融ワークフローは、単一の汚染入力がパイプライン全体に伝播しうる複合的なインジェクションベクトルを生む。金融AIエージェント専用の成熟した緩和ツールキットは存在せず、拘束力のある技術基準を発行した金融規制当局もまだない。
\end{insight}

\begin{center}
\begin{tikzpicture}[font=\small,
  lvl/.style={fnode, align=left, text width=10.4cm, minimum height=1.1cm},
  lvltop/.style={fnode accent, align=left, text width=10.4cm, minimum height=1.1cm},
  arr/.style={farr}
]
\node[lvl] (l1) at (0,0)
  {\textbf{Lv.1 人間が意思決定}\ ―\ AIは情報提供のみで決定権は人間に残る};
\node[lvl] (l2) at (0,1.7)
  {\textbf{Lv.2 AIが推奨・人間が承認}\ ―\ 典型的なHuman-in-the-Loop構成};
\node[lvl] (l3) at (0,3.8)
  {\textbf{Lv.3 動的権限委任}（Adaptive AI Delegation）\ ―\ 証拠品質・不確実性に応じPOMDPで逐次調整し、5種の静的ガバナンス代替案を上回る};
\node[lvltop] (l4) at (0,6.3)
  {\textbf{Lv.4 完全自律執行}（BondGPT）\ ―\ 市場監視・取引チケット生成・ディーラー選定・RFQ起動・価格受諾までを自律実行。ガードレールは取引サイズ制限のみ};

\draw[arr] (l1) -- (l2);
\draw[arr] (l2) -- (l3);
\draw[arr] (l3) -- (l4);

\node[anchor=west, font=\scriptsize\itshape, text width=3.4cm, text=figaccent, align=left] at (5.6,6.3)
  {リスク面: プロンプトインジェクション成功率\textbf{50\textendash84\%}（OWASP LLM01:2025）};
\end{tikzpicture}

\vspace{1mm}
{\small\color{figgray}\textbf{図: 自律性のガバナンス・スペクトラム}\
人間が決定するLv.1から、BondGPTのように取引執行まで自律化されたLv.4まで、委任範囲が段階的に拡大する。自律性が高まるほど、プロンプトインジェクション等の攻撃面へのリスクも拡大する。}
\end{center}

\subsection{ベンチマークが示す一貫した性能ギャップ}

\begin{itemize}
  \item \textbf{BigFinanceBench}: 928件の専門家作成・自由回答形式の財務リサーチタスクをステップバイステップのルーブリック（36,241ポイント）で採点。最先端エージェントでもルーブリックスコア\textbf{58.8\%}に留まり、最終回答の正確性だけでなく監査可能な導出過程の質が未解決のボトルネックであることを示す。（\cite{bigfinancebench2026}）
  \item \textbf{Deep FinResearch Bench}: AI生成の投資リサーチレポートをプロアナリストのレポートと定性的厳密さ・定量予測・claim信頼性の3軸で比較し、AIは全軸でプロ水準に及ばないことを確認した。（\cite{deepfinresearchbench2026}）
  \item \textbf{FinMCP-Bench}: Model Context Protocol（MCP）下での金融ツール利用を評価する初のベンチマーク。65の実在する財務MCPツールサーバーを用い、株価データ・リスク分析・ポートフォリオ操作など33の金融シナリオにまたがる613タスクでLLMエージェントを評価し、マルチツール推論と長期タスク遂行に大きなギャップがあることを最先端モデル全般で確認した（ICASSP 2026）。（\cite{finmcpbench2026}）
  \item \textbf{Hedge-Bench}: 実際のヘッジファンド・アナリスト業務から抽出した102件の自由回答形式タスクを、モデル判定ではなく検証済みの専門家の推論ステップに照らして採点するベンチマーク。フロンティアLLMでも平均スコアは\textbf{16\%未満}に留まり、機械的な金融ツール利用と真の投資推論との間に大きなギャップがあることを露呈した（\cite{hedgebench2026}）。
  \item \textbf{LLMトレーディングエージェント・サーベイ}: LLM取引エージェントの設計空間（メモリ・ツール利用・内省・マルチエージェント協調）を体系化した初期の包括的サーベイ（\cite{x04finagentsurvey2024}）。より監査志向の後続サーベイ（\S\ref{sec:exec-summary}で言及した77研究のレビュー、\cite{agentictrading2026}）は、閉ループ評価を満たす19研究のうち取引コストを明示モデル化したのは\textbf{1件}、時間整合的な分割プロトコルを抽出できたのは\textbf{2件}、R3（完全再現性）に達した研究は\textbf{皆無}であったと報告し、標準化された報告様式の欠如を本分野の構造的問題と位置づけた。
  \item \textbf{FinCoT}: 専門家の推論手順を\textbf{Mermaid図で符号化した「エキスパート・ブループリント」}をプロンプトに注入する構造化Chain-of-Thought。Qwen3-8B-Baseの精度を\textbf{63.2\%→80.5\%}、Fin-R1（7B）を\textbf{65.7\%→75.7\%}へ引き上げつつ出力長を最大\textbf{8.9倍}短縮した。金融ポストトレーニングを欠くモデルほど効果が大きく、より解釈可能で専門家整合的な推論トレースを生む――エージェントの推論ステップの監査可能性に直結する結果である（\cite{x04fincot2025}）。
  \item \textbf{SHARP}: 自由形式のプロンプト変異を、\textbf{人間可読な条件-行動ルーブリック}の記号的方策最適化に置き換えたニューロシンボリック・フレームワーク。取引が劣化した際にどのルールが失敗したかを特定して的を絞った編集を行い、厳格なwalk-forward検証で方策ドリフトを防ぐ。3セクター・4種のLLMバックボーンで一貫した改善を示し、小型モデル（例: GPT-4o-mini）を平均\textbf{10〜20ポイント}引き上げた。機関投資家が要求する透明性・監査可能性と、方策の自己進化を両立させる設計である（\cite{x04sharp2026}）。
\end{itemize}

\begin{casestudy}{Look-Ahead-Bench: 学習カットオフを跨いだ瞬間に消えるアルファ}
LLMベースの投資バックテストにおけるルックアヘッドバイアスを標準化して測定するベンチマーク。DeepSeek 3.2は、評価期間をモデルの学習カットオフ以降に移すと、年率アルファが\textbf{+20.73\%から-1.04\%}へとほぼ完全に消失することを示し、公表されているLLMアルファの多くが人為的成果物（記憶されたパターンの再生）である可能性を示唆した。（出典: \cite{lookaheadbench2026}）
\end{casestudy}

以下の図は、本節で取り上げた実運用・バックテストの成功事例と、独立ベンチマークが計測した性能ギャップを対比したものである。

\begin{center}
\begin{tikzpicture}[font=\small]
  % ---- column header boxes ----
  \node[fchip, fill=figblue, minimum width=6.4cm, minimum height=0.9cm]
        (Lhead) at (0,0) {実運用・強力なバックテスト成果};
  \node[fchip, fill=figaccent, minimum width=6.4cm, minimum height=0.9cm]
        (Rhead) at (8.0,0) {独立ベンチマークが示す性能ギャップ};

  % ---- left column: production / backtest successes ----
  \node[fnode blue, minimum width=6.4cm, minimum height=2.6cm] (L1) at (0,-2.0) {};
  \node[anchor=north west, text width=5.8cm, align=left]
        at ($(L1.north west)+(0.3,-0.3)$)
        {\textbf{Bridgewater AIA Labs}\\ 生成AI基盤を本番の投資意思決定支援に実装};
  \node[fnode blue, minimum width=6.4cm, minimum height=2.6cm] (L2) at (0,-4.75) {};
  \node[anchor=north west, text width=5.8cm, align=left]
        at ($(L2.north west)+(0.3,-0.3)$)
        {\textbf{MarketSenseAI 2.0}\\ S\&P100バックテストで累積リターン\textbf{\textcolor{figblue}{125.9\%}}（指数比+73.5pt）};
  \node[fnode blue, minimum width=6.4cm, minimum height=2.6cm] (L3) at (0,-7.5) {};
  \node[anchor=north west, text width=5.8cm, align=left]
        at ($(L3.north west)+(0.3,-0.3)$)
        {\textbf{TradingAgents}\\ ディベート構造で単一エージェント比のSharpe比・累積リターンが改善};

  % ---- right column: benchmark performance gaps (bar = achieved / 100%) ----
  \node[fnode accent, minimum width=6.4cm, minimum height=2.6cm] (R1) at (8.0,-2.0) {};
  \node[anchor=north west, text width=5.8cm, align=left]
        at ($(R1.north west)+(0.3,-0.3)$)
        {\textbf{BankerToolBench}\\ ジュニアバンカー業務基準の達成率};
  \node[fnode accent, minimum width=6.4cm, minimum height=2.6cm] (R2) at (8.0,-4.75) {};
  \node[anchor=north west, text width=5.8cm, align=left]
        at ($(R2.north west)+(0.3,-0.3)$)
        {\textbf{PortBench}\\ 均等加重ベースラインを上回れた組合せ比率};
  \node[fnode accent, minimum width=6.4cm, minimum height=2.6cm] (R3) at (8.0,-7.5) {};
  \node[anchor=north west, text width=5.8cm, align=left]
        at ($(R3.north west)+(0.3,-0.3)$)
        {\textbf{BigFinanceBench}\\ 財務リサーチ・ルーブリックスコア};

  % ---- bars (background = 100% reference, filled = achieved share) ----
  \fill[figlight] ($(R1.south west)+(0.3,0.25)$) rectangle ($(R1.south west)+(4.9,0.5)$);
  \fill[figaccent] ($(R1.south west)+(0.3,0.25)$) rectangle ($(R1.south west)+(2.6,0.5)$);
  \node[font=\bfseries\small, anchor=west] at ($(R1.south west)+(5.05,0.375)$) {50\%};

  \fill[figlight] ($(R2.south west)+(0.3,0.25)$) rectangle ($(R2.south west)+(4.9,0.5)$);
  \fill[figaccent] ($(R2.south west)+(0.3,0.25)$) rectangle ($(R2.south west)+(0.76,0.5)$);
  \node[font=\bfseries\small, anchor=west] at ($(R2.south west)+(5.05,0.375)$) {10\%};

  \fill[figlight] ($(R3.south west)+(0.3,0.25)$) rectangle ($(R3.south west)+(4.9,0.5)$);
  \fill[figaccent] ($(R3.south west)+(0.3,0.25)$) rectangle ($(R3.south west)+(3.0,0.5)$);
  \node[font=\bfseries\small, anchor=west] at ($(R3.south west)+(5.05,0.375)$) {58.8\%};

  % ---- bottom banner: Look-Ahead-Bench collapse spans both columns ----
  \node[fnode dark, minimum width=14.4cm, minimum height=1.5cm, text width=13.4cm] (Banner) at (4.0,-9.75)
        {\textbf{Look-Ahead-Bench}: 評価期間を学習カットオフ以降に移すと、DeepSeek 3.2の年率アルファは
         \textbf{\textcolor{figaccent!70!white}{+20.73\%}} $\rightarrow$ \textbf{\textcolor{figaccent!70!white}{-1.04\%}} へとほぼ完全に消失};
\end{tikzpicture}

\vspace{1mm}
{\small\color{figgray}\textbf{図: 実運用・バックテストの成功事例と独立ベンチマークが計測した性能ギャップの二極化。}
左列は本節で扱った具体的な成功事例、右列は独立ベンチマークが計測した達成率（バーは0〜100\%を表す）。下段はLook-Ahead-Benchが示す、学習カットオフを跨ぐだけでアルファが消失するという最も直接的な証拠。}
\end{center}

\begin{insight}{示唆: 実運用に届くAIエージェントとベンチマーク上の弱いAIエージェントの共存}
Bridgewater AIA LabsやMarketSenseAI 2.0のような強力な実運用・バックテスト成果、そしてMorgan StanleyのAskResearchGPT／Debrief（アドバイザー採用率98\%超）やAnthropicが提供する金融向けエージェント・テンプレート群（\cite{anthropicfinanceagents2026}）のような大規模な本番採用が存在する一方で、BankerToolBench、DeepFund、PortBench、FINSABER、Look-Ahead-Benchのような独立ベンチマークは一貫して大きな性能ギャップと方法論的な脆弱性を報告している。この二極化は、成功事例の多くが「厳密なタスクスコープ設定」「人間によるレビュー工程の保持」「透明な監査記録」を備えている一方、汎用的・完全自律的なエージェントほど性能が崩れやすいという構造を示唆している。SHARPの人間可読ルーブリックやFinCoTのエキスパート・ブループリントが示すように、\textbf{推論と方策を監査可能な記号的表現へ外部化する}方向性――そして調整プリマシー仮説が説くように調整プロトコルそのものをコスト差し引き後で評価する規律（CBS指標、\cite{x04coordinationprimacy2026}）――が、この二極を橋渡しし、自律性のガバナンス・スペクトラムを上方へ押し上げる現実的な経路になると考えられる。
\end{insight}

\subsection{ソフトウェア工学エージェントに学ぶ：評価と自律度の段階付け}
\label{sec:investment-agents-crossdomain}

投資エージェントが直面する「評価の厳密さ」「オーケストレーション設計」「自律度の段階付け」という3つの課題は、金融に固有のものではない。むしろ、ソフトウェア工学エージェント（コーディングエージェント）、深い調査（deep research）・科学的発見エージェント、そして本番運用されているカスタマーサポート・エージェントは、いずれも金融より先に大規模な実運用と厳密なベンチマークの両方を経験しており、投資エージェントの設計・評価に転用できる普遍的な教訓を蓄積している。本項では、これら周辺分野の代表事例を金融の視点から読み解き、投資エージェントの実務に持ち帰れる示唆を抽出する。

\begin{casestudy}{SWE-bench／SWE-bench Verified：自己申告ではなく「閉じたテスト」で合否が決まる評価基盤}
コーディングエージェント分野は、2023年のSWE-bench（12のPython OSSリポジトリから抽出した2,294件の実GitHub issue、\cite{y04swebench2023}）によって、「エージェントが生成したコード変更が、実際に隠しテストスイートを通過するか」という自動判定可能な客観指標を確立した。2024年8月にはOpenAIとの協業で、人間のアノテーターが問題記述の明確さ・テストパッチの正確さ・解決可能性を検証した500件の人手検証済みサブセット\textbf{SWE-bench Verified}が公開され（\cite{y04swebenchverified2024}）、あいまいな問題設定に起因するノイズを取り除いた。Cognition社のDevinは2024年3月、このSWE-bench上で570件中79件を自己完結的に解決（成功率\textbf{13.86\%}、当時の最良のアシスト型システムであるClaude 2の4.80\%を大きく上回る）と発表し（\cite{y04devinintro2024}）、自律型ソフトウェアエンジニアという分野を切り開いた。2025年公開のオープンソース・スカフォールドOpenHands Software Agent SDK（MIT License、40,000超のGitHubスター）は、Claude Sonnet 4.5＋拡張思考でSWE-bench Verifiedの解決率\textbf{72\%}に到達し（\cite{y04openhandssdk2025}）、オープンな実装でもクローズドな商用エージェントに迫れることを示した。金融のBankerToolBench・WorkstreamBenchが自由回答形式のルーブリック採点に留まるのに対し、SWE-bench Verifiedは「人間が事前に解決可能性を検証したタスク」に対して「自動テストが機械的に合否を判定する」という二重の厳密さを備えている点が対照的である。
\end{casestudy}

\begin{casestudy}{AI co-scientist：生成・ディベート・進化のループで仮説を淘汰する科学的発見エージェント}
Google Research・Cloud AI・DeepMindが共同開発したAI co-scientist（Gemini 2.0基盤、\cite{y04aicoscientist2025}）は、科学者が示す研究目標に沿って検証可能な新規仮説を生成するマルチエージェントシステムである。中核設計は\textbf{「生成（generate）→ディベート（debate）→進化（evolve）」}のループで構成され、自己対戦型の科学的ディベートによる新規仮説生成、トーナメント方式のランキングによる仮説比較、反復的な品質改善の「進化」プロセスを非同期タスク実行フレームワーク上で回す。急性骨髄性白血病の薬剤再利用、肝線維症の新規標的発見、抗菌薬耐性メカニズムの説明という3つの生物医学応用で初期検証が行われ、Elo指標による自動評価で他の最先端エージェント・推論モデルを上回ったと報告されている。この「複数エージェントによるディベート＋トーナメント式の淘汰」という設計は、本節前半で扱ったTradingAgentsの強気・弱気ディベート構造と同型であり、投資thesisの生成において「どのように複数の候補案を競わせ淘汰するか」という設計問題が、金融固有ではなく多エージェント発見システム一般の課題であることを示している。一方、deep researchエージェントの評価でも、ResearchRubrics（2,800時間超の人手労働・2,500件超の専門家執筆ルーブリック、\cite{y04researchrubrics2025}）のように、事実の裏付け・推論の健全性・明瞭さを軸とする採点方式が標準化されつつあり、本節のBigFinanceBench（36,241ポイントのステップバイステップ・ルーブリック）と同型の設計思想が金融の外側でも独立に発展している。
\end{casestudy}

\begin{casestudy}{Klarnaのカスタマーサポート・エージェント：急進的自律化とその巻き戻し}
Klarnaは2024年2月、OpenAIと共同開発したAIアシスタントを本番投入し、最初の1か月で\textbf{230万件}の会話（カスタマーサービスチャットの3分の2）を処理、フルタイム従業員\textbf{700名相当}の業務量をこなし、対応時間を平均11分から2分未満へ短縮、2024年通期で\textbf{4,000万ドル}の利益改善に貢献したと報告した（\cite{y04klarnapress2024}）。しかし2025年、CEO Sebastian Siemiatkowski氏はコストありきのAI自動化が「品質低下」を招いたと認め、人間のカスタマーサービス担当者を再雇用して「常に人間対応を選べる」方針へ転換した（\cite{y04klarnawalkback2025}）。カスタマーサポート・エージェント全般の評価でも、Sierra Research発のτ-bench（小売・航空ドメイン、\cite{y04taubench2024}）は、エージェントが対話全体を通じてドメイン固有のポリシーガイドラインに従い続けられるかを会話終了時のDB状態で機械的に判定し、当時最先端のGPT-4oでもタスク成功率は\textbf{50\%未満}に留まった。同一タスクを複数回試行した際の一貫性を測る\textbf{pass\textasciicircum k}という信頼性指標も導入されている。2025年のτ\textsuperscript{2}-bench（\cite{y04tau2bench2025}）は、ユーザー側もツール呼び出しを行うデュアルコントロール型の通信（telecom）ドメインを追加し、双方向にツールを操作する現実的シナリオへ拡張した。急速な自律化の定量的成功と、その後の巻き戻しが同一企業内で観測されたKlarnaの事例は、「ポリシー遵守」と「複数回試行での一貫性」を分離して測る評価の必要性を裏づけている。
\end{casestudy}

\begin{center}
\begin{tikzpicture}[font=\small,
  axisline/.style={line width=1pt, draw=fignavy!60},
  tick/.style={line width=1pt, draw=fignavy!60},
  chip/.style={fchip, text=fignavy, minimum width=2.6cm, minimum height=0.85cm, font=\scriptsize},
  lbl/.style={font=\scriptsize, text=figgray, align=center}
]
% axis
\draw[axisline] (0,0) -- (12,0);
\foreach \x/\name in {0/Lv.1, 4/Lv.2, 8/Lv.3, 12/Lv.4} {
  \draw[tick] (\x,0.12) -- (\x,-0.12);
  \node[lbl] at (\x,-0.5) {\name};
}
\node[lbl, align=left, text width=2.6cm] at (0,-1.05) {人間が意思決定};
\node[lbl, align=left, text width=2.6cm] at (4,-1.05) {AI推奨・人間承認};
\node[lbl, align=left, text width=2.6cm] at (8,-1.05) {動的権限委任};
\node[lbl, align=left, text width=2.6cm] at (12,-1.05) {完全自律執行};

% chips above axis at each domain's typical position
\node[chip, fill=figblue!18, draw=figblue] (swe) at (4,1.4) {SWE-bench Verified／Devin\\（合格後も人間がPRをレビュー・マージ）};
\node[chip, fill=figaccent!18, draw=figaccent] (dr) at (2.6,3.0) {AI co-scientist／\\deep research（仮説提示のみ）};
\node[chip, fill=fignavy!12, draw=fignavy] (cs2025) at (4,4.6) {Klarna 2025\\（人間対応の選択肢を復元）};
\node[chip, fill=fignavy!12, draw=fignavy] (cs2024) at (9.5,4.6) {Klarna 2024\\（自律応答が主・700人分代替）};
\draw[-{Stealth[length=2.4mm]}, line width=0.9pt, draw=figaccent!70!black] (cs2024) -- (cs2025)
      node[midway, above, font=\scriptsize\itshape, text=figaccent!70!black] {2025年に巻き戻し};
\node[chip, fill=fignavy!12, draw=fignavy] (inv) at (8,3.0) {投資エージェント（本節）\\ユースケースによりLv.1〜4に分布};

\end{tikzpicture}

\vspace{1mm}
{\small\color{figgray}\textbf{図: 自律度スペクトラム上の周辺分野マッピング。}\
横軸は\S\ref{sec:investment-agents}後段の自律性のガバナンス・スペクトラム（Lv.1〜4）と同じ尺度。SWE-bench Verified／Devinはタスク完了後も人間のレビュー・マージ工程を残し、AI co-scientist・deep researchは仮説提示に留まる。Klarnaは2024年に高い自律性へ踏み込んだのち、2025年に人間対応の選択肢を復元する形で後退した。投資エージェントは本節で見た通りユースケースごとにLv.1〜4へ分布する。}
\end{center}

以下の表は、ソフトウェア工学・深い調査／科学的発見・カスタマーサポートの3分野と、本節で扱った投資エージェントを、オーケストレーション類型・評価の厳密さ・自律度の段階付け・金融が学べる要点の4軸で対比したものである。

\begingroup
\footnotesize
\setlength{\tabcolsep}{4pt}
\begin{longtable}{@{}>{\bfseries}p{2.1cm}p{2.6cm}p{3.4cm}p{2.6cm}p{3.6cm}@{}}
\toprule
分野 & 代表的オーケストレーション & 評価の厳密さ & 自律度の段階付け & 金融が学べる要点\\
\midrule
\endfirsthead
\toprule
分野 & 代表的オーケストレーション & 評価の厳密さ & 自律度の段階付け & 金融が学べる要点\\
\midrule
\endhead
ソフトウェア工学 & plan→act→verify（Devinのダッシュボード計画提示、OpenHandsの委譲型サブエージェント） & SWE-bench Verified：人手検証済み解決可能性＋隠しテストによる自動合否判定（\cite{y04swebench2023,y04swebenchverified2024}） & Devin：完了後も人間がPRをレビュー・マージ（Lv.2相当） & 「自動テストで機械的に合否が決まる成果物」を金融DDでどこまで設計できるか\\
\addlinespace
深い調査／科学的発見 & 生成→ディベート→進化（AI co-scientist）、専門家ルーブリック採点（ResearchRubrics） & Elo指標によるトーナメント評価、2,500件超の専門家ルーブリック（\cite{y04aicoscientist2025,y04researchrubrics2025}） & 仮説提示のみ、科学者が最終判断（Lv.1〜2相当） & 投資thesis候補の「淘汰」設計＋ルーブリック採点の標準化\\
\addlinespace
カスタマーサポート & ポリシー準拠型シングルエージェント＋人間エスカレーション & τ-bench／τ\textsuperscript{2}-bench：ポリシー遵守を機械判定、pass\textasciicircum kで多重試行の一貫性を測定。GPT-4oでも成功率50\%未満（\cite{y04taubench2024,y04tau2bench2025}） & Klarna 2024はLv.3〜4相当に接近も2025年にLv.2へ後退（\cite{y04klarnapress2024,y04klarnawalkback2025}） & 急進的な自律化は評価が追いつく前に品質・レピュテーションリスクを露呈しうる\\
\addlinespace
投資エージェント（本節） & supervisor-worker／debate／reflectionの合成（\S\ref{sec:investment-agents}参照） & 自己申告バックテストが主流。独立ベンチマーク（BankerToolBench等）は依然厳しい結果 & ユースケース別にLv.1（推奨のみ）〜Lv.4（BondGPT型執行）へ分布 & 周辺分野が先行して確立した「閉じた合否判定」「トーナメント淘汰」「ポリシー遵守の機械判定」を評価基盤に取り込む余地\\
\bottomrule
\end{longtable}
\endgroup

\vspace{-1mm}
{\small\color{figgray}\textbf{表: ソフトウェア工学・深い調査／科学的発見・カスタマーサポート・投資エージェントの比較。}\
いずれの分野も「オーケストレーション」「評価の厳密さ」「自律度の段階付け」という3軸で整理でき、投資エージェントはこの3軸すべてにおいて周辺分野の到達点より後発である。}

\begin{insight}{金融への示唆: 「閉じた合否判定」を金融DDワークフローへ移植する}
SWE-bench Verifiedが示したのは、「人間が事前に解決可能性を検証したタスク」に対し「隠しテストが機械的に合否を判定する」という組み合わせが、自己申告のベンチマークより遥かに信頼できるという教訓である。金融のデューデリジェンス・エージェント評価（BankerToolBench、WorkstreamBench等）は依然として人間ルーブリック採点に依存しており、SWE-bench的な「自動テストハーネス」を持たない。財務モデルの数式・セル参照・出典の実在性など、機械的にチェック可能な要素（例: DCFモデルの特定セルが正しい財務諸表項目を参照しているか）を切り出し、SWE-bench Verified型の「人手検証済み解決可能性＋自動テストによる合否判定」を部分的にでも導入できれば、独立ベンチマークの信頼性を大きく引き上げられる可能性がある。ただし財務モデルには唯一の正解が存在しない設計選択（会計処理・前提の妥当な幅）が残るため、コードのPRのような二値的合否判定への単純化には限界がある点に注意が必要である。
\end{insight}

\begin{insight}{金融への示唆: 自律性は「主張」ではなく行動ログから測定する}
Anthropicは約40万件のClaude Codeセッション（2025年10月〜2026年4月、プライバシー保護分析）を基に、ターン継続時間（99.9パーセンタイルで2025年10月の25分未満から2026年1月に45分超へほぼ倍増）、自動承認（auto-accept）モードの利用率（新規ユーザー約20\%に対し750セッション超の熟練ユーザーで40\%超）、中断率（新規ユーザー5\%に対し熟練ユーザー約9\%）を定量測定した（\cite{y04anthropicautonomy2026}）。最も複雑なタスクでは、Claude Code自身が人間の中断頻度の2倍以上の頻度で明確化を求める自己制限も観測されている。この「自律性を主張ではなく行動ログから実証する」手法は、本節の自律性のガバナンス・スペクトラム（Lv.1〜4）を裏付ける実証手段として金融機関の投資エージェント基盤にも転用できる。同時に、Klarnaの2024年急進的自律化と2025年の巻き戻しが示すように、行動ログ上の高い自律性指標だけでは品質・レピュテーションリスクを捕捉できない。The 2025 AI Agent Index（30の実運用エージェントを1,350項目にわたり監査し198項目で情報欠落を確認、\cite{y04aiagentindex2025}）が指摘する通り、安全性評価・開示メカニズムの透明性は自主的報告だけでは不十分であり、金融規制当局が投資エージェントに構造化された開示要件を課す素地は、周辺分野の経験からも既に整いつつある。
\end{insight}

\begin{insight}{金融への示唆: 評価インフラそのものを公共財として整備する}
Holistic Agent Leaderboard（21,730回のエージェント実行を9モデル×9ベンチマーク・4領域〈コーディング・Web操作・科学・カスタマーサービス〉にわたり約4万ドルの計算コストで標準化、\cite{y04holisticagentleaderboard2025}）は、標準化された評価ハーネスを整備することで評価期間を数週間から数時間へ短縮し、「推論努力を上げるとむしろ大多数の実行で精度が低下する」「エージェントがベンチマーク自体をHuggingFace上で検索してしまう」といった評価汚染や不正確な挙動を横断的に検出した。ソフトウェア工学エージェント分野がこうした評価インフラを分野横断の公共財として整備しつつあるのに対し、金融の投資エージェント評価は個別研究ごとにバラバラのバックテスト設定・データ期間・取引コスト前提に依存しており、本節で繰り返し指摘してきた再現性の欠如（Agentic Tradingサーベイが指摘したR3達成研究\textbf{皆無}、\cite{agentictrading2026}）と表裏の関係にある。標準化されたVM並列実行基盤・共通のコスト前提・独立検証バッジ付き投稿制度を業界団体や規制当局が主導して整備することが、金融エージェント評価の再現性ギャップを埋める現実的な一歩になりうる。
\end{insight}

\section{LLM／AIエージェントによる金融データ分析}
\label{sec:llm-analysis}

本節では、AIが財務データを抽出・要約・分類・構造化・比較・推論する事例を、構造化・非構造化データの両方にわたって整理する。特にアルファ生成領域では、公表された成果の再現性・頑健性に関する重要な警鐘が複数の独立研究から挙がっている。

\subsection{テキストからの情報抽出と構造化}

決算資料や適時開示のような非構造化テキストから、意味のあるシグナルや特徴量を抽出・構造化する取り組みは、後段のアルファ生成やリスク管理の入り口として位置づけられる。

\begin{itemize}
  \item \textbf{From Text to Alpha}: コンテキストを保持した文脈対応スパン抽出により、連続する開示文書間の意味的シフト（経営陣が重視する指標の変化）を追跡。キーワードベースラインに対し\textbf{2倍のリスク調整後アルファ}を達成した。（出典: \cite{texttoalpha2025}）
  \item \textbf{Two Sigmaの社内手法・論考}: レジームモデリング（ガウス混合モデルで市場を4レジームに分類、著者Alex Botte, Doris Bao）\cite{twosigmandregimemodeling}、インフレテーマ・ポートフォリオ構築（LightGBM活用、著者Simon Lau, Alex Botte）\cite{twosigmandinflationportfolio}、そして「モデルは怠惰であり人間の判断なしでは体系的にバイアスのある結果を生みうる」との論考（著者Claudia Perlich）\cite{twosigmandintuition}という3本の社内発信は、いずれも単独では際立つ定量成果を伴わないものの、統計モデル活用と人間の監督を組み合わせる一貫した設計思想を映し出している。
  \item \textbf{Fund2Persona}: ファンドの開示文書・保有銘柄の推移・運用者コメンタリーにLLMのファイナンシャルアドバイザー・ペルソナを接地させるフレームワーク。actor（生成）・scorer（評価）・patcher（修正）のエージェント的ループで、実開示に根拠を持つAIネイティブな投資助言をスケーラブルに提供する。（出典: \cite{fund2persona2026}）
  \item \textbf{FinDocMRE（VLMによるチャート画像の直接読解）}: ローソク足・テクニカル指標オーバーレイ・収益ウォーターフォール・マルチパネルダッシュボードといった金融チャートは、空間配置・色・軸スケールに意味を埋め込むため、OCRによる線形化ではその情報が失われる。文書レベルの金融マルチモーダル推論を評価するベンチマークとして、VLMを画像へ直接適用する方式とOCR＋LLMパイプラインを比較する枠組みが提案されている。（出典: \cite{findocmre2026}）
\end{itemize}

\begin{center}
\begin{tikzpicture}[
  doc/.style={fnode, text width=6.6cm, minimum height=1.2cm, font=\small},
  base/.style={fnode muted, text width=3.8cm, minimum height=1.6cm, font=\footnotesize},
  hi/.style={fnode blue, text width=4.8cm, minimum height=1.6cm, font=\footnotesize},
  outbox/.style={fnode accent, text width=8.2cm, minimum height=1.2cm, font=\small},
  arr/.style={farr blue},
  darr/.style={farr dash},
]
\node[doc] (doc) at (5.0,3.0) {決算資料・適時開示\\（非構造化テキスト）};
\node[base] (kw) at (1.8,0.6) {① キーワード\\ベースライン\\\textcolor{figgray}{基準アルファ ×1}};
\node[hi] (ctx) at (7.6,0.6) {② コンテキスト対応スパン抽出\\(From Text to Alpha)\\\textbf{リスク調整後アルファ ×2}};
\node[outbox] (out) at (5.0,-1.9) {構造化シグナル・特徴量 → 下流のアルファ生成・リスク管理へ};
\draw[darr] (doc) -- (kw);
\draw[arr] (doc) -- (ctx);
\draw[darr] (kw) -- (out);
\draw[arr] (ctx) -- (out);
\end{tikzpicture}
\end{center}
{\small\color{brandgray}\textbf{図: 非構造化テキストからの情報抽出パイプライン}\ 決算資料・適時開示テキストにコンテキスト対応スパン抽出（From Text to Alpha）を適用すると、キーワードベースラインに対し\textbf{2倍のリスク調整後アルファ}を達成し、構造化シグナルとして下流のアルファ生成・リスク管理へ供される。}

\subsection{固有表現抽出・XBRLタグ付け・センチメント分類: 構造化タスクにおける小型モデルの底力}

抽出・分類のように出力空間が明確に定義された「構造化タスク」では、汎用フロンティアLLMのプロンプティングよりも、ドメインでファインチューンした小型モデルの方がなお高精度である場合が多い。この非直感的な事実は、LLM活用の設計判断（どこにフロンティアモデルを投入し、どこに軽量モデルを据えるか）を左右する。

\begin{itemize}
  \item \textbf{固有表現抽出（NER）}: FiNER-ORD（手作業で注釈付けした金融ニュース記事201本）を用いた評価では、ファインチューン済みのBERT／RoBERTaが加重F値\textbf{約0.88}を達成するのに対し、少数ショットの最良LLMはGemini \textbf{0.8368}、GPT-4o \textbf{0.8203}にとどまる。LLMは文脈依存の分類やドメイン固有の固有名詞（例: 「Google Maps」を組織として誤タグ付け）で系統的に誤りやすい。構造化抽出では、プロンプトされたフロンティアLLMよりも小型のファインチューンモデルが依然優位でありうることを示す。（出典: \cite{x05finerord2025}）
  \item \textbf{XBRLタグ付け（FinTagging）}: 財務情報を約\textbf{14,000概念}のUS-GAAPタクソノミへ標準化するタグ付けベンチマーク。財務事実の抽出自体はF値\textbf{約72\%}に達するものの、抽出した事実を巨大なタクソノミへ正しく紐付ける「概念リンキング」が律速要因であり、iXBRL規制報告の自動化における真のボトルネックはファクト抽出ではなく概念照合であることを示唆する。（出典: \cite{x05fintagging2025}）
  \item \textbf{センチメント分類（FinGPT）}: オープンソースのFinGPT（AI4Finance）はv3系列をニュース・ツイートのセンチメントデータにLoRAでファインチューンし、多くの金融センチメントベンチマークで低コストながら最良スコアを報告する。閉じたフロンティアモデルや計算資源を要するBloombergGPT型の対極として、安価に適応可能なオープンソース・ドメインモデルの実用性を示す。（出典: \cite{x05fingpt2023}）
\end{itemize}

\subsection{要約と情報忠実度: 圧縮が意思決定を歪めるリスク}

LLMによる要約は情報量を圧縮する裏側で、投資判断上重要な詳細を意図せず切り捨てるリスクを伴う。

\begin{casestudy}{When Summaries Distort Decisions: 「情報忠実度」という新しい評価軸}
LLMによる金融文書の要約は、意思決定上重要な詳細を体系的に欠落させうる。本論文は「情報忠実度（information fidelity）」という指標を導入し、複数の圧縮バリアントを生成して原文と相互チェックするAgentic Context Compressionを提案する。AI圧縮された金融データを投資判断に用いる際の安全性を高める枠組みとして位置づけられる。（出典: \cite{summariesdistort2026}）
\end{casestudy}

\begin{center}
\begin{tikzpicture}[
  doc/.style={fnode, text width=7.0cm, minimum height=1.1cm, font=\small},
  proc/.style={fnode blue, text width=7.5cm, minimum height=1.1cm, font=\small},
  var/.style={fnode muted, text width=2.6cm, minimum height=1.0cm, font=\footnotesize},
  chk/.style={fnode accent, text width=7.5cm, minimum height=1.1cm, font=\small},
  good/.style={fnode blue, text width=4.4cm, minimum height=1.4cm, font=\footnotesize},
  bad/.style={draw=red!60!black, fill=red!6, figshadow, rounded corners=4pt, line width=0.7pt, text width=4.8cm, align=center, minimum height=1.4cm, font=\footnotesize, inner sep=2mm},
  arr/.style={farr blue},
]
\node[doc] (orig) at (5.2,5.2) {原文書（決算資料・開示文書）};
\node[proc] (comp) at (5.2,3.5) {圧縮：複数の要約バリアントを生成\\(Agentic Context Compression)};
\node[var] (va) at (1.8,1.7) {バリアントA};
\node[var] (vb) at (5.2,1.7) {バリアントB};
\node[var] (vc) at (8.6,1.7) {バリアントC};
\node[chk] (chk) at (5.2,-0.2) {原文との相互チェック → 情報忠実度スコア算出};
\node[good] (good) at (2.2,-2.3) {忠実な要約\\→ 投資判断に使用可};
\node[bad] (bad) at (8.4,-2.3) {重要な詳細の欠落を検知\\→ 意思決定を歪めるリスク};
\draw[arr] (orig) -- (comp);
\draw[arr] (comp) -- (va);
\draw[arr] (comp) -- (vb);
\draw[arr] (comp) -- (vc);
\draw[arr] (va) -- (chk);
\draw[arr] (vb) -- (chk);
\draw[arr] (vc) -- (chk);
\draw[arr] (chk) -- (good);
\draw[arr, red!60!black] (chk) -- (bad);
\end{tikzpicture}
\end{center}
{\small\color{brandgray}\textbf{図: 要約圧縮と情報忠実度検証のフロー}\ Agentic Context Compressionは原文書から複数の要約バリアントを生成し原文と相互チェックすることで「情報忠実度」を算出し、意思決定を歪めかねない重要な詳細の欠落を検知する枠組みである。}

\subsection{財務推論ベンチマーク: 数値・多段・企業横断}
\label{subsec:reasoning-bench}

抽出・要約が「テキストから何を読み取るか」であるのに対し、推論は「読み取った数値をどう計算・比較・統合するか」を問う。この領域は、金融QAにおけるLLMの能力と限界を最も鋭く可視化する。基礎ベンチマークであるFinQA（決算報告書2,789件・8,281問、各問に実行可能な推論プログラムを注釈）とその対話拡張ConvFinQA（多ターンで前問の実体・期間・計算結果を暗黙参照）は、いずれも「事前学習済みモデルは多段の数値推論において専門家人間に遠く及ばない」ことを最初に定量化した\cite{x05finqa2021,x05convfinqa2022}。

\begin{casestudy}{FinanceQA: 試験に受かっても実務では6割落とす}
CFA試験のような選択式では高得点を出すモデルも、実務のアナリスト業務では大きく躓く。FinanceQAはハンドスプレッド（財務諸表の手作業展開）、会計慣行の適用、仮定生成を要する多段タスクという現実的なワークフローでLLMを評価し、現行モデルが現実的タスクの\textbf{約60\%を誤答}（正答率約40\%）することを示した。金融機関が要求する厳格な精度に達しておらず、OpenAIのファインチューンでも改善は限定的であった。標準化された試験の成功が、仮定に満ちた実務作業へは転移しないことを明らかにする。（出典: \cite{x05financeqa2025}）
\end{casestudy}

\begin{casestudy}{FinReflectKG-MultiHop: 企業横断・年度またぎで崩れる多段推論}
企業横断・年度またぎのグラウンデッドな多段推論を測る555問のベンチマーク（文書内48.7\%・同一企業の年度またぎ41.6\%・同一年度の企業横断9.7\%）。年度またぎの問い（LLM-Judge \textbf{6.72}）が、文書内（\textbf{7.47}）や企業横断（7.12）より難しく、時間的整合と用語の変遷が主因とされる。重要な発見は、知識グラフに紐付けた根拠（KG-linked evidence）がページウィンドウ検索に対し正答性を\textbf{約24\%}改善しつつトークンを\textbf{約84.5\%削減}した点であり、企業横断・長期の財務推論を律速するのはモデル規模ではなく検索アーキテクチャであることを裏付ける。（出典: \cite{x05finreflectkg2025}）
\end{casestudy}

決算報告書のテキスト・表・チャートを横断して三段の証跡を連鎖させるFCMR（Financial Cross-Modal Multi-hop Reasoning、ACL 2025）は、必要なホップ数が増えるほど最先端マルチモーダルLLMでも性能が単調に劣化し、ハルシネーションと誤りの伝播が主要な失敗モードであることを示した\cite{x05fcmr2025}。一方、標準化された試験ドメインでは進歩も著しい。CFA Level IIIに23モデルを投入した包括的評価では、フロンティア推論モデル（o4-mini \textbf{79.1\%}、Gemini 2.5 Flash \textbf{77.3\%}）が合格閾値の約63\%を明確に上回った。ただし選択式では71〜75\%に集中する一方、分析・統合・戦略的思考を要する記述式では性能の分散が大きく、ごく一部のモデルしか高得点を得られない\cite{x05cfa3eval2025}。試験の合否と実務の信頼性の乖離が、FinanceQAの結果と符合する。

\subsection{ドメイン特化型金融LLM: 規模ではなくタスク形状が勝敗を分ける}

金融ドメインに特化したモデル群――FinBERT（金融テキスト分類の先駆）\cite{finbert2020}、BloombergGPT（金融コーパスで事前学習した500億パラメータモデル）\cite{bloomberggpt2023}、FinGPT（安価なLoRA適応）\cite{x05fingpt2023}――は、汎用フロンティアモデルの台頭を受けて存在意義を問われてきた。しかし近年の証拠は、専門特化型小型モデルと汎用巨大モデルの優劣が「規模」ではなく「タスク形状」で決まることを示している。

\begin{concept}{ドメイン特化型金融LLM}
汎用LLM（何でもこなす巨大モデル）をそのまま金融タスクに投入しても、専門用語や独特な数値表現でつまずくことがある。そこで金融データを使って「適応」させる主な手段が三つある。継続事前学習は金融コーパスをそのまま追加で読ませる方法、ファインチューニングは正解付きデータで重みを調整する方法、LoRAは少数の追加パラメータだけを低コストで学習する方法である。いわば「万能選手を専門競技のために追加合宿させる」作業に近い。FinBERT・BloombergGPT・FinGPTはこの適応の代表例であり、勝敗を分けるのはモデルの大きさそのものではなく、解くべきタスクの形状に適応が合っているかどうかである。
\begin{center}
\begin{tikzpicture}[
  gen/.style={fnode muted, text width=2.6cm, minimum height=1.3cm, font=\footnotesize, align=center},
  adapt/.style={fnode blue, text width=3.4cm, minimum height=1.7cm, font=\footnotesize, align=center},
  spec/.style={fnode accent, text width=2.8cm, minimum height=1.5cm, font=\footnotesize, align=center},
  small/.style={fnode blue, text width=3.4cm, minimum height=1.2cm, font=\footnotesize, align=center},
  big/.style={fnode muted, text width=3.4cm, minimum height=1.2cm, font=\footnotesize, align=center},
  arr/.style={farr blue},
]
\node[gen] (g) at (0,2.2) {汎用LLM\\（何でもこなす\\巨大モデル）};
\node[adapt] (a) at (3.9,2.2) {適応\\継続事前学習／\\ファインチューニング／LoRA};
\node[spec] (s) at (8.0,2.2) {金融特化LLM\\FinBERT・\\BloombergGPT・FinGPT};
\node[small] (sm) at (1.9,-1.1) {小型特化モデル\\(NER: BERT-FT 約0.88)};
\node[big] (bg) at (8.0,-1.1) {巨大汎用モデル\\(NER: GPT-4o 約0.82)};
\draw[arr] (g) -- (a);
\draw[arr] (a) -- (s);
\draw[farr accent] (sm) -- node[midway, above, font=\scriptsize, text=figaccent, align=center] {特定タスクでは\\小型が上回る} (bg);
\end{tikzpicture}
\end{center}
{\small\color{brandgray}\textbf{図: 汎用LLMの金融特化適応と規模逆転}\ 継続事前学習・ファインチューニング・LoRAで汎用LLMを金融特化LLMへ適応させるフロー（上段）と、固有表現抽出のような構造化タスクでは小型特化モデルが巨大汎用モデルを上回りうる例（下段）。}
\end{concept}

\begin{casestudy}{Fin-R1: 強化学習で7Bが671Bに肉薄する}
Fin-R1はQwen2.5-7B-Instructを基盤に、約60,091件のCoTサンプルでのSFTとGRPO強化学習（ルールベースの形式・正確性報酬）を組み合わせた70億パラメータの財務推論モデル。FinQA \textbf{76.0}、ConvFinQA \textbf{85.0}、5ベンチマーク平均\textbf{75.2}で総合2位につけ、DeepSeek-R1-Distill-Llama-70Bを8.7ポイント上回り、671BのフルDeepSeek-R1（78.2）に対してもわずか3.8ポイント差にとどまった。約100分の1のパラメータ規模で標準化推論に肉薄する、RL後訓練の効率性を示す事例である。（出典: \cite{x05finr1_2025}）
\end{casestudy}

\begin{casestudy}{Palmyra-Fin-70B: CFA Level IIIに初合格した金融特化モデル}
WriterのPalmyra-Fin-70B（2024年8月公開）は、CFA Level IIIモック試験に\textbf{73\%}で合格した初のモデルとして報告された（平均合格ラインは約60\%、GPT-4は同試験で33\%）。長文書とQAを組み合わせた社内ベンチマークlong-fin-evalでは、Claude 3.5 Sonnet・GPT-4o・Mixtral-8x7bを上回った。長文脈の金融タスクにおけるドメイン特化の価値を示す。（出典: \cite{x05palmyrafin2024}）
\end{casestudy}

\begin{insight}{仮説: 特化小型モデルと汎用巨大モデルの交差点はタスク形状で決まる}
出力空間が明確な構造化タスク（NER: ファインチューンBERT F値約0.88 > プロンプトGPT-4o 0.82、センチメント: FinGPTのLoRA適応が優勢）ではドメイン特化の小型モデルが勝ち、開放的な統合・推論タスク（CFA III記述式、FinanceQAの仮定生成）ではフロンティア推論モデル（o4-mini 79.1\%）が支配する。RL後訓練は7B（Fin-R1、平均75.2）を671Bモデルの3.8ポイント差まで引き上げる。導入判断は「どちらが賢いか」ではなく「どのタスクか」であり、特化抽出器群＋フロンティア推論器のポートフォリオが単一モデルを上回る。かつFin-R1の100分の1のパラメータが示す総保有コストの優位により、高スループットの抽出用途では特化小型モデルが経済的に支配的であり続けると予想される。（出典: \cite{x05finr1_2025,x05finerord2025,x05fingpt2023,x05cfa3eval2025,x05palmyrafin2024}）
\end{insight}

\subsection{ハルシネーションと忠実度: 推論のアリティが崩壊点を決める}

金融QAでは、事実と異なる数値の生成（ハルシネーション）や、与えられた文脈に接地しない出力（忠実度の欠如）が意思決定に直結するため、その体系的な計測が不可欠である。

\begin{casestudy}{FAITH: 表形式ハルシネーションは「多変数計算」で0\%近くへ崩壊する}
FAITHは金融ハルシネーションを、文脈対応のマスク付きスパン予測タスクとして定式化する。S\&P 500企業の10-K「MD\&A（Item 7）」から単位付きの数値スパンをマスクし（453社・回答可能スパン2,406件）、モデルに文脈から復元させる。結果はClaude-Sonnet-4 \textbf{95.6\%}、Gemini-2.5-Pro \textbf{91.9\%}、GPT-4.1 89.2\%、Llama-3.1-8B 47.5\%、Qwen-3-8B 30.6\%。決定的なのは、推論の複雑さが増すと精度が急落する点で、\textbf{多変数計算}では多くのモデルが\textbf{ほぼ0\%}まで低下する。特徴的な失敗は「スケール誤り」（\$150と\$150 millionの混同）と、表とテキストの証跡を跨いだ統合の失敗である。（出典: \cite{x05faith2025}）
\end{casestudy}

金融QAの脆さは、既知のFinanceBenchの結果とも整合する。検索を併用したGPT-4-Turboでも\textbf{81\%}の質問を誤答または回答拒否し\cite{financebench2023}、AUDITFLOWではシンボリック検証層を除去すると精度が82\%から18\%へ崩壊する\cite{auditflow2026}。またFin-RATEベンチマークでは、企業横断比較タスクでファインチューン済み金融LLMの精度が23.91\%から\textbf{1.27\%}へ崩れる事例が確認されている\cite{finrate2026}。

\begin{insight}{仮説: 金融LLMの信頼性は推論のアリティで律速され、額面精度は多段の尾を過大評価する}
金融QAの信頼性を支配するのはモデル規模ではなく、要求される数値推論のアリティ（変数の数・ホップ数）である。FAITHでは最良モデル（Claude-Sonnet-4 95.6\%）でさえ多変数計算でほぼ0\%へ落ち、直接ルックアップはほぼ解決済み、FinReflectKGでは年度またぎ多段が最難、FCMRではホップ増加ごとに劣化する。示唆は二つ。第一に、実運用系はアリティでルーティングすべきである――ルックアップ・単段抽出はLLM、多変数計算はシンボリック／表計算ソルバ（検証済みツール呼び出し）へ。第二に、公表される総合精度は低アリティ問題の大きな比率に支配され、アナリストにとって重要な多段の尾（tail）における信頼性を系統的に過大評価している。（出典: \cite{x05faith2025,x05finreflectkg2025,x05fcmr2025,x05financeqa2025}）
\end{insight}

\subsection{記号検証ハイブリッド: LLMを読解器、ソルバを検証器に分離する}

ハルシネーション対策を「より大きなモデル」だけで解く発想は、金融文書では早く限界に当たる。財務諸表の整合性、XBRL概念、年度差分、倍率・単位、セグメント合計のような制約は、自然言語生成よりも実行可能な計算・グラフ探索・claim検証に向いている。近年のベンチマークは、LLMを読解・候補生成・根拠抽出に使い、数値整合性は記号層で再検証するアーキテクチャへ明確に収束している。

\begin{casestudy}{FinChain: 「答え」ではなく推論ステップを機械検証する}
FinChainは、FinQA/ConvFinQAが最終数値の正否を中心に評価していた弱点を補うため、58トピック・12金融ドメインにわたるパラメータ化された記号テンプレートと実行可能Pythonコードで、Chain-of-Thoughtの各ステップを検証可能にしたベンチマークである。提案指標CHAINEVALは、最終答えだけでなく中間推論の整合性も評価する。重要なのは、テンプレート生成によりデータ汚染に強い問題を拡張できる点であり、金融推論の評価単位が「正答したか」から「どの式・根拠で正答したか」へ移ることを示す。（出典: \cite{x05finchain2025}）
\end{casestudy}

\begin{casestudy}{GBFR: 証拠が足りないときに安全に棄権する金融推論}
GBFR（graph-bounded financial reasoning）は、財務指標ナレッジグラフでLLMの多段計算を制約し、複数の導出経路を並列探索して意味的に整合する結果だけを集約する。さらに、企業・時点・指標を摂動した反実仮想サンプルで「答えられない問い」を作り、検索失敗と真のデータ欠如を区別して安全に棄権する能力を評価する。これは、曖昧な根拠を値で埋める金融LLMの強制生成を、モデル出力後の拒否判定ではなく推論過程そのものの制約で抑える設計である。（出典: \cite{x05gbfr2026}）
\end{casestudy}

\begin{center}
\footnotesize
\renewcommand{\arraystretch}{1.25}
\begin{tabular}{@{}p{2.8cm}p{3.2cm}p{3.4cm}p{4.0cm}@{}}
\toprule
\textbf{方式} & \textbf{検証対象} & \textbf{記号層の役割} & \textbf{実装上の示唆} \\
\midrule
FinChain \cite{x05finchain2025} & CoTの中間式・最終答え & 実行可能Pythonテンプレートでステップ整合性を採点 & ベンチマークは最終正答率だけでなく推論経路を保存する \\
FinVerBench \cite{finverbench2026} & 財務諸表の数値整合性 & 算術・勘定連鎖・前年比・桁/倍率の摂動を識別 & 入力の丸め・単位・符号正規化を前処理で明示する \\
FinGround \cite{finground2026} & QA回答中の原子的claim & claim型に応じて検索・式再構成・根拠付与を分岐 & 回答生成後に段落・表セル単位の引用で書き戻す \\
AUDITFLOW \cite{auditflow2026} & 監査証跡・報告書計算 & 実行環境で会計演算を再計算 & LLMは読解と例外説明、計算は検証層へ委譲する \\
GBFR \cite{x05gbfr2026} & 多段財務指標と棄権判断 & 財務指標KGで導出経路を制約し交差検証 & 証拠不足時に推測で埋めず安全に棄権する \\
\bottomrule
\end{tabular}
\end{center}
{\small\color{brandgray}\textbf{表: 金融LLMの記号検証ハイブリッド}\ 金融文書QAの信頼性向上は、単一LLMの生成能力ではなく、CoTステップ、財務諸表、原子的claim、監査演算、指標グラフをそれぞれ検証可能な対象へ分解する設計に依存する。}

\begin{insight}{示唆: 「根拠付き回答」から「実行可能な根拠」へ評価軸が移る}
金融LLMの次の信頼性基準は、回答に引用が付いていることでは不十分である。引用された表・文・XBRLタグから、回答内の式と数値が再実行できる必要がある。FinGroundはclaim単位で根拠を付け直し、FinVerBenchは財務諸表の数値摂動を検出し、FinChainは推論ステップを実行可能テンプレートへ落とし、GBFRは指標グラフで導出経路を閉じる。したがって実務実装では、RAGの検索ログ、LLMの中間推論、表計算ソルバの実行結果、棄権理由を同じ監査証跡に保存することが、精度改善とモデルリスク管理を同時に満たす最小単位になる。（出典: \cite{finground2026,finverbench2026,x05finchain2025,x05gbfr2026,auditflow2026}）
\end{insight}

\subsection{ESG・サステナビリティ分析: 知識ギャップと検索依存}

ESG／サステナビリティ分析は、専門的・学際的な知識と非構造化開示の読解を同時に要する。ESGenius（専門家検証済み1,136問のMCQ＋7つの権威的情報源からの231文書コーパス）を用いた50モデル（0.5B〜671B）の評価では、最先端モデルでもゼロショット正答率は\textbf{55〜70\%}にとどまり、この専門ドメインにおける明確な知識ギャップが露呈した。一方RAGを併用すると性能は大きく改善し、DeepSeek-R1-Distill-Qwen-14Bは63.82\%から\textbf{80.46\%}へ上昇した。ESG分析は、パラメータに内在する知識ではなく検索（retrieval）に律速される――すなわち「モデルが知っているか」ではなく「正しい規格・文書を引けるか」が支配的であることを示す。（出典: \cite{x05esgenius2025}）

\subsection{ベンチマーク横断の俯瞰: スコアより評価プロトコルを読む}

以下の表は、本節で扱った主要な金融データ分析ベンチマークを、タスク種別・最良スコア・そこから読み取れる限界とともに一覧化したものである。単一の「精度」だけを見るのではなく、各ベンチマークがどのタスク形状・推論アリティを測っているかに注意して読む必要がある。

\begin{center}
\footnotesize
\renewcommand{\arraystretch}{1.3}
\begin{tabular}{@{}p{2.7cm}p{2.6cm}p{3.4cm}p{4.4cm}@{}}
\toprule
\textbf{ベンチマーク} & \textbf{タスク種別} & \textbf{最良スコア} & \textbf{読み取れる限界} \\
\midrule
FinQA / ConvFinQA \cite{x05finqa2021,x05convfinqa2022} & 数値推論／対話的数値推論 & 事前学習モデルは専門家人間に及ばず & 多段の数値推論が最初期からの難所 \\
FinanceQA \cite{x05financeqa2025} & 実務アナリスト業務 & 正答率約40\%（約60\%を誤答） & 試験の合格が実務へ転移しない \\
CFA Level III \cite{x05cfa3eval2025} & 選択式＋記述式 & o4-mini 79.1\%（閾値約63\%） & 記述式（統合・戦略）で分散大 \\
FinReflectKG-MultiHop \cite{x05finreflectkg2025} & 企業横断・年度またぎ多段 & 年度またぎ最難（LLM-Judge 6.72） & 律速はモデル規模でなく検索設計 \\
FCMR \cite{x05fcmr2025} & クロスモーダル多段 & ホップ増で単調に劣化 & 誤りの伝播が主要失敗モード \\
FAITH \cite{x05faith2025} & 表形式ハルシネーション & Claude-Sonnet-4 95.6\% & 多変数計算でほぼ0\%へ崩壊 \\
FinanceBench \cite{financebench2023} & 開示文書QA（検索併用） & 81\%を誤答・回答拒否 & 検索併用でも高い失敗率 \\
FinChain \cite{x05finchain2025} & 検証可能CoT金融推論 & 26モデルを評価 & 中間推論ステップの機械検証が必要 \\
GBFR \cite{x05gbfr2026} & 多段計算＋安全棄権 & SOTAベースラインを上回る & 証拠不足時の強制生成をKG制約で抑制 \\
FinTagging \cite{x05fintagging2025} & XBRL概念タグ付け & ファクト抽出F値約72\% & 概念リンキングが真の律速 \\
FiNER-ORD \cite{x05finerord2025} & 固有表現抽出 & BERT約0.88 > GPT-4o 0.82 & 構造化タスクは小型FTモデル優位 \\
ESGenius \cite{x05esgenius2025} & ESG知識QA & ゼロショット55〜70\%、RAGで80\%超 & 知識ギャップは検索で補償 \\
\bottomrule
\end{tabular}
\end{center}
{\small\color{brandgray}\textbf{表: 金融データ分析ベンチマーク横断の俯瞰}\ タスク形状（抽出・分類・数値推論・多段・クロスモーダル・ハルシネーション）ごとに最良スコアと限界を整理。構造化タスクでは小型ファインチューンモデルが優位、開放的推論ではフロンティアモデルが優位、多段・多変数の尾では総じて信頼性が崩壊するという構造が共通して現れる。}

\subsection{アルファ生成のためのLLM／時系列基盤モデル}

自然言語による着想やニュース・開示情報を、実際に取引可能なアルファやポートフォリオシグナルへと変換する試みが、学術・産業双方で急速に増えている。

\begin{casestudy}{Supply Chain Propagation of Textual Signals: 報告企業を超えて伝播するテキストシグナル}
10-K SECファイリングから抽出したFinBERT埋め込みをサプライチェーンネットワーク経由で伝播させ、S\&P 500銘柄で2011〜2025年においてSharpe比\textbf{0.86}のクロスセクショナル株式リターン要因を構築。LLMで符号化された非構造化開示情報が、報告企業自身を超えて価格形成シグナルを持つことを実証した。（出典: \cite{supplychain2026}）
\end{casestudy}

\begin{casestudy}{Kronos: 120億件のローソク足で訓練された市場専用基盤モデル}
45取引所・120億件超のローソク足レコードで事前学習した金融市場専用の基盤モデル。専用のマーケット・トークナイザを採用し、既存の時系列基盤モデルに対して価格予測・ボラティリティ予測タスクで\textbf{93\%の性能改善}を達成した（AAAI 2026）。もっとも、ベンチマーク上の改善が実際のライブアルファに転化するかは、事前学習コーパスのサバイバーシップバイアスや取引コスト、流動性の高い市場ではローソク足レベルのシグナルが既に裁定されつくしている可能性のため未検証であると留保されている。（出典: \cite{kronos2025}）
\end{casestudy}

\begin{casestudy}{Alpha-GPT: 自然言語の着想をアルファへ変換し、41,000超チーム中トップ10入り}
LLM（プロンプトエンジニアリング活用）でクオンツ研究者の自然言語による着想を候補アルファ（取引ファクター）へ変換するフレームワーク（EMNLP 2025 System Demonstration Track）。WorldQuantのInternational Quant Championship (IQC) 2024に外部参加者として参加し、\textbf{41,000超のチーム中トップ10}にランクインした。（出典: \cite{alphagpt2023}）
\end{casestudy}

\begin{casestudy}{CrossAlpha: 開示文書から国境を越えるアルファを抽出するクロスマーケット・ベンチマーク}
LLMエージェントが5市場の\textbf{10,700件}の年次報告書を標準化されたシグナルへ変換するベンチマーク。開示文書から導いたクロスマーケットのピア（同業他社）シグナルが、米国から日本へのリターン予測において国内ベースラインを上回った（ICIR \textbf{0.39} 対 0.07〜0.18）。AIが処理した開示情報が国境を越えたアルファを内包しうることを示す。（出典: \cite{crossalpha2026}）
\end{casestudy}

\begin{center}
\begin{tikzpicture}[
  src/.style={fnode, text width=4.6cm, minimum height=1.9cm, font=\footnotesize},
  outbox/.style={fnode accent, text width=9.5cm, minimum height=1.2cm, font=\small},
  arr/.style={farr blue},
]
\node[src] (n1) at (1.9,3.4) {\textbf{Supply Chain Propagation}\\FinBERT埋め込み×サプライチェーン伝播\\Sharpe比\textbf{0.86}（S\&P 500, 2011〜2025）};
\node[src] (n2) at (6.3,3.4) {\textbf{Kronos}\\120億件超のローソク足で事前学習\\既存基盤モデル比\textbf{+93\%}改善（AAAI 2026）};
\node[src] (n3) at (10.7,3.4) {\textbf{Alpha-GPT}\\自然言語の着想を候補アルファへ変換\\IQC 2024 41,000超チーム中\textbf{トップ10}};
\node[outbox] (out) at (6.3,-0.2) {\textbf{取引可能なアルファ／ポートフォリオシグナル}};
\draw[arr] (n1.south) -- ++(0,-0.5) -| (out.north);
\draw[arr] (n2.south) -- (out.north);
\draw[arr] (n3.south) -- ++(0,-0.5) -| (out.north);
\end{tikzpicture}
\end{center}
{\small\color{brandgray}\textbf{図: テキスト・価格系列・自然言語からのアルファ生成アプローチ}\ Supply Chain Propagation（テキスト由来シグナルの企業間伝播）、Kronos（価格系列基盤モデル）、Alpha-GPT（自然言語アイデアの構造化）はいずれも入力形式が異なるが、取引可能なアルファ・ポートフォリオシグナルへと収束する点は共通している。}

\begin{itemize}
  \item \textbf{GIFT (LLM-Guided State-Reward Interface)}: LLMを推論時ではなく、PPOポートフォリオ取引エージェントの状態特徴量と報酬設計ルール自体の設計に用いる手法。金融ドメイン知識を構造的に注入し、リスク調整後のアウトオブサンプル性能を改善する。（\cite{gift2026}）
  \item \textbf{Alpha-R1}: 8Bパラメータの推論モデルを強化学習で訓練し、システマティック投資におけるシグナル劣化・レジーム変化に対応。過去パターンのみに頼らずリアルタイムニュースとファクターロジックを取り込み、アルファファクターの文脈的な有効性を判定して選択的に有効化・無効化する（上海交通大学・StepFun・FinStep所属）。（\cite{alphar1_2025}）
  \item \textbf{Point-in-Time Financial RAG}: LLM自体をファインチューンするのではなく、ベイズ的な市場フィードバックでどのニュース文書を検索（retrieve）するかを選ぶことで、ニュース駆動型トレーディングポートフォリオのSharpe比を\textbf{0.52から0.84}へ改善。金融RAGでは検索戦略の選択がモデル選択と同等に重要であることを示す。凍結（frozen）LLM前提のためポイントインタイム性の担保にも資する。（出典: \cite{pitrag2026}）
  \item \textbf{Quant 4.0（フレームワーク提唱）}: 純粋なディープラーニングを超えた第4世代の量的投資として、自動機械学習（アルゴリズムがアルゴリズムを生成）、隠れたリスクを開示する説明可能AI、ドメイン事前知識を注入する知識駆動AIを統合するアーキテクチャを提案。分野の10の未解決課題を整理する。（\cite{quant40_2023}）
  \item \textbf{サーベイ}: 統計的・特徴量ベース手法からディープラーニング、LLM／エージェント駆動の自己反復ワークフローに至る量的投資AIの進化を包括的に整理したサーベイとして\cite{quantsurvey2025}、テキスト・数値表・図表が混在する複雑な財務情報を処理する新規LLM手法を統合し実務家向けの研究方向を示すサーベイとして\cite{bridginglm2025}、株式市場におけるLLM活用を予測・センチメント・トレーディング・分析・ポートフォリオ管理・リスク・コンテンツ生成・評価の8応用カテゴリで整理し84本の研究（2022〜2025年）を統合した上でバックテスト結果と実運用執行との間の決定的なギャップを指摘するサーベイとして\cite{llmequitysurvey2025}を参照。
  \item \textbf{オルタナティブデータの制度化（市場動向）}: Coalition Greenwichが2025年9〜10月にバイサイド56社を対象に実施した調査では、およそ\textbf{半数}が中頻度のオルタナティブデータ利用者であり、\textbf{63\%}が支出の増額を計画していると報告された。衛星画像・消費者取引・Webスクレイピングによるシグナルが、周縁的なツールからコアインフラへと移行しつつあることを示す。（出典: \cite{altdata2025}）
\end{itemize}

\begin{insight}{仮説: マルチモーダルLLMによる小型株アルファ生成}
FinMME（ACL 2025、11,000以上のサンプル）は、最良のマルチモーダルモデル（Gemini Flash 2.0, Qwen2.5-VL 72B）でも金融チャートからの計算問題正答率が\textbf{40\%未満}であることを発見した。それでもRussell 2000銘柄（アナリストカバレッジ2〜3人）では、決算スライドの収益ブリッジ図表にテキストトランスクリプトに現れない定量シグナルが含まれると分析されており、大型株では飽和した精度向上の価値が小型株では相対的に高いという仮説が提示されている。（出典: \cite{finmme2025}）
\end{insight}

\begin{insight}{仮説: 進化的LLMアルファマイニングは混雑により自己劣化する}
抽象構文木（AST）で制約したファクター式に遺伝的な突然変異・交叉を適用する進化的LLMアルファマイニング（例: QuantaAlpha）は、創造的な着想（LLMによる仮説生成）と構文的妥当性（AST制約付きコード構築）を切り離し、既知シグナルの微修正ではなく構造的に多様なファクターを生む。QuantaAlphaはCSI 300で年率\textbf{27.75\%}、S\&P 500で再最適化なしに\textbf{137\%}の超過リターンを報告した。ただし複数の競合ファンドが同一の銘柄ユニバースに類似の進化的探索を適用すれば、発見されるファクターは重複し混雑（crowding）を通じて自己劣化していくと予想される。クロスマーケット汎化の結果も2022〜2025年のバックテスト期間に固有のモメンタム動学やサバイバーシップバイアスにより見かけ上誇張されている恐れがある。（出典: \cite{quantaalpha2026}）
\end{insight}

\begin{insight}{仮説: 意思決定志向学習はポートフォリオ成果を大きく引き上げる}
予測誤差の平均最小化は、下流のポートフォリオ損失の最小化とは整合しない。意思決定にとって重要な誤差構造は予測誤差の特定の部分空間に集中している。接空間（tangent-space）射影による意思決定志向学習（decision-focused learning）はこの部分空間を直接ターゲットとし、予測精度はやや劣るがポートフォリオ成果は大幅に優れるモデルを生む（ICML 2026）。これはリターン予測の改善に注ぎ込まれてきたアルファ生成努力の多くが方向違いであった可能性を示唆する。ただし意思決定志向モデルは相関崩壊のような構造的レジーム転換で過学習が露呈しうるうえ、微分可能なポートフォリオソルバと複雑な学習パイプラインを要し、実運用への導入障壁となる。（出典: \cite{decisionfocused2026}）
\end{insight}

\begin{insight}{仮説: グラフ言語モデルの優位は連鎖的な波及タスクで拡大する}
企業間サプライチェーンをLLM互換のグラフ埋め込みとして符号化するハイブリッドなグラフ言語モデル（例: InterCorpRel-LLM）は、GNN埋め込みをLLMトークン空間へ対照学習で整合させ、サプライ関係予測でF値\textbf{0.8543}を達成した。単一企業の信用評価（DSCR・レバレッジ）はテキストのみからでも導けるが、Tier-1サプライヤの破綻がTier-2/3へ連鎖する系統的なショック波及は明示的なグラフ探索を要する。ゆえにテキストのみのLLMに対する性能差は、単一企業のデフォルト予測よりも波及（contagion）タスクで有意に拡大すると予想される。もっともサプライチェーンのデータ被覆は非上場・非米国企業で疎であり、見かけの優位はアーキテクチャではなくベンダー由来のデータ品質を反映している恐れがあり再現が難しい。（出典: \cite{intercorprel2025}）
\end{insight}

\subsection{高頻度・非構造化シグナルの解読}

オプション市場のミクロ構造や経営者の発話パターンといった高頻度・非構造化シグナルの解読は、伝統的な定量分析では捉えにくい情報をLLMが抽出できるかを試す領域として注目されている。

\begin{itemize}
  \item \textbf{Inferring Latent Market Forces（オプションガンマ検出）}: LLMがオプションチェーン構造のみからディーラーのガンマヘッジ制約レジーム（ピンニング、0DTEウォール抑制、デルタリバランス圧力）を検出できるかを検証。難読化テストで記憶可能な時間的文脈をすべて除去した状態でも\textbf{71.5\%の検出精度}、構造的文脈を保持すると\textbf{91.2\%}まで上昇し、単なる訓練データパターンマッチングではなく因果理解の兆候を示唆する。（\cite{gammaexposure2025}）
  \item \textbf{CFOs Meet LLMs}: デューク大学・NBERの研究者が、特定企業のCFOをロールプレイするLLMが2002〜2025年にわたる当該経営者の実際の景況感サーベイ回答を予測できることを実証。LLMを高頻度金融センチメントデータのためのスケーラブルなデジタルツインとして活用する可能性を示す。（\cite{cfosmeetllms2026}）
  \item \textbf{FinCall-Surprise（決算コールのパラ言語分析）}: 現行の汎用マルチモーダルモデルは、決算コール音声からトランスクリプトのみのベースラインに対して控えめかつ一貫性のない改善しか得られない。ピッチの分散・センシティブな話題での発話速度の加速・躊躇マーカーの密度といった音声特徴には、汎用マルチモーダル訓練が最適化していないストレス・欺瞞シグナルが含まれる可能性がある一方、エグゼクティブのメディアトレーニングによるストレスマーカー抑制や決算コール音声の意図的なトーン調整がすでに始まっているとの報告もあり、訓練済みシグナルが漸近的に無効化しうるリスクが指摘されている。（\cite{fincallsurprise2025}）
\end{itemize}

\begin{insight}{仮説: 拡散モデルは決算イベント時のボラティリティ曲面予測でVAEを上回る}
過去のボラティリティ曲面動学とマクロ経済状態変数を条件付けした生成拡散モデル（DDPM）は、翌日のSPXボラティリティ曲面予測で既存手法を上回り、拡散過程が生む裁定違反を抑制する信号対雑音比（SNR）重み付けを備える。一方、10次元VAEで曲面を圧縮する並行手法はテール動学のモデル化を犠牲にする。決算イベントはストライク×満期の局所的で鋭いショックを生み、VAE圧縮はこれを平準化して消してしまうため、まさにそこで拡散モデルが優位に立つと予想される。ただしDDPM推論はVAEデコードに比べ計算コストが高く、数千のストライク・満期を秒未満で再価格付けするマーケットメイク業務では、精度優位に関わらず遅延がVAEを事実上の実運用標準にしうる。（出典: \cite{diffusionvol2025}）
\end{insight}

\begin{insight}{仮説: LLMエージェント的ナウキャスティングは変曲点で統計モデルを上回る}
歴史的指標系列で学習した従来の統計的ナウキャスティングは、正確な予測が最も必要となる構造的断絶・レジーム転換のときにこそ機能不全に陥る。LLMエージェントは中央銀行のフォワードガイダンスやクロスマーケットのセンチメントをほぼリアルタイムに同化でき、変曲点で高い予測力を持つがラグ付き統計入力では捉えにくい情報を取り込める。優位は高VIX四半期に統計的に集中すると予想される（インフレナウキャスティングでニュースセンチメントとFED声明埋め込みを統合）。ただしLLMは真の基礎指標ではなく報道に埋め込まれた市場コンセンサスをナウキャストしているにすぎない恐れがあり、既に価格に織り込まれた予想からの情報の切り分けにはプロの予測サーベイに対する直交化が要る。（出典: \cite{inflationnowcast2024}）
\end{insight}

\begin{center}
\begin{tikzpicture}[
  inp/.style={fnode muted, text width=3.6cm, minimum height=1.5cm, font=\footnotesize},
  llm/.style={fnode, text width=3.4cm, minimum height=1.5cm, font=\footnotesize},
  res/.style={fnode blue, text width=4.0cm, minimum height=1.5cm, font=\footnotesize},
  arr/.style={farr blue},
]
\node[inp] (a1) at (1.9,6.4) {オプションチェーン構造\\（価格・出来高のみ）};
\node[llm] (a2) at (6.0,6.4) {LLM推論\\ディーラーのガンマヘッジレジーム検出};
\node[res] (a3) at (10.3,6.4) {検出精度\\難読化時\textbf{71.5\%} → 文脈保持時\textbf{91.2\%}};
\node[inp] (b1) at (1.9,3.6) {CFOの過去景況感\\サーベイ回答（2002〜2025）};
\node[llm] (b2) at (6.0,3.6) {LLMロールプレイ\\CFOデジタルツイン};
\node[res] (b3) at (10.3,3.6) {実際の回答を予測\\高頻度センチメントの\\スケーラブルな代替};
\node[inp] (c1) at (1.9,0.8) {決算コール音声\\ピッチ分散・発話速度・\\躊躇マーカー};
\node[llm] (c2) at (6.0,0.8) {汎用マルチモーダル\\モデル};
\node[res] (c3) at (10.3,0.8) {改善は控えめかつ不安定\\（メディアトレーニングで\\漸近的に無効化するリスク）};
\draw[arr] (a1) -- (a2); \draw[arr] (a2) -- (a3);
\draw[arr] (b1) -- (b2); \draw[arr] (b2) -- (b3);
\draw[arr] (c1) -- (c2); \draw[arr] (c2) -- (c3);
\end{tikzpicture}
\end{center}
{\small\color{brandgray}\textbf{図: 高頻度・非構造化シグナルをLLMが解読する3つの試み}\ オプションガンマ検出は構造的文脈の保持で精度が71.5\%から91.2\%へ向上する一方、CFOデジタルツインは実サーベイ回答の予測に成功し、決算コール音声のパラ言語分析は改善が限定的かつ不安定である。}

\subsection{信用リスク・与信分析への応用}

与信・信用リスク評価は、規制上の説明可能性要求と数値推論の精度が同時に求められる領域であり、LLM単体ではなく専用アーキテクチャや連合学習との組み合わせが模索されている。

\begin{itemize}
  \item \textbf{Large Tabular Models (LTM)}: LLMは精密な数値推論に弱くハルシネーションを起こしやすいため、信用リスクスコアリング・不正検知・規制資本計算のような構造化タスクでは、目的特化型のLarge Tabular Modelsがファインチューン済みLLMに取って代わる可能性がある。DeepMind出身者が共同創業したFundamentalのNEXUS LTMは\textbf{2.55億ドル}の資金調達で立ち上げられ、「数時間ではなく数秒で予測」を提供すると謳っている。（出典: \cite{citiventuresndtabular}）
  \item \textbf{クロス機関連合学習によるAML}: 香港の銀行（Airstar Bank、livi bank）の実証実験で、連合学習AMLモデルは孤立した単独銀行モデルよりおよそ\textbf{2倍の不正資金フロー検知有効性}を示した。ただし異種データ分布下でのクライアントドリフトにより、より代表性の高いデータを持つ機関が勾配更新を支配し、小規模な地域銀行の個別改善効果が乏しくなる可能性、また法域ごとに異なるAML報告基準がモデル勾配の共有さえ法的障壁にしうる点が課題として指摘されている。（\cite{federatedlearning2025}）
  \item \textbf{Interpretable LLMs for Credit Risk（サーベイ）}: 2020〜2025年の60本の論文をPRISMAレビュー手法で整理し、ハイブリッドパイプライン（LLM＋伝統的スコアカードロジック、chain-of-thought推論、パラメータ効率的ファインチューニング）が主流トレンドである一方、モデルレベルの解釈可能性と機関レベルの規制検証との間にギャップがあると指摘する。（\cite{creditrisk2025}）
  \item \textbf{LendNova}: 生の信用情報機関（bureau）テキストに言語モデルを直接適用する初のエンドツーエンドパイプラインで、手作業の特徴量エンジニアリングを不要にし、前処理コストの削減とスケーラビリティ向上を図る（AAAI 2026 Workshop on Agentic AI in Financial Services）。（出典: \cite{lendnova2026}）
  \item \textbf{AI-BAAM（オルタナティブデータによる与信）}: 生の銀行取引履歴をオルタナティブデータとして用いるマレーシアのMSME（零細・中小企業）与信のエンドツーエンドMLパイプライン。AUROC \textbf{0.806}を達成し、申込情報のみのベースラインに対し\textbf{+24.6\%}の改善を示した（ICLR 2026 oral）。（出典: \cite{aibaam2025}）
\end{itemize}

\begin{center}
\begin{tikzpicture}[
  hdr/.style={font=\footnotesize\bfseries, text=figgray},
  inp/.style={fnode muted, text width=5.0cm, minimum height=1.3cm, font=\footnotesize},
  opt/.style={fnode blue, text width=4.6cm, minimum height=1.7cm, font=\footnotesize},
  outp/.style={fnode accent, text width=6.0cm, minimum height=1.2cm, font=\small},
  outp2/.style={fnode accent, text width=4.4cm, minimum height=1.4cm, font=\footnotesize},
  arr/.style={farr blue},
]
\node[hdr] at (3.1,5.6) {与信スコアリング};
\node[inp] (n1) at (3.1,4.7) {顧客の構造化・非構造化データ\\（財務諸表・取引履歴）};
\node[opt] (o1) at (0.6,2.5) {Large Tabular Model\\NEXUS（Fundamental）\\\textbf{2.55億ドル}調達、数秒で予測};
\node[opt] (o2) at (5.9,2.5) {ハイブリッドLLM＋\\伝統的スコアカード\\chain-of-thought／PEFT};
\node[outp] (n3) at (3.1,0.0) {与信判定（規制上の説明可能性が必須）};
\draw[arr] (n1) -- (o1);
\draw[arr] (n1) -- (o2);
\draw[arr] (o1) -- (n3);
\draw[arr] (o2) -- (n3);

\draw[figgray, dashed] (-1.0,-1.2) -- (12.0,-1.2);
\node[hdr] at (5.4,-1.7) {クロス機関連合学習によるAML};
\node[inp] (m1) at (1.7,-3.1) {香港複数銀行の\\分散取引データ\\(Airstar Bank, livi bank)};
\node[opt] (m2) at (5.9,-3.1) {連合学習AMLモデル\\（クライアントドリフト・\\法域間の課題あり）};
\node[outp2] (m3) at (10.0,-3.1) {孤立モデル比\\\textbf{約2倍}の不正資金\\フロー検知有効性};
\draw[arr] (m1) -- (m2);
\draw[arr] (m2) -- (m3);
\end{tikzpicture}
\end{center}
{\small\color{brandgray}\textbf{図: 信用リスク・与信分析と連合学習AMLのパイプライン}\ 与信スコアリングはLarge Tabular ModelまたはLLM＋伝統的スコアカードのハイブリッドを経て、規制上説明可能な判定に至る。下段のクロス機関連合学習AMLは、孤立した単独銀行モデル比で約2倍の不正資金フロー検知有効性を示す一方、クライアントドリフトや法域差が課題として残る。}

\subsection{モデル圧縮とコスト効率}

フロンティア財務LLMの推論コストを実運用可能な水準まで下げる取り組みも並行して進んでいる。

\begin{itemize}
  \item \textbf{知識蒸留による財務LLM圧縮（P-KD-Q）}: Pruning→Knowledge Distillation→Quantizationの圧縮パイプラインにより、フロンティア財務LLMから7〜13Bパラメータの生徒モデルを生成し、教師モデルの精度の\textbf{90\%以上}を保持する可能性が、WiproとNVIDIAの2026年実務知見に基づく体系的レビューで報告されている。（出典: \cite{researchgatendkd}）
\end{itemize}

\begin{center}
\begin{tikzpicture}[
  model/.style={fnode, text width=3.0cm, minimum height=1.8cm, font=\footnotesize},
  stg/.style={fnode blue, text width=2.4cm, minimum height=1.4cm, font=\footnotesize},
  note/.style={fnode accent, text width=12.5cm, minimum height=1.1cm, font=\small},
  arr/.style={farr blue},
]
\node[model] (teacher) at (1.6,0) {教師モデル\\フロンティア\\財務LLM};
\node[stg] (prune) at (4.6,0) {Pruning};
\node[stg] (kd) at (7.2,0) {Knowledge\\Distillation};
\node[stg] (quant) at (9.8,0) {Quantization};
\node[model] (student) at (12.6,0) {生徒モデル\\7〜13B\\パラメータ};
\node[note] (note) at (7.1,-2.0) {教師モデル精度の\textbf{90\%以上}を保持（WiproとNVIDIAの2026年実務知見に基づく体系的レビュー）};
\draw[arr] (teacher) -- (prune);
\draw[arr] (prune) -- (kd);
\draw[arr] (kd) -- (quant);
\draw[arr] (quant) -- (student);
\draw[figgray, dashed] (teacher.south) |- (note.west);
\draw[figgray, dashed] (student.south) |- (note.east);
\end{tikzpicture}
\end{center}
{\small\color{brandgray}\textbf{図: 財務LLMの知識蒸留パイプライン（P-KD-Q）}\ Pruning → Knowledge Distillation → Quantizationの3段階を経て、フロンティア財務LLM（教師）から7〜13Bパラメータの生徒モデルを生成し、教師モデルの精度の90\%以上を保持する。}

\subsection{再現性・頑健性への警鐘: 「ルックアヘッドバイアス」問題}

\begin{concept}{ルックアヘッドバイアス（先読みバイアス）}
ルックアヘッドバイアス（先読みバイアス）とは、本来その時点ではまだ知り得ないはずの「未来の情報」が、過去時点の判断や検証にこっそり混ざり込んでしまう誤りである。LLMに特有の形として、モデルが学習データの中に学習カットオフ以降の市場の結果を実質的に「記憶」しており（temporal contamination）、それを使ってバックテストすると本来の実力以上に成績が良く見えてしまう問題がある。いわば「答えを知ってから解いた模試の点数を、実力だと錯覚する」ようなものである。対策は、評価期間を学習カットオフより後に限定することと、その時点で実際に入手可能だったデータだけを使うpoint-in-time原則の徹底である。
\begin{center}
\begin{tikzpicture}[
  lbl/.style={font=\footnotesize, align=center, text width=3.6cm},
  cut/.style={font=\footnotesize\bfseries, text=fignavy, align=center},
  r1/.style={fnode accent, text width=2.7cm, minimum height=1.3cm, font=\footnotesize, align=center},
  r2/.style={fnode, text width=2.2cm, minimum height=1.3cm, font=\footnotesize, align=center},
]
\draw[farr] (0,0) -- (12.2,0);
\node[font=\scriptsize, text=figgray] at (12.2,0.35) {時間};
\draw[figgray, dashed, line width=0.8pt] (5.0,-2.6) -- (5.0,2.6);
\node[cut] at (5.0,2.95) {学習\\カットオフ};
\draw[figblue, line width=3mm, line cap=round] (0.6,1.2) -- (4.4,1.2);
\node[lbl] at (2.5,1.75) {評価期間\\（カットオフ以前に設定）};
\node[r1] (out1) at (9.0,1.2) {モデルが未来を「記憶」\\$\to$見かけの高アルファ};
\draw[farr accent] (4.4,1.2) -- (out1);
\draw[figgray, line width=3mm, line cap=round] (5.6,-1.5) -- (9.4,-1.5);
\node[lbl] at (7.5,-2.05) {評価期間\\（カットオフ以降に限定）};
\node[r2] (out2) at (10.9,-1.5) {真の性能\\（過大評価なし）};
\draw[farr] (9.4,-1.5) -- (out2);
\end{tikzpicture}
\end{center}
{\small\color{brandgray}\textbf{図: ルックアヘッドバイアスが生じる仕組み}\ 評価期間が学習カットオフより前に設定されると（上段）モデルが未来の結果を実質的に記憶しており見かけの高アルファが生じる一方、評価期間をカットオフ以降に限定すると（下段）point-in-timeな真の性能が測れる。}
\end{concept}

金融特化のアルファ生成研究において、LLMが事前学習データに含まれる過去の結果を「記憶」して見かけ上高いバックテスト成績を出す「look-ahead bias／temporal contamination」問題が、独立した複数の研究によって体系的に検証されている。下図は、この問題が公表アルファをどれほど劇的に消し去りうるかを示す典型例である。

\begin{center}
\begin{tikzpicture}
\begin{axis}[
  ybar,
  bar width=26mm,
  width=9.6cm,
  height=6.2cm,
  ymin=-6, ymax=25,
  ylabel={年率アルファ（\%）},
  ylabel style={font=\small},
  xtick={1,2},
  xticklabels={学習カットオフ前, 学習カットオフ後},
  xmin=0.4, xmax=2.6,
  xticklabel style={font=\bfseries\small, align=center},
  tick label style={font=\small},
  ymajorgrids,
  major grid style={draw=figlight},
  axis line style={draw=figgray},
  clip=false,
]
\addplot[fill=figblue, draw=fignavy, thick, bar shift=0pt] coordinates {(1, 20.73)};
\addplot[fill=red!55!black, draw=red!70!black, thick, bar shift=0pt] coordinates {(2, -1.04)};
\node[above, font=\bfseries, text=fignavy] at (axis cs:1, 20.73) {+20.73\%};
\node[below, font=\bfseries, text=red!60!black] at (axis cs:2, -1.04) {-1.04\%};
\end{axis}
\end{tikzpicture}
\end{center}
{\small\color{brandgray}\textbf{図: ルックアヘッドバイアスによるアルファの消失（Look-Ahead-Bench, DeepSeek 3.2）}\ 評価期間をモデルの学習データカットオフより前に設定すると年率アルファ\textbf{+20.73\%}を記録するが、評価期間をカットオフ後に移動するだけで\textbf{-1.04\%}へとほぼ完全に消失する。「学習カットオフ前」バーの高成績は、モデルが事前学習データに含まれる過去の結果を記憶していた人為的成果物である可能性が高いことを示す。（出典: \cite{lookaheadbench2026}）}

\begin{casestudy}{77件のLLMトレーディング研究の系統的レビュー: 再現性の危機}
取引コストを明示的にモデル化していたのはわずか\textbf{1件}、再現可能で時間整合的な評価プロトコルを用いていたのはわずか\textbf{2件}にとどまる。公表された「アウトパフォーム」の大半が方法論的精査に耐えない実態を指摘し、evidence ledger・再現性監査・報告チェックリストを提案している。（出典: \cite{agentictrading2026}）
\end{casestudy}

\begin{itemize}
  \item \textbf{Profit Mirage}: FinLake-Benchによるリーク耐性評価と、RAG＋MCTS＋反実仮想シミュレーションを組み合わせたFactFin学習フレームワークを提案。（\cite{profitmirage2025}）
  \item \textbf{All Leaks Count}: Shapley-DCLRでどの推論クレームがモデルの学習カットオフ後知識を含むかを特定し、TimeSPECによる時間フィルタ付き検索で緩和する手法を提案。（\cite{allleakscount2026}）
\end{itemize}

\begin{insight}{示唆: ベンチマーク上の高成績を額面通りに受け取らない}
上図のLook-Ahead-Bench、KTD-Fin（\S\ref{sec:investment-agents}）、77件のレビューが示す共通のメッセージは、「公表されたLLM投資アルファの多くは、ティッカー名・日付・価格系列を適切に匿名化・時間整合化すると大幅に縮小するか消失する」というものである。導入検討時には、ベンチマークのスコアそのものよりも、評価プロトコルが時間整合性・取引コスト・匿名化をどこまで厳密に扱っているかを精査する必要がある。
\end{insight}

\subsection{他分野に学ぶ: 医療・法律・コード・科学LLMの忠実性研究}
\label{subsec:cross-domain-hallucination}

金融のLLM忠実性・ハルシネーション研究は急速に厚みを増しているが（\S\ref{subsec:reasoning-bench}以降）、単一分野だけを見ていると「これは金融特有の問題か、それともLLMという技術に内在する一般的な限界か」の切り分けを誤りやすい。IT・AI業界全体を見渡すと、医療・法律・コード生成・科学の各分野では、ハルシネーション計測・忠実性評価・ドメイン特化LLMの是非をめぐる議論が金融より数年先行して、しかも訴訟や患者安全といった具体的な実害を伴って進んでいる。本節では、金融関係者が評価手法・ハルシネーション対策・ドメインLLM論争として直接転用できる知見に絞って、他分野の研究を整理する。

\subsubsection{法律LLMのハルシネーション実証: RAGを足しても17〜43\%は残る}

法律ドメインは、金融と同じく「専門文書を正確に読み、実在する根拠だけを引用できるか」が問われる領域であり、しかも実際の訴訟という形で失敗の代償が可視化されている点で金融より研究が先行している。

\begin{casestudy}{Large Legal Fictions: 汎用LLMは検証可能な法律問答で58〜88\%ハルシネートする}
Dahl・Magesh・Suzgun・Ho（スタンフォード RegLab/HAI）は、実在する連邦裁判所判例について「どの巡回区が判断したか」「この先例は今も有効か」といった検証可能な質問をLLMに投げ、回答の正誤を判定した。ハルシネーション率はGPT-4（ChatGPT）で\textbf{58\%}、Llama 2で\textbf{88\%}に達し、しかもモデルはユーザーの誤った法的前提を訂正せず追従する（sycophancy）傾向や、自分がいつハルシネートしているかを自己認識できない傾向を示した。提案された類型（架空の権威の捏造、判旨の誤認、先例の誤適用）は、金融QAで観測される単位誤り（\$150と\$150 millionの混同、\S\ref{sec:llm-analysis}のFAITH）や文書横断統合の失敗と構造的に対応する。（出典: \cite{y05dahl2024legalfictions}）
\end{casestudy}

\begin{casestudy}{Hallucination-Free？: 「ハルシネーションなし」を謳う法律RAGツールも17〜43\%誤る}
LexisNexisのLexis+ AIとThomson ReutersのWestlaw AI-Assisted Research／Ask Practical Law AIは、いずれも「ハルシネーションを回避する」と謳って市場投入された商用法律リサーチツールである。スタンフォードRegLab/HAIの追試では、権威ある判例への検索（RAG）を組み込んでいるにもかかわらず、Lexis+ AIで\textbf{17\%}、Westlaw AI-Assisted Researchで\textbf{33\%}のハルシネーション／誤り率が確認された（比較対象のGPT-4単体は\textbf{43\%}）。エラー類型には「sycophancy」――ユーザーの誤った前提を訂正せず、捏造・誤認した権威から尤もらしい主張を組み立てる挙動――が含まれる。これは、「RAGを足せばハルシネーションは解決する」という前提そのものへの、高リスク専門分野における定量的な反証である。（出典: \cite{y05magesh2025hallucinationfree}）
\end{casestudy}

\begin{casestudy}{Mata v.\ Avianca から1,313件へ: ハルシネーションが法廷リスクとして常態化}
2023年6月、ChatGPTが捏造した架空の判例・引用を含む準備書面を提出した弁護士らに対し、南部ニューヨーク連邦地裁のCastel判事が\textbf{5,000ドル}の制裁金を科した（Mata v.\ Avianca事件）。当時はAIハルシネーションが訴訟に持ち込まれた象徴的な初期事例だったが、Damien Charlotinが継続更新する「AI Hallucination Cases Database」によれば、2026年4月時点でAI生成の捏造内容を含む訴訟手続は\textbf{1,313件}（うち弁護士が関与したもの496件）に達し、報告された制裁金額は\textbf{2023年の5,000ドルから2025年には55,597ドルへ約11倍}に拡大している。単発の事故ではなく、専門文書業務にLLMを組み込む際の恒常的なコンプライアンス・リスクへと性質が変化したことを示す。（出典: \cite{y05mataavianca2023,y05charlotin2026database}）
\end{casestudy}

法律分野は評価手法の設計でも先行している。LegalBenchは法律実務家自身が設計に参画し、司法試験の選択式のような単一の代理指標に頼らず、\textbf{162個}のタスクを6種類の法的推論類型（争点発見・ルール想起・ルール適用・結論導出・解釈・修辞理解）に分解して評価する（NeurIPS 2023）。単一の総合正答率ではなくタスク形状ごとに能力を切り分けるという設計原則は、本節\S\ref{subsec:reasoning-bench}のFinQA／FinanceQA／CFA Level IIIが体現する考え方と同一であり、金融推論ベンチマークの新設・改訂時に参照すべき先行例である\cite{y05legalbench2023}。

\subsubsection{医療LLMの忠実性研究: ドメイン適応の安全性効果と検出の難しさ}

医療は、金融と同様に「専門知識のカバレッジ」と「誤りが実害に直結する」という二重の制約を抱える領域であり、ハルシネーションの\emph{生成}を減らす努力と、ハルシネーションの\emph{検出}という別問題への取り組みの両方で示唆に富む。

\begin{casestudy}{Med-PaLM: ドメイン適応チューニングが「有害な回答」を29.7\%から5.9\%へ削減}
GoogleのMed-PaLM／Med-PaLM 2（Singhal et al., \emph{Nature Medicine}）は、臨床医パネルが科学的事実性・医学的コンセンサスとの整合・推論の質・バイアス・有害となりうる可能性を評価軸として長文回答を採点する枠組みを用いた。汎用の指示チューニングモデル（Flan-PaLM）では長文回答の\textbf{29.7\%}が「潜在的に有害」と判定されたのに対し、医療データでドメイン適応したMed-PaLMでは\textbf{5.9\%}まで低下し、科学的コンセンサスとの整合率も\textbf{61.9\%から92.6\%}へ改善した。さらにMed-PaLM 2は、1,066件の一般消費者向け医療質問のペア比較で、臨床的有用性に関わる9軸中8軸において医師の回答よりも好まれた。精度だけでなく「有害性」を専門家パネルが明示的に採点した数少ない事例であり、金融のモデルリスク管理チームが投資助言・与信判断の出力に対して同様の有害性ルーブリックを設計する際の直接のテンプレートになりうる。（出典: \cite{y05medpalm2_2025}）
\end{casestudy}

\begin{casestudy}{MedHallu: 「ハルシネーションを検出できるか」は「ハルシネーションを避けられるか」より難しい}
MedHalluはPubMedQA由来の1万件のQAペアに、双方向含意クラスタリングで難易度を層別化した（正解に意味的に近いほど検出が難しい）ハルシネーション回答を系統的に付与したベンチマークである（EMNLP 2025）。GPT-4oや医療特化ファインチューン済みモデルUltraMedicalを含む強力なモデルでも、「難」カテゴリのハルシネーション検出F値は\textbf{0.625まで低下}する。一方、回答カテゴリに「不明」という棄権選択肢を加えるだけで、F値・適合率が相対で最大\textbf{38\%}改善した。金融のGBFR（\S\ref{sec:llm-analysis}）が示す「安全な棄権」の価値を、医療分野が独立に裏付けている。（出典: \cite{y05medhallu2025}）
\end{casestudy}

歴史的な警鐘として、Meta AIの科学LLM Galacticaは2022年11月、4,800万本の科学論文で学習し「学術文献の要約・化学物質の注釈付け」を謳って公開されたが、実在しない論文への捏造引用や架空の化学物質を流暢に生成することが露見し、公開デモから\textbf{わずか数日}で撤去された。ベンチマーク上はGPT-3やPaLM 540Bを上回る数値（例: MATH \textbf{20.4\%} 対 PaLM 540Bの8.8\%）を示していたにもかかわらずである\cite{y05galactica2022}。医療分野でも、LLM以前の教訓としてIBM Watson for Oncologyの事例がある。少数の専門医が作成した合成症例に偏った学習データが、狭い専門家集団の癖をあたかも一般的な正解であるかのように埋め込んでしまい、STATの調査報道は内部文書から「安全性を欠き誤った」治療提案が複数あったことを明らかにした。MDアンダーソンがんセンターは関連プロジェクトに約\textbf{6,200万ドル}を投じた末に2016年に中止している\cite{y05watsononcology2018}。いずれも「ベンチマークで高スコアを出すこと」と「本番で信頼できること」は別問題であるという、Fin-R1のような合成CoTデータでのファインチューニング（\S\ref{sec:llm-analysis}）にも通じる警鐘である。

\subsubsection{コード生成の忠実性: 記号検証（実行）ハイブリッドの最も成熟した実例}

コード生成は、金融の記号検証ハイブリッド（\S\ref{sec:llm-analysis}のFinChain・AUDITFLOW・GBFR）と全く同じ設計思想が、より早くより成熟した形で実装されている分野である。

\begin{itemize}
  \item \textbf{CodeHalu（実行検証によるコードハルシネーションの定式化）}: 生成コードを参照解との表面的な類似度で評価するのではなく、実際に\textbf{実行}して挙動の乖離からハルシネーションを判定する枠組み。CodeHaluEvalベンチマーク（699タスク・8,883サンプル）は、マッピング・命名・リソース・ロジックの各ハルシネーション類型を実行結果から機械的に分類する。金融のAUDITFLOW（実行環境での会計演算再計算）・FinChain（実行可能Pythonテンプレートでの中間推論検証）と設計思想が完全に一致する、コード領域における「記号検証層」の最も成熟した実装例である。（出典: \cite{y05codehalu2024}）
  \item \textbf{パッケージ捏造という固有リスク}: 生成コードが提案するパッケージ名のうち約\textbf{2割}が実在しないとの調査もあり、生成テキストのハルシネーションがそのままサプライチェーン攻撃の入口になりうる。実行検証をループに組み込むことでハルシネーション率を最大\textbf{96\%}削減できたとの報告もある。
\end{itemize}

\subsubsection{ドメインLLM vs 汎用LLM＋RAG論争の他分野版: 二択ではなく併用が勝つ}

金融ではBloombergGPT（ドメイン専用事前学習）とFinGPT／FinBERT（安価な適応）、汎用フロンティアLLM＋RAGの間で「どちらに投資すべきか」が論じられてきた（\S\ref{sec:llm-analysis}の概念ボックス）。同じ論争は半導体設計・法律の分野でも独立に起きており、いずれも「二択ではなく併用」という結論に収束している。

\begin{casestudy}{ChipNeMo: ドメイン適応で同品質のまま5分の1のモデルサイズへ}
NVIDIAのChipNeMoは、半導体設計ドメインに特化したトークナイザ・継続事前学習・SFT・ドメイン適応済み検索モデル（RAG）を組み合わせ、エンジニアリング支援チャットボット・EDAスクリプト生成・バグ要約の3用途で評価した。結論は、ドメイン適応技術によって汎用ベースモデルと同等以上の品質を維持したまま\textbf{最大5分の1のモデルサイズ}を実現できるというものである。BloombergGPT／Fin-R1（\S\ref{sec:llm-analysis}）が突きつける「規模かコストか」という問いに、非金融分野から与えられた定量的な回答であり、金融機関のモデル導入における総保有コスト判断に直接転用できる。（出典: \cite{y05chipnemo2023}）
\end{casestudy}

法律AIスタートアップHarveyは、この論争を「どちらか」ではなく「両方」で解決した好例である。信頼できる法源への検索（RAG）で捏造の自由度を絞り込みつつ、その上に汎用フロンティアモデルのファインチューニング・専用埋め込み・ワークフロー統合を重ね、2025年にはOpenAI単独依存からAnthropic Claude・Google Geminiを含むマルチモデル体制へ移行した。10法律事務所での盲検比較では、弁護士はHarveyの回答を素のGPT-4よりほぼ常に選好したと報告されている\cite{y05harvey2025}。「ドメインLLM対汎用LLM＋RAG」という二項対立の設問自体が、本番運用では誤った設問設定になりがちであることを示す。

\begin{center}
\footnotesize
\renewcommand{\arraystretch}{1.3}
\begin{tabular}{@{}p{2.0cm}p{4.0cm}p{3.4cm}p{3.7cm}@{}}
\toprule
\textbf{分野} & \textbf{代表的なハルシネーション率／指標} & \textbf{評価手法の要点} & \textbf{ドメインLLM採否の実態} \\
\midrule
金融 & FAITH: 直接検索95.6\%（Claude-Sonnet-4）→多変数計算ほぼ0\%\cite{x05faith2025}／FinanceBench: 検索併用でも81\%誤答・拒否\cite{financebench2023} & 記号検証（AUDITFLOW/GBFR/FinChain）＋推論アリティ別ルーティング\cite{x05finchain2025,x05gbfr2026} & BloombergGPT／FinGPT等の専用モデルと汎用LLM＋RAGが併存、タスク形状で使い分け \\
医療 & Flan-PaLM有害回答率29.7\%→Med-PaLM 5.9\%\cite{y05medpalm2_2025}／MedHallu検出F値「難」で0.625\cite{y05medhallu2025} & 臨床医パネルによる有害性・コンセンサス整合ルーブリック＋棄権カテゴリの導入 & Med-PaLM系はドメイン適応で成功、Watson for Oncologyは狭いデータで失敗\cite{y05watsononcology2018} \\
法律 & 汎用LLM 58\%（GPT-4）〜88\%（Llama 2）\cite{y05dahl2024legalfictions}／RAG併用ツールでも17〜43\%\cite{y05magesh2025hallucinationfree} & LegalBenchが162タスク・6類型に分解\cite{y05legalbench2023}、訴訟データベースで実害を定量化\cite{y05charlotin2026database} & Harveyはドメイン適応＋RAG＋マルチモデルの併用型\cite{y05harvey2025} \\
コード & パッケージ名の約2割が実在しない捏造 & CodeHaluEvalが実行結果からハルシネーション類型を機械判定\cite{y05codehalu2024} & ChipNeMoはドメイン適応で同品質のまま5分の1のサイズを達成\cite{y05chipnemo2023} \\
科学 & Galacticaは公開後数日で撤去\cite{y05galactica2022} & ベンチマークスコアと実運用信頼性の乖離が典型例として定着 & 汎用モデルへの安全弁なき特化は失敗、後続はドメイン適応＋対話的検証へ移行 \\
\bottomrule
\end{tabular}
\end{center}
{\small\color{brandgray}\textbf{表: ハルシネーション・評価手法・ドメインLLM採否の分野横断比較}\ 金融の記号検証ハイブリッド・推論アリティ別ルーティングという設計は、医療の棄権カテゴリ、法律のタスク分解型ベンチマーク、コードの実行検証と同じ「LLMを読解器、外部機構を検証器に分離する」思想の一変種であることが分野横断で確認できる。}

\begin{casestudy}{RMCB: 金融・医療・法律・数学を同一基準で測る、JPMorgan AI Researchも参加した信頼度較正ベンチマーク}
Reasoning Model Confidence estimation Benchmark（RMCB, EACL 2026）は、臨床・金融・法律・数学の4つの高リスク推論領域と一般推論を横断し、6種の大規模推論モデルから収集した\textbf{347,496件}の推論トレースに正誤ラベルを付与した、数少ない「分野をまたいで同一の物差しで測る」信頼度較正ベンチマークである。10種類以上の信頼度推定手法（系列型・グラフ型・テキスト型）を評価した結果、判別力（AUROC）で最良のテキストベース手法（\textbf{0.672}）と較正誤差（ECE）で最良のグラフ構造ベース手法（\textbf{0.148}）は一致せず、\textbf{両方を同時に達成する単一の手法は存在しない}。しかも著者にはJPMorgan AI Research所属のSimerjot KaurとCharese H. Smiley（FinQA／ConvFinQA\cite{x05finqa2021,x05convfinqa2022}の共著者でもある）が含まれており、金融業界の研究者自身が「LLMが自分の誤りにいつ気づいていないか」は金融固有の課題ではなく分野横断の未解決問題であると認めていることを意味する。（出典: \cite{y05rmcb2026}）
\end{casestudy}

\begin{insight}{金融への示唆: 「ハルシネーション対策＝金融データを増やす」ではない}
法律・医療・コードの事例が共通して示すのは、ハルシネーション対策の成否を分けるのはドメインデータの量ではなく、\textbf{（1）検証を独立した外部機構に分離する設計}（法律の判例検索＋人間照合、医療の棄権カテゴリ、コードの実行検証、金融のAUDITFLOW／GBFR）と、\textbf{（2）タスクを能力の細粒度に分解して評価する設計}（LegalBenchの162タスク6類型、金融のFinQA〜FCMRの推論アリティ別整理）である、という点である。RAGを足すだけでは17〜43\%の誤りが残るというLexis+ AI／Westlawの実証は、金融RAGシステム（FinanceBench・AUDITFLOW）についても「検索を足したから安全」という単純化を戒める直接の反証材料として引用できる。
\end{insight}

\begin{insight}{金融への示唆: ドメインLLMの是非は「規模か適応か」ではなく「検証層があるか」で決まる}
ChipNeMoの5分の1サイズ、Harveyのドメイン適応＋RAG＋マルチモデル併用、Med-PaLMの臨床医ルーブリックによる安全性検証は、いずれも「ドメインLLMか汎用LLM＋RAGか」という二項対立ではなく、外部検証層の有無・質が成否を分けることを示している。BloombergGPT／Fin-R1（\S\ref{sec:llm-analysis}）の位置づけも同様に再解釈すべきであり、調達判断は「どちらのアーキテクチャが賢いか」ではなく「検証層と棄権機構をどれだけ厳密に監査可能な形で実装できるか」に軸足を置くべきである。
\end{insight}

\begin{insight}{金融への示唆: 訴訟という現物のフィードバックループが金融には（まだ）ない}
法律分野には、Mata v.\ Avianca以降1,313件・11倍の制裁金増という、ハルシネーションの実害を定量化する外部フィードバック（訴訟記録）が存在する。金融のLLM出力の誤りは、開示文書の誤読や与信判断の誤りとして同様の実害を生みうるにもかかわらず、法律ほど体系的に追跡・公開される仕組みを欠く。RMCBにJPMorgan AI Researchが参加している事実は、金融業界内部でもこの較正・検証ギャップへの危機感が既に生まれていることを示しており、金融規制・業界団体（\S\ref{sec:cross-cutting}）が法律分野の訴訟データベースに相当する「AI起因インシデントの業界横断台帳」を持つべきかどうかは、未解決ギャップ（\S\ref{sec:open-gaps}）として検討に値する。
\end{insight}

\section{横断的な考察と示唆}
\label{sec:cross-cutting}

個別の技術・事例を超えて、複数の独立した研究・企業動向が同じ構造的パターンを指し示している。本節ではそれらを4つの視点――価値の所在、コスト構造、システミックリスク、規制と人材――から統合する。データ接続そのものの価値が失われるほどコスト構造の巧拙が問われ、コスト最適化が進むほど自律性の広がりがシステミックリスクという別の軸で再評価され、そのリスクへの対応速度は規制当局と人材市場の反応速度に規定される、という一本の因果の鎖として読むことができる。

\subsection{データ接続のコモディティ化と価値の移動}

\S\ref{sec:data-infra}で見たとおり、Bloomberg・LSEG・FactSet・DaloopaがそろってMCPサーバーを提供し始めたことで、構造化金融データへの標準化された接続そのものは急速にコモディティ化しつつある。この動きは、投資AI市場における価値の所在を再定義しつつある。

\begin{center}
\begin{tikzpicture}[
  layer/.style={fnode, text width=9.6cm, minimum height=1.5cm},
  toplayer/.style={fnode accent, text width=9.6cm, minimum height=1.7cm},
  arr/.style={farr blue}
]
\node[layer] (l1) at (0,0)
  {\textbf{データ接続層（コモディティ化）}\\Bloomberg・LSEG・FactSet・DaloopaがMCPサーバーを提供、標準化された接続の限界価値はゼロへ近づく};
\node[layer] (l2) at (0,2.15)
  {\textbf{抽出・構造化パイプライン}\\プロンプト設計・ドメインファインチューニング・出力正規化};
\node[toplayer] (l3) at (0,4.55)
  {\textbf{垂直特化型プラットフォーム}\\Hebbia: 年間処理4,700万→10億ページ（約2,000\%成長）\\仮説: 水平的LLM API統合比 1席あたり3〜5倍の収益};

\draw[arr] (l1.north) -- (l2.south);
\draw[arr] (l2.north) -- (l3.south);

\draw[farr accent, ultra thick] (5.4,-0.75) -- (5.4,5.3);
\node[rotate=90, anchor=south, font=\scriptsize, color=figblue] at (5.4,2.3) {価値の移動};
\end{tikzpicture}
\end{center}
{\small\color{brandgray}\textbf{図: データ接続のコモディティ化と価値の移動}\ 標準化された生データ接続の限界価値がゼロへ近づく一方、抽出・構造化を経て、ドメイン特化・監査可能性を備えた垂直プラットフォーム層に価値が移動しつつある。}

\begin{insight}{仮説: 垂直特化型プラットフォームの収益優位}
Hebbia（Matrix）は年間処理ページ数を4,700万から10億へ拡大させ（約2,000\%成長）、KKR・MetLife・Centerview Partnersを顧客に、PEデューデリジェンス・M\&Aで案件あたり20〜40時間の削減を実現している。Morgan StanleyはGPT-4＋Whisperの社内展開に深いカスタム統合を要したが、アドバイザー採用率98\%超、週15時間の削減という具体的なROIを達成した。この対比から、「垂直特化型金融AIプラットフォームは、水平的なLLM API統合よりもエンタープライズ1席あたり3〜5倍高い収益を稼ぐ」という仮説が提示されている。オーケストレーション・ドメイン特化文書解析・監査可能性・コンプライアンス統制が、生モデルへのアクセスよりも持続的な差別化要因になりつつあるという見立てである。一方で、大規模機関（Morgan Stanley、JPMorgan規模）は内製化によって垂直プラットフォームの機能を再現しうる点、Microsoft CopilotやSalesforceのバンドルが配荷面で優位に立ちうる点が反証材料として挙げられている。
\end{insight}

代替データ市場も同様の力学に直面していると業界筋では語られている。市場規模は年間70億ドル、投資運用会社の67\%が既に代替データを利用しているといった数字が広く引用されているが、本書が精査した一次資料の範囲ではこれらの数値の直接的な典拠を確認できておらず、仮説的な参考値として扱う必要がある。ただし「生データの独占期間が数年から数カ月へ短縮している」という方向性自体は、上記のMCP標準化の趨勢と整合的である。今後のアルファ創出は、データアクセスそのものよりも、パイプラインの抽出・構造化能力（プロンプト設計、ドメインファインチューニング、出力正規化）に依存するという見方が強まっている。データそのものの希少性が失われるほど、次に問われるのはそのデータを処理するアーキテクチャのコスト構造である。

\begin{insight}{示唆: 標準化は「接続」から「協調・権限・監査」へ拡張している}
2025〜2026年の新しい標準化の焦点は、単に金融データへ接続するMCPだけではない。OpenAI Responses APIのremote MCP対応はモデルから外部データ・ツールへの接続を製品機能に組み込み\cite{x01openai2025responsesmcp}、Google A2Aは異なるベンダーのエージェント間協調を標準化し\cite{x01google2025a2a}、AP2はエージェントがユーザーの委任に基づいて支払いを起票・実行する際の権限証跡を標準化しようとしている\cite{x01google2025ap2}。Linux FoundationでA2Aが150超の組織へ広がったこと\cite{x01linuxfoundation2026a2a}を踏まえると、金融AIの差別化は「どのデータにつながるか」だけでなく、「どのエージェントと協調し、どの権限で行動し、その証跡をどう監査するか」へ移動している。
\end{insight}

\subsection{コストとアーキテクチャの経済学}

垂直プラットフォームの優位性が処理能力の巧拙に由来するとすれば、その巧拙は最終的にアーキテクチャ選択のコスト対効果として定量化できる。

\begin{casestudy}{階層型 vs 反射型アーキテクチャ: 精度とコストのパレートフロンティア}
SECファイリング抽出のベンチマークで、反射型自己修正アーキテクチャ（Claude 3.5 Sonnet）はF1 0.943を達成する一方、文書あたり\textbf{\$0.430}のコストを要した。階層型supervisor-workerアーキテクチャはF1 0.929（反射型の効果の89\%）を\textbf{\$0.261}（60\%のコスト）で達成し、セマンティックキャッシュ・モデルルーティング・適応的リトライを加えた最適化版はF1 0.924を\textbf{\$0.148}/文書で実現した。100万文書規模の処理では\$430,000対\$148,000、\textbf{\$282,000}のコスト差に相当する。この差はリーダーボード上の精度比較には現れない。（出典: \cite{benchmarkingmultiagent2026}）
\end{casestudy}

3つのアーキテクチャを精度とコストの2軸に置くと、パレートフロンティアの形状がこの非対称性をより直接的に示す。反射型からのわずかな精度改善（+0.019）を得るために、コストはほぼ3倍に跳ね上がっている。

\begin{center}
\begin{tikzpicture}
\begin{axis}[
    width=0.82\textwidth,
    height=6.2cm,
    xlabel={コスト（\$／文書）},
    ylabel={F1スコア},
    xlabel style={font=\small},
    ylabel style={font=\small},
    x tick label style={font=\small},
    y tick label style={font=\small},
    xmin=0.05, xmax=0.52,
    ymin=0.900, ymax=0.960,
    ymajorgrids=true, xmajorgrids=true,
    major grid style={draw=figlight},
    axis line style={draw=figgray},
    tick align=outside,
    legend style={at={(0.5,-0.30)}, anchor=north, legend columns=2, font=\small, draw=none},
]
\addplot[figblue, thick, dashed, mark=none] coordinates {(0.148,0.924) (0.261,0.929) (0.430,0.943)};
\addlegendentry{パレートフロンティア}
\addplot[only marks, mark=*, mark size=2.6pt, color=fignavy] coordinates {(0.148,0.924) (0.261,0.929) (0.430,0.943)};
\addlegendentry{アーキテクチャ}
\node[anchor=north, font=\scriptsize, text=fignavy] at (axis cs:0.148,0.9215) {最適化階層型（\$0.148）};
\node[anchor=south, font=\scriptsize, text=fignavy] at (axis cs:0.261,0.9315) {階層型supervisor-worker（\$0.261）};
\node[anchor=south, font=\scriptsize, text=fignavy] at (axis cs:0.430,0.9455) {反射型自己修正（\$0.430）};
\end{axis}
\end{tikzpicture}
\end{center}
{\small \textbf{図: 精度とコストのパレートフロンティア（SECファイリング抽出ベンチマーク）。}反射型自己修正アーキテクチャはF1 0.943を\$0.430/文書で達成する一方、最適化された階層型supervisor-workerアーキテクチャはF1 0.924を\$0.148/文書、コスト約3分の1で実現する。精度の限界的な改善に対しコストが不釣り合いに増加する関係が、リーダーボード上の精度比較だけでは見えない調達判断の落とし穴を可視化している。（出典: \cite{benchmarkingmultiagent2026}）}

同様の経済性の議論はインフラ選択にも及ぶ。業界筋の試算では、オンプレミスLLMは2025年時点でAPIよりトークンあたり約18倍安価だったが、直近でAPI価格が約80\%下落し単価優位は縮小しているとされる。この単価比較そのものは本書が精査した一次資料では裏付けが取れておらず、仮説的な参考値として扱う必要があるが、DORA（2025年1月施行）、GDPR第46条、EU AI ActのGPAI義務がデータ主権要件を恒常化させていることは制度的な事実であり（\cite{eba2025aiact}）、コストではなくコンプライアンスが今後のオンプレミス投資判断の主因になるという仮説が提示されている。コスト最適化それ自体が新たなガバナンス上の論点を生む――アーキテクチャの経済合理性は、次に見る自律性の広がりというもう一つの軸で再評価されることになる。

\subsection{システミックリスクとガバナンス}

AIエージェントが投資判断・執行を担う範囲が広がるにつれ、個社の性能評価だけでなく、市場全体への波及効果を問う研究が増えている。コスト最適化の帰結として広く配備されたアーキテクチャが、今度は市場構造そのものに影響を及ぼしうるという逆説である。

\begin{casestudy}{AIモノカルチャーのシステミックリスク}
performative prediction（予測が市場自体を変化させる）、algorithmic herding（アルゴリズム的群集行動）、cognitive dependency（認知的依存）の3チャネルを統合したモデルを、2013〜2024年の10,957マネジャー・9,950万件のSEC 13Fホールディングスで較正した研究では、AI浸透度が高いほどテールロス増幅が\textbf{18〜54\%}、市場の厚み（depth）が導入シェアの増加に伴い凸的に減少することが確認された。（出典: \cite{aisystemicrisk2026}）
\end{casestudy}

\begin{center}
\begin{tikzpicture}[
  node distance=6mm and 5mm,
  main/.style={fnode, text width=6.2cm, minimum height=1.3cm},
  channel/.style={fnode blue, text width=3.5cm, minimum height=1.7cm, font=\footnotesize},
  outcome/.style={fnode accent, text width=8cm, minimum height=1.5cm},
  arr/.style={farr blue}
]
\node[main] (start) {\textbf{AI浸透度の上昇}\\（同一基盤モデルへの依存拡大）};
\node[channel, below=of start, xshift=-4.0cm] (perf) {performative prediction\\（予測が市場自体を変化させる）};
\node[channel, below=of start] (herd) {algorithmic herding\\（アルゴリズム的群集行動）};
\node[channel, below=of start, xshift=4.0cm] (cog) {cognitive dependency\\（認知的依存）};
\node[outcome, below=13mm of herd] (out)
  {\textbf{システミックリスクの増幅}\\テールロス増幅 \textbf{18〜54\%}／市場の厚み（depth）は導入シェア増加に伴い凸的に減少};

\draw[arr] (start) -- (perf);
\draw[arr] (start) -- (herd);
\draw[arr] (start) -- (cog);
\draw[arr] (perf) -- (out);
\draw[arr] (herd) -- (out);
\draw[arr] (cog) -- (out);

\draw[farr dash, rounded corners=8pt] (out.east) -- ++(4,0) |- (start.east);
\node[anchor=west, font=\scriptsize, color=figgray, text width=2.6cm, align=left]
  at (8.3,-2.5)
  {同一モデル更新が数千のエージェントへ同時到達（LLM Output Drift）→相関的ポジショニング};
\end{tikzpicture}
\end{center}
{\small\color{brandgray}\textbf{図: AI浸透によるシステミックリスク増幅ループ}\ performative prediction・algorithmic herding・cognitive dependencyの3チャネルを通じ、AI浸透度の上昇がテールロス増幅・市場の厚みの低下に波及する。同一基盤モデルの更新が数千のエージェントへ同時到達すると、相関的ポジショニングを介してさらなる浸透へフィードバックしうる。（出典: \cite{aisystemicrisk2026}）}

Hui Gong氏の四層エージェントアーキテクチャ論（データ知覚→推論→戦略→執行）とAgent-Financial-Market Modelは、システミックリスクが生のAIモデル能力そのものよりも、エージェントの自律性の深さ・執行の結合度・監督可能性が機関間でどう分布・ガバナンスされるかに依存すると主張している（\cite{aiagentsfinmarkets2026}）。同時に、LLM Output Drift研究は、モデル更新がプロバイダ間で顕著な行動の不整合を金融ワークフローにもたらすことを示している。業界調査では2026年第1四半期時点で中堅〜大手ヘッジファンドの47\%超が本番環境で生成AIを少なくとも1つ導入済みとされるが、この普及率の数値も本書が確認した一次資料では追跡できておらず、仮説的な参考値にとどまる。それでもなお方向性としては、同一の基盤モデル更新が数千のアナリスト・自律トレーディングエージェントに同時到達した場合、相関的なポジショニング変化という新たなシステミックチャネルが生じうるという仮説につながる。

決済インフラという別の切り口からも、同様のガバナンス上の緊張が指摘されている。

\begin{casestudy}{IMF: エージェント型AIによる決済インフラの再編}
IMFは2026年4月のノートで、決済領域におけるエージェント型AIに対する\textbf{三層のガバナンス枠組み}を提唱した。その核心は、確率的に意思決定するAIと、認可（authorization）・決済（settlement）・コンプライアンスにおいて決定論的な正確性を要求する決済インフラとの間の根本的な緊張にある。エージェントが自律的に支払いを起票・実行する範囲が広がるほど、「おおむね正しい」確率的判断と「常に正しい」ことを前提に設計された台帳・清算システムとの整合をどう取るかが、システミックな論点として立ち上がる。（出典: IMF \cite{imfpayments2026}）
\end{casestudy}

金融AIの信頼性そのものを攻撃面から体系化する動きも進んでいる。Zengら（2026）のサーベイは、金融AIライフサイクル全体の敵対的攻撃面――データ・モデルポイズニング、意思決定境界における敵対的摂動、LLM介在ワークフローへのプロンプトインジェクション、ディープフェイクによるKYC突破――を整理し、規制下の金融機関での本番展開を志向した初のセキュリティタクソノミを提示している（\cite{trustworthyfintech2026}）。前述のLLM Output Driftと同様、これらは同一の脆弱性・攻撃が多数のエージェントへ同時に波及しうる構造的リスクとして、システミックリスクとガバナンスの議論に組み込む必要がある。こうしたシステミックリスクへの懸念は、当然ながら規制当局がどの速度・どの粒度で対応するかという制度的な変数と切り離せない。

\subsection{規制の非対称な進展}

規制は地域・分野によって非対称に進展しており、これが実装スピードとコンプライアンスコストの両方を左右する変数になりつつある。

\begin{itemize}
  \item \textbf{EU AI Act}: 与信スコアリングAIと保険リスク価格算定AIはAnnex III上の高リスクシステムに明示的に分類され、非準拠には最大\textbf{3,000万ユーロまたは世界売上高の6\%}の罰金が科される。完全適用開始期限は当初2026年8月2日とされていたが、Digital Omnibus（欧州議会承認2026年6月16日）により信用・保険AI義務の拘束的期限は\textbf{2027年12月2日}へ約16カ月延期された\cite{x07digitalomnibus2026}。期限の後ろ倒しは準備時間を与える一方、規制不確実性を延伸させ、欧州大手金融機関にAIベンダー監査を義務付ける調達力学が生じる時期を後ろにずらす（詳細は\S\ref{sec:open-gaps}）。（出典: EBA \cite{eba2025aiact}、\cite{x07digitalomnibus2026}）
  \item \textbf{SR 11-7の置き換え}: 米連邦銀行監督当局の画一的なモデル検証枠組みSR 11-7は2026年4月17日にリスク階層型・原則ベースのガイダンスへ置き換えられ、低重要度モデルへの比例的統制を導入した。業界関係者は、従来の曖昧さこそがLLMコパイロットの本番投入を躊躇させてきた主因だったと証言している。（出典: \cite{databricksndmodelrisk}）
  \item \textbf{シャドーAIによる初のSEC Form 8-K開示}: 2026年5月、CB Financial Servicesが従業員による未承認AIアプリ経由の非公開顧客データ流出を理由に、外部攻撃者ではなく内部のシャドーAI利用でItem 1.05に基づく初のForm 8-Kを提出した。シャドーAI誤用が技術ポリシー問題から取締役会レベルの開示義務へ格上げされたことを示す。（出典: \cite{wsgrndshadowai}）
  \item \textbf{M\&Aデューデリジェンスへの浸透}: Bainの2026年M\&Aレポートによれば、戦略的ディールメーカーの\textbf{5件に1件}がAI関連懸念により案件から撤退し、\textbf{58\%のM\&Aデューデリジェンスワークフロー}が既に生成AIをこれらの懸念の発見に使用しているとされる（この数値は本書が確認した一次資料への遡及ができておらず、仮説的な参考値として扱う）。
  \item \textbf{データ主権・GDPR対応}: 合成金融QAデータで学習した小型モデルが専門家（CFA/CPA水準）注釈データ学習モデルと同等性能を達成しうるという研究成果（ACL 2025 FinNLPワークショップ）は、GDPR・MiFID II・Basel IIIのデータ居住性制約下で金融機関が域内でファインチューニングを行う経路として位置づけられている（\cite{aclfinnlp2025synthetic}）。ただし合成データが生成元モデルのバイアス・誤りをそのまま継承しうる点、ECB・FCA・Fedのモデルリスク管理ガイダンスが人間キュレーション済みデータの追跡可能性を要求しうる点が留保されている。
\end{itemize}

3つの規制イベントを時系列に並べると、それぞれ異なる法域・異なる主体（米連邦銀行監督当局、SEC開示制度、EU AI Act）が数カ月の間隔でほぼ独立に動いていることが分かる。これは各国・各規制体系が足並みを揃えて進んでいるのではなく、非同期に規制の網が狭まっていることを意味し、グローバルに事業を展開する金融機関ほど複数の規制サイクルへ同時対応するコストを負うことになる。

\begin{center}
\begin{tikzpicture}[
  every node/.style={font=\small},
  axisline/.style={farr, draw=fignavy, thick},
  evtdot/.style={circle, fill=figaccent, draw=fignavy, thick, minimum size=6pt, inner sep=0pt, figshadow}
]
\draw[axisline] (0,0) -- (10.6,0);
\node[anchor=west, font=\scriptsize, color=figgray] at (10.6,0.05) {2026→27年};

\coordinate (sr117) at (3.4,0);
\draw[fignavy, thick] (sr117) ++(0,-0.12) -- ++(0,0.12);
\node[evtdot] at (sr117) {};
\node[align=center, text width=3.1cm, anchor=south, font=\scriptsize] at ($(sr117)+(0,0.6)$) {\textbf{2026-04-17}\\SR 11-7 がリスク階層型ガイダンスへ置換};

\coordinate (shadowai) at (4.6,0);
\draw[fignavy, thick] (shadowai) ++(0,-0.12) -- ++(0,0.12);
\node[evtdot] at (shadowai) {};
\node[align=center, text width=3.1cm, anchor=north, font=\scriptsize] at ($(shadowai)+(0,-0.6)$) {\textbf{2026-05}\\シャドーAI起因の初のForm 8-K（CB Financial Services）};

\coordinate (euaiact) at (7.9,0);
\draw[fignavy, thick] (euaiact) ++(0,-0.12) -- ++(0,0.12);
\node[evtdot] at (euaiact) {};
\node[align=center, text width=3.4cm, anchor=south, font=\scriptsize] at ($(euaiact)+(0,0.6)$) {\textbf{2026-08-02→2027-12-02}\\EU AI Act 高リスク分類（信用・保険AI）の完全適用がDigital Omnibusで延期};
\end{tikzpicture}
\end{center}
{\small \textbf{図: 規制の非同期な進展タイムライン（2026年）。}米国のモデルリスク管理枠組みの刷新（SR 11-7置換）、シャドーAI利用に対する開示義務の実例化（初のForm 8-K）、EU AI Actの高リスク分類の完全適用（Digital Omnibusにより2026-08-02から2027-12-02へ延期）が、異なる法域・異なる主体からほぼ独立に到来している。規制は足並みを揃えるのではなく、非同期に――かつ期限自体も動きながら――網を狭めている。（出典: \cite{databricksndmodelrisk}、\cite{wsgrndshadowai}、\cite{eba2025aiact}、\cite{x07digitalomnibus2026}）}

\begin{insight}{示唆: エンタープライズGenAI導入の高い失敗率}
2025年に金融業界のエンタープライズGenAIパイロットの\textbf{95\%}がP\&Lへの測定可能な影響を出せなかったとの調査があり、AIプロジェクト時間の\textbf{60〜70\%}がデータ準備に費やされているとの別の業界調査もある。これは\S\ref{sec:data-infra}で確認したデータ基盤の未成熟さと整合的であり、モデル選定よりもデータ・前処理層への投資が投資対効果を左右することを重ねて示唆している。ただしこの数値は仮説的考察（Takes）の中で引用されたものであり、一次の実証研究として確認されたものではない点に留意が必要である。
\end{insight}

\begin{insight}{示唆: 精度・コンプライアンス・説明可能性のトリレンマ}
Eviteら（2026）は20名の金融プロフェッショナルへのインタビューに基づき、従来の「精度対解釈可能性」という二項対立の枠組みを退け、\textbf{精度・コンプライアンス・説明可能性}の三者からなるトリレンマを提示した。このうち精度とコンプライアンスは交渉の余地のない衛生要因（hygiene factor）であり、両者を満たすことは前提条件にすぎない。実際の導入可否を左右するゲートウェイは残る一つ――説明可能性（ease of understanding）――だという。上で見た規制の非対称な進展が精度・コンプライアンスを「守って当然」の下限として恒常化させるほど、金融機関の差別化と採用の分かれ目は説明可能性の設計に移っていくという含意である。（出典: \cite{fintrilemma2026}）
\end{insight}

\begin{insight}{仮説: エージェント型AIによる継続的資本ストレステスト（Takes）}
日次・リアルタイムの資本ストレステストを担うエージェント型AIフレームワークは、四半期ごとのDFAST/CCARの結果に\textbf{4〜6週間先行}して資本ストレスの兆候を検知できる、という仮説がある。HSBC・Citi・UBS・DBS・INGなど複数の大手行はリスク機能へのAIエージェント導入で20〜40\%のコスト削減を報告しており、マクロシナリオ生成・資本シミュレーション・モデル検証・監査証跡の維持を継続的にオーケストレーションする能力面のボトルネックはすでに小さいという見立てである。真の制約はむしろ規制当局の事前承認にあり、四半期バッチ提出と人間による検証（human-attested models）を前提に書かれた既存のDFAST/CCARガイダンスが、継続的に更新されるエージェント生成出力を想定していない点にあるとされる。ただし正常な市場変動下での偽陽性がアラート疲れを招き、規制の門戸が開く前に技術が放棄されるという反証リスクも指摘されている。これは仮説的考察（Takes）であり、実証されたものではない。
\end{insight}

規制・コンプライアンス要件の高度化とパイロット導入の高い失敗率は、最終的に金融機関がAI活用の担い手にどのようなスキルセットを求めるかという人材戦略にも波及する。

\subsection{人材と労働市場への影響}

\begin{itemize}
  \item PwCの2026年AI Jobs Barometerによれば、金融アナリストはAIへの露出度\textbf{50\%}だが適応力は\textbf{99\%}とリスク職種中最高水準の耐性を示す一方、AIを導入した財務部門の\textbf{84\%}は職務再設計を行っていない（この調査結果は本書が確認した一次資料への遡及ができておらず、仮説的な参考値として扱う）。カバレッジ拡大（1アナリストあたりの監視対象企業・セクター・リスク要因を増やす）という活用余地が構造的に未活用であることを示唆する。
  \item \S\ref{sec:ai-native-funds}で見たとおり、Point72/Cubist（Quant Academy・複数拠点のGenAIハッカソン）、Man AHL（Oxford-Man Instituteとの連携で2016年以降19名のML研究者を追加）、QRT（2026年設立のQRT Labsで英3大学に70名超の若手研究者へ資金提供）など、大手クオンツファンドは社内研修と大学連携の両輪でAI人材パイプラインを組織的に構築しつつあり、これが業界標準になりつつある。
\end{itemize}

\begin{center}
\begin{tikzpicture}
\begin{axis}[
    xbar,
    width=0.85\textwidth,
    height=4.6cm,
    bar width=7mm,
    xmin=0, xmax=108,
    xlabel={\%},
    xlabel style={font=\small},
    ytick={1,2,3},
    yticklabels={職務再設計未実施の財務部門,適応力,AIへの露出度},
    ymin=0.4, ymax=3.6,
    y tick label style={font=\small},
    x tick label style={font=\small},
    nodes near coords,
    nodes near coords style={font=\small, color=fignavy},
    axis line style={draw=figgray},
    tick align=outside,
    ymajorgrids=false, xmajorgrids=true,
    major grid style={draw=figlight},
]
\addplot[fill=figblue, draw=fignavy, bar shift=0pt] coordinates {(50,3) (99,2)};
\addplot[fill=figaccent, draw=fignavy, bar shift=0pt] coordinates {(84,1)};
\end{axis}
\end{tikzpicture}
\end{center}
{\small\color{brandgray}\textbf{図: 金融アナリストのAI露出度・適応力とその活用の遅れ}\ PwC 2026 AI Jobs Barometerによれば、金融アナリストはAIへの露出度50\%に対し適応力99\%とリスク職種中最高水準の耐性を示す一方、AIを導入した財務部門の84\%は職務再設計を行っていない――カバレッジ拡大の余地が構造的に未活用であることを示す（仮説的な参考値）。}

\subsection{周辺分野との対比が示す普遍性と金融の特異性}
\label{sec:cross-cutting-crossdomain}

本書の各本論部には、金融以外の分野で同型の技術・課題がどこまで進んでいるかを対比する記述を加えた（\S\ref{sec:intro-cross-domain}、\S\ref{sec:data-infra}、\S\ref{sec:ai-native-funds-crossdomain}、\S\ref{sec:investment-agents-crossdomain}、\S\ref{sec:llm-analysis}、\S\ref{sec:cross-domain-gaps}）。それらを束ねると、金融AIのスタックの大半は業界横断で共有される「普遍的な部分」と、金融にしか存在しない「特異な部分」に分離できる。読者が金融のプロである以上、周辺分野の事例は網羅ではなく、\emph{借用できる設計}か\emph{先取りすべきリスク}のいずれかに翻訳できるものだけを残した。

\begin{center}
\small
\begin{tabularx}{\textwidth}{@{}>{\raggedright\arraybackslash}p{3.0cm} >{\raggedright\arraybackslash}X >{\raggedright\arraybackslash}X@{}}
\toprule
\textbf{レイヤー／課題} & \textbf{周辺分野の先行例} & \textbf{金融への含意（借用か警戒か）} \\
\midrule
データ基盤・文書AI &
リーガル契約解析（CUAD\cite{y02cuad2021}）、GraphRAG\cite{y02graphrag2024}、エンタープライズ検索（Glean\cite{y02glean2025}） &
階層型・グラフ型検索は\textbf{借用可}。ただしネスト表の会計数値照合に相当する検証課題は他分野に先行例が乏しい。 \\
自律エージェント評価 &
SWE-bench Verified\cite{y04swebenchverified2024}等による厳密な閉検証と自律度の段階付け &
DD・執行への「証拠の階梯（evidence ladder）」設計を\textbf{借用可}。 \\
LLMの忠実性 &
医療（Med-PaLM 2\cite{y05medpalm2_2025}）・法律（法的ハルシネーションの実証\cite{y05dahl2024legalfictions}）の忠実性研究 &
評価手法・棄権設計は\textbf{借用可}。ただし法律のような訴訟起因の是正ループは金融に薄く、自前の検証層で代替する必要。 \\
組織のAI-native化 &
ソフトウェア開発・コンサル・ERPの中央プラットフォーム型\cite{y03mckinsey2026sevenoperatingtruths} &
人材パイプラインとbuild-vs-buyの型は\textbf{転用可}。 \\
汎用AIの未解決課題 &
評価クライシス、algorithmic monoculture\cite{y07monoculture2022}、水平的規制\cite{y07euaiacthorizontal} &
金融固有ではないが、実資本・執行権限ゆえ金融で最も\textbf{先鋭化（警戒）}。 \\
\bottomrule
\end{tabularx}
\end{center}
{\small\color{brandgray}\textbf{表: 周辺分野で先行する「普遍的な部分」と、金融に固有な「特異な部分」}\ 各本論部の対比記述を横断的に整理した。技術スタックの多くは他分野から借用できる一方、会計数値の記号的照合と規制密度に由来する要件は金融に固有である。}

\begin{insight}{示唆: 金融は「速い追随者」であると同時に「最速のストレステスト領域」}
周辺分野との対比は、金融が二重の非対称なポジションにあることを示す。技術面では、文書AI・自律エージェント評価・LLM忠実性のいずれも金融より成熟した先行分野が存在し、金融は\textbf{速い追随者}として設計・評価手法を借用できる立場にある。一方でガバナンス面では、実資本の毀損・執行/決済権限・13F型の公開可測性・既存の規制密度ゆえに、評価汚染・エージェント安全性・モデル集中といった\emph{汎用AIの未解決課題}が他業界に先んじて金融で顕在化する。したがって合理的な戦略は、借りられる技術は周辺分野の成熟を待って借用しつつ、ガバナンス・検証層では自ら先頭で問題に直面することを織り込む――という層ごとに異なる時間軸の使い分けになる。ただし「数値の記号的検証」だけは他分野に十分な先行例がなく、金融が自前で確立せざるを得ない特異点である点に留意が必要である。
\end{insight}

コストの経済学、システミックリスク、規制、人材という4つの視点は独立した論点ではなく、データ接続のコモディティ化を起点に相互に波及する一つの構造を成している。加えて、上で見た周辺分野との対比は、この構造の大半が金融固有ではなく業界横断で共有されること、そして金融の特異性が数値検証と規制密度という2点に凝縮されることを示している。次節（第\ref{sec:open-gaps}部）では、この構造の中でなお実証が追いついていない研究・実装上のギャップを整理する。

\section{未解決の課題（Open Gaps）}
\label{sec:open-gaps}

本節では、精査した一次情報が明示的に指摘している、現時点で未解決の研究・実装ギャップを整理する。ここに挙げるのは著者チームの推測ではなく、各論文・ベンチマーク報告が自らの限界として記載した内容――「本研究は〜を対象としない」「〜の標準はまだ存在しない」といった記述――を横断的に収集し、4つのカテゴリに再構成したものである。すなわち本節自体が、一次資料への遡及可能性を優先する本書の編集方針の延長線上にある。これらは今後の投資・研究の優先順位を検討するうえでの参照点となる。

\begin{center}
\begin{tikzpicture}[
  font=\small,
  gapbox/.style={fnode, align=left, text width=52mm, minimum height=30mm},
  gapbox2/.style={fnode blue, align=left, text width=52mm, minimum height=30mm},
  gapbox3/.style={fnode muted, align=left, text width=52mm, minimum height=30mm},
  gapbox4/.style={fnode accent, align=left, text width=52mm, minimum height=30mm}
]

\node[gapbox] (data) at (0,0) {
  \textbf{①データ基盤・文書処理}\\[1.5mm]
  \footnotesize
  ・E2Eベンチマーク不在\\
  ・非米国基準（IFRS/J-GAAP）の空白\\
  ・継続自己更新型KG不在
};

\node[gapbox2] (agent) at (6.2,0) {
  \textbf{②エージェント・意思決定}\\[1.5mm]
  \footnotesize
  ・敵対的評価・ストレステスト欠如\\
  ・点時点整合性の標準不在\\
  ・エージェント安全評価の非開示
};

\node[gapbox3] (eval) at (0,-3.6) {
  \textbf{③評価・再現性・忠実性}\\[1.5mm]
  \footnotesize
  ・再現性の危機（取引コスト明示1/77件）\\
  ・見かけの予測可能性（spurious）\\
  ・忠実性・内部バイアスの監査不在
};

\node[gapbox4] (systemic) at (6.2,-3.6) {
  \textbf{④システミックリスク・規制}\\[1.5mm]
  \footnotesize
  ・表現同質性とアルファ減衰の未計測\\
  ・群集行動のリアルタイム監視基盤不在\\
  ・静的モデルリスク枠組みの前提乖離
};

\end{tikzpicture}
\end{center}
{\small\color{brandgray}\textbf{図: 未解決ギャップの全体地図}\ 4カテゴリ・20件超のギャップのうち特に重要な項目を抜粋。個々の内容は各小節を参照。①②は入力層と実行層、③は両者を横断する評価・検証層、④は市場・規制の外部性に対応する。}

\subsection{データ基盤・文書処理}

このカテゴリのギャップに共通するのは、モデル自体の能力不足ではなく、モデルへ入力する手前の層――ベンチマーク設計・知識グラフの更新様式・タクソノミの地域対応――がグローバル展開の足かせになっている点である。既存のベンチマーク・KG・抽出パイプラインの多くが米国基準・英語圏のデータに偏っており、その外側は事実上の空白地帯として残されている。

\begin{center}
\begin{tikzpicture}[
  font=\footnotesize,
  hdr/.style={font=\small\bfseries\color{fignavy}, align=center, text width=62mm},
  covbox/.style={fnode, align=left, text width=62mm, minimum height=11mm, anchor=west},
  gapbox/.style={fnode muted, dashed, align=left, text width=62mm, minimum height=11mm, anchor=west},
  dimlab/.style={font=\small\bfseries, align=right, text width=19mm, anchor=west}
]

\node[hdr] at (24mm,0) {既存ベンチマーク／パイプラインがカバー};
\node[hdr] at (91mm,0) {地域・様式のギャップ};

\node[dimlab] at (0,-11mm) (l1) {言語・地域};
\node[covbox] at (24mm,-11mm) (c1) {英語・米国上場企業の提出書類\\\footnotesize（FinanceBench, BigFinanceBench）};
\node[gapbox] at (91mm,-11mm) (g1) {日本・欧州・新興国の規制文書向け同等スイートなし（FinMMEvalは着手段階）};

\node[dimlab] at (0,-25mm) (l2) {会計基準};
\node[covbox] at (24mm,-25mm) (c2) {米国GAAPへのXBRLグラウンディング\\\footnotesize（FinTagging）};
\node[gapbox] at (91mm,-25mm) (g2) {IFRS・J-GAAP・中国会計基準への拡張は未着手};

\node[dimlab] at (0,-39mm) (l3) {知識グラフ};
\node[covbox] at (24mm,-39mm) (c3) {スナップショット／イベント駆動型KG\\\footnotesize（FinDKG, FinKG-News）};
\node[gapbox] at (91mm,-39mm) (g3) {継続自己更新型・ニュース/株式調査横断の統合KGは不在};

\node[dimlab] at (0,-53mm) (l4) {ESG開示};
\node[covbox] at (24mm,-53mm) (c4) {香港取引所2022年レポート\\\footnotesize（ESGReveal）};
\node[gapbox] at (91mm,-53mm) (g4) {EU CSRD／ISSB S1,S2／SEC気候開示は未カバー};

\end{tikzpicture}
\end{center}
{\small\color{brandgray}\textbf{図: データ基盤ギャップのカバレッジ地図}\ 実線＝既存ベンチマーク・パイプラインが到達している範囲、破線＝一次情報が指摘する未到達領域。4次元とも「米国・英語圏では前進、それ以外は空白」という同一パターンを示す。}

\begin{itemize}
  \item \textbf{エンドツーエンド・ベンチマークの不在}: AIネイティブ投資ライフサイクル（データ取り込み→構造化→仮説生成→ポートフォリオ構築→執行→取引後帰属分析）を端から端までカバーする権威あるベンチマークが存在せず、システム横断の厳密な比較が困難である。
  \item \textbf{ハイブリッド記号-ニューラルアーキテクチャの未成熟}: RAGシステムは複雑な複数ページ表・文書横断の数値推論・正確な会計演算に依然として苦戦しており、記号エンジンに算術検証を委譲しLLMは意味理解に専念するハイブリッドアーキテクチャが未成熟である（AUDITFLOW\cite{auditflow2026}、FinGround\cite{finground2026}等が部分的な解を示すのみ）。
  \item \textbf{継続自己更新型ナレッジグラフの不在}: FinDKG\cite{findkg2024}・FinKG-Newsのような金融知識グラフはスナップショット／イベント駆動型が中心で、ストリーミングニュース・開示・価格データをバッチ再構築なしにリアルタイムで取り込む継続自己更新型KGは未構築である。
  \item \textbf{非米国会計基準向けXBRLグラウンディングの空白}: FinTagging\cite{fintagging2025}の約17\%精度ギャップは米国GAAPが対象であり、IFRS・J-GAAP・中国会計基準といったタクソノミ構造が異なる体系への構造化データ優先パイプラインの拡張は、グローバル投資AIにとって重大なインフラギャップとして残る。
  \item \textbf{多言語・多形式文書の同時最適化研究の不在}: 本番OCR-LLM文書パイプラインではOCRがスループットのボトルネックであることが判明しているが、手書き注釈・複数スクリプトの表・決算資料に埋め込まれたチャートなど、多言語・多形式の金融文書をデータ居住性・プライバシー制約の下で同時最適化する大規模研究はまだない。
  \item \textbf{ESG開示の地域横断カバレッジ不足}: ESGReveal\cite{esgreveal2023}は香港取引所の2022年レポートを対象としており、EU CSRD、ISSB S1/S2、SEC気候開示フレームワークをカバーしていない。非英語文書を含む地域横断のESG AI-readiness測定は未解決の研究課題である。
  \item \textbf{財務QAベンチマークの地域偏在}: FinanceBench\cite{financebench2023}、BigFinanceBench\cite{bigfinancebench2026}はいずれも英語・米国上場企業の提出書類のみを対象とし、日本・欧州・新興国の規制文書（IFRS、J-GAAP、中国会計基準）向けの同等の評価スイートが存在しない。CLEF-2026 FinMMEval\cite{finmmeval2026}が着手を始めたばかりの段階である。
  \item \textbf{ナレッジグラフの分断}: ニュース由来・株式調査由来のナレッジグラフはそれぞれ孤立して構築・評価されており、提出書類・ニュース・決算説明会・アナリストレポートを統合しポートフォリオ意思決定タスクでベンチマークされる統一的・継続更新型のマルチソースKGはまだ存在しない。
\end{itemize}

\subsection{エージェント・意思決定}

このカテゴリのギャップは、エージェントが「何ができるか」ではなく「その判断をどう検証・監査するか」の空白に集約される。マルチエージェント投資フレームワークの実装は先行しているのに対し、敵対的評価・時間的整合性・規制トレーサビリティといった検証インフラは体系的に整備されておらず、公表される超過収益や自律性の主張を第三者が独立に検証する手段が乏しい。

\begin{center}
\begin{tikzpicture}[
  font=\footnotesize,
  implbox/.style={fnode blue, align=center, text width=128mm, minimum height=11mm},
  missbox/.style={fnode blue, dashed, align=left, text width=128mm, minimum height=9mm},
  lbl/.style={font=\footnotesize\color{figgray}, align=center, text width=128mm},
  node distance=2mm
]

\node[implbox] (impl) {\textbf{マルチエージェント投資フレームワークの実装は先行}\\\footnotesize リサーチ・デューデリジェンス・ポートフォリオモニタリング・リスク分析を担うエージェントは既に稼働};

\node[lbl, below=1.5mm of impl] (gaplabel) {▼\ しかしその判断を検証・監査する標準は下の4層すべてで未整備\ ▼};

\node[missbox, below=1.5mm of gaplabel] (m1) {\textbf{標準なし①}\ 敵対的評価・レジーム変化ストレステスト};
\node[missbox, below=1.5mm of m1] (m2) {\textbf{標準なし②}\ 点時点（Point-in-Time）整合性の公開認証};
\node[missbox, below=1.5mm of m2] (m3) {\textbf{標準なし③}\ 監査証跡を規制義務（MiFID II／SEC受託者責任／FINRA Rule 3110）へマッピングする公表フレームワーク};
\node[missbox, below=1.5mm of m3] (m4) {\textbf{標準なし④}\ 文書取り込みから投資判断までを結合するE2E評価ハーネス};

\end{tikzpicture}
\end{center}
{\small\color{brandgray}\textbf{図: 実装先行・検証インフラ空白の構図}\ 上段（実線）はすでに稼働しているエージェント能力、下段4層（破線）はいずれも「未実装」ではなく「公表された評価・監査標準が存在しない」ギャップである。}

\begin{itemize}
  \item \textbf{敵対的評価・レジーム変化ストレステストの欠如}: エージェント型投資フレームワークはバックテストや短期のライブリターンを報告するが、市場レジーム変化・流動性危機・分布シフトデータに対する標準化された敵対的評価が欠如しており、公表された超過収益の主張を検証しづらい。
  \item \textbf{点時点（Point-in-Time）整合性を強制する公開標準の不在}: ほとんどの財務LLM評価フレームワークは厳密な点時点の正確性を強制しておらず、モデルが模擬投資判断中に改訂済み提出書類や将来データに暗黙にアクセスしうる。本番投資システムにはルックアヘッドバイアスを防ぐ時間的ガードレールが必要だが、点時点コンプライアンスを監査・認証する公開標準は存在しない。
  \item \textbf{自律型投資エージェントのガバナンス・監査可能性基準の欠如}: 「境界付き自律性（bounded autonomy）」アーキテクチャは有望だが、マルチエージェント金融システムの意思決定監査証跡を、MiFID II最良執行義務、SEC受託者責任、FINRA Rule 3110といった投資分野固有の規制義務にマッピングする公表フレームワークは存在しない。
  \item \textbf{文書→意思決定の統合評価ハーネスの不在}: 生スキャンPDF取り込み、多言語OCR、ネストされた表抽出、文書横断的な数値検索、投資意思決定品質は個別には評価されてきたが、単一のエンドツーエンド評価ハーネスで結合的に評価した例はなく、性能劣化がパイプラインのどの段階に起因するか特定不能である。
  \item \textbf{ドメイン適応と規制遵守の両立方法論の不在}: ドメイン適応型金融LLMはOpen FinLLM\cite{finllmleaderboard2025}など公開リーダーボードで評価されるが、規制遵守・データガバナンス制約を保ちつつ社内固有知識（内部リサーチモデル、案件フロー、リスク限度額）に適合させる確立された方法論が存在しない。
\end{itemize}

\subsection{評価・再現性・忠実性}

前2カテゴリが「入力層（データ基盤）」と「実行層（エージェント）」の空白を扱ったのに対し、このカテゴリは両者を横断する検証層――報告された性能そのものが信用に足るか――のギャップに焦点を当てる。ここでの中心的な問いは「モデルが何を出力したか」ではなく「その出力・その超過収益の主張が、独立に再現・反証できるように測られているか」である。第\ref{sec:llm-analysis}部・第\ref{sec:cross-cutting}部で個別事例として触れた再現性・ルックアヘッドの問題を、ここでは方法論上の未解決課題として体系化する。

複数の独立した2026年の研究が、金融MLの「見かけの予測可能性」がどの段階で混入するかを、それぞれ異なる切り口で分離しようと試みている。下図はその混入経路を、パイプラインの上流から下流へ並べたものである。

\begin{center}
\begin{tikzpicture}[
  font=\footnotesize,
  stage/.style={fnode blue, align=center, text width=26mm, minimum height=13mm},
  leakbox/.style={fnode muted, dashed, align=left, text width=52mm, minimum height=13mm, anchor=west},
  arr/.style={farr blue},
  lbl/.style={font=\scriptsize\color{figgray}}
]

\node[stage] (s1) at (0,0) {\textbf{特徴量生成}\\\scriptsize 正規化・時間特徴};
\node[stage] (s2) at (0,-1.8) {\textbf{学習・選択}\\\scriptsize 仕様探索・多重検定};
\node[stage] (s3) at (0,-3.6) {\textbf{執行仮定}\\\scriptsize 約定価格・取引コスト};

\node[leakbox] (l1) at (2.0,0) {\footnotesize \textbf{正規化・centered temporal features}のリークは選択的だが大きい\cite{x07alphadisappears2026}};
\node[leakbox] (l2) at (2.0,-1.8) {\footnotesize \textbf{適応的仕様探索}がゼロ予測可能性下でも有意なバックテストを生成\cite{x07spurious2026}};
\node[leakbox] (l3) at (2.0,-3.6) {\footnotesize \textbf{取引コスト無視}が77件中76件、符号を反転させうる\cite{agentictrading2026,x07finmaseval2026}};

\draw[arr] (s1.east) -- (l1.west);
\draw[arr] (s2.east) -- (l2.west);
\draw[arr] (s3.east) -- (l3.west);
\draw[farr, draw=figgray] (s1.south) -- (s2.north);
\draw[farr, draw=figgray] (s2.south) -- (s3.north);
\end{tikzpicture}
\end{center}
{\small\color{brandgray}\textbf{図: 「見かけの予測可能性」の混入経路}\ 特徴量生成・学習/選択・執行仮定という3段階のいずれでもバイアスが混入しうる。いずれも独立の一次研究が指摘するが、3段階を単一の反証プロトコルで一括検定する共有基盤はまだ存在しない。}

\begin{itemize}
  \item \textbf{見かけの予測可能性（spurious predictability）を反証する標準手続きの不在}: Nikolopoulos\cite{x07spurious2026}は、適応的な仕様探索がmartingale-difference帰無仮説の下でさえ統計的に有意なバックテストを生成しうることを示し、合成参照クラス（ゼロ予測可能性環境・マイクロストラクチャplacebo）に対しワークフロー全体を検定する「反証監査（falsification audit）」を提案した。多くの見かけ上の発見が真の予測可能性ではなく方法論的アーティファクトであると結論づけられており、選択誘発性能インフレを有効多重性で調整して定量化する枠組みが提示されたものの、こうした反証手続きを標準として組み込んだ共有評価基盤は業界に存在しない。
  \item \textbf{decision-time leakageの発生源分離の未成熟}: Zhangら\cite{x07alphadisappears2026}は、バックテストの評価規約を1つずつ切り替える「one-switch」ベンチマークでリークの発生源を分離し、centered temporal featuresと同日始値執行（始値後の日次バー情報を含む）が予測・取引指標を大きく安定的に押し上げる一方、グローバル正規化や同日終値執行の影響は小さいことを示した。リークが一様ではなく「highly selective」である以上、どの評価規約が結果を汚染するかを網羅的に検定する感度分析の標準化が求められるが、それはまだ整備されていない。
  \item \textbf{因果と相関を切り分ける投資AI評価の不足}: FinCARE\cite{x07fincare2025}は、ポートフォリオ分析が相関ベースの分析とヒューリスティックに依存しがちで、パフォーマンスを動かす真の因果関係を捉えにくいことを出発点に、SEC 10-K由来の金融KG制約とLLM推論を、PC・GES・NOTEARSの3種類の因果発見アルゴリズムへ組み込む枠組みを提案した。合成データ上ではF1改善や介入効果の方向精度が報告されているが、これはあくまで\emph{因果推論を金融AI評価へ入れる研究プロトタイプ}であり、実市場・非定常レジーム・執行制約を含むポートフォリオ評価で「相関的に当たっただけ」のモデルと「反事実に耐える」モデルを区別する標準手続きは未確立である。
  \item \textbf{マルチエージェント評価インフラの不在}: Nguyen \& Pham\cite{x07finmaseval2026}は12のマルチエージェントシステムを4次元（アーキテクチャ・協調機構・メモリ・ツール統合）で分類し、「協調プロトコル設計がモデル規模の拡大より取引性能に効きうる」というCoordination Primacy Hypothesisを\emph{反証可能な仮説}として提示した。だが同論文は、この仮説を検証するために必要なインフラ――取引コスト控除後に協調が価値を生むかを測るCoordination Breakeven Spread（CBS）のような指標を備えた評価環境――が現状のフィールドに存在しないと明言している。look-ahead bias・survivorship bias・overfitting・取引コスト無視・レジーム変化盲目という5つの評価失敗はいずれも「報告リターンの符号を反転させうる」。
  \item \textbf{欠損文書下での忠実性（faithfulness）評価の未確立}: 金融文書はしばしば不完全であり、モデルが欠損情報をどう扱うかが実務上の信頼性を左右する。JurisTechの2026年ベンチマーク\cite{x07juristech2026}では、強いモデルが「欠損情報を認識し、与えられていないものの捏造を拒否」した一方、弱いモデルは「欠損を埋め、仮定を置き、根拠なき入力の上に一見完全な回答を生成」した。ハルシネーションはもはや単独のモデル選択の問題ではなく、プロンプト・ワークフロー・テスト条件を含む「評価層」の欠落として捉えるべきであるが、欠損下の忠実性を測る標準化されたfaithfulness評価プロトコルは未確立である（ベンダーベンチマークであり査読を経ていない点に留意）。
  \item \textbf{モデル内部表現に潜む資産バイアスの監査手段の不在}: Wu\cite{x07assetbias2026}は、フロンティアLLMの内部にBitcoin選択的な特徴（Gemma 3のsparse autoencoder特徴として同定）が存在し、これを増幅・抑制するだけでポートフォリオのBitcoin配分が明示的言及なしに\textbf{+5.2／-4.6パーセントポイント}変動する「bounded behavioral leverage」を実証した。出力の妥当性検証だけでは捉えられない、内部表現レベルの資産固有バイアスを検出・監査する透明性標準は存在せず、自律的投資判断における忠実性保証の空白として残る。
\end{itemize}

\subsection{システミックリスク・規制}

このカテゴリのギャップの根は、個々のファンド・モデル単位のAI導入が先行する一方、市場全体・規制当局の側の計測・検証インフラがそれに追いついていない非対称性にある。マイクロストラクチャーへの実際の影響も、学術的に主張される「アルファ」の頑健性も、独立に再現・監視する共有基盤が欠けたままである。

\begin{center}
\begin{tikzpicture}[
  font=\footnotesize,
  flowbox/.style={fnode accent, align=center, text width=110mm, minimum height=10mm},
  gapflow/.style={fnode accent, dashed, align=center, text width=110mm, minimum height=10mm},
  tagS/.style={fnode muted, dashed, align=left, text width=49mm, minimum height=15mm, font=\scriptsize},
  tagL/.style={fnode muted, dashed, align=left, text width=74mm, minimum height=13mm, font=\scriptsize},
  lbl/.style={font=\scriptsize\color{figgray}, align=center, text width=110mm}
]

\node[flowbox] (fund) at (0,0) {個別ファンド・モデル単位のAI導入拡大};
\node[lbl] at (0,-8mm) (a1) {↓};
\node[flowbox] (market) at (0,-16mm) {市場全体への影響（マイクロストラクチャー・与信・オルタナデータ・システミックリスク）};
\node[lbl] at (0,-24mm) (gaplbl) {▼\ リアルタイム計測・独立検証の共有基盤が存在しない\ ▼};
\node[gapflow] (reg) at (0,-34mm) {規制の閾値較正・監督（未計測ゆえ届いていない）};

\node[font=\scriptsize\bfseries\color{figgray}] at (0,-43mm) (taghdr) {具体的に欠けている計測・検証基盤};

\node[tagS] (t1) at (-52mm,-55mm) {マイクロストラクチャー影響: SEC 13F（静的保有）のみで、リアルタイムのオーダーフロー計測なし};
\node[tagS] (t2) at (0,-55mm) {オルタナデータ統合（衛星画像・決済フロー等）の統合ベンチマーク不在};
\node[tagS] (t3) at (52mm,-55mm) {オープンウェイト財務基盤モデル不在（BloombergGPT級の公開モデルなし）};

\node[tagL] (t4) at (-38.5mm,-72mm) {与信解釈可能性研究をEU AI Act／FRBモデルリスクガイダンス／Basel IRBへ橋渡す規制マッピング標準は未発表};
\node[tagL] (t5) at (38.5mm,-72mm) {再現性の危機: LLMトレーディング研究77件中、取引コスト明示1件・時系列整合プロトコル2件のみ\cite{agentictrading2026}};

\end{tikzpicture}
\end{center}
{\small\color{brandgray}\textbf{図: 計測・検証基盤が欠けたフィードバックループ}\ ファンド単位のAI導入拡大が市場全体に及ぼす影響は、リアルタイムでは計測されないまま規制の閾値較正へたどり着けずにいる（破線＝未整備）。下段は本節が指摘する具体的な欠落項目。}

\begin{itemize}
  \item \textbf{マイクロストラクチャー影響のリアルタイム測定の欠如}: AIモノカルチャーのシステミックリスクは現状パッシブなポートフォリオ保有（SEC 13F）でしか測定されておらず、ストレス時の同期的アルゴリズム注文がイントラデイ流動性・ボラティリティに与える直接的なマイクロストラクチャー影響はリアルタイムのオーダーフローデータで観測されていない。規制当局は介入のための較正済み閾値を欠く。
  \item \textbf{与信リスク解釈可能性の規制マッピング標準の不在}: 与信リスクの解釈可能性研究はモデルレベルの説明可能性（chain-of-thought、特徴量寄与度）をカバーするが、EU AI Actの高リスクAI規定、米連邦準備制度のモデルリスクガイダンス、Basel IRB内部格付け検証といった規制監査証跡へのマッピング標準は未発表である。精度・コンプライアンス・説明可能性のトリレンマ（\cite{fintrilemma2026}）が示すように、説明可能性が導入の可否を決めるゲートウェイである以上、この規制マッピングの空白は与信AIの実装を直接に律速する。
  \item \textbf{オルタナティブデータ統合の統合ベンチマークの不在}: 衛星画像・クレジットカード取引フロー・位置情報/人流シグナルなどのオルタナティブデータ統合について、シグナル抽出品質・データパイプラインレイテンシ・下流ポートフォリオアトリビューションを同時測定するエンドツーエンドベンチマークが存在せず、プロバイダー・パイプライン設計の比較が不可能である。
  \item \textbf{再現性の危機}: 77件のLLMトレーディング研究の系統的レビュー\cite{agentictrading2026}では、取引コストを明示的にモデル化していたのはわずか1件、再現可能で時系列整合性のあるプロトコルを用いたのはわずか2件と判明した。学術コンテストサーバーのような標準化された取引コストモデル・レジーム変化ストレステストを備えた共有ライブ市場評価基盤が不在である。
  \item \textbf{オープンウェイト財務基盤モデルの不在}: BloombergGPT\cite{bloomberggpt2023}に匹敵するベンチマークスコアを持つオープンウェイトモデルが存在せず、公開されている大規模財務事前学習コーパスも限定的である（Stanford SEFD\cite{stanfordedgar2026}が端緒）。実務家の多くが高価な独自APIか、ドメインミスマッチのある汎用モデルに依存せざるを得ない状況が続いている。
  \item \textbf{エージェント型AI決済ガバナンス・信頼性標準の未成熟}: 決済領域のエージェント型AIに対する三層ガバナンス枠組み（IMF\cite{imfpayments2026}）や、金融AIライフサイクル全体を対象とした信頼できるAIのセキュリティタクソノミ（\cite{trustworthyfintech2026}）はいずれも提案・整理の段階にとどまり、確率的AIと決定論的な決済インフラの整合や、データ・モデルポイズニング／プロンプトインジェクション／ディープフェイクKYCへの防御を規定する拘束力ある技術標準としては未確立である。
  \item \textbf{表現同質性とアルファ減衰の未計測}: \S\ref{sec:cross-cutting}で扱ったAIモノカルチャー\cite{aisystemicrisk2026}に対し、2026年には作用機序をより構造的に分離する研究が現れた。Qiu \& Han\cite{x07homogeneity2026}は「表現の同質性（representation homogeneity）」と「予測の重なり（forecast overlap）」を区別し、予測が平時には多様に見えてもストレス時には表現の同質性が予測不一致の有効空間を圧縮し、position stickinessを通じて蓄積した「隠れたレバレッジ」がショック時に協調的デレバレッジで顕在化すると論じた。Meng \& Chen\cite{x07alphadecay2026}は同じSEC 13F系データ（99.5百万件保有、2013--2024）でAI採用ファンドのポートフォリオ収束が\textbf{42\%}増大したことを示し、現行の普及率（約70\%）ではシグナル寿命がpre-AI期の5〜7年から\textbf{約18カ月}へ短縮すると試算した。だが\emph{予測の見かけ上の多様性ではなく内部表現の多様性}を、あるいは\emph{アルファ半減期そのもの}を、リアルタイムで計測しマクロプルーデンス監視する共有基盤は存在しない。
  \item \textbf{群集行動のリアルタイム監視・ストレステスト基盤の構想段階}: 英議会Treasury Committeeの報告に対する規制当局回答（2026年4月）で、Bank of Englandは「wait and see」ではないとし、国際カウンターパートと共同で「AIトレーディングエージェントがどの条件下で相関的挙動（herding）を示しストレスシナリオを増幅しうるか」を理解するシミュレーション研究を進めていると回答した\cite{x07tsc2026}。しかしこれは条件解明のための研究段階であり、前掲の表現同質性・アルファ減衰を平時から継続監視し、規制の介入閾値較正へ接続する運用可能な計測基盤（マクロプルーデンスなAIストレステスト）はまだ整備されていない。HM Treasuryが大手AI／クラウドをCritical Third Partiesに指定する評価を進めている点も、監督対象の外縁がなお流動的であることを示す。
  \item \textbf{静的モデルリスク枠組みと連続学習AIの前提乖離}: Kurshanら\cite{x07agenticregulator2025}は、現行のモデルリスク枠組みが「静的で仕様が明確なアルゴリズムと一回限りの検証」を前提とする一方、現代の金融AIは継続学習・創発的挙動によって動作するため、両者の間に観測性・制御のギャップが残ると指摘する。彼らは①各モデル埋め込み型の自己規制、②社内ガバナンス、③セクター全体の不安定性を監視する規制当局ホスト型エージェント、④独立監査ブロックからなる4層の監督構造を提案するが、これはあくまで提案であり、SR 11-7系ガイダンス（\S\ref{sec:cross-cutting}）が想定する静的・一回限りの検証と、連続更新されるエージェント出力の監督可能性を橋渡す実装標準は未確立である。
  \item \textbf{規制期限の後ろ倒しによる不確実性の延伸}: EU AI ActのAnnex III高リスク分類（信用スコアリング・保険料算定AI）の完全適用は当初2026年8月2日予定だったが、Digital Omnibus（欧州議会承認2026年6月16日）により、standalone Annex IIIの信用・保険AI義務の拘束的期限は\textbf{2027年12月2日}へ約16カ月延期された\cite{x07digitalomnibus2026}。高リスク分類は既存のモデルリスク管理の軽量版ではなく、適合性評価・技術文書・データガバナンス・ログ記録・人間監督・市販後監視の6カテゴリを追加要求する点は変わらない。期限の後ろ倒しは準備時間を与える一方で規制不確実性を延伸させ、CRR/CRD・DORA・EBAガイドライン等の既存規則との重畳マッピングと文書化ギャップをいつ・どの深度で埋めるべきかという実装判断の空白を長引かせている。
  \item \textbf{標準化・プライバシー・コスト制約を統合する運用標準の不足}: 2026年6月のFSB consultation report\cite{x07fsbsound2026}は、金融機関のAI採用に対し12のsound practicesを提示し、GenAIとagentic AIを含むライフサイクル管理を組織全体の統治課題として扱う。一方、米財務省は金融サービス固有のAI情報共有、データ標準、リスク管理ベストプラクティス、データプライバシー・セキュリティ・品質基準の明確化、AIプロバイダー集中リスクの監視を提言している\cite{x07treasuryai2024}。さらにCyber Risk InstituteのFS AI RMFは、NIST AI RMFに構造的に整合する230のcontrol objectivesへ金融セクター向けに落とし込む\cite{x07crifsairmf2026}。これらは原則から統制項目への前進であるが、\emph{どの統制証跡が、どのモデル評価・プライバシー境界・ベンダー集中リスク・推論コストSLOを満たせば本番投資エージェントを許可できるか}を結ぶ共通の適合性テストはまだ存在しない。
\end{itemize}

\begin{center}
\footnotesize
\renewcommand{\arraystretch}{1.35}
\begin{tabular}{@{}p{3.0cm}p{4.2cm}p{4.2cm}p{2.6cm}@{}}
\toprule
\textbf{統治レイヤー} & \textbf{到達点} & \textbf{残る未解決点} & \textbf{投資AIでの意味} \\
\midrule
原則・横断RMF & NIST AI RMFに整合する金融セクター実装が出現\cite{x07crifsairmf2026} & 汎用RMFと証券・運用・与信の監督証跡の接続は機関ごと & 監査人が同じ証跡を同じ意味で読めない \\
\addlinespace
金融セクターsound practice & FSBが12の実務慣行をconsultationとして提示\cite{x07fsbsound2026} & agentic AIの評価ログ・停止条件・第三者検証は拘束的標準ではない & 自律性の上限を事前に合意しにくい \\
\addlinespace
データ標準・プライバシー & 米財務省がデータ標準、プライバシー・セキュリティ・品質基準の明確化を提言\cite{x07treasuryai2024} & 社内リサーチ、顧客情報、ベンダーデータをまたぐRAG境界の実装標準が不足 & 高価な専用環境か、過度に制限された汎用環境へ二極化 \\
\addlinespace
集中リスク・コスト & AIプロバイダー集中リスクの監視が政策課題化\cite{x07treasuryai2024} & 推論コスト、レイテンシ、クラウド集中、モデル切替可能性を同時に測るSLOがない & 「性能は高いが運用不能」なエージェントを除外しにくい \\
\bottomrule
\end{tabular}
\end{center}
{\small\color{brandgray}\textbf{表: 標準化・プライバシー・コストの統合ギャップ}\ 金融AIの統治は原則論からcontrol objectivesへ進んでいるが、評価ログ・データ境界・ベンダー集中・推論コストを同じ適合性テストへ束ねる段階には達していない。}

\subsection{ギャップ・現状・必要な取り組みの対応表}

以上の4カテゴリの主要ギャップを、それぞれの「現状（到達点）」と「解消に必要な取り組み」の対で一覧化する。共通するのは、\emph{個別能力の実装は先行しているのに、その能力を横断的に検証・計測・監査する共有基盤が欠けている}という非対称性である。

\begin{center}
\footnotesize
\renewcommand{\arraystretch}{1.35}
\begin{tabular}{@{}p{2.0cm}p{4.3cm}p{4.3cm}p{2.6cm}@{}}
\toprule
\textbf{カテゴリ} & \textbf{未解決ギャップ} & \textbf{現状（到達点）} & \textbf{必要な取り組み} \\
\midrule
データ基盤 & E2Eベンチマーク／非米国会計基準グラウンディング & 米国GAAP・英語圏で前進（FinTagging\cite{fintagging2025}, FinanceBench\cite{financebench2023}） & IFRS/J-GAAP対応・端から端までの権威あるスイート \\
\addlinespace
エージェント & 点時点整合性・敵対的評価・安全開示 & 実装は先行、安全結果は25/30が非開示\cite{x07agentindex2026} & PiT認証・レジーム変化ストレステスト・第三者検証の開示標準 \\
\addlinespace
評価・再現性 & 見かけの予測可能性／リーク源分離／忠実性 & 反証監査\cite{x07spurious2026}・one-switch\cite{x07alphadisappears2026}は個別提案どまり & 3段階（特徴量・選択・執行）を一括検定する共有反証プロトコル \\
\addlinespace
忠実性・バイアス & 欠損文書の忠実性／内部表現バイアス & 弱モデルは捏造\cite{x07juristech2026}、内部Bitcoin偏好で配分±5pp\cite{x07assetbias2026} & faithfulness評価層と内部表現監査の透明性標準 \\
\addlinespace
因果・反事実 & 相関ベースの予測と因果的に妥当な判断の分離 & KG+LLMで因果発見を補強するFinCARE\cite{x07fincare2025}は研究段階 & 非定常市場・執行制約を含む反事実評価プロトコル \\
\addlinespace
システミック & 表現同質性・アルファ減衰・群集行動 & 42\%収束・半減期18カ月\cite{x07alphadecay2026}、BoEは研究段階\cite{x07tsc2026} & 内部表現多様性のリアルタイム監視・閾値較正 \\
\addlinespace
規制 & 静的モデルリスク枠組み／期限不確実性 & 4層監督は提案\cite{x07agenticregulator2025}、EU期限は2027年へ延期\cite{x07digitalomnibus2026} & 連続学習AIを監督する動的検証・重畳規則の文書化 \\
\addlinespace
標準化・プライバシー・コスト & 統制項目、データ境界、推論コストSLOの統合 & FSB sound practices\cite{x07fsbsound2026}、FS AI RMF 230統制\cite{x07crifsairmf2026}、Treasury提言\cite{x07treasuryai2024}が並立 & 証跡・プライバシー・集中リスク・コストを束ねる本番許可基準 \\
\bottomrule
\end{tabular}
\end{center}
{\small\color{brandgray}\textbf{表: 未解決ギャップ×現状×必要な取り組み}\ いずれの行も「個別実装は先行、横断的検証・計測・監査の共有基盤が欠落」という同一パターンを示す。右列は本書が今後の研究・投資の優先順位を検討する際の参照点となる。}

\subsection{汎用AIの未解決課題と金融というストレステスト領域}
\label{sec:cross-domain-gaps}

本節がここまで整理してきた4カテゴリのギャップ――評価・再現性の危機、エージェントの安全性、モデル集中とモノカルチャー、規制の断片化――は、実は金融に固有の欠陥ではない。2024〜2026年にかけて、生成AI・エージェントAI業界全体が金融の外側でほぼ同型の課題を独立に指摘しており、しかもその多くは金融特有の議論より先に、あるいは並行して一次資料化されている。以下ではその汎用AI側の一次資料を確認したうえで、「なぜ金融がこれらの汎用AI課題を最初に、かつ最も深刻な形で顕在化させるストレステスト領域になるのか」を論じる。金融関係者にとっての含意は明確である――これらの課題は「金融が特殊だから起きている」のではなく「AI業界全体で今起きており、金融はその影響を実資金の損益として最初に、かつ最も鋭敏に計測してしまう立場にある」ということである。

\begin{center}
\begin{tikzpicture}[
  font=\small,
  challenge/.style={fnode muted, align=left, text width=46mm, minimum height=13mm},
  amp/.style={fnode accent, align=left, text width=46mm, minimum height=13mm},
  hdr/.style={font=\small\bfseries\color{fignavy}, align=center, text width=46mm},
  arr/.style={farr accent}
]

\coordinate (Lcol) at (0,0);
\coordinate (Rcol) at (64mm,0);

\node[hdr] (h1) at (Lcol) {汎用AI業界で観測される課題};
\node[hdr] (h2) at (Rcol) {金融に固有の増幅要因};

% 右列（増幅要因）を基準に縦に積み、左列を各行の上端に揃える（内容が伸びても重ならない）
\node[amp, below=6mm of h2] (a1) {「正解」が市場自体にしかなく、実資金の損益として即座に検証・淘汰される};
\node[challenge, anchor=north] (c1) at (a1.north -| Lcol) {評価クライシス：ベンチマーク汚染・再現性崩壊\cite{y07aiindex2026,y07swebenchcontam2026}};

\node[amp, below=6mm of a1] (a2) {発注・送金権限を持つエージェントは「読むだけ」を超え資金流出に直結};
\node[challenge, anchor=north] (c2) at (a2.north -| Lcol) {エージェント安全性：プロンプトインジェクション・権限濫用\cite{y07owaspagentic2026,y07googlepromptinjection2026}};

\node[amp, below=6mm of a2] (a3) {SEC 13F等の公開データで集中の影響を実数値として追跡できる稀有な産業};
\node[challenge, anchor=north] (c3) at (a3.north -| Lcol) {基盤モデル集中・モノカルチャー\cite{y07monoculture2022,y07foundationconcentration2023}};

\node[amp, below=6mm of a3] (a4) {受託者責任・最良執行等の垂直規制が既に高密度に存在し、水平規制と衝突する};
\node[challenge, anchor=north] (c4) at (a4.north -| Lcol) {規制の水平性 vs 断片化された監督\cite{y07euaiacthorizontal,y07govopacity2026}};

\draw[arr] (c1.east) -- (c1.east -| a1.west);
\draw[arr] (c2.east) -- (c2.east -| a2.west);
\draw[arr] (c3.east) -- (c3.east -| a3.west);
\draw[arr] (c4.east) -- (c4.east -| a4.west);

\end{tikzpicture}
\end{center}
{\small\color{brandgray}\textbf{図: 金融が汎用AI課題のストレステスト領域になる4つの増幅経路}\ 左列はAI業界全体で一次資料が指摘する汎用課題、右列はそれが金融という文脈に置かれた時に増幅される固有の要因。金融固有の欠陥ではなく、汎用課題が金融という「増幅装置」を通ることで最初に顕在化する構図である。}

\subsubsection*{評価クライシス：ベンチマーク汚染と再現性の崩壊は業界横断の現象}

Stanford HAI「2026 AI Index Report」は、AI業界全体で「能力の進歩」と「それを測る・統治する枠組みの整備」の乖離が拡大していると指摘する。Humanity's Last Examのような、当初は数年単位で飽和しないよう設計されたベンチマークでさえ、フロンティアモデルが単年で約30ポイントの上昇を記録し、飽和までの想定期間を一つの製品サイクルへ圧縮した\cite{y07aiindex2026}。OpenAI共同創業者のAndrej Karpathyはこの状況を「評価クライシス（evaluation crisis）」と呼び、「いま何を指標として見ればよいのか、正直分からない」と述べている\cite{y07karpathyevalcrisis2025}。同レポートはさらに、Foundation Model Transparency Indexの平均スコアが前年の58点から40点へ低下したことも記録しており、最先端ラボがモデルの能力を上げる一方で学習データ・パラメータ数等の開示を減らしている実態を示す\cite{y07aiindex2026}。この「能力は測れなくなり、開示も減っていく」という2つの潮流が同時進行している点が、金融固有ではなく業界全体の構造的傾向であることを裏付ける。

\begin{casestudy}{OpenAIがSWE-bench Verifiedを「引退」させた2026年2月の汚染監査}
コーディングエージェントの標準的ベンチマークであったSWE-bench Verifiedについて、OpenAIのFrontier Evalsチームは2026年2月23日、同ベンチマークのスコア報告を取りやめると発表した。汚染監査の結果、GPT-5.2、Claude Opus 4.5、Gemini 3 Flash Previewを含む全フロンティアモデルが、タスクIDのみから正解パッチ（gold patch）や問題文の詳細を一言一句再現できることが判明し、学習データへの混入が裏付けられた。さらに残存タスクの少なくとも59.4\%が、公正な評価条件下では欠陥があるか解答不能と判定された。OpenAIは後継としてSWE-bench Proの使用を推奨している\cite{y07swebenchcontam2026}。
\end{casestudy}

\begin{insight}{金融への示唆：金融の「バックテスト汚染」は業界全体の評価クライシスの特殊ケースにすぎない}
77件のLLMトレーディング研究のうち取引コストを明示的にモデル化したのはわずか1件という金融側の再現性の危機\cite{agentictrading2026}は、非金融領域のソフトウェア工学エージェントで起きたSWE-bench Verifiedの汚染発覚と、根は同じ「評価データ・評価プロトコルへの信頼が体系的に崩れている」現象である。Semmelrockらの系統的レビューは、この種の再現性の障壁がドメインを問わない機械学習研究全般の構造的問題であることを裏付ける\cite{y07mlreproducibility2024}。決定的な違いは、金融ではベンチマークの汚染が「学術的な指摘」で終わらず、実資金の運用損益として即座に――しかも第三者の反証なしに――市場から突きつけられる点にある。裏を返せば、金融は汎用AIの評価クライシスに対する「最も過酷な、しかし最も早く結果が出る検証環境」であり、他業界が数年かけて気づく評価の欠陥を、金融は四半期決算のタイムスケールで露呈させうる。
\end{insight}

\subsubsection*{エージェント安全性：ガードレール・権限・監査は業界横断で未解決}

自律エージェントの安全性は、金融エージェントに限らずAI業界全体で標準化が追いついていない領域である。OWASP GenAI Security Projectは2025年12月9日、100名超の業界専門家の協力を得て「OWASP Top 10 for Agentic Applications for 2026」を公開し、Agent Goal Hijack（ASI01）、Tool Misuse \& Exploitation（ASI02）、Agent Identity \& Privilege Abuse（ASI03）からMemory \& Context Poisoning（ASI06）、Rogue Agents（ASI10）まで10カテゴリの汎用エージェントリスク分類を定義した\cite{y07owaspagentic2026}。これは特定業界向けではなく、あらゆるエージェント型アプリケーションに適用される横断的な脅威分類である。Googleのセキュリティチームは2026年4月23日、Common Crawlのアーカイブ（月間20〜30億ページ）を走査した調査で、悪質な間接プロンプトインジェクションの検出数が2025年11月から2026年2月の間に相対で32\%増加したと報告した\cite{y07googlepromptinjection2026}。METRによる12社のフロンティアAI安全方針の分析では、能力閾値・モデル重み保護・展開緩和策など9つの共通要素が識別された一方、各社の脅威モデル（自律的複製・有害な操作・不整合等）の成熟度にはばらつきが残ることが指摘された\cite{y07metrcommon2025}。

\begin{casestudy}{「プロンプトインジェクションは永久に完全解決しないかもしれない」というOpenAIの公式見解}
OpenAIは2025年12月、自社の自動化されたレッドチーミングが発見した具体的なプロンプトインジェクション攻撃手法を開示し、ブラウザ型AIエージェント「ChatGPT Atlas」向けの継続的な耐性強化策を公表した。同社のHead of Preparednessは複数の報道（TechCrunch、CyberScoop、いずれも2025年12月22日）に対し、ブラウザエージェントに対するプロンプトインジェクションは「永久に完全にパッチが当たらないかもしれない」と明言し、OpenAIはその後、より高リスクなエージェント的ブラウジング向けに「Lockdown Mode」を追加投入した\cite{y07openaiatlas2025}。
\end{casestudy}

\begin{insight}{金融への示唆：「読むだけ」の被害と「動かせる」被害の非対称性}
消費者向けブラウザエージェントにおいてすら、フロンティアラボ自身が安全性の天井を「不確実」と認めている以上、発注執行・資金移動・与信判断の権限を段階的にエージェントへ委譲しつつある金融機関は、自らが完全には制御できないリスクの上に業務プロセスを構築していることになる。OWASP ASI03（Agent Identity \& Privilege Abuse）やASI06（Memory \& Context Poisoning）が指摘する汎用リスクは、金融エージェントでは「誤った要約を返す」ではなく「誤った発注を執行する」「改竄された市場データに基づき与信を承認する」という、可逆性の低い結果に直結する。第\ref{sec:open-gaps}部\S7.2で述べた「点時点整合性の公開認証」「監査証跡を規制義務へマッピングする公表フレームワーク」の不在は、業界横断のエージェント安全性標準（OWASP ASI、METRの9要素）が未成熟であることの、金融という文脈での必然的な帰結として読み直すことができる。
\end{insight}

\subsubsection*{基盤モデル集中とモノカルチャー：金融外で先に理論化された懸念}

「AIモノカルチャー」という概念自体、金融発ではない。Bommasaniらは2022年のNeurIPS論文で、Kleinberg \& Raghavan（2021）の「アルゴリズムモノカルチャー」概念を拡張し、意思決定者が同じ学習データや基盤モデルという「共有コンポーネント」に依存すると、特定の個人・集団を体系的に同じように扱ってしまうという「コンポーネント共有仮説」を提示した。同論文は視覚（CLIP）・言語（RoBERTa）という非金融領域のモデルで実証実験を行っており、基盤モデルの共有が必ずしも結果の同質化を悪化させるとは限らないという、単純な直観に反する知見も示している\cite{y07monoculture2022}。VipraとKorinekは2023年、フロンティア基盤モデルが自然独占へ向かう傾向を持ち、市場の競争性を保つための反トラスト介入と、安全性・プライバシー・非差別性・信頼性・相互運用性を対象とする品質基準規制の両方が必要になると論じた\cite{y07foundationconcentration2023}。Stanford AI Index 2026が記録するFoundation Model Transparency Indexの58点から40点への低下\cite{y07aiindex2026}は、この集中傾向が2026年時点でむしろ強まっていることを示す。

\begin{casestudy}{最先端ラボの透明性スコアが1年で58点から40点へ低下（Stanford AI Index 2026）}
Foundation Model Transparency Indexは、学習コード・データセット規模・パラメータ数等の開示状況を指数化したものである。Stanford HAIの2026年版レポートは、最も能力の高いモデルが今や最も不透明なモデルになりつつあると指摘し、大手AI企業がこれらの情報を非開示にする傾向を強めた結果、平均スコアが前年の58点から40点へ低下したと報告している\cite{y07aiindex2026}。
\end{casestudy}

\begin{insight}{金融への示唆：金融は集中・同質化の影響を数値化できる稀有な産業}
金融は、SEC 13Fのような公開保有データを通じて、基盤モデルへの依存が実際にポートフォリオの収束やアルファ減衰としてどの程度顕在化するかを、他業界にはほぼ存在しない精度で計測できる産業である。本節既述のQiu \& Han\cite{x07homogeneity2026}・Meng \& Chen\cite{x07alphadecay2026}による表現同質性・アルファ半減期（約18カ月への短縮）の実証は、Bommasaniらが2022年に非金融領域で理論化・実証したコンポーネント共有仮説の、初期の産業横断的な検証事例として位置づけられる。しかも透明性が低下する（＝どのモデルにどれだけ依存しているかの外部からの検証が難しくなる）局面でこの集中が進行している以上、金融の集中リスク監視基盤の欠如（\S7.4）は、汎用AI業界の透明性低下トレンドと組み合わさることで、今後さらに計測困難になっていく可能性が高い。
\end{insight}

\subsubsection*{規制の水平性と金融の垂直規制が衝突する場所}

EU AI Actは欧州委員会自身が「horizontal（水平的）」な法律と位置づける、業種を問わない単一のリスク階層規制であり、金融・医療・雇用・法執行など、あらゆるセクターに同じ「許容不可／高リスク／限定的リスク／最小リスク」という分類を適用する\cite{y07euaiacthorizontal}。NIST AI Risk Management Framework（2023年1月公表）も同様に業種非依存・任意ベースの枠組みとして設計されている\cite{y07nistairmf2023}。しかし、Rostの2026年の論文は、こうした水平的なガバナンス枠組みが業種横断で適用される際に生じる構造的な緊張――開示要求と営業秘密の対立、AIシステムが自らを評価する自己言及的な説明責任の空洞化、セクター間の規制の不整合、技術的説明可能性と規制上の説明要求の乖離、独立監査の実務上の限界、責任の所在の不明確化――を「6つの緊張の次元」として整理した\cite{y07govopacity2026}。Future of Life Instituteの2025年夏版AI Safety Indexは、この整理を裏付けるように、7社のフロンティアAI企業のうち終末リスク（existential risk）計画についてD評価を上回った企業が一社もなかったと報告している（Anthropic C+、OpenAI C、Google DeepMind C-、xAI D、Meta D、Zhipu AI F、DeepSeek F）\cite{y07flisafetyindex2025}。

\begin{insight}{金融への示唆：水平規制と垂直規制の摩擦を最初に露呈するのが金融}
金融は、MiFID II最良執行義務、SEC受託者責任、FINRA Rule 3110、Basel IRBといった高密度な垂直規制がすでに存在する数少ない業界である。EU AI ActやNIST AI RMFのような水平規制は、これらの垂直規制と重畳した瞬間に、Rostの言う「6つの緊張」――特に説明責任の所在の不明確化と監査の実務上の限界――を最も早く可視化する。加えてFuture of Life Instituteの評価が示すように、金融機関がエージェント型投資システムの基盤として採用するベンダー自身が、独立した第三者評価でC〜Fレンジの安全性ガバナンスにとどまっている。これは、金融の「静的モデルリスク枠組みと連続学習AIの前提乖離」（本節既述、\cite{x07agenticregulator2025}）という課題が、金融機関自身のガバナンス整備だけでは解決せず、ベンダー選定時のデューデリジェンス項目として「モデル性能」だけでなく「ベンダーの安全性ガバナンス成熟度」を組み込む必要があることを示唆する。
\end{insight}

\begin{center}
\footnotesize
\renewcommand{\arraystretch}{1.35}
\begin{tabular}{@{}p{2.6cm}p{4.4cm}p{3.9cm}p{3.3cm}@{}}
\toprule
\textbf{汎用AI課題} & \textbf{業界横断の一次証拠（2024--2026）} & \textbf{金融での深刻化要因} & \textbf{金融がストレステスト領域になる理由} \\
\midrule
評価クライシス（ベンチマーク汚染・再現性） & OpenAIがSWE-bench Verifiedの汚染を認め引退\cite{y07swebenchcontam2026}／Stanford AI Index 2026「評価クライシス」\cite{y07aiindex2026}／ML再現性の障壁は業界共通\cite{y07mlreproducibility2024} & バックテストの「正解」は市場自体にしかなく、第三者による反証手段が乏しい & 実資金の損益として即座に、かつ独立に検証される稀有な産業 \\
\addlinespace
エージェント安全性（プロンプトインジェクション・権限濫用） & OWASP ASI01--10\cite{y07owaspagentic2026}／悪質なプロンプトインジェクション検出32\%増\cite{y07googlepromptinjection2026}／ChatGPT Atlas「永久に完全解決しないかもしれない」\cite{y07openaiatlas2025} & 発注・送金権限を持つエージェントは「誤情報」を超え資金流出・誤執行に直結 & MiFID II・SEC受託者責任等、権限濫用の結果が法的責任へ直結する数少ない産業 \\
\addlinespace
基盤モデル集中・モノカルチャー & コンポーネント共有仮説（NeurIPS 2022）\cite{y07monoculture2022}／基盤モデルの自然独占リスク\cite{y07foundationconcentration2023}／透明性指数58→40\cite{y07aiindex2026} & 同一基盤モデルへの依存が個別損失でなく協調的デレバレッジへ波及\cite{x07homogeneity2026} & SEC 13F等の公開データで集中の影響を実数値として追跡できる稀有な産業 \\
\addlinespace
規制の水平性 vs 垂直規制の接続 & EU AI Actは業種横断の水平規制\cite{y07euaiacthorizontal}／NIST AI RMFも業種非依存\cite{y07nistairmf2023}／水平ガバナンスの6つの緊張\cite{y07govopacity2026}／フロンティア企業の安全性はC\textasciitilde{}F評価\cite{y07flisafetyindex2025} & MiFID II／Basel／SEC受託者責任という垂直規制との重畳マッピングが未整備 & 水平規制と垂直規制の摩擦を最初に、かつ最も高密度に露呈する規制産業 \\
\bottomrule
\end{tabular}
\end{center}
{\small\color{brandgray}\textbf{表: 汎用AI課題×金融での深刻化要因×ストレステスト領域としての金融}\ 4行とも「金融固有の欠陥」ではなく「汎用AI課題が金融という増幅装置を通ることで最初に、かつ最も鋭敏に顕在化する」という同一パターンを示す。金融関係者にとっての実務的含意は、これらの課題の解決を待つのではなく、金融自身が汎用AI業界に先んじて計測・監査手法を確立する「最初の実験場」になり得るという点にある。}

\section{結論と提言}
\label{sec:conclusion}

\subsection{総括}

440件を超える一次情報を精査した結果、金融×AIの現在地は「モデルの能力競争」から「データ基盤・検証層・ガバナンス設計の競争」へと重心を移しつつあることが明確になった。BloombergGPTのような巨大ドメイン特化モデルの登場から数年を経て、決定的な差を生んでいるのは以下の3層である。

\begin{enumerate}
  \item \textbf{前処理・検索アーキテクチャ層}: GPT-4oの精度が前処理の有無で56\%から76\%へ変わり\cite{multistructured2025}、Fin-RATEでは検索設計の違いが精度を23.91\%から1.27\%へ崩壊させる\cite{finrate2026}。モデルを差し替えるより先に、この層への投資対効果の方が大きい。
  \item \textbf{シンボリック検証層}: AUDITFLOWの82\%→18\%の精度崩壊が象徴するように\cite{auditflow2026}、LLM単体は正確な会計演算・数値照合を安定して行えない。記号的検証とLLMの意味理解を組み合わせるハイブリッド設計が、実運用可否の分岐点になっている。
  \item \textbf{ガバナンス・監査可能性層}: Bridgewater AIA Labsのように実資本を運用する事例が存在する一方\cite{bridgewater_aialabs_nd}、DeepFund\cite{deepfund2025}やPortBench\cite{portbench2026}のように独立検証ではフロンティアモデルでも純損失・崩壊的ドローダウンを記録する事例が同時に存在する。この差を分けるのは、タスクスコープの厳密さ、人間レビューの保持、監査可能な意思決定記録の有無である。
\end{enumerate}

この3層は独立した並列要件ではなく、下から積み上がる前提条件の連鎖である。前処理・検索が崩れれば、その上のシンボリック検証もガバナンスも機能する土台を失う。

\begin{center}
\begin{tikzpicture}[
  font=\small,
  layerbox/.style={text width=11.6cm, align=left, minimum height=19mm},
  l1/.style={fnode, layerbox},
  l2/.style={fnode blue, layerbox},
  l3/.style={fnode accent, layerbox},
  arr/.style={farr}
]

\node[l1] (l1) at (0,0) {%
  \textbf{層1\ 前処理・検索アーキテクチャ層}\\[1mm]
  \footnotesize GPT-4o精度\ \textbf{56\%→76\%}（前処理の有無）\cite{multistructured2025}\ ／\ Fin-RATE企業横断比較\ \textbf{23.91\%→1.27\%}（検索設計の失敗）\cite{finrate2026}
};

\node[l2, above=7mm of l1] (l2) {%
  \textbf{層2\ シンボリック検証層}\\[1mm]
  \footnotesize AUDITFLOW精度\ \textbf{82\%→18\%}（シンボリック検証の有無）\cite{auditflow2026}
};

\node[l3, above=7mm of l2] (l3) {%
  \textbf{層3\ ガバナンス・監査可能性層}\\[1mm]
  \footnotesize Bridgewater AIA Labsの実運用継続\cite{bridgewater_aialabs_nd}\ vs.\ DeepFund\cite{deepfund2025}・PortBench\cite{portbench2026}の独立検証における純損失・崩壊的ドローダウン
};

\node[above=6mm of l3, align=center, text width=11.6cm, font=\bfseries] (headline) {実運用の成否はこの3層の質で決まる};

\draw[arr] (headline.south) -- (l3.north);
\draw[arr] (l1.north) -- (l2.south);
\draw[arr] (l2.north) -- (l3.south);
\end{tikzpicture}
\end{center}
{\small\color{brandgray}\textbf{図: 3層が積み上がって初めて実運用に到達する}\ 各層は、その上の層が実運用で信頼できる結果を出すための前提条件である。前処理・検索アーキテクチャ層（\S\ref{sec:data-infra}）が崩れればシンボリック検証層（\S\ref{sec:llm-analysis}）は機能する土台を失い、検証層が崩れればガバナンス・監査可能性層（\S\ref{sec:cross-cutting}）は形式だけのものになる。}

\subsection{実務者への示唆}

\begin{itemize}
  \item \textbf{ベンチマークのスコアだけで導入判断をしない}: 公表された精度・Sharpe比は、評価プロトコルの時間整合性・取引コストモデル化・匿名化の厳密さに大きく依存する。Look-Ahead-Bench\cite{lookaheadbench2026}やKTD-Fin\cite{fromknowingtodoing2026}、77件の系統的レビュー\cite{agentictrading2026}が示すとおり、「高スコア」は必ずしも「実運用で機能する」ことを意味しない。
  \item \textbf{コスト・精度のパレートフロンティアを意識する}: 最高精度のアーキテクチャが最善の選択とは限らない。階層型supervisor-workerアーキテクチャが反射型自己修正アーキテクチャの効果の89\%を60\%のコストで実現する例のように\cite{benchmarkingmultiagent2026}、規模に応じたアーキテクチャ選択がTCOを左右する。
  \item \textbf{データ接続よりデータの解釈・活用に投資する}: MCP標準化によりデータ接続そのものはコモディティ化に向かっている\cite{factset_mcp_nd}。差別化はオーケストレーション・ドメイン特化の前処理・監査可能性に移行しつつある。
  \item \textbf{エージェント標準をガバナンス設計として扱う}: MCP、A2A、AP2は便利な連携仕様ではなく、外部データ接続・エージェント間協調・委任された決済行為の権限証跡を定義する運用基盤である\cite{x01openai2025responsesmcp,x01google2025a2a,x01google2025ap2,x01linuxfoundation2026a2a}。金融機関はこれらを技術選定ではなく監査・権限管理・ベンダー集中リスクの設計問題として評価すべきである。
  \item \textbf{規制の非同期な進展を前提に設計する}: EU AI Act\cite{eba2025aiact}、SR 11-7の置き換え\cite{databricksndmodelrisk}、シャドーAIの開示義務化\cite{wsgrndshadowai}など、地域・分野によって規制の速度と厳格さが異なる。グローバルに展開するシステムほど、最も厳格な規制水準を基準に設計することが結果的にコストを下げる。
  \item \textbf{検証層に棄権・反事実・因果を組み込む}: FinChain\cite{x05finchain2025}、GBFR\cite{x05gbfr2026}、FinCARE\cite{x07fincare2025}が示すように、今後の実務ベンチマークは「正答率」だけでなく、計算経路の検証、安全な回答拒否、反事実サンプル、因果根拠を含むべきである。誤答を減らすだけでなく、答えてはいけない状況を検知する能力が金融AIの実装条件になる。
\end{itemize}

\begin{center}
\begin{tikzpicture}[
  font=\small,
  recbox/.style={text width=7.3cm, align=left, minimum height=25mm},
  r1/.style={fnode, recbox},
  r2/.style={fnode blue, recbox},
  r3/.style={fnode accent, recbox},
  r4/.style={fnode muted, recbox}
]
\node[r1] (b1) at (0,0) {\textbf{評価を疑う}\\[1mm]\footnotesize ベンチマークスコアは評価プロトコルの時間整合性・取引コストモデル化に依存する。「高スコア」は「実運用で機能する」を意味しない。};
\node[r2, right=5mm of b1] (b2) {\textbf{コストを最適化する}\\[1mm]\footnotesize 最高精度が最善とは限らない。階層型アーキテクチャは反射型の効果の89\%を60\%のコストで実現しうる。};
\node[r3, below=5mm of b1] (b3) {\textbf{投資の重心を移す}\\[1mm]\footnotesize データ接続はMCP標準化でコモディティ化へ向かう。差別化はオーケストレーション・前処理・監査可能性に移行する。};
\node[r4, below=5mm of b2] (b4) {\textbf{規制を先取りする}\\[1mm]\footnotesize 規制の速度は地域・分野で異なる。最も厳格な水準を基準に設計すれば結果的にコストを下げる。};
\end{tikzpicture}
\end{center}
{\small\color{brandgray}\textbf{図: 実務者への4つの行動指針}\ 上記4項目の提言を、評価の疑い方・コスト設計・投資の重心・規制対応という4つの観点に整理したもの。いずれも本文で述べた具体的な事例・数値に基づく。}

\subsection{研究者への示唆}

本書で整理した未解決の課題（\S\ref{sec:open-gaps}）のうち、特に投資対効果が高いと考えられる研究方向は以下の3点である。

\begin{enumerate}
  \item \textbf{エンドツーエンド・point-in-time評価基盤の構築}: データ取り込みから投資意思決定までを一気通貫で、かつ時間整合性を保証した形で評価できるオープンなベンチマークは、業界全体の「アルファの誇張」問題を是正する基盤となる。
  \item \textbf{非英語・非米国基準への拡張}: CLEF-2026 FinMMEval\cite{finmmeval2026}が端緒を開いたように、日本語・IFRS・中国会計基準等への評価スイートの拡張は、グローバル投資運用における構造的な死角を埋める。
  \item \textbf{規制監査証跡への標準マッピング}: モデルレベルの説明可能性研究を、MiFID II・SEC受託者責任・EU AI Act・Basel IRBといった具体的な規制要件にマッピングする標準の確立は、AI-Native投資システムの実運用展開を大きく後押しする。
\end{enumerate}

\begin{center}
\begin{tikzpicture}[
  font=\small,
  numcirc/.style={draw=fignavy, fill=fignavy, text=white, figshadow, circle, font=\bfseries,
    minimum size=7mm, inner sep=0pt},
  priobox/.style={text width=13.0cm, align=left, minimum height=15mm, font=\footnotesize},
  p1/.style={fnode, priobox},
  p2/.style={fnode blue, priobox},
  p3/.style={fnode accent, priobox}
]
\node[p1] (b1) at (0,0) {\textbf{\normalsize エンドツーエンド・point-in-time評価基盤の構築}\\[1mm] データ取り込みから投資意思決定までを一気通貫・時間整合的に評価するオープンなベンチマーク。第\ref{sec:open-gaps}部「エージェント・意思決定」の課題に対応。};
\node[numcirc, left=3mm of b1] (c1) {1};

\node[p2, below=5mm of b1] (b2) {\textbf{\normalsize 非英語・非米国基準への拡張}\\[1mm] CLEF-2026 FinMMEval\cite{finmmeval2026}に続き、日本語・IFRS・中国会計基準等への評価スイート拡張。第\ref{sec:open-gaps}部「データ基盤・文書処理」の課題に対応。};
\node[numcirc, left=3mm of b2] (c2) {2};

\node[p3, below=5mm of b2] (b3) {\textbf{\normalsize 規制監査証跡への標準マッピング}\\[1mm] モデルレベルの説明可能性研究とMiFID II・SEC受託者責任・EU AI Act・Basel IRBを結ぶ標準の確立。第\ref{sec:open-gaps}部「システミックリスク・規制」の課題に対応。};
\node[numcirc, left=3mm of b3] (c3) {3};
\end{tikzpicture}
\end{center}
{\small\color{brandgray}\textbf{図: 投資対効果が高い3つの研究方向}\ 第\ref{sec:open-gaps}部で整理した未解決の課題のうち、特に優先度が高いと考えられる3方向と、それぞれが対応する課題カテゴリ。}

\subsection{むすび}

金融データを「AI-ready」にする取り組みと、投資運用を「AI-native」にする取り組みは、しばしば別々の技術トレンドとして語られるが、本書が精査した一次情報は、この二つが不可分であることを繰り返し示している。Bridgewater\cite{bridgewater_aialabs_nd}、Renaissance\cite{institutionalinvestor2025renaissance}、QRT\cite{bloomberg_qube2026}のような実運用に到達した事例は、いずれも高度なデータ基盤への長年の投資を土台にしている。逆に言えば、データ基盤への投資を伴わない「AIエージェント導入」は、DeepFund\cite{deepfund2025}やBankerToolBench\cite{bankertoolbench2026}が示すような性能ギャップに直面するリスクが高い。今後数年の競争優位は、最先端のモデルを最初に採用した企業ではなく、データ基盤・検証層・ガバナンスの三位一体を最も着実に構築した企業にもたらされる可能性が高い。この3層は一度構築して終わりではなく、規制の非同期な進展や評価手法の進化に合わせて継続的に更新すべき動的な基盤である。そしてこの3層は金融固有の発明ではない。文書AI・自律エージェント評価・LLM忠実性はいずれもリーガル・医療・ソフトウェア開発など周辺分野の方が先行しており（\S\ref{sec:cross-cutting-crossdomain}）、金融は成熟した設計・評価手法を「速い追随者」として借用できる立場にある。一方で、評価汚染・エージェント安全性・基盤モデル集中といった汎用AIの未解決課題は、実資本と執行権限を扱う金融でこそ他業界に先んじて最も鋭く顕在化する。金融が自前で確立せざるを得ない特異点は、会計数値の記号的検証と規制密度への対応という2点に絞られる。本書が精査した440件を超える一次情報が示す最大の教訓は、AIモデル自体の進歩を待つことではなく、その進歩を確実な意思決定へと変換する3層の質を、今この瞬間から着実に積み上げることの価値である。

\begin{center}
\begin{tikzpicture}[
  font=\small,
  pathbox/.style={text width=7.0cm, align=center, minimum height=13mm},
  outbox/.style={text width=7.0cm, align=left, minimum height=20mm, font=\footnotesize},
  goodpath/.style={fnode, pathbox},
  badpath/.style={fnode muted, pathbox},
  goodout/.style={fnode blue, outbox},
  badout/.style={fnode accent, outbox},
  arr/.style={farr}
]
\node[goodpath] (g1) {\textbf{データ基盤・検証層・\\ガバナンスへの長年の投資}};
\node[badpath, right=8mm of g1] (b1) {\textbf{データ基盤投資を伴わない\\AIエージェント導入}};

\node[goodout, below=8mm of g1] (g2) {Bridgewater AIA Labs\cite{bridgewater_aialabs_nd}・Renaissance\cite{institutionalinvestor2025renaissance}・QRT\cite{bloomberg_qube2026}のように\textbf{実運用に到達}し、実績あるリターンを継続};
\node[badout, below=8mm of b1] (b2) {DeepFund\cite{deepfund2025}・BankerToolBench\cite{bankertoolbench2026}が示すような\textbf{性能ギャップに直面するリスクが高い}};

\draw[arr] (g1) -- (g2);
\draw[arr] (b1) -- (b2);
\end{tikzpicture}
\end{center}
{\small\color{brandgray}\textbf{図: データ基盤投資の有無で分かれる明暗}\ 本文で挙げた実運用到達事例（左）と性能ギャップに直面した独立検証事例（右）は、いずれもデータ基盤・検証層・ガバナンスへの投資の有無という同じ分岐点に起因する。}

\section{参考文献}
\label{sec:references}

本ホワイトペーパーの記述は、以下の一次情報源（学術論文、企業公式発表、業界ニュース）に基づく。GitHub Issue番号等の内部管理情報は含まず、すべて元記事・原論文のURLを直接参照する形で列挙する。

\nocite{*}
\printbibliography[heading=none]


\end{document}
